\documentclass[11pt,a4paper]{article}
\usepackage[margin=25mm]{geometry}
\usepackage{fontspec}
\setmainfont{Latin Modern Roman}
\setsansfont{Latin Modern Sans}
\setmonofont[Scale=0.87]{Latin Modern Mono}
\usepackage{amsmath,amssymb,amsthm,mathtools}
\providecommand{\mathscr}{\mathcal}
\usepackage{array,longtable,booktabs}
\usepackage{microtype}
\usepackage[dvipsnames]{xcolor}
\usepackage{hyperref}
\hypersetup{colorlinks=true,linkcolor=MidnightBlue,citecolor=MidnightBlue,urlcolor=MidnightBlue,pdftitle={AN GRH: Anabelian point separation and effective Galois bounds},pdfauthor={AN GRH research audit}}
\usepackage{enumitem}
\setlist{itemsep=2pt,topsep=5pt}
\setlength{\parindent}{1.2em}
\setlength{\parskip}{3pt}
\setlength{\emergencystretch}{2em}
\renewcommand{\arraystretch}{1.15}
\newcommand{\et}{\mathrm{\acute et}}
\title{\textbf{AN\_GRH}\\[7pt]Anabelian point separation, cyclotomic certificates,\\and effective Galois bounds}
\author{Research audit of the supplied programme sources\\GPT-6 Astra Pro; run of 21 November 2026}
\date{21 November 2026 (Europe/Berlin)}
\begin{document}
\maketitle
\begin{abstract}
We reconstruct and check an unconditional proof of the curve point-separation
assertion A1 from the recovered later run. The requested Hooley route has a
separate geometric premise: we prove that its universal primitive-root theta
condition is false by a six-prime obstruction on an explicit genus-two curve.
We prove weaker prime-power and composite primitive-factor sufficient
conditions, including an all-genus, all-characteristic rational primitive-factor criterion
for odd cyclic degrees prime to the characteristic,
with exact rank-one examples of degrees $13$ and $15$ and an
ordinary example of degree $4159$. A Lang--Weil test and its fixed-prime
limitation give a curve-dependent numerical route. Under M, a Hooley--Lenstra
calculation supplies a positive-density family of composite candidate orders.
Generic curves on both positive-rank strata avoid every nonzero odd torsion
order. For each fixed finite prime support, all geometric theta torsion is
confined to an explicit finite field; one rank-one example has no nonzero
three-primary theta point. No fixed finite prime support works uniformly over
finite fields, even allowing all exponents and products. A separate ordinary
example defeats the older large-prime quotient test over every constant
extension, despite having a theta witness. Frobenius-orbit and Weil-pairing
arguments give further exact restrictions on candidate primes and successful
numerical tests. A parity improvement gives a separate rational cover test,
and a rank-one curve defeats that size test over every extension. A
prime-power canonical-class criterion supplies prescribed-point rational
primitive factors without size inequalities, and strictly extends the
recovered prime-level detector on an actual curve. A genus-three example
refutes a proposed universal rule for detecting quotients. An exact
degree-seven geometric witness exists on the rank-one numerical-obstruction
curve even though its original rational points all miss the canonical detector.
A quintic reflection construction reduces the genus-two point-section
assertion to ordinary genus-two Jacobians with Rosati-invariant integral real
multiplication by $\mathbf Q(\sqrt5)$. A separate cyclic-cover criterion permits
nonzero rational two-primary Prym torsion of exponent $M$ when
$n>(2M-1)^2$; an exact degree-nineteen example realizes $M=2$.
Every purely two-primary canonical point lies over the field fixing geometric
four-torsion, of degree at most six, and over the degree-six extension in the
real multiplication target case. An ordinary degree-thirteen new factor
separately illustrates the distinction between rational units and geometric
rank zero. The real multiplication class is shown to be nonempty and proper.
Finite-support theta avoidance persists densely on it, including the
ramified prime five, while an actual quintic reflection target has no
rational five-primary canonical component.
An inverse dihedral construction realizes every ordinary real multiplication
curve as a reflection quotient of an etale quintic genus-two cover and
realizes its arbitrary points as actual split transferred images.
A quantitative refinement supplies fixed-field examples where every
lift and every reflection of one quintic cover miss the sufficient
canonical-prime test, including both Weierstrass points on rank-one
sources. A fixed Lang cover gives explicit equidistribution for prescribed
finite primary components and a dense set of canonically detectable target
points, without forcing a prescribed finite transfer fibre to meet that set.
In a tower of characteristic-power constant extensions, an explicit finite
set contains every signed-doubling-eligible torsion prime; all its primary
exponents stabilize in one computable field, although odd theta points
remain infinite. The finite set can contain successful witnesses. In an odd auxiliary-prime
tower, compatible orders instead have at most one new prime and a cofactor
from an explicit finite list; their theta incidence is not supplied.
The bounded-exponent criterion works over the original field with nonconstant
deck group, in every characteristic for genus at most that characteristic,
and otherwise leaves an explicit totally ramified wild elliptic map.
A pointed degree-thirty-nine example has whole Prym exponent two before
deck constantization and sixteen afterwards. Integral semilinear descent
proves equality for each primitive conductor over its own minimal deck field,
while identifying proper-conductor growth as the source of this difference.
An exact Riemann--Roch certificate separately gives three-primary
detectors at both Weierstrass points of a particular rank-one source.
An integral divisor-sign argument eliminates the earlier elliptic exception,
even with nonconstant geometric deck group. The remaining sufficient input
is a pointed rational primitive-factor unit over its specified finite field.
A weighted integral-sign argument further permits a nonzero rational
primitive exponent $M_n$ when $n>(2^{\omega(n)}M_n-1)^2$, in the
range $2\le g\le p$ and with a precise wild elliptic obstruction
in larger genus. This can ignore large proper-conductor torsion. A count of constant twists
gives good unpointed primitive factors, but an exact descent calculation
shows why pointing a different twist restores any original pointed torsion.
Actual norm compatibility and semilinear tower descent bound rational
generator numbers, but leave their orders uncontrolled. Ordinary principally
polarized Weil restrictions show that the listed weights and Frobenius
polynomial shape alone allow exponentially large exponents. A general small
integral-annihilator criterion and a fixed-field lattice construction give
another sufficient bound, with the same wild elliptic obstruction. The bound
requires an estimate on the actual annihilator index; trace relations
factor through earlier levels and do not improve the degree comparison.
Universal existence of such covers, including the residual ordinary real
multiplication case and universal literal primitive-cover existence, remains
open within the requested route.
Under the stipulated Alp\"oge theorem, Bach's applicable general unit-generator
constant is $3$, while exact unconditional coefficient-$2$ certificates are
provided for $53$ moduli and two infinite families. The Masullo draft claim
supplies specified Dedekind-GRH effective Chebotarev consequences, including
explicit full adic semisimplification tests. The audit covers every supplied
flag and all $800$ recovered full register rows, with $591$ available proof
blocks. The missing $390$ full register statements and $209$ further LM7
proof blocks prevent an exhaustive audit of the stated $1190$-row source set.
No proof through Hooley, formal verification of this paper, or verification
of Masullo's unfinished main theorem is claimed.
\end{abstract}
\tableofcontents
\clearpage
\part{Checked findings and precisely scoped cited consequences}
\section{Scope, analytic premises, and what has been established}
\label{sec:premises}

The principal point-separation assertion A1 is proved here without
number-field GRH, by a complete checked reconstruction of the later
ANABELIAN\_A1 argument. The originally proposed passage from Hooley's
primitive-root density to theta-divisor incidence is not a theorem:
A04 identifies the logical gap, and A06 gives an explicit geometric
counterexample to its universal primitive-root version. This changes
the starting diagnosis. It does not turn the counterexample into a
counterexample to A1.

The exhaustive audit remains in progress; the
requested proof through Hooley is also not established. A1 itself
is proved by the different route just specified. In source coverage,
800 of the 1,190 full register statements were recovered, and 591
statement/proof blocks are available. The 210-row
ANABELIAN\_A1\_SUCCESSOR and 180-row GT\_TOWER\_CHECK registers were not
recovered; their visible excerpts are audited separately. Availability
of a proof is not a claim that its whole content was checked. Each
actual check is recorded in \texttt{CHECK\_EVENTS.tsv}. The positive-rank
theta-incidence variant and unrelated open hypotheses in programme
rows remain recorded as open, rather than being assigned consequences
of GRH.

\subsection{The two inputs are deliberately different}

\paragraph{Input A: the stipulated Alp\"oge theorem.}
We accept the cited Dirichlet GRH theorem as instructed. The exact
supplied declarations in \texttt{ALPOGE\_Main.lean} are:
\begin{quote}\small\raggedright
\texttt{riemannHypothesis}\\
\texttt{riemannZeta\_eq\_zero\_re}\\
\texttt{LFunction\_ne\_zero\_of\_half\_lt\_re}\\
\texttt{LFunction\_eq\_zero\_re}\\
\texttt{LFunction\_eq\_zero\_re\_of\_isPrimitive}.
\end{quote}
The definitions in \texttt{ALPOGE\_Statements.lean} quantify over
Dirichlet characters of modulus $N$; the critical-strip assertion
has $0<\Re(s)<1$ and concludes $\Re(s)=1/2$.
The note's summary GRH claim is step~20, p.~4 \cite{Alpoge}.
No other intermediate step of the four-page note is needed in the
proofs below. In particular we do not convert its phrases about
purity into new geometric theorems.

This run inspected those declarations, definitions, and the attached
scan report. It did not build the Lean archive or rescan 175,187
source files. Acceptance as a cited theorem is the convention of
this task, not a newly reported verification of its formal proof.

\paragraph{Input M: the Masullo draft claim.}
Write M, or $\mathrm H_M$, for the draft's claimed GRH for the Hecke
$L$-functions of every number field, as described in \S1.2
\cite{Masullo}. Every use is explicitly conditional on that claim.
The trivial Hecke character over $F$ gives $\zeta_F$, so M includes
Dedekind GRH for every field used below.

\paragraph{Finding S01 (the exact numbered-source boundary).}
The supplied draft does not contain a completed numbered GRH
statement at its intended culminating locator. On printed/PDF
page~522, Chapter~13, Theorem~13.3.1.1 is a trace formula, referring
back to Theorem~12.5.1.1. The next heading is
Theorem~13.3.1.2, \emph{Divisor of $L(s,\chi)$}, but it has no
statement. Its preceding cross-reference is unresolved, and the
footnotes ask for insertion of the statement and proof of the main
theorems. Therefore this report cites M as the claim of \S1.2, with
13.3.1.2 as its unfinished intended location; it never cites the
latter as an already stated and proved GRH theorem.

\begin{proof}
The page was inspected both by direct PDF text extraction and by
rendering the page. The displayed material under 13.3.1.1 is the
trace formula. After the heading 13.3.1.2 the page has only the
editorial notes, and Chapter~14 starts on the following page.
This proves the source-status assertion. It neither supplies a
mathematical refutation of M nor claims a review of all 940 pages.
\end{proof}

\paragraph{Finding S02 (the analytic dependency rule).}
Input A discharges Dedekind GRH for finite abelian extensions of
$\mathbb Q$. M discharges it for arbitrary number fields. Neither
input, as used here, supplies a missing geometric existence theorem,
a missing finite-ramification hypothesis, or a description of an
unspecified finite quotient. No Artin holomorphy conjecture is needed
for the Chebotarev bounds chosen in this report.

\begin{proof}
For an abelian extension $F/\mathbb Q$ the Dedekind zeta function
is the product of the primitive Dirichlet $L$-functions corresponding to the
characters of $\operatorname{Gal}(F/\mathbb Q)$. At an unramified
prime with Frobenius order $f$, the character product of Euler
factors is $(1-p^{-fs})^{-[F:\mathbb Q]/f}$, exactly the Dedekind
factor; the inertia-invariant character factors give the ramified
version. Their nontrivial
zeros therefore satisfy the asserted line condition under A.
For arbitrary $F$, the trivial character specialization of M is
exactly $\zeta_F$. The remaining sentences specify hypotheses not
included in either premise; the concrete failures to make the
corresponding inferences are proved in A04, A06, T02 and G04.
The exact Chebotarev theorem G-CHEB assumes only the indicated
Dedekind zeta RH statements.
\end{proof}

There is a necessary fourth status in the audit: \emph{already
unconditional}. A qualitative Chebotarev theorem, a finite-field
Weil bound, an elementary modular computation, or a previously proved
geometric result does not become unconditional for the first time
because A is added. Calling such a result ``unconditional by Alp\"oge''
would misstate its dependency. The row tables retain the requested
A-only/M-only/beyond-both distinctions for actual analytic upgrades,
and this fourth status for existing statements.

\subsection{The geometric framework and exact locators}

All finite number-field monodromy quotients below are finite etale
covers of an open subscheme of $\operatorname{Spec}\mathbb Z$ (or
$\operatorname{Spec}\mathcal O_K$). The cyclotomic examples are the
covers obtained by adjoining roots of unity. The underlying
fibre-functor and generic-point correspondences are SGA~1,
Expos\'e~V, Th\'eor\`eme~4.1, \S7 and Propositions~8.1--8.2.
For a geometrically connected curve, its arithmetic fundamental-group
exact sequence and point sections are SGA~1, IX, Th\'eor\`eme~6.1
and Corollaire~6.4. Purely inseparable invariance is IX,
Th\'eor\`eme~4.10. These are used with the arithmetic Galois quotient
unchanged.

The configuration argument in A05 uses smooth proper curve lifting
(SGA~1, III, Th\'eor\`eme~7.3 and Corollaire~7.4), with marked
points then lifted by smoothness over the complete henselian base,
proper smooth
specialization (X, Th\'eor\`eme~3.8 and Corollaire~3.9), and the
prime-to-characteristic comparison and inertia of smooth proper pairs
(XIII, Corollaire~2.9, \S2.10, Lemme~2.11). Its surface relation is
XIII, Corollaire~2.12, and its tame Kummer extension at a boundary
uses XIII, Proposition~5.2 \cite{SGA1}. The exact archive context is
cote~149, scans~33--36 for Kummer separation and its covering
criterion, and scans~37--38 for the surrounding geometric reductions.

For compactified curves over finite fields, A05 uses the trace
formula and purity of $H^1$, with Frobenius eigenvalues of weight one:
Deligne, \emph{La conjecture de Weil. I}, formula~(1.5.1) and
Th\'eor\`eme~1.6 \cite{WeilI}. The equivalent action on the
Tate module and its dual must retain the full integral adic module,
not just its reduction modulo a prime. The curve cohomology and
orientation used there are also explicitly given in
\cite[Proposition~14.2 and Remark~14.10]{MilneLEC}.

The coherent Frobenius calculation in A03 and A06 concerns the
Artin--Schreier sequence (\cite[Example~7.9(b)]{MilneLEC}) and the stable part of $H^1(C,\mathcal O_C)$.
It is a characteristic-$p$ calculation of $p$-rank. No theorem about
complex zeros of number-field $L$-functions changes that rank.
Finite quotient representations in G01--G03 are weight-zero Artin
realizations; the original Tate modules may have positive weight.
No new purity theorem, motivic standard conjecture, crystalline
comparison, divided-power identity, or $\lambda$-ring theorem is
asserted. Those additional structures are not inputs to the proofs
in this report.


Findings A07--A51 develop the covering criterion of cote~149,
scans~35--36, inside the Picard variety and finite \'etale covers of
curves (SGA~1, V, Th\'eor\`eme~4.1 and IX, Th\'eor\`eme~6.1).
A08's finite-type spreading uses EGA~IV$_3$, Theorems~8.8.2 and
9.7.7, and EGA~III$_2$, Theorems~7.7.5 and 7.9.4.
A12--A14 and A19--A20 combine the finite \'etale Lang quotient,
coherent Frobenius and the weight-one Jacobian polynomial; the
point-count and absolute-value formulas are
\cite[Theorem~19.1(a)--(c)]{MilneAV} and
\cite[Theorem~11.1]{MilneJV}, with the cohomological formulation
in Weil~I, formula~(1.5.1) and Theorem~1.6.
A15--A18 use cyclic character factors of the Jacobian, coherent
Frobenius, and the characteristic-two \'etale Tate module.
A17 remains a number-field statement about Frobenius conditions
on Kummer covers of an open subscheme of $\operatorname{Spec}\mathbf Z$;
it does not turn those covers into a theta-incidence theorem.

A21--A22 concern the finite \'etale symplectic level torsors of the
relative Picard scheme and their arithmetic descent (SGA~1, V,
Th\'eor\`eme~4.1 and IX, Th\'eor\`eme~6.1). The finite-field
specialization uses effective affine descent and the finite-field
trace estimates specified in A22; it is unrelated to Dedekind GRH.
A23 uses the same weight-one curve and Jacobian Frobenius polynomials
as A13 to test the numerical covering criterion of cote~149,
scans~35--36, in the exact recovered form A1 B.3.

A24 combines that integral Frobenius action with the genus-two
intersection bound; its separate density calculation lies on the
Kummer and quadratic covers of an open subscheme of
$\operatorname{Spec}\mathbf Z$. A25 uses the polarized finite
\'etale torsion sheaf, its Galois-equivariant Weil pairing, and
the same weight-one point bounds. The precise pairing locator is
\cite[\S16 and Lemma~16.2(e)]{MilneAV}.

A26--A29 concern the same pointed finite etale coverings of cote~149,
scans~35--36. Their group-ring sums are character Euler factors of
unramified finite-order sheaves. Genus-two Riemann--Roch truncates their
reduction modulo the characteristic to a point-character sum; the
canonical class controls its inverse-character symmetry. A29 uses
the Picard functor, Lang isogeny and cyclic prime-power norm projector,
with the fibre-functor and arithmetic exact-sequence framework of
SGA~1, V, Theorem~4.1 and IX, Theorem~6.1. A28 is the explicitly
credited recovered successor A.26, not a new prime-degree theorem.

A30 tests the exact genus boundary of that character argument by a
ramified double map to a supersingular elliptic curve and local Euler
factors. Its proposed preliminary double cover is excluded from the
prime-to-characteristic fibre functor. A31 gives an exact principal
divisor, its Frobenius-twisted theta class, and the resulting tame cyclic
cover on the same curve as A27. The geometry remains within the
pointed-cover question of cote~149, scans~35--36.

A32 distinguishes a rational class-number unit from geometric ordinary rank,
using exact character values and the SGA~$4\frac12$ and SGA~7~II comparisons.
A33 makes the characteristic-two genus-two section problem reduce to its
ordinary real multiplication locus by a tame quintic cover and two reflection
maps. A34 admits bounded rational two-primary Prym exponent and realizes
strictness on an exact degree-nineteen cover. A35 confines every purely
two-primary canonical residual point to an explicit small finite field.
These refinements remain in the pointed-cover framework of cote~149,
scans~35--36; none asserts universal literal primitive-cover existence.

A36 gives an actual ordinary quintic reflection target whose rational
Jacobian has no five-primary part, and proves the exact inhomogeneous
Picard correspondence identities. A37 proves generic and dense
finite-field fixed-support avoidance on the ordinary real multiplication
locus itself, using full integral Hilbert--Blumenthal level structure.
The supplements to A33 prove that this locus is both nonempty and a
proper subclass of ordinary genus-two curves. They do not solve the
remaining variable-prime or prescribed-image separation problem.

A38 exhibits $126$ non-Weierstrass points on an ordinary integral
real multiplication curve for which the canonical class has exact
order $4159$, while two has odd order modulo that prime. Thus real
multiplication alone does not force the A29 canonical detector.
Membership in a particular prescribed quintic-transfer image remains
a separate condition, not deduced from this example.

A39 supplies an actual rational three-primary detector for both
Weierstrass points of the A27 rank-one curve, with complete finite
Riemann--Roch certificates on its quintic reflection quotient.
A40 constructs an inverse dihedral cover for every ordinary genus-two
curve with the stated integral real multiplication. It realizes the
bad canonical points of A38 as actual split transferred images.
The sufficient canonical test therefore cannot be forced merely from
membership in that image class; choosing a suitable cover remains open.

A41 bounds the exceptional inverse-construction parameters by thirty-four.
For the A13 curve over $\mathbf F_{128}$, at least ninety-five parameters
give covers, each with at least eleven split source points whose five
reflection images all have canonical class of exact order $4159$.
Changing the lift or reflection within that cover cannot force A29;
choosing a different cover remains an open possibility.

A42 gives an explicit fixed-cover count for each prescribed finite set
of primary components, including the characteristic-primary component.
On the actual Raynaud theta curve it supplies a dense set of odd-order
points and a dense set meeting A29 after rationalizing a useful odd
torsion component. It does not impose whole-order orbit coverage or
incidence with a specified finite transfer fibre. A43 proves that the
sources in the special A41 family have two-rank one and that both
Weierstrass points are rational and split, while every reflection image
over either point has canonical order $4159$. The existential choice
of a different cover is still unresolved.

The A43 supplement proves that its order-three automorphism is absent
on some ordinary real multiplication targets, using A37 with the prime
three and functoriality of the actual Raynaud line bundle. A44 treats a
different attempt: the tower of constant extensions of degrees $2^m$
has an explicit finite set of eligible odd torsion primes, and all their
powers occur in a single computable field. Odd theta points are still
infinite in that tower. A direct divisor identity supplies four
order-three theta points on A13 over $\mathbf F_{16}$, illustrating
that the finite candidate set can be nonempty.

Finding A46 removes the rational primitive-factor qualification in A15/A18.
For every characteristic and every genus at least two, a pointed cyclic
etale cover of odd degree prime to the characteristic suffices if its
actual rational primitive factor has no characteristic-primary torsion.
The deck group need not be constant. Existence of such a cover remains
a separate geometric and arithmetic input.

A45 extends the actual-exponent test to nonconstant deck groups over the
original field, in every characteristic for $2\le g\le p$, and identifies
the precise wild elliptic obstruction in larger genus. Integral semilinear
descent proves equality over each conductor factor's own minimal deck field.
The degree-$39$ twist gives actual whole exponents $2$ and $16$ before and
after full deck constantization. The A44 odd-tower supplement leaves one
new prime and a finite list of possible cofactors; it proves a necessary
restriction, not the existence of theta incidences of those orders.

A47 permits bounded nonzero rational torsion on the actual primitive
factor, using a weighted integral-sign comparison. Its inequality
$n>(2^{\omega(n)}M_n-1)^2$ separates in the range $2\le g\le p$;
in larger genus the remaining elliptic map has degree dividing $M_n$
and at most $g-1$. It recovers A46 when $M_n=1$ in every genus.

A48 gives at least $n-f_n$ good constant twists, where $f_n$ is the
primitive geometric $p$-rank. In genus two this is at least
$n-\varphi(n)>0$. The prescribed fibre chooses one twist;
pointing any other twist over an extension identifies it with the
original pointed cover and restores its rational primitive torsion.
The same cancellation is proved for nonconstant deck groups.

A49 gives an exact original-field Fourier criterion on an odd abelian
quotient which kills the prescribed canonical class. When the quotient's
unit group is generated by the characteristic and inversion, nonuniformity
of the point-image augmentation element is equivalent to a successful
primitive rational-unit cyclic conductor. This tests an actual quotient;
universal existence of the required quotient remains open.


\section{A1: the question, the older route, and the recovered proof}
\label{sec:a1}

\subsection{What A1 says}
The archive contains no statement labelled ``A1'': this is the programme's
label. Cote~149, scan~35, distinguishes the full prime-to-characteristic
fundamental group from its abelian quotient; scan~36 asks whether the
point-separation assertion survives in characteristic $p$, and whether its
covering condition forces $P=Q$. The exact displayed maps denoted $(*)$ and
$(**)$ occur before the supplied scans. We do not reconstruct their omitted
wording. The following is the curve statement of the runs, with all
quantifiers restored from the recovered full manuscript
\cite[A.1, A.31--A.33]{A1run}.

\paragraph{Question A1.}
Let $K$ be finitely generated over $\mathbb F_p$ and let $X/K$ be a smooth
geometrically connected hyperbolic curve. Put
\[
 \Delta_X=\pi_1(X_{\overline K})^{(p')},\qquad
 \Pi_X=\pi_1(X)/\ker\bigl(\pi_1(X_{\overline K})\to\Delta_X\bigr).
\]
The arithmetic quotient is the \emph{whole} $G_K$; only the geometric
kernel is replaced. Is
\[
 X(K)\longrightarrow
 \operatorname{Sec}(\Pi_X\to G_K)/\Delta_X,
 \qquad P\longmapsto[s_P]
 \tag{A1}
\]
injective? In characteristic zero the corresponding assertion uses every
geometric prime. Hyperbolic means $2g-2+r>0$ for the geometric smooth
compactification and its $r$ missing points.

For a proper curve of genus at least two over a finite field, the equivalent
covering formulation is this: for every finite constant extension $L/k$
and every connected geometrically prime-to-$p$ Galois cover $D\to C_L$
with a chosen $P'\in D(L)$ above $P$, there is $Q'\in D(L)$ above $Q$
such that $[Q'-P']\in J_D(L)[p^\infty]$. The term ``$p$-torsion'' in
scan~36 has this primary meaning, as its preceding phrase about order a
power of $p$ confirms. The phrase ``a cover of $X$'' in that scan must be
read as a cover of $Y$ in its local notation; see the transcription table.

\paragraph{Finding A01 (Kummer and the covering formulation).}
The preceding covering condition is equivalent to equality of point-section
classes over a finite field. Equality of sections always implies the
stated primary torsion condition. No number-field GRH occurs here.

\begin{proof}
For a finite field $L$, the prime-to-$p$ Kummer kernel of $J_D(L)$ is its
$p$-primary subgroup: finite Kummer maps inject $J_D(L)/nJ_D(L)$ into
$H^1(L,J_D[n])$, and $J_D(L)$ is finite. Equal sections give the same
action on the fibre functor, hence compatible rational lifts to any
pointed cover, with equal abelian Kummer classes.

For the converse we first establish uniqueness of the specified lift.
If two lifts differ by a nonidentity deck element $a$ of order $m$ prime
to $p$, then $p^e(aR-R)$ is principal. Choose a function with this
divisor. Its norm for $\langle a\rangle$ is constant; over the algebraic
closure scale it to norm one. Hilbert~90 writes it as $a(u)/u$.
Thus $\operatorname{div}(u)-p^eR$ is $a$-invariant. The free deck action
has every orbit of length $m$, so the degree $-p^e$ of this divisor is
divisible by $m$, a contradiction.

Take the cofinal characteristic open normal subgroups $N$ of the finitely
generated geometric group $\Delta_C$. The section $s_P$ makes
$N\rtimes s_P(G_k)$ a pointed cover. Make its geometric deck group
constant by a finite extension and apply the assumed condition.
Uniqueness descends the lift of $Q$ to $k$; norm maps and uniqueness
make these lifts compatible as $N$ varies. They form a fixed point of
the profinite path torsor from $P$ to $Q$, which is exactly a geometric
conjugacy of the two sections. The prime-to-$p$ topology on each cover
is the full one: normal cores embed the additional quotient in a finite
product of prime-to-$p$ groups. This also justifies the use of its
Jacobian Kummer map. This checks the form in \cite[A.31]{A1run}.
\end{proof}

\subsection{The exact contribution of Hooley}

\paragraph{Finding A02 (the analytic consequence, conditional on M).}
For every integer $a\ne-1$ which is not a square, write
$K_n(a)=\mathbb Q(\zeta_n,a^{1/n})$. Under M, Hooley's primitive-root
theorem applies, since the trivial Hecke character of $K_n(a)$ has
$L$-function $\zeta_{K_n(a)}$. In particular,
\[
 \#\{\ell\le x:\operatorname{ord}_{\ell}(2)=\ell-1\}
 = A\operatorname{Li}(x)
   +O\!\left(\frac{x\log\log x}{\log^2x}\right),\qquad
 A=\prod_{r\ {
m prime}}\left(1-\frac1{r(r-1)}\right)>0.
 \tag{A02}
\]
This is a cited consequence of \cite[\S7, concluding theorem]{Hooley}; the precise modern
numbered formulation used for the asymptotic and error is
\cite[Theorem~1.1, specialized to $a=2$]{FanPollack}. We do not re-prove
Hooley's sieve. The needed fields have squarefree $n$; already $K_3(2)$
has nonabelian Galois group $S_3$. The stipulated Alp\"oge Dirichlet
theorems therefore do not supply this application of Hooley.

\begin{proof}
Apply M to the trivial character over each $K_n(2)$, and then the
cited theorem. Its density series is
$\sum_{n\ge1}\mu(n)/[K_n(2):\mathbb Q]=A$. The polynomial $T^3-2$
is irreducible by Eisenstein and has nonsquare discriminant $-108$;
its splitting field is $K_3(2)$ and has group $S_3$. This identifies
a specific nonabelian analytic input absent from the Dirichlet scope.
It does not assert that Dirichlet GRH logically contradicts, or can
never imply by some different argument, the needed Dedekind GRH.
\end{proof}

\paragraph{Finding A03 (what the theta route actually needs).}
For a genus-two curve $C/\overline{\mathbb F}_2$, a sufficient geometric
input is the following. The theta divisor is defined by the bundle of
locally exact differentials on the Frobenius twist $C^{(1)}$; we
transport its geometric line bundles to $C$ by the absolute
Frobenius-twist identification, and denote the resulting divisor by
$\Theta_B$. The required condition is
\[
 \begin{gathered}
 \exists\ell>2\text{ prime},\qquad \operatorname{ord}_{\ell}(2)=\ell-1,\\
 \exists L\in\operatorname{Pic}^0(C)[\ell]\setminus\{0\}
                  \quad\text{with }L\in\Theta_B.
 \end{gathered}
 \tag{T(C)}
\]
Together with the necessary covering implication of A01, this implies
point separation for any arithmetic model of $C$. Hooley supplies the
first condition on $\ell$, not the joint assertion $T(C)$.

\begin{proof}
We spell out the geometric implication, so that the exact missing input
is visible. The line bundle $L$ defines a connected cyclic etale cover
$D\to C$ of degree $\ell$. Its nontrivial character spaces
$H^1(C,L^i)$, $1\le i<\ell$, are one-dimensional by Riemann--Roch.
Absolute Frobenius sends the $i$-space to the $2i$-space semilinearly.
The locally exact differential sequence
$0\to\mathcal O_{C^{(1)}}\to F_*\mathcal O_C\to B\to0$,
tensored with $L^i$, identifies its kernel with
$H^0(C^{(1)},B\otimes L^i)$, with the usual Frobenius-twist
identifications. A theta incidence therefore makes one arrow zero.
The primitive-root assumption makes all nontrivial characters one
cycle, so Frobenius on their sum is nilpotent. The Artin--Schreier
sequence identifies its stable part with etale $\mathbb F_2$
cohomology. Thus the new factor $\operatorname{im}(\sigma-1)$ of
$J_D$ has two-rank zero.

Suppose distinct points collide. After a finite extension A01 gives
lifts $P',Q'$ with $T=[Q'-P']$ two-primary. Since
$(\sigma-1)T$ lies in a two-rank-zero abelian variety, it is zero.
Hence $\sigma Q'+P'\sim Q'+\sigma P'$. Equality of the divisors
is impossible by freeness of the deck action and $P'\ne Q'$.
They therefore give a degree-two pencil on $D$. But $D$ cannot be
hyperelliptic: an odd-order free deck automorphism would induce a
nontrivial automorphism of $\mathbb P^1$, which has a fixed point;
its fibre has at most two geometric points, each fixed by the
odd-order element. This contradicts freeness. The implication follows.
The same proof is the prime-degree instance of the older intake's
A.C2.7 and A.C2.19, without requiring its general composite-conductor
criterion.
\end{proof}

\paragraph{Finding A04 (the density inference is false).}
Infinitude, or positive density, of theta-supporting primes would not
by itself combine with Hooley to prove $T(C)$. Indeed the primes
$\ell\equiv1\pmod8$ have density $1/4$, while both $2$ and $-1$
are squares modulo every such prime. Consequently
$\langle2,-1\rangle\ne\mathbb F_\ell^\times$ throughout that set.
This refutes the proposed inference between properties of sets of
primes; it is not a curve counterexample to A1.

\begin{proof}
Quadratic reciprocity gives $(2/\ell)=1$ and $(-1/\ell)=1$ in the
specified progression. Their subgroup lies in the proper subgroup
of squares. Dirichlet's prime-distribution theorem gives density
$1/\varphi(8)=1/4$. The same obstruction persists at $\ell^a$ by
reduction modulo $\ell$. A fixed finite Chebotarev condition can be
combined with primitive roots using \cite[Theorem~3.1, and
Theorems~4.1--4.2 for positivity]{Lenstra1977}; no such encoding,
or positivity for the curve-dependent torsion condition, is supplied
by M. Goldfeld's average theorem \cite[Theorem, p.~223]{Goldfeld}
averages the integer base, and does not supply it for the fixed base~2.
\end{proof}

There is a second limitation in higher genus. The spaces $H^1(C,L^i)$
then have dimension $g-1$. One singular arrow round a Frobenius orbit
need not make its composite nilpotent: take every arrow to be the
identity except one $\operatorname{diag}(0,1,\ldots,1)$.
Its stable subspace has dimension $g-2$. Thus replacing ``ordinary''
by ``not ordinary'' cannot supply zero primitive $p$-rank. The prime-to-$p$ cover statements in
\cite[Propositions~2.3--2.4]{PopSaidi} do not repair this
step. In the latter source ``pairwise prime orders'' is justified by
a proof of pairwise \emph{coprime} orders, not of prime integer orders.

\subsection{A1 does have a different proof in the later run}

\paragraph{Finding A05 (A1, independently of A and M).}
The assertion A1 holds, using the full arithmetic quotient. The
following proof is a checked reconstruction of
\cite[A.2--A.5, A.14--A.19, A.30--A.33]{A1run}, recovered from the
exact checkpoint identified by the supplied summary. It is not a
proof through Hooley, and no originality or formal verification is
claimed. Its decisive input is cuspidal transport with integral
Frobenius information, together with the already known finite
exceptional-set theorem of Sa\"idi--Tamagawa.

\subsubsection*{Step 1. The relative configuration fibre is retained}
For the finite-field assertion let $k$ be finite, let $C/k$ be proper
of genus at least two, and let $P,Q\in C(k)$. To prove the required
geometric group comparison first work with a complex proper curve
$C$ of genus $g\ge2$, and write
\[
 \Lambda=\pi_1^{\rm top}(C-P),\quad
 \Gamma=\pi_1^{\rm top}(C^2-\operatorname{diag}),\quad
 B_0=\pi_1^{\rm top}(C),\quad \Delta=B_0^{(p')}.
\]
The second projection is a locally trivial fibre bundle, obtained
locally by moving the missing point with an isotopy. Since
$\pi_2(C)=0$, it gives $1\to\Lambda\to\Gamma\to B_0\to1$.
Put $F=\Lambda^{(p')}$, $E=\Gamma^{(p')}$,
$K_F=\ker(F\to\Delta)$ and
$K_E=\ker(E\to\Delta\times\Delta)$. Fix $r\ne p$ and replace
each $K$ by its maximal pro-$r$ quotient. Denote the resulting
groups by $F_c,E_c$. We prove
\[
 1\longrightarrow F_c\longrightarrow E_c\longrightarrow\Delta
 \longrightarrow1. \tag{A05.1}
\]

We need a free-group lemma. If $1\to L\to H\to B\to1$ has $L$
free of finite rank at least two, and $H$ acts on
$L/[L,L]L^r$ through an $r$-group, then $L^{(r)}\to H^{(r)}$
is injective. For a finite $r$-group $Q$, the kernel of
$\operatorname{Aut}(Q)\to\operatorname{Aut}(Q/\Phi(Q))$ is an
$r$-group: it acts freely on the $|\Phi(Q)|^d$ ordered generating
lifts of a fixed basis. Characteristic Frattini quotients therefore
show that the closure of the action on $L^{(r)}$ is pro-$r$.
It extends to $H^{(r)}$ and restricts on $L^{(r)}$ to the inner
action. This action is faithful, because a nonabelian free pro-$r$
group has trivial centre. One direct verification uses its Magnus
embedding into $1+(X_1,\ldots,X_d)$ in the noncommutative formal
power-series algebra over $\mathbb F_r$. Injectivity follows by
substitution into each finite $r$-group algebra, whose augmentation
ideal is nilpotent. A series commuting with both $X_1$ and $X_2$
is scalar, so a central Magnus unit of augmentation one is one.
Compactness makes the injection a closed embedding; hence every
finite $r$-quotient of $L$ is dominated by one of $H$.

Now take any finite quotient $Q$ of $F_c$, with normal $r$-subgroup
$R$ the image of $K_F$, and put $A=Q/R$. The corresponding proper
Galois cover $C'\to C$ gives
$Y=(C'\times C')-(a\times a)^{-1}(\operatorname{diag})$.
The removed components are disjoint labelled graphs. Its fibre
$V'=C'-a^{-1}(P)$ has free fundamental group $L$. In
$H_1(V',\mathbb F_r)$, monodromy is trivial both on boundary
classes and on the quotient $H_1(C',\mathbb F_r)$, since the
compactified family is a product. It is therefore an $r$-group of
upper block unipotent transformations. The lemma supplies a finite
$r$-quotient of $H=\pi_1^{\rm top}(Y)$ dominating the restriction
of $Q$ to $L$. Taking its core under the finite quotient
$\Gamma/H=A\times A$ gives a quotient of $\Gamma$ which is an
$r$-group extended by $A\times A$, still prime to $p$.
Its image of $K_E$ is an $r$-group, so it factors through $E_c$.
On the fibre its kernel is contained in $\ker(F_c\to Q)$.
Every finite quotient of $F_c$ is thus detected in $E_c$, proving
injectivity. Right exactness follows from the universal property of
pro-prime-to-$p$ completion; the further killed subgroup lies in
the fibre kernel. This proves (A05.1).

Lift the smooth proper curve in characteristic $p$ to characteristic
zero (SGA~1, III, Corollary~7.4); its marked points lift by smoothness
over the complete henselian base. Apply the
prime-to-$p$ comparison for smooth proper pairs to the pointed
curve and to its square with diagonal boundary (SGA~1, XIII,
Corollary~2.9, \S2.10 and Lemma~2.11). Functoriality identifies
the groups, the relative kernels, and boundary inertia. Thus
(A05.1) holds for the etale groups. Arithmetic fibre comparison
over a finite field is an isomorphism on geometric kernels and
on the unchanged $G_k$, hence on the whole extensions.

\subsubsection*{Step 2. Equal proper sections transport the cusp}
In the arithmetic configuration quotient the diagonal boundary has
a common decomposition group $\mathcal D$ with
\[
 1\to\mathbb Z_r(1)\to\mathcal D\to\Pi_C\to1,
 \tag{A05.2}
\]
mapping diagonally to $\Pi_C\times_{G_k}\Pi_C$. Here is why no
boundary injection is assumed without proof. The boundary of the
$n$th root construction along the diagonal is the gerbe of roots
of its normal line bundle. Finite tame covers extend on its
Kummer charts by SGA~1, XIII, Proposition~5.2. The inverse-limit
boundary group therefore has kernel the image of tame inertia.
After the relative quotient, this is a quotient of $\mathbb Z_r(1)$.
It is the full group: the boundary word of the punctured fibre is
the product of handle commutators, and sending one handle pair
to $1+e_{12},1+e_{23}$ in $\operatorname{UT}_3(\mathbb Z_r)$
sends it to $1+e_{13}$ of infinite order. Step~1 retains this
inertia. The transverse fibre has multiplicity one, so its cusp
group is the full inverse image of $s_P(G_k)$ in $\mathcal D$.

If $s_P,s_Q$ are geometrically conjugate, lift their conjugator
$\gamma\in\Delta_C$ through $\mathcal D$. Its conjugation
identifies the two fibre groups and cusps, and acts on their first
proper coordinate by $\operatorname{Inn}(\gamma)$. A geometric
element of the target fibre mapping to $\gamma^{-1}$ cancels this
inner automorphism. Consequently we obtain
\[
 \beta:G_P\xrightarrow\sim G_Q
 \quad\text{over }\operatorname{id}_{\Pi_C},
 \tag{A05.3}
\]
where $G_P,G_Q$ are the maximal cuspidally pro-$r$ quotients of
the actual punctured curves. It preserves cusp conjugacy classes.

\subsubsection*{Step 3. Corresponding open covers have the same counts}
If $K\triangleleft F$ is closed and
$N=\ker(K\to K^{(r)})$, every open $I\le F$ containing $N$
satisfies $I^{(r)}\simeq(I/N)^{(r)}$. Indeed take a finite
$r$-quotient $I/J$, put $I_0=\operatorname{core}_F(I)$, and take
the core $M$ of $J\cap I_0$. Then $I_0/M$ is an $r$-group.
The image of $K$ in $F/I_0$ is also an $r$-group because
$N\subset I_0$. Thus its image in $F/M$ is an $r$-group, and
the defining property of $N$ gives $N\subset M\subset J$.

Apply this to a geometric open cover of $G_P$ or $G_Q$. Killing
its actual cusp inertias and taking its pro-$r$ quotient recovers
the full pro-$r$ group of its smooth compactification. This is
equivariant for arithmetic Frobenius. Thus $\beta$ identifies
the integral $r$-adic $H^1$ of corresponding compactifications.
The trace formula gives identical compactified point counts over
every finite extension; the cusp double cosets give identical
degrees of the missing points. Subtracting them gives identical
counts for the open covers. The inputs are the usual curve trace
formula and purity, not analytic RH over $\mathbb Q$.

\subsubsection*{Step 4. Compactness recognizes actual point sections}
Fix $R\in(C-P)(L)$ and let $S=\beta(D_R)$. For every open
subgroup $H$ of $(G_Q)_L$ containing $S$, the source cover for
$\beta^{-1}H$ has an $L$-point above $R$. Step~3 shows that the
corresponding target cover $X_H$ has an $L$-point. Its finite
sets $X_H(L)$ have nonempty inverse limit: any finite collection
of compatibility conditions is satisfied by a point on the cover
for the intersection of the corresponding subgroups. Compactness
then supplies a compatible family, and a lift of its base point
$T\in(C-Q)(L)$ to the full profinite fibre. Its decomposition
group lies in every $H$ (rationality at each level, and injectivity
of the decomposition group's projection to $G_L$), hence in
$\bigcap_{H\supset S}H=S$. Both project isomorphically to $G_L$,
so they are equal. Thus $\beta$ transports every point section
to a genuine point section.

Sa\"idi--Tamagawa \cite[Definitions~1.3, 1.5, 1.10 and
Proposition~1.8(vi)]{ST2009} give one finite exceptional image
$E_C\subset C^{\rm cl}$ for the cofinite prime set $p'$.
Outside it the proper decomposition group uniquely determines
the geometric point, even up to commensurability. Remark~1.2
makes this one exceptional set valid over every constant extension.
Since $\beta$ is over the identity of $\Pi_C$, the transported
point is $R$ whenever $R\notin E_C\cup\{P,Q\}$.

\subsubsection*{Step 5. The Kummer difference is impossible}
The proper abelianization gives $[P-Q]\in J_C(k)[p^\infty]$.
Choose $f\in k(C)^\times$ with
$\operatorname{div}(f)=p^e(P-Q)$; descent of a geometric
generator follows from Hilbert~90. Put $M=\mathbb Z_r(1)$.
The geometrically cuspidally central quotients $E_P,E_Q$ have
kernel $M$. Infinite inertia follows from the same Heisenberg
test, and their geometric extension classes are the degree-one
point classes in $H^2(\Delta_C,M)=\mathbb Z_r$.
The map induced by $\beta$ acts on $M$ by a unit $u$.
Since it fixes the proper quotient and the degree-one class,
$u=1$.

For completeness the orientation identification uses the
punctured surface presentation (SGA~1, XIII, Corollary~2.12).
Equivalently it is the Kummer first Chern class of
$\mathcal O_C(P)$, with residue $+1$. The group/etale
cohomology comparison follows from the tower of all geometric
prime-to-$p$ covers: cyclic torsors kill $H^1$ classes, covers
of degree divisible by $r^a$ kill $H^2(-,\mu_{r^a})$, and a
smooth proper curve has no higher such cohomology. Taking the
associated spectral sequence and then inverse limits gives
the asserted $\mathbb Z_r(1)$ identification.

On $E_P\times_{\Pi_C}E_Q$ the difference
$\chi(u,v)=\overline\beta(u)v^{-1}$ is an $M$-valued
crossed homomorphism. Pull it back to the group of
$C-\{P,Q\}$. Its residues are $(1,-1)$, so
$\kappa_r(f)-p^e\chi$ has zero residues and descends to
$H^1(\Pi_C,M)$. Inflation-restriction identifies this group
with $H^1(G_k,M)$: the Frobenius eigenvalues on
$\operatorname{Hom}(T_rJ_C,\mathbb Z_r(1))\otimes\mathbb Q_r$
have absolute value $\sqrt{\#k}$, hence no invariant, and its
integral module is torsion free. Finite-field Kummer theory
therefore gives $a\in k^\times$ with
\[
 \kappa_r(f)=p^e\chi+\kappa_r(a). \tag{A05.4}
\]

At a good closed point $R$ put $L=k(R)$. Step~4 makes
$\beta(D_R^P)$ conjugate to $D_R^Q$. Its conjugator projects
to the normalizer in $\Pi_{C_L}$ of the proper section $D_R$, which is $D_R$
itself by \cite[Lemma~1.7]{ST2009}. Multiplying by an element
of $D_R^Q$ puts it in the cuspidal kernel; in $E_Q$ it lies
in $M$. Thus $\chi$ evaluates to a coboundary, and (A05.4)
gives $\kappa_{r,L}(f(R)/a)=0$.

If $P\ne Q$, the function $f$ is nonconstant. The excluded
points and their $f$-values form a finite set. Choose a
nontrivial $r$-power root of unity $\zeta\in\overline k$
with $a\zeta$ outside that set, and a point $R$ above it.
Then $\zeta=f(R)/a\in k(R)^\times$. The $r$-adic Kummer
map of this finite cyclic group is injective on its $r$-primary
subgroup, contradicting its asserted zero value. This proves
injectivity over every finite field, including every constant
extension.

\subsubsection*{Step 6. Finitely generated fields and open curves}
For a finitely generated field $K/\mathbb F_p$, spread a proper
curve and two distinct colliding points to a smooth connected
finite-type base, shrinking so that the sections are disjoint.
At a closed point use its henselian local ring $R$ and fraction
field $L$. Smooth proper specialization (SGA~1, X,
Theorem~3.8 and Corollary~3.9) identifies the geometric
prime-to-$p$ groups. The arithmetic map $G_L\to G_{k(s)}$
is onto, since finite etale covers of the normal henselian
domain have connected generic fibres. The same conjugator
would therefore give colliding distinct points on the finite
residue-field curve, a contradiction.

We record also the Kummer fact needed for open curves. For a
semiabelian $G$ over $K$ the prime-to-$p$ Kummer kernel is
$G(K)[p^\infty]$; in characteristic zero, using all primes,
it is zero. After a finite splitting extension of its torus,
the divisor sequence on a normal proper model makes
$K^\times$ an extension of constant units by a free abelian
group. Constant units are finite over finite constants, and
finite torsion plus free abelian over number-field constants.
Lang--Neron (\cite[Theorem~2.1]{Conrad}, with the perfect constant
field algebraically closed in $K$), together with Mordell--Weil on the trace over
number fields, makes the abelian quotient's rational points
finitely generated. Lifting its finitely many generators and
relations shows that $G(K)$ is free abelian plus a finitely
generated group; the intersections of its prime-to-$p$
multiples and of all its multiples are respectively its
$p$-primary torsion and zero. This is the argument of
\cite[A.30]{A1run}; it does not falsely assume $K^\times$
finitely generated.

Make a finite algebraic extension to split the compactification
and boundary; purely inseparable base change leaves the etale
category and Galois group unchanged (SGA~1, IX,
Theorem~4.10). If the compactified genus is at least two,
map a collision to the proper curve. In genus zero, choose
three boundary points as $0,1,\infty$; their coordinate embeds
the curve into $\mathbb G_m$, and equal Kummer classes have
ratio of $p$-power order in $K^\times$, hence ratio one.
In genus one use a rational point as origin of the compact
elliptic curve $E$. A collision difference has order $p^a$.
Pull back the open curve under $[m]$, with $(m,p)=1$ and
$m^2>p^a$, and split its boundary and deck group by finite
extension. Lift the collision with first point zero. Its class
in the generalized Jacobian has exact order $p^a$: the torus
kernel has no nonzero geometric $p$-primary points. Thus a
function with divisor $p^a(Q'-P')$ equals one at all points
of the pulled-back reduced boundary. There are at least
$m^2$ of them, whereas $f-1$ has pole degree $p^a$, impossible.
In characteristic zero the all-prime Kummer kernel is zero
and the Abel map to the generalized Jacobian separates points.
This proves A1 with the stated scope. It does not assert the
full higher-dimensional or section-surjectivity conjecture.

\paragraph{Consequence for the two inputs.}
A05 needs neither A nor M. Therefore it remains true if only
Alp\"oge's stated theorems are added to the existing background.
Hooley's number-field application is still conditional on M;
the theta-support and primitive-cover questions are separate
geometric assertions. The older intake's unresolved A1 status
has been superseded by the recovered later proof, not by an
unproved implication from primitive-root density.


\subsection{An explicit obstruction to the proposed Hooley bridge}

\paragraph{Finding A06 (the universal primitive-root theta condition is false).}
The condition $T(C)$ of A03 does not hold for every genus-two curve in
characteristic two, even after arbitrary finite extension of its field of
definition. Take the smooth proper model of
\[
 \begin{split}
 k&=\mathbb F_2[t]/(M(t)),\quad
 M(Z)=Z^{38}+Z^{14}+Z^{10}+Z^9+Z^8+Z^5+Z^2+Z+1,\\
 d&=t^3+1,\qquad
 C:\ y^2+y=A(u)^3,\qquad
 A(u)=(u^2+t^2u+t)/d.
 \end{split} \tag{A06.1}
\]
For its locally exact differential theta divisor $\Theta$, there is no
odd prime $\ell$ with $\operatorname{ord}_\ell(2)=\ell-1$ and
$\Theta\cap J_C[\ell](\bar k)\ne\{0\}$. This holds geometrically;
number-field GRH cannot change it.

\begin{proof}
The curve and scalar Frobenius are those of \cite[B.23]{A1run}.
We recall the algebraic data used here and give the new small
primitive-root obstruction in full. Rabin's irreducibility test gives
$Z^{2^{38}}=Z$ modulo $M$ and
$\gcd(M,Z^{2^2}-Z)=\gcd(M,Z^{2^{19}}-Z)=1$.
Thus $k$ is the stated field and $d\ne0$. Reduction of the
Artin--Schreier polynomial by even powers leaves leading odd term
$(t^2/d^3)u^5$. There is a unique point $P$ at infinity, the genus
is two, and $K_C\sim2P$. The one-pole Artin--Schreier curve has
two-rank zero (equivalently, Cartier on the basis $du,u\,du$ has
stable rank zero). Hence $J_C[2](\bar k)=0$ and the unique geometric
theta characteristic is $\mathcal O(P)$. In characteristic two
the bundle of locally exact differentials is this theta characteristic,
so its theta divisor is $\Theta=C-P$. With relative Frobenius this
statement is on $C^{(1)}$; the same equations with coefficients squared
and the same scalar Frobenius give exactly the argument below.
In particular it is symmetric
and contains zero.

Put $B(u)=(tu^2+u+t^2)/d$. There are two maps to
$E:y^2+y=x^3$, given by $(A(u),y)$ and
$(B(u),y+(u^3+1)/d)$. Their equations follow from
\[
 (u^2+t^2u+t)^3+(tu^2+u+t^2)^3
 =d(u^3+1)^2+d^2(u^3+1).
\]
Their joint map is birational onto its image, since $u=B+tA$.
The image, whose normalization has genus two, cannot lie in an
elliptic translate in $E^2$; thus $J_C\to E^2$ is an isogeny.
Both maps take $P$ to zero. The duplication formula on $E$ is
$[2](x,y)=(x^4,y^4+1)$, so its fourth-power Frobenius is $[-2]$.
It follows that the $k$-Frobenius of $J_C$ is $[-2^{19}]$.
Indeed composing the difference with the isogeny gives zero,
and a homomorphism from a connected abelian variety into a finite
kernel is zero. These geometric identities and the irreducibility
calculation were separately recomputed in this run.

A nonzero difference in $J_C$ has at most two ordered representations
as a difference of points on $\Theta$. For if
$R_1-S_1=R_2-S_2\ne0$, then
$R_1+S_2\sim R_2+S_1$. Either the effective divisors coincide,
which gives the same representation, or they belong to the unique
degree-two hyperelliptic pencil. The latter determines the sole
alternative representation $(\iota S_1,\iota R_1)$.

Suppose a primitive-root prime $\ell$ and a nonzero
$L\in\Theta[\ell]$ existed. Since the theta divisor is defined over
$k$ and symmetric, it contains $hL$ for
$h\in H=\langle2^{19}\rangle\subset\mathbb F_\ell^\times$.
Write $e=\gcd(\ell-1,19)$ and $m=(\ell-1)/e$.
It contains these $m$ distinct points and zero. Counting their
ordered nonzero differences in the cyclic subgroup gives
\[
 m(m+1)\le2(\ell-1). \tag{A06.2}
\]
If $e=1$, this forces $\ell\le2$, impossible. If $e=19$,
it forces $m\le37$, hence $\ell\le704$ and
$\ell\equiv1\pmod{19}$. Exhaustive trial division of the eighteen
integers $19m+1$ with $m=2,4,\ldots,36$ gives exactly
\[
\begin{array}{c|rrrrrr}
 \ell&191&229&419&457&571&647\\ \hline
 \operatorname{ord}_\ell(2)&95&76&418&76&114&323.
\end{array}
\]
Only $419$ remains. Here $g=2^{19}\bmod419=119$, and $H$ contains
\[
 (g^0,g^{11})=(1,418),\quad
 (g^{19},g^{20})=(71,69),\quad
 (g^9,g^8)=(350,348).
\]
These are three distinct ordered pairs with the same nonzero
difference $2$ modulo $419$, contradicting the two-representation
bound. The included standard-library certificate
\texttt{check\_theta\_obstruction.py} verifies every finite calculation.
\end{proof}

\paragraph{Scope of the refutation.}
The stronger source statement \cite[B.23]{A1run} excludes every
nontrivial odd cyclic etale cover with zero primitive two-rank on this
curve. Its exhaustive conductor reduction and its complete polynomial
B\'ezout certificates at $191$ and $229$ were rerun successfully;
the verifier, data and logs are included. The shorter A06 proof does
not depend on those large certificates. In particular the weaker
arithmetic condition $\langle2,-1\rangle=\mathbb F_\ell^\times$
cannot repair the universal geometric assertion on this example,
by that checked B.23 result.

This curve already has two-rank zero, so ordinary Kummer theory
separates its points. It is not a counterexample to A1, nor a
refutation of a strategy which first removes all zero-rank curves
and then seeks a different theorem for the remaining curves.
The positive-two-rank theta-incidence problem has not been solved
by this argument. Its unresolved status does not affect A05.


% Follow-up supplement. Standalone finding IDs reserved by root: A07,A09,A10,A11.
% A08 is supplied separately by effective_galois/finite_prime_avoidance.tex.
% Bibliography additions are in POSITIVE_RANK_BIB.tex.

\section{Positive rank and a weaker geometric bridge}
\label{sec:positive-rank-followup}

\paragraph{Finding A07 (the generic positive-rank bridge fails).}
For each $f\in\{1,2\}$ there is a smooth proper genus-two curve over a
field finitely generated over $\mathbb F_2$, of geometric $2$-rank $f$,
whose Raynaud theta divisor has no nonzero point of any odd prime order.
Consequently the assertion $T(C)$ used in A03 is false on this larger
class of ground fields, even after imposing positive $2$-rank.

\begin{proof}
Choose an irreducible component $Z$ of the nonempty stratum
$\mathcal M_2^f$ and a curve $C_\Omega$ representing its geometric
generic point. Ozman--Pries, Theorem~1.1(2), equivalently
Theorem~8.3, applies with $p=2$, $g=2$, and either $f=1$ or $f=2$:
for every odd prime $\ell$, its theta divisor contains no point of
order $\ell$ \cite[pp.~456,473]{OzmanPries2019}.
These assertions apply to the same geometric generic curve for all
$\ell$; they do not require choosing a different curve for each prime.
The curve and a projective embedding have finitely many coefficients,
so they descend to a finitely generated subfield $K\subset\Omega$.
Smoothness and geometric connectedness descend under this field
extension. The $2$-rank and the absence of geometric theta torsion
are preserved by algebraically closed extension, so this curve over
$K$ has the claimed properties. Any required rational base point can
also be included among the finitely many coefficients.
\end{proof}

This conclusion does not give a counterexample over
$\overline{\mathbb F}_2$: the generic curve has transcendental moduli.
Finding A08 gives the separate finite-field consequence with its
correct quantifiers, namely avoidance of each prescribed finite set
of primes by a curve that may depend on that set.
For the ordinary generic curve, the stronger statement that all
prime-to-$2$ cyclic etale covers are ordinary is the Nakajima result
recorded as Proposition~4 in Yang's Theorem~1.1
\cite[Theorem~1.1]{Yang2021}. No restriction $p>2$ occurs there.
For a generic positive-rank stratum, Yang's Section~3.6.2 and
Corollary~3.6 explicitly retain prime degree; they do not supply a
prime-power version for rank one.

\paragraph{Source qualification.}
The original Nakajima paper is \emph{On generalized Hasse--Witt
invariants of an algebraic curve}, Adv. Stud. Pure Math.
\textbf{2} (1983), 69--88. Its original Proposition~4/Theorem~2 was
not recovered in this follow-up. The generic assertion above is
verified in Yang's paper and, for prime degree, directly in
Ozman--Pries. The latter's Theorem~8.2(3) is not used: its summary of
Pop--Saidi omits their finite-field hypothesis and describes prime
degrees, whereas the cited argument establishes pairwise coprime
orders, not prescribed prime orders. The finite-primary projection
argument in Pop--Saidi, Section~2, is the evidence used below.
\cite[Propositions~2.1,2.3 and their proofs, pp.~112--114]{PopSaidi}
In particular, Yang's stated all-cyclic-degree result also rules
out the sufficient prime-power incidence condition of A11 on a
generic ordinary curve. This stronger consequence is attributed to
the verified statement in Yang, not to an inspection of Nakajima's
unavailable original text. No corresponding all-prime-power claim
is made for the generic rank-one stratum.
There is also a direct modern theorem to cite: Yang's
Theorem~3.5 allows any integer $m$ prime to $p$ and says that a
cyclic cover of a geometric generic pointed curve is new-ordinary
exactly when the associated ramification divisors are Frobenius
stable. In the unpointed etale case every such divisor is zero,
which satisfies Definition~3.3 and Lemma~3.2 automatically.
This verifies the all-odd-degree conclusion at $p=2$ without
claiming access to Nakajima's original paper
\cite[Definition~3.3, Theorem~3.5]{Yang2021}.



\subsection{No uniform finite list of primes in positive 2-rank}
\label{sec:finite-prime-avoidance}

\paragraph{Published input (Ozman--Pries).}
Let $\ell\ne p$ be prime, let $g\ge2$, and let $0\le f\le g$, with
$f\ne0$ when $g=2$.  For each irreducible component $Z$ of
$\mathcal M_g^f$, the Raynaud theta divisor of the geometric generic
curve of $Z$ contains no point of order $\ell$.  This is
Ozman--Pries, \emph{Ordinary and almost ordinary Prym varieties},
Asian J. Math. \textbf{23} (2019), 455--478,
Theorem~1.1(2), p.~456, equivalently Theorem~8.3, p.~473 \cite{OzmanPries2019}.
In particular its hypotheses allow $p=2$, $g=2$, and either $f=1$ or
$f=2$.  The assertion concerns every geometric point of order $\ell$,
not just a selected cyclic cover.

\paragraph{Finding A08 (finite prime avoidance).}
For either $f\in\{1,2\}$ and every finite set $S$ of odd primes, there
exist an integer $m\ge1$ and a smooth projective geometrically connected
curve $C/\mathbf F_{2^m}$ of genus two and $2$-rank $f$ such that
\[
 \Theta_C(\overline{\mathbf F}_2)\cap
 \bigl(J_{C^{(1)}}[\ell](\overline{\mathbf F}_2)\setminus\{0\}\bigr)
 =\varnothing\qquad(\ell\in S).
 \tag{FP.1}
\]
Here $C^{(1)}$ is the relative Frobenius twist, and $\Theta_C$ is the
theta support of the sheaf of locally exact differentials on $C^{(1)}$.
This is the convention of Ozman--Pries, Section~8.1, p.~472.
One can equivalently use the usual Frobenius-semilinear relabelling of
this support; the existence and order of its geometric torsion points
are unchanged.  For $f=2$ the curve is ordinary.

\paragraph{Recovered source scope.}
The ordinary finite-list obstruction already appears in recovered
Successor finding A.16 \cite[A.16]{ASrun}, which allows arbitrary
characteristic and genus.  Its full recovered proof uses generic
ordinary Pryms and finite cover moduli.  The argument below states the
geometric theta-support conclusion explicitly, gives a scheme-family
spreading proof without coarse-moduli descent, and also treats the
fixed $2$-rank-one stratum.

\begin{proof}
\emph{1. A curve at a geometric generic point.}
The strata $\mathcal M_2^1$ and $\mathcal M_2^2$ are nonempty.  For
example, the smooth projective models of
\[
 y^2+y=x^3+x^{-1},\qquad
 y^2+y=x+x^{-1}+(x-1)^{-1}
\]
have genus two and $2$-ranks one and two, respectively.  For an
Artin--Schreier equation with $r+1$ poles of odd orders $d_j$, the
genus is $\bigl(\sum_j(d_j+1)-2\bigr)/2$ and the $2$-rank is $r$;
see Pries--Zhu, \emph{The $p$-rank stratification of Artin--Schreier
curves}, Ann. Inst. Fourier \textbf{62} (2012), 707--726,
Lemma~2.6, p.~712 \cite{PriesZhu2012}.

Fix an irreducible component of $\mathcal M_2^f$ over
$\overline{\mathbf F}_2$ and a curve $C_\Omega$ at its geometric generic
point, where $\Omega$ is an algebraically closed field of characteristic
two.  Apply the published input to each $\ell\in S$, on this same
component and this same geometric generic curve.  Thus $C_\Omega$
satisfies all the required avoidance conditions simultaneously.

\emph{2. Spreading the curve while preserving its $2$-rank.}
Choose a point $P\in C_\Omega(\Omega)$ and a projective embedding.
The equations and the point involve finitely many coefficients.
They therefore descend to an integral finitely generated
$\mathbf F_2$-subalgebra $A\subset\Omega$; the precise limit reference
for the scheme and the marked-point morphism is EGA~IV$_3$,
Theorem~8.8.2(ii) and (i), printed p.~28, respectively;
see also the Stacks Project, Lemma~32.10.1(1)--(3),
Tag~\texttt{01ZM} \cite{EGAIV3,Stacks}.
After replacing
$T=\operatorname{Spec}A$ by a nonempty affine open, they give a smooth
projective family $\pi:\mathcal C\to T$ of geometrically connected
genus-two curves with a section, whose base change to $\Omega$ is
$(C_\Omega,P)$.  This uses the usual finite-presentation spreading
argument: the projective model is proper, the nonsmooth locus has
closed image missing the generic point, and we remove that image.
The geometric-connectedness locus is constructible by EGA~IV$_3$,
Theorem~9.7.7(ii), p.~79; since it contains the generic point of
the integral base, it contains a nonempty open.  The genus is
preserved by constancy of the Euler characteristic in this smooth
proper family, EGA~III$_2$, Theorem~7.9.4, pp.~76--77.
There is no need to choose a universal curve over a coarse moduli
space or to descend a generic curve to that coarse space's function
field.

We may also ensure that every geometric fibre has $2$-rank $f$.
For completeness, locally trivialize the rank-two bundle
$E=R^1\pi_*\mathcal O_{\mathcal C}$ and let $H$ be the matrix of its
Frobenius-semilinear endomorphism.  On a geometric fibre, the stable
rank of Frobenius is the $2$-rank, and in dimension two the image has
stabilized by the second iterate.  Hence the $2$-rank is
$\operatorname{rank}(H H^{(2)})$, where $H^{(2)}$ denotes entrywise
squaring.  All minors of size $f+1$ vanish at the generic point and
thus vanish identically on the integral base; some minor of size $f$
is nonzero at the generic point.  Removing its zero locus gives the
assertion.  Equivalently one can use the locally closed $p$-rank
strata recalled in Ozman--Pries, Section~2.5, p.~458.

\emph{3. The relative theta incidence is closed.}
Put $\mathcal C^{(1)}=\mathcal C\times_{T,F_T}T$, and let
$F:\mathcal C\to\mathcal C^{(1)}$ be relative Frobenius.
For a smooth relative curve in characteristic two it is finite locally
free of degree two, and
\[
 0\longrightarrow\mathcal O_{\mathcal C^{(1)}}
 \longrightarrow F_*\mathcal O_{\mathcal C}
 \longrightarrow\mathcal B\longrightarrow0
 \tag{FP.2}
\]
has a line-bundle quotient $\mathcal B$, compatible with base change.
These claims can be checked in a relative etale coordinate, where
$F_*\mathcal O$ has basis $1,t$ over the Frobenius-twisted coordinate
ring.  On each fibre, $\mathcal B$ is the sheaf used in the definition
of the Raynaud theta divisor.

Let $\mathcal J=\operatorname{Pic}^0(\mathcal C^{(1)}/T)$, an abelian
scheme.  The chosen section gives a Poincare line bundle $\mathcal P$
on $\mathcal C^{(1)}\times_T\mathcal J$, normalized along that section.
Upper semicontinuity of cohomology for its proper projection to
$\mathcal J$ shows that
\[
 Z_\Theta=
 \{z\in\mathcal J:
 h^0((\mathcal C^{(1)}\times_T\mathcal J)_{\overline z},
       (\operatorname{pr}_1^*\mathcal B\otimes\mathcal P)_{\overline z})>0\}
 \tag{FP.3}
\]
is a closed subset, where $\overline z$ denotes a geometric point
over $z$.  Its geometric fibres over $T$ are
exactly the supports of the relevant theta divisors.  Thus the
argument requires only this closed incidence locus, and does not
require a separate global assertion that a chosen relative theta
divisor is Cartier.  The precise semicontinuity theorem is
Grothendieck, EGA~III$_2$, Theorem~7.7.5(I), printed p.~67;
equivalently, the Stacks Project, Lemma~36.32.1, Tag~\texttt{0BDN} \cite{EGAIII2,Stacks}.
Its hypotheses hold because the base is noetherian, the projection
is proper and smooth, and
$\operatorname{pr}_1^*\mathcal B\otimes\mathcal P$ is a line bundle,
hence coherent and flat over $\mathcal J$.

\emph{4. All nonzero torsion points are removed at once.}
For each odd prime $\ell$, the group scheme $\mathcal J[\ell]$ is
finite etale over $T$.  Its zero section is both open and closed, so
\[
 \mathcal J[\ell]^\times:=\mathcal J[\ell]\setminus0_T
\]
is still finite etale over $T$.  The intersection
$Z_\ell=Z_\Theta\cap\mathcal J[\ell]^\times$ is closed in this finite
scheme.  Consequently its image $D_\ell\subset T$ is closed, because
finite morphisms are closed.  A geometric fibre over $t$ contains a
nonzero $\ell$-torsion point on its theta support if and only if
$t\in D_\ell$.

By Step~1, $D_\ell$ misses the generic point of $T$ for every
$\ell\in S$: otherwise its nonempty generic fibre would remain
nonempty after the field extension to $\Omega$.  Therefore
\[
 U=T\setminus\bigcup_{\ell\in S}D_\ell
\]
is a nonempty open subscheme, containing the generic point.
It is of finite type over $\mathbf F_2$, so it has a closed point $u$
with finite residue field $\kappa(u)=\mathbf F_{2^m}$, by the
Nullstellensatz.  The curve $C=\mathcal C_u$ has genus two and
$2$-rank $f$, and satisfies (FP.1) over an algebraic closure of its
residue field.  This proves the claim.
\end{proof}

\paragraph{Ordinary strengthening of A08 (any finite list of odd orders).}
For $f=2$, the finite set $S$ may be replaced by any finite set
$\mathcal N$ of odd integers $n\ge3$: one can require
\[
 \Theta_C(\overline{\mathbf F}_2)\cap
 \bigl(J_{C^{(1)}}[n](\overline{\mathbf F}_2)\setminus\{0\}\bigr)
 =\varnothing\qquad(n\in\mathcal N).
 \tag{FP.4}
\]
In particular there is no uniform finite list of odd prime-power
orders for the sufficient incidence condition in A11.

\begin{proof}
Yang, \emph{Raynaud--Tamagawa theta divisors and new-ordinariness of
ramified coverings of curves}, J. Algebra \textbf{587} (2021),
263--294, Theorem~3.5 \cite{Yang2021}, applies to arbitrary cyclic degree $m$
prime to the characteristic over a generic pointed stable curve.
For an unpointed smooth curve every ramification divisor $D(j)$
is zero, which is Frobenius stable by Definition~3.3 and
Lemma~3.2.  Hence every connected cyclic etale cover of the
generic genus-two curve, of any odd degree, has ordinary new
factor.  The generic curve itself is ordinary, as also recorded
in Yang's Theorem~1.1.  These statements concern this same
geometric generic curve.

To verify that this excludes odd-order theta points, suppose $L$
has exact odd order $d>1$ on its Raynaud theta support, viewed on
the Frobenius twist.  Choose $e$ with $2e\equiv1\pmod d$ and set
$M=(F^*L)^e$ on the original curve.  Then $M$ has order $d$,
$M^2=F^*L$, and $M^{(1)}=L$; the last equality follows by applying
$F^*$, which is injective on odd-order torsion.
The twisted Frobenius sequence
\[
 0\longrightarrow L\longrightarrow F_*F^*L
 \longrightarrow B\otimes L\longrightarrow0
\]
therefore identifies $H^0(B\otimes L)$ with the kernel of the
semilinear Frobenius arrow
$H^1(C,M)\to H^1(C,M^2)$.
The two spaces are lines, since $C$ has genus two and the line
bundles are nontrivial of degree zero.  This arrow is zero.
In the connected degree-$d$ Kummer cover defined by $M$, it is
one arrow in the nontrivial-character sum of $H^1(\mathcal O)$.
That sum is Frobenius stable, and its Frobenius is not injective,
so its stable rank is smaller than its dimension.  Its new
factor is therefore not ordinary, contradicting Yang's theorem.
Thus the geometric generic theta support contains no nonzero
point of any odd order.

Now use this generic ordinary curve in Steps~2--4 above.  For each
$n\in\mathcal N$, the scheme $\mathcal J[n]\setminus0_T$ is finite
etale.  Its theta incidence has closed image missing the generic
point, because every nonzero element has odd order dividing $n$.
Remove the finitely many images and choose a closed point of the
remaining open.  This gives (FP.4).
\end{proof}

This stronger odd-order assertion is stated only for the ordinary
case $f=2$.  Yang's Section~3.6.2 and Corollary~3.6 retain a prime
degree hypothesis on a generic fixed lower-$p$-rank stratum; no
prime-power extension for $f=1$ is inferred here.

\paragraph{The precise nonuniformity conclusion.}
Taking $S$ to be all odd primes at most an arbitrary real number $B$
shows that no bound depending only on characteristic two, genus two
and the fixed positive $2$-rank can guarantee a nonzero prime-order
theta point below that bound for every curve over finite extensions
of $\mathbf F_2$.  More generally, no fixed finite list of candidate
odd primes works for all those curves.  The extension degree $m$ is
allowed to depend on the list; this proof gives no numerical bound
for $m$ or for the defining equations.

This statement has quantifiers
\[
 \text{for every finite }S\quad\text{there exists }C_S
 \quad\text{avoiding every }\ell\in S.
\]
It does not assert that one finite-field curve avoids every odd prime,
and it does not obstruct a curve-dependent prime of unbounded size.
An infinite intersection of the open avoidance loci need not contain
a closed point.  No analytic hypothesis is used in this deduction.


\paragraph{Finding A09 (an explicit rank-one theta point of order thirteen).}
Let
\[
 C/\mathbb F_2:\quad Y^2+X^2Y=X^5+X,\qquad P=\infty.
\]
This curve has genus two and $2$-rank one. Put
\[
 k=\mathbb F_2[u,v]/(u^2+u+1,\ v^2+u^2v+1)\simeq\mathbb F_{16},
 \qquad R=(u,v)\in C(k).
\]
Then $L=\mathcal O_C(R-P)$ has exact order $13$ and lies on the
Raynaud theta divisor. Since $2$ has order $12$ modulo $13$,
this is a positive-$2$-rank instance of $T(C)$.

\begin{proof}
The substitution $x=X^{-1}$, $y=Y/X^2$ gives
$y^2+y=x^3+x^{-1}$. Its two Artin--Schreier poles have odd orders
$3$ and $1$. Riemann--Hurwitz gives $2g-2=-4+4+2=2$, and
Deuring--Shafarevich gives $f_C-1=2(-1)+2=0$, so $g=2$ and $f_C=1$.
There is one point $P$ above $X=\infty$, with
$\operatorname{ord}_P(X)=-2$ and $\operatorname{ord}_P(Y)=-5$.
The differential $dX/X^2=d(1/X)$ has divisor $2P$.
It is locally exact, and hence its corresponding section of the
rank-one sheaf $B^1$ has divisor $P$. Thus
$B^1\simeq\mathcal O_C(P)$ and the theta support is $C-P$.
Here the Frobenius twist is identified with the same equation over
$\mathbb F_2$; no choice of a twist over a nonperfect field is used.

The field presentation is explicit. Over $\mathbb F_4=\mathbb F_2(u)$,
putting $w=v/u^2$ gives $w^2+w=u^2$. Its trace to $\mathbb F_2$ is
$u^2+u=1$, so the second quadratic is irreducible. The relation
$u^3=1$ gives $u^5+u=1$, proving $R\in C(k)$.

Define
\begin{align*}
 A(X)={}&u(v+1)(X^6+X^3)+u^2vX^5+uX^4
                 +(uv+1)X^2+uX+u^2,\\
 B(X)={}&X^4+uX^3+X+1,\qquad F=A(X)+B(X)Y.
\end{align*}
Direct polynomial reduction in the displayed field gives
\[
 A(X)^2+X^2A(X)B(X)+(X^5+X)B(X)^2=(X+u)^{13},
 \tag{PR.1}
\]
and
\[
 F(R)=0,\qquad F(\iota R)=u\ne0,
 \qquad\iota(X,Y)=(X,Y+X^2).
 \tag{PR.2}
\]
The left side of (PR.1) is the norm $F\,\iota(F)$.
As $u\ne0$, the two points over $X=u$ are unramified and distinct.
Equations (PR.1)--(PR.2) imply that every finite zero of $F$ is $R$
and that its multiplicity is exactly $13$. There are no finite poles.
At $P$, the term $X^4Y$ has pole order $13$, while every other term
has smaller pole order. Therefore
\[
 \operatorname{div}(F)=13R-13P.\tag{PR.3}
\]
The class $R-P$ is nonzero: otherwise a rational function with that
divisor would give a degree-one map from a genus-two curve to
$\mathbb P^1$. Since $13$ is prime, (PR.3) proves its exact order.
Also $B^1\otimes L\simeq\mathcal O_C(R)$, so $L\in\Theta_{B^1}$.
Finally, the successive powers of $2$ modulo $13$ are
\[
 1,2,4,8,3,6,12,11,9,5,10,7,
\]
which are all the nonzero residues.
\end{proof}

\paragraph{An explicit cover and the arithmetic check.}
The normalization of the Kummer equation
\[
 W^{13}=F
\]
is a geometrically connected etale $\mu_{13}$-torsor over $C_k$.
Indeed, (PR.3) makes every valuation divisible by $13$, so it is
etale; it is geometrically connected because $F$ cannot be a
thirteenth power, again by nontriviality of $R-P$. With the local
parameter $t=X^2/Y$ at $P$, the unit $t^{13}F$ has residue $1$.
The root $tW=1$ therefore gives a $k$-rational point above $P$.
After extending constants to $\mathbb F_{4096}$, which contains
$k$ and $\mu_{13}$, its deck group is the constant group
$\mathbb Z/13$. Its genus is $14$ and its new Jacobian factor has
$2$-rank zero by A11 below (or by the prime case A03).
For this assertion about the displayed Kummer model, note that
$C$ and its theta support are defined over $\mathbb F_2$, so
Frobenius relabelling preserves the theta incidence; the explicit
twist normalization in A11 applies to its defining torsion class.
The direct torsor construction here is not asserted to be the same
arithmetic model as the Lang cover over $\mathbb F_{64}$ in the
recovered A1 row B.39, working identifier B.LIT7.

The file \texttt{theta13\_direct.py} constructs the identity using
Hasse--Taylor linear algebra over
$\mathbb F_2[z]/(z^4+z+1)$, with $u=z^2+z$, $v=z^3$.
Its certificate lists every coefficient. The separate verifier
\texttt{verify\_theta13\_norm.py} imports no producer code and
checks (PR.1)--(PR.2) by polynomial reduction over $\mathbb F_2$.
These are exact finite computations; they do not extrapolate to a
classification of positive-rank curves.

\paragraph{Nonabelian content of the example.}
The same curve has
\[
 \#C(\mathbb F_2)=4,\qquad \#C(\mathbb F_4)=4,\qquad
 H_2(U)=U^4+U^3+2U+4,\qquad \#J_C(\mathbb F_2)=8.
 \tag{PR.10}
\]
Indeed, the Artin--Schreier model has its two rational branch
points, and $x=1$ contributes two points. For
$x\in\mathbb F_4\setminus\mathbb F_2$, its right side is
$x^3+x^{-1}=1+x^2=x$,
whose absolute trace is $1$, so no additional points occur.
The two counts and the genus-two functional equation give the
displayed characteristic polynomial; evaluating it at $1$ gives
$8$. Thus all four rational points have the same prime-to-$2$
abelian Kummer class, because their divisor differences lie in
the $2$-group $J_C(\mathbb F_2)$. Nevertheless A09 and A11 show
that their full geometrically prime-to-$2$ point-section classes
are pairwise distinct. This is a concrete positive-rank instance
where the nonabelian covering argument supplies information
absent from the abelian quotient.

\paragraph{Finding A10 (three is excluded for every nonordinary genus-two base).}
Let $C$ be a smooth proper genus-two curve over an algebraically
closed field of characteristic two, with $f_C<2$. Then its Raynaud
theta divisor contains no nonzero point of order $3$.
Consequently every connected cyclic etale degree-three cover has
ordinary new factor of dimension two. In particular, this conclusion
also holds when $f_C=1$.

\begin{proof}
Work on the Frobenius twist, temporarily denoting that curve by
$C$ as well, and write $B=B^1$ there. In
characteristic two this is a line bundle with $B^{\otimes2}\simeq
\omega_C$ and degree one. The square relation follows in an etale
coordinate by comparing transition functions of locally exact
differentials with those of $\omega_C$.
From the relative Frobenius exact sequence,
$H^0(B)$ is the kernel of Frobenius on $H^1(\mathcal O_C)$.
Nonordinarity implies that Frobenius is not bijective, so $H^0(B)\ne0$.
Thus $B\simeq\mathcal O(W)$ for a point $W$, and
$2W\sim K_C$. Its theta support is $C-W$.

If $L\ne0$ had order $3$ on this support, write
$L=\mathcal O(R-W)$ with $R\ne W$. Then $3(R-W)$ would be
principal, so there would be a function with unique pole of exact
order $3$ at $W$. But Riemann--Roch and $2W\sim K_C$ give
\[
 h^0(2W)=2=h^0(3W).
\]
Every function with pole divisor at most $3W$ therefore already has
pole divisor at most $2W$, a contradiction.
For a cyclic degree-three cover, every nontrivial eigenspace
$H^1(C,L^i)$ is one-dimensional. The Frobenius maps between the
two eigenspaces have zero kernel precisely when the associated
theta incidences are absent, by the same exact sequence used in
A11. They are both bijective, so the new factor is ordinary.
\end{proof}

\subsection{A prime-power bridge and the remaining order obstruction}
\label{sec:prime-power-bridge}

\paragraph{Finding A11 (a sufficient prime-power incidence condition).}
Let $C$ be a smooth proper genus-two curve in characteristic two.
Let $n=\ell^r$ with $\ell$ an odd prime and $r\ge1$. Suppose that
\[
 \langle 2,-1\rangle=(\mathbb Z/n\mathbb Z)^\times
 \tag{PR.4}
\]
and that the Raynaud theta divisor contains a class of exact order
$n$. Then over an algebraic closure there is a connected cyclic etale
degree-$n$ cover $D_n\to C$ whose primitive conductor-$n$ factor has $2$-rank zero.
Writing $D_{n/\ell}$ for its intermediate degree-$n/\ell$ cover of
$C$, this factor is precisely
\[
 \operatorname{Prym}(D_n/D_{n/\ell})
\]
up to isogeny. In the finite-field covering formulation of A1, this
single cover separates every proposed pair of distinct colliding
points. Thus prime-power incidence with (PR.4) is sufficient; prime
order and the stronger equality $\langle2\rangle=(\mathbb Z/n)^\times$
are not required.

\begin{proof}
\emph{1. Symmetry and a zero Frobenius map.}
Let $B$ be the locally exact differential bundle on $C^{(1)}$ and
$L\in\operatorname{Pic}^0(C^{(1)})$ the given class. Since
$B^{\otimes2}\simeq\omega_{C^{(1)}}$ and $\chi(B\otimes L)=0$,
Serre duality gives
\[
 h^0(B\otimes L)=h^0(B\otimes L^{-1}).
 \tag{PR.5}
\]
Thus both $L$ and $L^{-1}$ lie on theta. Tensoring the relative
Frobenius sequence by $L$ gives
\[
 0\longrightarrow L\longrightarrow F_*F^*L
       \longrightarrow B\otimes L\longrightarrow0.
 \tag{PR.6}
\]
The two degree-zero line bundles $L$ and $F^*L$ are nontrivial:
$n$ is odd and Frobenius pullback preserves its exact order.
They have no sections, and each has one-dimensional $H^1$ by
Riemann--Roch. Therefore $H^0(B\otimes L)\ne0$ says that the
induced map from $H^1(L)$ to $H^1(F^*L)$ has nonzero kernel;
being a map of lines, it is zero. The same applies to $L^{-1}$.
To make the character labels precise, choose $e$ with $2e\equiv1
\pmod n$ and put $M=(F^*L)^e$ on $C$. This class still has exact
order $n$, and $M^2=F^*L$. If $M^{(1)}$ denotes its Frobenius
twist on $C^{(1)}$, then $F^*(M^{(1)})=M^2=F^*L$.
The pullback $F^*$ is injective on odd-order torsion, so
$M^{(1)}=L$. Consequently (PR.6) gives literally the
Frobenius-semilinear arrow from $H^1(C,M)$ to $H^1(C,M^2)$;
the inverse class gives the arrow from $H^1(C,M^{-1})$ to
$H^1(C,M^{-2})$.

\emph{2. Every primitive orbit is nilpotent.}
Over an algebraic closure, the Kummer correspondence gives the
connected cyclic etale cover associated to $M$, choosing its
character labels so that the eigensheaves are $M^i$.
When $C$ is over a finite field, all these data descend after a
finite constant extension. Its cohomology decomposes as
\[
 H^1(D_n,\mathcal O_{D_n})=
   H^1(C,\mathcal O_C)\oplus
   \bigoplus_{i=1}^{n-1} V_i,\qquad
   V_i=H^1(C,M^i),\quad \dim V_i=1.
 \tag{PR.7}
\]
Absolute Frobenius sends $V_i$ to $V_{2i}$. The primitive part is
the sum over $i\in(\mathbb Z/n)^\times$.
Condition (PR.4) means that the action of multiplication by $2$
on this unit group has at most two orbits, represented by $1$ and
$-1$. By Step~1, each orbit contains a zero arrow. A sufficiently
large iterate of Frobenius is therefore zero on every $V_i$ in
the primitive sum: traverse the finite orbit until the zero arrow
is reached. The stable Frobenius part of that sum is zero, which
is the assertion that the primitive factor has $2$-rank zero.
The one-dimensional spaces are those on the original genus-two
base; no inference from a single singular higher-dimensional
matrix is being made.

\emph{3. Prime powers identify the entire relative factor.}
Choose a generator $\sigma$ of the order-$n$ deck group and put
$a=\sigma^{n/\ell}$, of order $\ell$.
The quotient by $\langle a\rangle$ is $D_{n/\ell}$.
On the character with index $i$, the element $a$ acts trivially
if and only if $\ell\mid i$. For $n=\ell^r$, its nontrivial
characters are exactly the primitive indices $\gcd(i,n)=1$.
Thus the primitive factor is exactly the new factor for this
degree-$\ell$ quotient, proving the stated Prym identification.
The dimension is $(n-n/\ell)(g(C)-1)=\varphi(n)$, as a check.

\emph{4. The point-separation consequence.}
Now work in the finite-field situation of the covering formulation
of A1. Enlarge constants to define the cover, make its deck group
constant, and choose a rational point $P_n$ above the given base
point $P$. If a distinct point $Q$ collided with $P$, the covering
condition would provide a rational $Q_n$ above $Q$ with
\[
 \delta=[Q_n-P_n]\in J_{D_n}[2^\infty].
\]
The point $(a-1)\delta$ is still $2$-primary and belongs to
$\operatorname{im}(a-1)=\ker(\operatorname{Nm})^0
=\operatorname{Prym}(D_n/D_{n/\ell})$.
An abelian variety of $2$-rank zero has no nonzero geometric
$2$-power torsion points. Thus $(a-1)\delta=0$.
Equivalently, the divisor
\[
 aQ_n+P_n-Q_n-aP_n
 \tag{PR.8}
\]
is principal. Its positive and negative parts both have degree two
and are disjoint: the deck action is free, and the orbits above
the distinct base points $P$ and $Q$ are disjoint. It therefore
defines a degree-two map $D_n\to\mathbb P^1$.

Such a map is separable, since a purely inseparable map of degree
two from a smooth curve over an algebraically closed field to
$\mathbb P^1$ would make the curve rational. Hence $D_n$ would be
hyperelliptic. Its genus is $n+1\ge4$, so its hyperelliptic pencil
is unique. The order-$\ell$ automorphism $a$ acts on its base
$\mathbb P^1$; this action is nontrivial, because the kernel of
the action on a separable degree-two function field has order
at most two. A nontrivial odd-order automorphism of $\mathbb P^1$
has a fixed point. The fibre over that point contains at most two
geometric points, and an odd-prime cyclic group acts trivially on
a set of size at most two. This supplies a fixed point of $a$ on
$D_n$, contradicting the freeness of the etale deck action.
Thus the proposed collision is impossible.
\end{proof}

Condition (PR.4) is a finite group computation for the actual integer
$n$ in question. No non-Wieferich assumption is needed for the
statement or proof. Hooley's theorem supplies prime-level
primitive-root information under its analytic hypotheses; it does
not supply a theta point of any of these prime or prime-power
orders, nor does it imply (PR.4) at an arbitrary higher power.

\paragraph{Attempted extension using primary projections.}
For every genus-two curve over $\overline{\mathbb F}_2$, the
Raynaud theta support is a translate of its Abel curve, since
$\deg B=1$. It therefore generates its Jacobian: a proper translate
of a subgroup containing theta would contain all differences of
points of the Abel curve, whereas those differences generate the
Jacobian.
Poonen's Theorem~1.8 consequently applies to this theta curve:
given a finite set $S_0$ of primes, there is a density-one set
$S\supseteq S_0$ for which projection of theta points onto the
$S$-primary torsion of $J$ is onto \cite[Theorem~1.8]{Poonen2005}.
In particular, take $S_0$ to contain $2$ and a selected odd prime
$\ell$, prescribe an exact-order-$\ell^r$ point in the $\ell$-part,
and prescribe zero at every other prime of $S$. A resulting theta
point has order
\[
 n=\ell^r m,\qquad (\ell,m)=1,\qquad
 \text{every prime dividing }m\text{ lies outside }S.
 \tag{PR.9}
\]
This gives odd-order points, and iterating with all previously
occurring primes in $S_0$ gives infinitely many of pairwise coprime
orders. It applies also to ordinary curves with split Jacobians.
It does not require or prove that $m=1$.
Indeed, for a nonordinary genus-two curve take $\ell=3$ and $r=1$.
Finding A10 proves that the resulting point cannot have order $3$,
so its extra factor $m$ is necessarily greater than one. The
uncontrolled factors cannot simply be omitted from this projection
argument.

Two attempted ways to remove $m$ fail at definite algebraic steps.
First, multiplying the point by $m$ removes the extra primary
components but places it on $[m](\Theta)$; this operation does not
preserve $\Theta$. Second, pass from a cyclic degree-$\ell m$
cover to its intermediate degree-$m$ base. Its relative
degree-$\ell$ factor contains the characters of every conductor
$\ell d$ with $d\mid m$, including conductor $\ell$ itself.
Theta incidence at exact order $\ell m$ controls only the
primitive conductor-$\ell m$ block. On the intermediate base the
character spaces have dimension $m$, so one zero constituent does
not make those whole spaces nilpotent.

Even if $\langle2,-1\rangle=(\mathbb Z/\ell m)^\times$, symmetry
and that group equality kill only the primitive block. They leave
the proper-conductor blocks untouched. More intrinsically, when
$n$ has at least two distinct prime factors, the least common
multiple of its proper divisors is $n$. A deck element that acts
trivially on all proper-conductor characters must therefore be
the identity; there is no nontrivial element to substitute for
$a$ in Step~4 using only the primitive-zero conclusion. For a prime
power, that least common multiple is $n/\ell$, precisely the
reason Step~3 works.

These observations limit the single-deck-element and intermediate-cover
strategies. They do not rule out a different primitive-factor argument:
A15 below uses a product of deck operators and a gonality bound to bypass
the proper-conductor factors on every rank-one genus-two base, and more
generally on a base with no separable elliptic double quotient. A16
exhibits this bypass with conductor $15$ and positive $2$-rank in the
proper-conductor factors. The remaining universal question is incidence
or orbit coverage on a single cyclic torsion subgroup; it is not a
requirement to kill every proper-conductor block by the same method.


\subsection{An explicit Lang--Weil test for pure primary theta torsion}
\label{sec:lang-weil-theta}

\paragraph{Source scope.}
The Lang quotient construction used here is the one in Poonen,
\emph{Multiples of subvarieties in algebraic groups over finite fields},
IMRN 2005, no.~24, 1487--1498, Lemma~5.1 (author PDF, p.~6).
The recovered \texttt{ANABELIAN\_SUCCESSOR}, A.10--A.11, also uses
finite-field character sums on Abel--Jacobi images.  The following
statement gives explicit constants for the genus-two theta translate
and a sufficient test for a point whose \emph{entire order} is an odd
prime power.  It is unconditional.

\paragraph{Finding A12 (a curve-dependent sufficient test).}
Let $C/\mathbf F_Q$ be a smooth projective geometrically connected
genus-two curve in characteristic two, and assume that $C$ has an
$\mathbf F_Q$-rational point.  Write $X=C^{(1)}$ for its one-fold
relative Frobenius twist,
$J=\operatorname{Pic}^0(X)$, and
\[
 0\longrightarrow\mathcal O_X\longrightarrow F_*\mathcal O_C
 \longrightarrow B\longrightarrow0
\]
for relative Frobenius and its rank-one quotient.  Let
$\Theta=\{L\in J:h^0(X,B\otimes L)>0\}$ denote the Raynaud theta
support.  For any subgroup $H\subseteq G:=J(\mathbf F_Q)$, put
$a=[G:H]$ and
$N_H=\#(\Theta(\mathbf F_Q)\cap H)$.  Then
\[
 N_H\ \ge\
 \max\left\{0,\
 \left\lceil
 \frac{Q+1-\left\lfloor2(a+1)\sqrt Q\right\rfloor}{a}
 \right\rceil\right\}.
 \tag{LW.1}
\]
In particular,
\[
       Q+1-2(a+1)\sqrt Q>a
 \tag{LW.2}
\]
ensures that $\Theta(\mathbf F_Q)\cap H$ contains a nonzero point.
The simpler condition
\[
                       Q>(2a+3)^2
 \tag{LW.3}
\]
implies (LW.2).  If $C$ is ordinary, the origin is not in $\Theta$,
and positivity of the numerator in (LW.1) already suffices.

For an odd prime $\ell$ and $r\ge1$, take
$H=J(\mathbf F_Q)[\ell^r]$.  If (LW.2) holds and
\[
             \langle 2,-1\rangle=(\mathbf Z/\ell^r)^\times,
 \tag{LW.4}
\]
then the prime-power bridge A11 applies: after a finite extension of
constants there is a connected cyclic cover of degree $\ell^s$,
for some $1\le s\le r$, whose final degree-$\ell$ Prym has $2$-rank zero.
Consequently the point-separation conclusion of A11 holds for $C$.
This last conclusion uses A11, whereas (LW.1)--(LW.3) are purely
finite-field counting statements.

\begin{proof}
\emph{1. The theta translate and the Lang quotient.}
Choose $P_0\in X(\mathbf F_Q)$.  The line bundle $B$ has degree one.
A degree-one line bundle on a genus-two curve has a nonzero section
exactly when it is $\mathcal O_X(R)$ for a unique geometric point $R$.
Thus the theta support is the image of the closed embedding
\[
 i_B:X\longrightarrow J,\qquad
 R\longmapsto\mathcal O_X(R)\otimes B^{-1}.
\]
It is the usual Abel--Jacobi embedding
$i_{P_0}(R)=\mathcal O_X(R-P_0)$ translated by the
$\mathbf F_Q$-rational point
$\tau=\mathcal O_X(P_0)\otimes B^{-1}$.
No rational effective divisor representing $B$ is required.

Let $\Phi$ denote the $Q$-power Frobenius endomorphism of $J$.
The homomorphism $\Phi-1$ has differential $-1$, so is an
\'etale isogeny.  Its kernel is the constant group scheme
$G=J(\mathbf F_Q)$.  Factoring through the constant subgroup $H$
gives \'etale isogenies
\[
  J\xrightarrow{\pi_H} J/H\xrightarrow{\psi}J,
  \qquad \psi\pi_H=\Phi-1,
  \qquad \ker\psi=G/H.
 \tag{LW.5}
\]
In particular $\deg\psi=a$.  Form the pullback
\[
                  Y=X\times_{i_B,J,\psi}(J/H).
 \tag{LW.6}
\]
It is an \'etale cover of $X$ of degree $a$.

\emph{2. Geometric connectedness, including the characteristic part.}
We give the argument to make clear that $a$ need not be odd.
Over an algebraic closure, if (LW.6) were disconnected, its
monodromy would be a proper subgroup $A_0$ of the constant abelian
deck group $A=G/H$.  Quotienting by $A_0$ would give a nontrivial
intermediate \'etale isogeny
$\psi':J'\to J$ such that $i_B$ lifts to a morphism $v:X\to J'$.
Indeed the resulting $A/A_0$-torsor on $X$ has trivial geometric
monodromy and hence has a section.
Replace $v$ by $v-v(P_0)$.  It now sends $P_0$ to zero and lifts
$i_{P_0}$.  The universal property of the Jacobian supplies a
homomorphism $u:J\to J'$ with
$u i_{P_0}=v-v(P_0)$.  It follows, again by that universal property,
that $\psi'u=\operatorname{id}_J$.  The two abelian varieties have
the same dimension, so $u$ is an isogeny and
$1=\deg(\psi'u)=\deg\psi'\deg u$, a contradiction.

This is also the full abelian \'etale fundamental-group property
of the Abel--Jacobi map; it applies to characteristic-primary
\'etale covers as well as prime-to-characteristic covers.
See Milne, \emph{Jacobian Varieties}, author version of
12 June 2021, Section~9, Proposition~9.1, p.~30;
its characteristic-primary argument is completed on pp.~31--32 \cite{MilneJV}.
Thus $Y$ is smooth, projective and geometrically connected.
Unramified Riemann--Hurwitz gives
\[
                         g(Y)=a(g(X)-1)+1=a+1.
 \tag{LW.7}
\]

\emph{3. Rational fibres and the exact count.}
For $R\in X(\mathbf F_Q)$, put $x=i_B(R)\in G$.
Choose $y\in J(\overline{\mathbf F}_Q)$ with
$(\Phi-1)y=x$, and put $z=\pi_H(y)\in\psi^{-1}(x)$.
Since $\pi_H$ is defined over $\mathbf F_Q$,
\[
                 \Phi(z)-z=\pi_H(x).
 \tag{LW.8}
\]
Thus $z$ is rational if and only if $x\in H$.
The deck group $G/H$ is constant, so in this case all $a$ points
in the fibre are rational; otherwise none is rational.
Consequently
\[
                           \#Y(\mathbf F_Q)=aN_H.
 \tag{LW.9}
\]
The Weil bound for the curve $Y$ and (LW.7) give
$\#Y(\mathbf F_Q)\ge Q+1-2(a+1)\sqrt Q$.
For the precise cohomological input, see Deligne,
\emph{La conjecture de Weil. I}, formula~(1.5.1), p.~275, and
Theorem~1.6, p.~276 \cite{WeilI}.  Since $\#Y(\mathbf F_Q)$ is an integer,
(LW.9) yields (LW.1).

There is at most one point of $X$ mapping to the origin of $J$.
Condition (LW.2) gives $N_H>1$ and hence a nonzero point of $H$ on
$\Theta$.  If $\sqrt Q>2a+3$, then
$Q-2(a+1)\sqrt Q>\sqrt Q>a-1$, proving (LW.3).
If $C$ is ordinary, Frobenius on $H^1(\mathcal O_C)$ is invertible.
The cohomology sequence defining $B$ therefore gives
$H^0(X,B)=0$, so the origin is not on $\Theta$.

\emph{4. The prime-power consequence.}
A nonzero point in $H=G[\ell^r]$ has exact order $\ell^s$ for some
$1\le s\le r$.  Reduction
$(\mathbf Z/\ell^r)^\times\to(\mathbf Z/\ell^s)^\times$ is surjective.
Thus (LW.4) implies
$\langle2,-1\rangle=(\mathbf Z/\ell^s)^\times$.
Apply A11 to this actual order, allowing the finite constant
extension needed to define the deck action and a chosen point
above the prescribed base point.
\end{proof}

\paragraph{Finding A13 (an ordinary example with order $4159$).}
Let $C/\mathbf F_2$ be the smooth projective model of
\[
                  y^2+y=x+x^{-1}+(x-1)^{-1}.
 \tag{LW.10}
\]
It is ordinary of genus two.  Its Raynaud theta support contains
at least four $\mathbf F_{128}$-rational points of exact order
$4159$.  Moreover
$\langle2,-1\rangle=(\mathbf Z/4159)^\times$.
Thus A11 applies to this ordinary curve with a cyclic cover of
prime degree $4159$.

\begin{proof}
There are exactly three simple poles in (LW.10), at $0,1,\infty$.
The Artin--Schreier genus and $p$-rank formulas give genus two and
$2$-rank two; see Pries--Zhu,
\emph{The $p$-rank stratification of Artin--Schreier curves},
Lemma~2.6, p.~712 \cite{PriesZhu2012}.  The three totally ramified points are rational.
For a polynomial model one may put $h=x(x+1)$ and $Y=hy$, obtaining
\[
               Y^2+(x^2+x)Y=x^5+x^3+x^2+x.
 \tag{LW.11}
\]

There are three points over $\mathbf F_2$, precisely the three
ramified points.  Over
$\mathbf F_4=\mathbf F_2(u)$ with $u^2+u+1=0$, the only nonpole
$x$-values are $u,u^2$.  The right side of (LW.10) at these
values is respectively $u^2,u$, both of trace one to $\mathbf F_2$.
Neither has an Artin--Schreier preimage, so
\[
                      \#C(\mathbf F_2)=\#C(\mathbf F_4)=3.
 \tag{LW.12}
\]
If $\alpha_1,\ldots,\alpha_4$ are the Frobenius eigenvalues, then
$\sum\alpha_i=0$ and $\sum\alpha_i^2=2$.
The functional equation of the genus-two Weil polynomial
therefore gives
\[
                          P(T)=T^4-T^2+4.
 \tag{LW.13}
\]
Write $\beta_1,\beta_2$ for the roots of $T^2-T+4$.
The roots of (LW.13) occur in opposite pairs with squares
$\beta_1,\beta_2$.  For odd $m$ the Jacobian point-count formula
\cite[Theorem~19.1(b)]{MilneAV} gives
\[
 \#J(\mathbf F_{2^m})=(1-\beta_1^m)(1-\beta_2^m)
                   =1-t_m+4^m,
 \quad t_m=\beta_1^m+\beta_2^m.
 \tag{LW.14}
\]
The exact recurrence
$t_0=2,\ t_1=1,\ t_m=t_{m-1}-4t_{m-2}$ yields
\[
                  (t_0,\ldots,t_7)=(2,1,-7,-11,17,61,-7,-251).
\]
Consequently
\[
                    \#J(\mathbf F_{128})=16636=4\cdot4159.
 \tag{LW.15}
\]
Here the relative Frobenius twist has the same model, since the
equation is defined over $\mathbf F_2$.

Trial division by the integers from $2$ through
$\lfloor\sqrt{4159}\rfloor=64$ proves that $4159$ is prime.
The following exact modular calculations certify its order condition:
\[
\begin{gathered}
 2079=3^3\cdot7\cdot11,\qquad 4158=2\cdot2079,\\
 2^{2079}\equiv1,\quad
 2^{693}\equiv2554,\quad
 2^{297}\equiv1245,\quad
 2^{189}\equiv102\pmod{4159}.
\end{gathered}
 \tag{LW.16}
\]
They imply $\operatorname{ord}_{4159}(2)=2079$.
This order is odd, so $-1\notin\langle2\rangle$ and
$\langle2,-1\rangle$ has order $4158$.

In (LW.1), now take $Q=128$ and
$H=J(\mathbf F_{128})[4159]$.  By (LW.15), $|H|=4159$ and $a=4$.
The coarse test already holds because $128>11^2$.
For the integer count,
$113^2<100\cdot128<114^2$ gives
\[
       N_H\ge\left\lceil\frac{129-113}{4}\right\rceil=4.
 \tag{LW.17}
\]
The curve is ordinary, so none of these theta points is the origin;
all have exact order $4159$.  This proves the assertion and
verifies both hypotheses of A11.
\end{proof}

\paragraph{Arithmetic evidence for A13.}
The standard-library certificate
\path{lang_weil_4159_certificate.py} checks the finite-field
counts, recurrence, primality, modular orders and integer bound.
An independent script
\path{verify_ordinary_4159_independent.py} obtains (LW.15) instead
as the determinant of the $11$-by-$11$ Sylvester matrix for
$P(T)$ and $T^7-1$.  These exact arithmetic checks complement
the proof above; they do not constitute formal proof-assistant
verification.

\paragraph{Finding A14 (a fixed-prime barrier to the sufficient test).}
Fix a genus-two curve over $\mathbf F_q$ and an odd prime
$\ell\ne\operatorname{char}\mathbf F_q$.  Let $\Phi$ be Frobenius on
the rank-four module $T_\ell J$, and choose $d\ge1$ with
$\Phi^d\equiv1\pmod\ell$.  One may take its order modulo $\ell$,
or the explicit integer
$d=|\operatorname{GL}_4(\mathbf F_\ell)|
 =\prod_{j=0}^3(\ell^4-\ell^j)$.
Set
\[
                 C_\ell=\ell^{v_\ell(\#J(\mathbf F_{q^d}))}.
 \tag{LW.18}
\]
For every $m\ge1$,
\[
       |J(\mathbf F_{q^m})[\ell^\infty]|
                    \le C_\ell\,\ell^{4v_\ell(m)}
                    \le C_\ell m^4.
 \tag{LW.19}
\]
Let $M$ be the first integer $m\ge4$ such that
\[
                         q^{3m}\ge4C_\ell^2m^8.
 \tag{LW.20}
\]
This is a finite, computable integer.  For every $m\ge M$ and
every $r\ge1$, condition (LW.2) fails for
$Q=q^m$ and $H=J(\mathbf F_{q^m})[\ell^r]$.
The assertion concerns the sufficient numerical test, and does
not assert absence of theta torsion.

\begin{proof}
Put $U=\Phi^d-1\in\ell M_4(\mathbf Z_\ell)$.
For every positive integer $t$, the binomial identity gives
\[
 \Phi^{dt}-1=tU V_t,\qquad
 V_t=1+\sum_{j=2}^t\frac{\binom tj}{t}U^{j-1}.
 \tag{LW.21}
\]
Since
$\binom tj/t=\binom{t-1}{j-1}/j$ and
$j-1-v_\ell(j)\ge1$ for odd $\ell$ and $j\ge2$,
every summand after $1$ lies in $\ell M_4(\mathbf Z_\ell)$.
Thus $V_t\equiv1\pmod\ell$ and its determinant is a unit.
No eigenvalue of $\Phi^d$ is one, by the Weil absolute values.
Taking determinants and valuations in (LW.21) proves the exact formula
\[
 v_\ell\#J(\mathbf F_{q^{dt}})
   =v_\ell\#J(\mathbf F_{q^d})+4v_\ell(t).
 \tag{LW.22}
\]
The inclusion
$J(\mathbf F_{q^m})\subseteq J(\mathbf F_{q^{dm}})$
and (LW.22) with $t=m$ prove (LW.19).

Write $s=q^{m/2}$ and
$a=[J(\mathbf F_{q^m}):J(\mathbf F_{q^m})[\ell^r]]$.
The Weil lower bound for the abelian surface and (LW.19) imply
\[
                     a\ge\frac{(s-1)^4}{C_\ell m^4}.
 \tag{LW.23}
\]
Condition (LW.2), after rearrangement, would require
\[
                        a<\frac{(s-1)^2}{2s+1}.
 \tag{LW.24}
\]
It is therefore impossible whenever
$(s-1)^2(2s+1)\ge C_\ell m^4$.
For $m\ge4$ we have $s\ge2$, and
$(s-1)^2(2s+1)\ge s^3/2$.
Condition (LW.20) gives $s^3/2\ge C_\ell m^4$,
proving failure of (LW.2) at that $m$.
Finally $q^{3m}/m^8$ tends to infinity and is strictly increasing
for integral $m\ge4$: its consecutive ratio is at least
$8(4/5)^8>1$.  Hence $M$ exists and the same conclusion holds
for every $m\ge M$.
\end{proof}

\paragraph{What an analytic prime theorem still does not supply.}
A12 allows a curve-dependent successful test and A13 proves that
it can succeed on an ordinary curve.  A14 rules out making this
particular test succeed automatically by enlarging constants at one
fixed odd prime.  Allowing the prime to vary would require additional
information: a sufficiently large pure prime-power subgroup of
$J(\mathbf F_{q^m})$ together with the unit-group condition (LW.4).
For example, (LW.2) requires its order to dominate the cofactor on
the scale specified by (LW.24).
The Jacobian orders are the integers
$\det(1-\Phi^m)$; a density theorem for primes at which $2$ has a
specified multiplicative order does not by itself control large
prime-power divisors of this sequence.
This is not asserted to be the strongest numerical criterion for
point separation.  The recovered \texttt{ANABELIAN\_SUCCESSOR},
A.10, gives a different rational criterion: for a quotient
$\chi:J_C(\mathbf F_Q)\twoheadrightarrow\mathbf Z/\ell$ with
kernel $K$, writing $N=\#C(\mathbf F_Q)$, the conditions
$N\ge2|K|$ and $\ell>N$, together with $2$ primitive modulo $\ell$,
give a suitable rational prime-degree cover.  That criterion can
permit larger cofactors and does not require pure-primary theta
incidence.  Its order hypothesis is stronger than (LW.4):
the prime $4159$ of A13 has $2$ of index two, not a primitive root.
The factorization and order conditions must therefore be stated
for the particular criterion being used.
Once the finite Lang quotient is specified, its rational points
are already controlled by the unconditional finite-field Weil bound.
Thus no Dedekind-zeta GRH, Artin holomorphy assumption, or
number-field least-prime theorem is a premise of A12--A14.
An inference that such an analytic theorem forces the missing
geometric incidence or the needed factorization would require a
separate argument.


% Finding A15. Prepared independently of final/.  One-fold Frobenius twist
% notation C^{(1)} agrees with the root paper's revised convention.
\subsection{Rank one permits a primitive factor of arbitrary odd conductor}
\label{sec:rank-one-composite-bridge}

\paragraph{Finding A15 (the composite primitive-factor criterion).}
Let $C/k$ be a smooth projective geometrically connected curve of genus
two over a finite field of characteristic two, with geometric $2$-rank
$f_C=1$. Let $n>1$ be odd. The following conditions suffice for the
point-section separation conclusion of A1.
First, there is a geometrically connected cyclic etale cover
$D\to C$ of degree $n$ whose primitive conductor-$n$ Jacobian factor
has $2$-rank zero. More generally, for a specified $P\in C(k)$ it
suffices to have a finite extension $L/k$, such a cover over $L$ with
constant deck group $G=\langle\sigma\rangle$, and a point
$P'\in D(L)$ over $P$, for which
\[
 A=\prod_{\ell\mid n}(1-\sigma^{n/\ell}),\qquad
 P_n=\operatorname{im}(A:J_D\longrightarrow J_D),\qquad
 P_n(L)[2^\infty]=0.                                      \tag{CB.1}
\]
The product is over distinct prime divisors of $n$.
In particular, the first condition holds if the Raynaud theta divisor
on $\operatorname{Pic}^0(C^{(1)})$ contains a point of exact odd order
$n$ and
\[
                    \langle 2,-1\rangle=(\mathbf Z/n)^\times.
                                                               \tag{CB.2}
\]
Here $C^{(1)}$ is the one-fold relative Frobenius twist. Thus for
$2$-rank one the order $n$ need not be a prime power. No assertion of
universal theta incidence is included.

\paragraph{Exact source scope.}
The geometric criterion underlying (CB.1), with the extra hypothesis
that the base is not geometrically a separable double cover of a
genus-one curve, is the recovered \texttt{ANABELIAN\_SUCCESSOR}
\emph{final row A.14}, working identifier AC.11, in
\path{agent_criteria/criteria_section.tex}. The refinement here
removes that extra hypothesis for the rank-one genus-two stratum and
combines the criterion with theta incidence. Final row A.35 is a
different, finite Fourier counting criterion; it is not the source of
the composite projector argument. We give the full argument below.

\begin{proof}
\emph{1. Rank one excludes the exceptional double quotient.}
Work first over an algebraic closure. Suppose that a genus-two curve
$C$ admits a separable degree-two map to a genus-one curve $E$.
The function-field extension is quadratic and Galois, with group
$\mathbf Z/2$. Riemann--Hurwitz says that its different has degree
$2g_C-2=2$. Every ramified point in this characteristic-two quadratic
extension has different exponent at least two, so there is exactly
one ramified point. The Deuring--Shafarevich formula consequently gives
\[
 f_C-1=2(f_E-1)+1,\qquad f_C=2f_E\in\{0,2\}.              \tag{CB.3}
\]
This contradicts $f_C=1$. We use the geometric $p$-rank form of that
formula, stated for $p$-group covers, for example, in
Anbar--Beelen--Nguyen, Theorem~1, equation~(1)
\cite{AnbarBeelenNguyenP}; no assumption of ordinarity is needed.

\emph{2. The primitive projector produces a small separable function.}
Suppose that distinct points $P,Q$ collided in the covering formulation
of A1. The pointed cover in (CB.1), after restricting the collision
to $L$, supplies $Q'\in D(L)$ above $Q$ with
$\delta=[Q'-P']\in J_D(L)[2^\infty]$. Hence $A\delta=0$.
Put $r=\omega(n)$ and $d=2^r$. Each element
$\sigma^{n/\ell}$ has order $\ell$, and these elements belong to
distinct prime-primary factors of the cyclic group. Their $2^r$
subset products are therefore all distinct. Expanding $A(Q'-P')$
gives a nonzero principal divisor with exactly $d$ simple zeros and
$d$ simple poles. They lie in the two disjoint free deck orbits of
$P'$ and $Q'$. Let $f$ be the resulting rational function on
$D_{\overline L}$. It has degree $d$ and is separable: its simple
zeros (or poles) rule out a square in the function field.

The relevant degree bound is
\[
                         n>(2^{\omega(n)}-1)^2=(d-1)^2.
                                                               \tag{CB.4}
\]
It holds for every odd $n>1$, with no additional size assumption.
For $r=1$ use $n\ge3>1$; for $r=2$ use $n\ge15>9$.
For $r\ge3$, the distinct odd prime factors give
$n\ge3\cdot5^{r-1}>4^r>(2^r-1)^2$; the middle inequality holds
at $r=3$ and its ratio increases by $5/4$ thereafter.

\emph{3. The field generated by the conjugate functions.}
Write $K=\overline L(D)$,
$M=\overline L(\gamma f:\gamma\in G)$, and $e=[K:M]$.
Poonen's Castelnuovo--Severi consequence gives
\[
                  e\mid d,\qquad g(M)\le(d/e-1)^2.
                                                               \tag{CB.5}
\]
Here is the entire induction needed from that consequence. Start with
one degree-$d$ function, whose rational field has index $d$ and genus
zero. If an intermediate field has index $d/s$ and genus at most
$(s-1)^2$, adjoining another degree-$d$ function with relative degree
$a$ gives index $d/(as)$. Castelnuovo--Severi bounds its genus by
\[
 a(s-1)^2+(a-1)(as-1)
  =(as-1)^2-as(a-1)(s-1)\le(as-1)^2.
\]
Induction over the finitely many conjugates proves (CB.5).
This is exactly the use of Proposition~2.1, Lemma~2.2, and
Corollary~2.3 of \cite{Poonen2007Gonality}, over the perfect field
$\overline L$.

Since $D\to C$ is etale, $g(D)=n+1$. Equations (CB.4)--(CB.5)
exclude $e=1$. Thus $e>1$ is a power of two. Let $Z$ be the smooth
projective curve with function field $M$, and let $v:D\to Z$ have
degree $e$. The group $G$ preserves $M$. Its induced action on $Z$
is faithful, since a kernel $H$ would satisfy
$|H|\mid[K:M]=e$ as well as $|H|\mid n$.
The action on $Z$ is also free. Indeed, a stabilizer of a geometric
point of $Z$ would contain an element of odd prime order $\ell$.
It would partition the fibre divisor of $v$ into orbits of length
$\ell$, since the deck action on $D$ is free; fibre multiplicities
are constant on these orbits. This forces $\ell\mid e$, impossible.

Riemann--Hurwitz for the free action now yields
\[
                    g(Z)-1=n\bigl(g(Z/G)-1\bigr).
\]
The bound $0\le g(Z)\le(d/e-1)^2<n$ forces
$g(Z)=g(Z/G)=1$: quotient genus zero would give $g(Z)=1-n<0$,
and quotient genus at least two would give $g(Z)\ge n+1$.

\emph{4. A reduced fibre forces the forbidden degree-two map.}
Write $f=h\circ v$, with $h$ a nonconstant function on $Z$.
Choose a zero $z$ of $h$. The fibre divisor $v^*z$ is reduced,
because the zero divisor of $f$ is reduced. It has $e$ distinct
geometric points, all contained in the two original deck orbits.
Each such orbit maps injectively to $Z$, by equivariance and the
free action on $Z$. Hence the fibre contains at most one point from
each orbit, giving $e\le2$. Thus $e=2$.

The extension $K/\overline L(f)$ is separable, so $K/M$ is separable.
Its quadratic involution $\tau$ is normalized by $G$, because $G$
preserves its fixed field $M$; it is centralized by $G$ since
$\operatorname{Aut}(\mathbf Z/2)=1$. Therefore taking the quotient
by $G$ yields a separable degree-two map
\[
                  C=D/G\longrightarrow Z/G
\]
to a genus-one curve. Its degree is two explicitly because
$[K:M^G]=2n$ and $[K:K^G]=n$.
This contradicts Step~1 and proves the
criterion (CB.1). A primitive factor of geometric $2$-rank zero
has no nonzero geometric $2$-power torsion points. Its rational
$2$-primary group is therefore zero after every finite constant
extension, so this geometric condition implies the required
pointed version of (CB.1) for every base point.

\emph{5. Theta incidence annihilates every primitive orbit.}
Let $L$ have exact order $n$ on the Raynaud theta divisor of
$C^{(1)}$, and let $B$ denote its locally exact differential bundle.
In characteristic two $B$ is a line bundle with $B^2\simeq\omega$.
Serre duality therefore gives
$h^0(B\otimes L)=h^0(B\otimes L^{-1})$, so both signs lie on theta.
For relative Frobenius $F:C\to C^{(1)}$, tensor its exact sequence
by $L$. The induced map
\[
 H^1(C^{(1)},L)\longrightarrow H^1(C,F^*L)
\]
has a nonzero kernel, since $H^0(B\otimes L)\ne0$.
Both spaces are lines: the two line bundles are nontrivial of
degree zero on genus-two curves. The map is consequently zero.
The same holds for $L^{-1}$.

Choose $b$ with $2b\equiv1\pmod n$ and put $M_0=(F^*L)^b$ on $C$.
Then $M_0$ has exact order $n$ and
$F^*(M_0^{(1)})=M_0^2=F^*L$. Since Frobenius pullback is
injective on odd-order torsion, $M_0^{(1)}=L$. On the connected
cyclic cover associated to $M_0$, label the eigensheaves by $M_0^i$.
For $i\ne0\bmod n$ the cohomology spaces
$V_i=H^1(C,M_0^i)$ are one-dimensional, and Frobenius maps
$V_i$ to $V_{2i}$. The arrows $1\to2$ and $-1\to-2$ just shown
are zero. Condition (CB.2) makes every doubling orbit in
$(\mathbf Z/n)^\times$ contain $1$ or $-1$. Thus every primitive
orbit has nilpotent Frobenius, and its stable part is zero.

Finally $A$ acts on character $i$ by
$\prod_{\ell\mid n}(1-\zeta_\ell^i)$, which is nonzero precisely
when $\gcd(i,n)=1$. Since $n$ is odd, the character projectors are
defined after an extension of $\mathbf Z_2$ and split the
$2$-primary representation. The image $P_n$ is exactly the
primitive factor up to isogeny; its stable Frobenius part is the
primitive sum just considered. It therefore has $2$-rank zero.
All the geometric cover data descend after a finite constant
extension, proving the stated theta consequence.
\end{proof}

A18 below removes this geometric exclusion for the geometric
primitive-rank-zero criterion in characteristic two, in every genus
at least two. The rational-only criterion retains the distinction.

Steps~2--5 apply to any genus-two base known not to admit a geometric
separable double cover of a genus-one curve; the rank-one hypothesis
is what makes this geometric exclusion automatic in Step~1.

\paragraph{How many prime factors can the single-point condition allow?}
For odd $n$, the Chinese remainder theorem and the cyclicity of units
modulo an odd prime power give
\[
 (\mathbf Z/n)^\times/\bigl((\mathbf Z/n)^\times\bigr)^2
                  \simeq (\mathbf Z/2)^{\omega(n)}.
\]
If (CB.2) holds, this quotient is generated by the images of $2$ and
$-1$, so necessarily $\omega(n)\le2$. This is an arithmetic
restriction on the one-point sufficient condition, not on the
geometric criterion (CB.1). Exact finite examples are
\[
\begin{array}{c|rrrrrr}
 n&15&21&35&45&33&105\\ \hline
 \varphi(n)&8&12&24&24&20&48\\
 \operatorname{ord}_n(2)&4&6&12&12&10&12\\
 |\langle2,-1\rangle|&8&12&24&24&10&24
\end{array}
\]
Thus the first four integers satisfy (CB.2); having at most two
distinct prime factors is not by itself sufficient, as $n=33$ shows.
These computations certify only the unit groups. Theta incidence at
these orders is a separate condition, illustrated by A16 for $n=15$.

\paragraph{Several theta points: the exact coverage condition.}
For any fixed cyclic torsion class $M_0$ of odd order $n$, define
\[
 Z(M_0)=\{i\in(\mathbf Z/n)^\times:
                    M_0^{(1)i}\in\Theta_C\}.
\]
The primitive factor of its cyclic cover has $2$-rank zero if and
only if $Z(M_0)$ meets every orbit of multiplication by $2$ on
$(\mathbf Z/n)^\times$. Indeed, the arrow $V_i\to V_{2i}$ is
zero exactly at the displayed theta incidence. A cyclic sequence
of maps of lines has zero stable part exactly when at least one
arrow is zero; if every arrow is nonzero, their composition is a
bijective semilinear map of a line. This proves the equivalence.
Consequently, on the rank-one base, sufficiently many theta points
\emph{in this one cyclic subgroup} suffice, even when $\omega(n)>2$.
Condition (CB.2) and symmetry give this coverage from a single point.

\paragraph{What the primary-projection theorem still leaves open.}
Applying Poonen's theorem with $S_0\supset\{2,\ell\}$ can prescribe
an $\ell^a$ component, but gives order $n=\ell^a m$, where all primes
of $m$ lie outside a subsequently chosen density-one set $S$.
In the proof, Lemma~6.4 first finds a point $h$ in a bounded-index
subgroup supported on a density-zero prime set and then takes
$U$ to be \emph{all} prime divisors of its order. It bounds the
density quantity $\eta(U)$, not $|U|$ or $\omega(\operatorname{ord}h)$.
Enlarging finite $S_0$ therefore does not give the extra restriction
$\omega(m)\le1$ needed even to qualify for (CB.2), or the unit-group
equality itself \cite[Theorem~1.8 and Lemmas~6.4--6.6]{Poonen2005}.
For $\ell=3$, A10 forces $m>1$ on the rank-one stratum; A15 now
permits such a cofactor, but does not manufacture a suitable one.
Theta points of pairwise coprime orders likewise do not establish
the displayed orbit-coverage condition: on their fibre-product
cover they control proper-conductor characters, while the primitive
characters are tensor products whose theta incidence is not supplied.
This identifies the additional incidence statement that would be
needed to complete this proposed route.

\paragraph{Subsequent strengthening.}
Finding A46 in Section~\ref{sec:primitive-integral-sign} removes the
rational primitive-factor elliptic qualification by a separate integral
sign argument. The earlier implications and conductor-rank calculations
remain valid, but no extra ordinarity or elliptic-quotient assumption
is needed in the final rational primitive-factor criterion.


\subsection{A checked composite-degree witness}

\paragraph{Finding A16 (exact order fifteen, with ordinary proper factors).}
Let $k=\mathbb F_2[z]/(z^3+z+1)\simeq\mathbb F_8$, and let
\[
 C/k:\quad Y^2+X^2Y=X^5+(z+1)X^3+X,\qquad
 P=\infty,\qquad R=(z^2,z).
 \tag{CE.1}
\]
Then $C$ has genus two and $2$-rank one. The class
$M=\mathcal O_C(R-P)$ has exact order $15$, and its Frobenius twist
$M^{(1)}$ lies on the Raynaud theta divisor of $C$.
There is a pointed connected cyclic etale cover of degree $15$
after a finite constant extension, of genus $16$ and $2$-rank $7$.
Its primitive conductor-$15$ factor has dimension $8$ and $2$-rank
zero, whereas its conductor-$3$ and conductor-$5$ factors are
ordinary of dimensions $2$ and $4$. Thus the whole new factor has
positive $2$-rank $6$. Finding A15 nevertheless separates every
proposed pair of distinct colliding point sections on this base.

\begin{proof}
The reciprocal coordinate $x=1/X$, $y=Y/X^2$ gives
$y^2+y=x^3+(z+1)x+x^{-1}$. Its two poles have orders $3$ and $1$,
so the genus is two and the $2$-rank is one, by the same
Artin--Schreier calculation used in A09
\cite[Lemma~2.6]{PriesZhu2012}. The point $P$ is the unique point
above $X=\infty$, with pole orders $2$ and $5$ for $X$ and $Y$.
As in A09, $\operatorname{div}(dX/X^2)=2P$; consequently
\[
 B^1_C\simeq\mathcal O_{C^{(1)}}(P^{(1)}),\qquad
 \Theta_C=C^{(1)}-P^{(1)}.
 \tag{CE.2}
\]
This twist is retained explicitly: the coefficients of (CE.1) are
in $\mathbb F_8$, not all in $\mathbb F_2$.

Define the following polynomials over $k$:
\begin{align*}
 A(X)={}&z+(z+1)X+(z^2+z)X^2+(z^2+z+1)X^3\\
        &+(z^2+z)X^4+(z^2+1)X^6+(z^2+z+1)X^7,\\
 B(X)={}&1+(z+1)X+z^2X^2+(z^2+z)X^3\\
        &+(z^2+z+1)X^4+X^5,\qquad F=A+BY.
\end{align*}
Reduction modulo $z^3+z+1$ gives the exact identities
\begin{align}
 A^2+X^2AB+(X^5+(z+1)X^3+X)B^2&=(X+z^2)^{15},
 \tag{CE.3}\\
 F(R)=0,\qquad F(\iota R)&=z\ne0,
 \quad\iota(X,Y)=(X,Y+X^2).
 \tag{CE.4}
\end{align}
The independent symbolic checker verifies every coefficient of
(CE.3) and both evaluations. Since $X=z^2$ is unramified, the norm
identity proves that the only finite zero of $F$ is $R$, with
multiplicity $15$. There are no finite poles. The leading term
$X^5Y$ has pole order $15$ at $P$, while $A$ has pole order at
most $14$. Thus
\[
                       \operatorname{div}(F)=15(R-P).
 \tag{CE.5}
\]

The class is nonzero, since a principal divisor $R-P$ would make
$C$ rational. Order $3$ is excluded by
$h^0(3P)=h^0(2P)=2$, because $2P$ is canonical. To exclude order
$5$, a function with divisor $5(R-P)$ would lie in
\[
 H^0(C,\mathcal O(5P))=\langle1,X,X^2,Y\rangle.
\]
Its $Y$ coefficient would be nonzero; after scaling it would be
$Y+a_2X^2+a_1X+a_0$. Its norm has coefficient $1$ at $X$, whereas
the monic polynomial $(X+r)^5$, $r=X(R)=z^2$, has coefficient
$r^4=z\ne1$ there. This is a contradiction. The only remaining
divisor of $15$ is $15$, proving the exact order. Equation (CE.2)
now shows that $M^{(1)}$ is on theta, as claimed. In the normalization
of A11 one obtains the original class $M$, since
$F_C^*(M^{(1)})=M^2$ and $2^{-1}\equiv8\pmod{15}$.

The normalization of $W^{15}=F$ is etale because every valuation
of $F$ is divisible by $15$. It is geometrically connected:
being a third or fifth power would make respectively
$5(R-P)$ or $3(R-P)$ principal, contrary to the exact order.
For the parameter $t=X^2/Y$ at $P$, the unit $t^{15}F$ has residue
$1$, so $tW=1$ gives a $k$-point above $P$. Over
$\mathbb F_{4096}$ the deck group is the constant group
$\mathbb Z/15$, and Riemann--Hurwitz gives genus $16$.

The units modulo $15$ consist of the two multiplication-by-$2$
orbits
\[
                \{1,2,4,8\},\qquad\{14,13,11,7\}.
\]
The theta point and its inverse give one zero Frobenius arrow on
each orbit, so the primitive factor has $2$-rank zero by the
character calculation of A11. It has dimension $\varphi(15)=8$.
The conductor-$3$ factor is ordinary of dimension $2$, by A10.

For the conductor-$5$ factor, the same norm-coefficient argument
above shows that every nonzero order-$5$ point on $C-P$ has
$X$ coordinate $1$. But at $X=1$ equation (CE.1) is
$Y^2+Y=z+1$, and $\operatorname{Tr}_{\mathbb F_8/\mathbb F_2}(z+1)=1$.
Thus there is no such $k$-point. All four nonzero powers of
$M^3$ are $k$-rational. The Abel embedding and Frobenius twisting
preserve this rationality, so none lies on the relevant theta
support. All four character arrows are therefore bijective, and
the conductor-$5$ factor is ordinary of dimension $4$.
The character decomposition, whose idempotents have odd
denominators, now gives
$f_D=1+2+4+0=7$. Finally A15 applies to the rank-one base and its
rank-zero primitive factor. This proves the point-separation claim
without asserting that the whole new factor has rank zero.
\end{proof}

\paragraph{Computation scope.}
The retained producer \texttt{construct\_theta15.py} obtains the
coefficients by Taylor linear algebra in $\mathbb F_8$.
The separate \texttt{verify\_theta15\_symbolic.py} imports no
producer arithmetic and reduces the polynomial identity using
SymPy over $\mathbb F_2$. Both pass. The example was found in a
finite search of $5488$ nonbranch rational points on $784$
displayed rank-one models over $\mathbb F_8$; this search is not
claimed to classify geometric curves or all torsion orders.
No claim is made that this curve lacks a prime-power witness:
the example verifies the composite-cover conclusion itself.

\paragraph{A smaller prime-degree comparison.}
On the rank-one curve
\[
 C_5/\mathbb F_2:\quad
 Y^2+X^2Y=X^5+X^4+X^3+X,
 \qquad R_5=(1,0),\quad P_5=\infty,
\]
the function $F_5=Y+X+1$ has norm $(X+1)^5$, vanishes at $R_5$,
has value $1$ at its conjugate, and has unique pole of order $5$
at $P_5$. Hence $\operatorname{div}(F_5)=5(R_5-P_5)$, and the
nonzero theta class has exact order $5$. Since $2$ is primitive
modulo $5$, A03 supplies a degree-$5$ separating cover. Meanwhile
$Q_5=(0,0)$ satisfies $\operatorname{div}(X)=2(Q_5-P_5)$, so
$P_5$ and $Q_5$ have the same prime-to-$2$ abelian Kummer class.
The pole calculation and theta-characteristic calculation are
identical to those above; this comparison is a direct proved
special case, not a claim of minimality among all covers.


\subsection{A positive-density arithmetic supply for composite conductors}
\label{sec:three-times-artin-prime}

\paragraph{Finding A17 (an explicit family of candidate orders).}
Assume the Dedekind-GRH premise M. Let
\[
 \mathcal Q=\{q\text{ prime}:q\equiv5\pmod8,
                         \operatorname{ord}_q(2)=q-1\}.
\]
Then $\mathcal Q$ has natural density $\mathcal A/2$ among the
primes, where
\[
 \mathcal A=\prod_{\ell\ {
m prime}}
                      \left(1-\frac1{\ell(\ell-1)}\right)>0
\]
is Artin's constant. For every $q\in\mathcal Q$,
\[
                  \langle2,-1\rangle=(\mathbf Z/3q)^\times.
                                                               \tag{AF.1}
\]
The analytic input used here is Lenstra's fixed finite Frobenius
condition version of Hooley's theorem. Its sufficient GRH family is
\[
 \zeta_{K_n}(s),\qquad
 K_n=\mathbf Q(\zeta_n,2^{1/n}),\qquad
 n\ge1\text{ squarefree}.                                    \tag{AF.2}
\]
In particular, this is a number-field Kummer Dedekind-zeta family,
which is contained in M. The density is an arithmetic candidate
supply; it is not a theta-incidence assertion.

\begin{proof}
\emph{1. Degrees and the finite congruence field.}
For odd squarefree $m>1$, the prime $2$ is unramified in
$\mathbf Q(\zeta_m)$. At any prime above $2$, the polynomial
$T^m-2$ is Eisenstein, because the valuation of $2$ is one.
It is irreducible there and hence over $\mathbf Q(\zeta_m)$.
Consequently
\[
                         [K_m:\mathbf Q]=m\varphi(m).
                                                               \tag{AF.3}
\]
The Galois group is the affine semidirect product
$C_m\rtimes(\mathbf Z/m)^\times$. For clarity, a generator
$\tau$ of $C_m$ sends $2^{1/m}$ to $\zeta_m2^{1/m}$, while the
cyclotomic automorphism indexed by $b$ may be chosen to fix
$2^{1/m}$ and sends $\zeta_m$ to $\zeta_m^b$.
Its conjugation sends $\tau$ to $\tau^b$. Since $m$ is odd,
$b=2$ is available, and its commutator with $\tau$ equals
$\tau$. The commutator subgroup therefore contains $C_m$,
and equals it because the quotient by $C_m$ is abelian.
Every abelian subextension of $K_m/\mathbf Q$ thus lies in
$\mathbf Q(\zeta_m)$.

Let $F=\mathbf Q(\zeta_8)$. Its intersection with $K_m$ is an
abelian extension of $\mathbf Q$, so lies in
$\mathbf Q(\zeta_m)\cap\mathbf Q(\zeta_8)=\mathbf Q$.
The latter equality follows from the coprime cyclotomic conductors.
Thus
\[
                       [FK_m:\mathbf Q]=4m\varphi(m).
                                                               \tag{AF.4}
\]
Equations (AF.3)--(AF.4) also hold for $m=1$ with $K_1=\mathbf Q$.

\emph{2. The density from finite inclusion-exclusion.}
The condition $q\equiv5\pmod8$ prescribes the automorphism
$\sigma_5:\zeta_8\mapsto\zeta_8^5$ on $F$. This automorphism
acts nontrivially on $\sqrt2=\zeta_8+\zeta_8^{-1}$, so $2$ is a
quadratic nonresidue at these primes. The prime-$2$ obstruction
to primitive order is already excluded.
For an odd prime $\ell$, the obstruction that
$\ell\mid(q-1)/\operatorname{ord}_q(2)$ is exactly complete
splitting in $K_\ell$, outside the finitely many ramified primes.
For a finite set $E$ of odd primes, the compositum of their
$K_\ell$ is $K_{\prod_{\ell\in E}\ell}$.
By (AF.4), the simultaneous condition of splitting in $K_m$ and
having Frobenius $\sigma_5$ on $F$ has density
$1/(4m\varphi(m))$ for odd squarefree $m$.
Finite inclusion-exclusion therefore gives the density
\[
 \sum_{m\mid\prod_{\ell\in E}\ell}
               \frac{\mu(m)}{4m\varphi(m)}.
                                                               \tag{AF.5}
\]
Lenstra's Theorem~3.1, with $K=\mathbf Q$, $W=\langle2\rangle$,
and the fixed Frobenius condition $(F,\{\sigma_5\})$, identifies
the limiting finite Chebotarev density with the natural density
under (AF.2) \cite[\S2 and Theorem~3.1]{Lenstra1977}.
Thus the analytic infinite-limit step is an application of that
theorem, not an inference from finite inclusion-exclusion alone.
The limit of (AF.5), an absolutely convergent Euler product, is
\[
 \frac14\prod_{\ell\ {
m odd}}
                \left(1-\frac1{\ell(\ell-1)}\right)
                         =\frac{\mathcal A}{2}>0.
\]

\emph{3. The composite unit group.}
For $q\in\mathcal Q$, the Chinese remainder theorem gives
\[
 \operatorname{ord}_{3q}(2)
       =\operatorname{lcm}(2,q-1)=q-1=\frac{\varphi(3q)}2.
\]
A power of $2$ equal to $-1$ modulo $q$ has exponent
$(q-1)/2$ modulo $q-1$, which is even because $q\equiv5\pmod8$.
A power equal to $-1$ modulo $3$ has odd exponent. Therefore
$-1\notin\langle2\rangle$ modulo $3q$. Adjoining it doubles
the subgroup of order $\varphi(3q)/2$, proving (AF.1).
\end{proof}

\paragraph{The exact remaining geometric gate.}
For a fixed genus-two curve over a finite field of
characteristic two, a theta point of exact order $3q$ for even one
$q\in\mathcal Q$ would satisfy A15 on its original locus, and
A18 on every genus-two base. Finding A16 exhibits this for
one particular curve and $q=5$. A17 supplies infinitely many
candidate primes for each fixed curve, but does not establish
such a point for every curve. Poonen's projection theorem with
$S_0\supseteq\{2,3\}$ supplies order $3m$ after prescribing an
exact order-three component; A10 forces $m>1$. Its proof neither
forces $m$ to be a single prime nor places that prime in
$\mathcal Q$. Those additional incidence requirements remain
unproved; the positive density in A17 does not remove them.


\subsection{A primitive rank-zero cover removes the elliptic exception}
\label{sec:geometric-primitive-all-genera}

\paragraph{Finding A18 (geometric primitive criterion in all genera).}
Let $C/k$ be a smooth projective geometrically connected curve of
genus $g\ge2$ over a finite field of characteristic two.
Suppose that, over a finite extension of constants, there is a
geometrically connected cyclic \'etale cover $D\to C$ of odd degree
$n>1$ whose primitive conductor-$n$ Jacobian factor has geometric
$2$-rank zero.  Then the finite-field point-section separation
conclusion in the covering formulation of A1 holds for $C$.
There is no hypothesis excluding a separable double map from $C$
to a genus-one curve, and no restriction that $n$ be a prime power.

More precisely, with constant deck group $G=\langle\sigma\rangle$,
write
\[
 A=\prod_{\ell\mid n}(1-\sigma^{n/\ell}),\qquad
 P_n=\operatorname{im}(A:J_D\longrightarrow J_D),
 \tag{GP.1}
\]
where the product runs over the distinct prime divisors of $n$.
Under a putative collision, if $P_n(L)[2^\infty]=0$ on a pointed
model over a finite field $L$, the exceptional outcome of the
projector argument necessarily satisfies
\[
                         f(P_n)=r\varphi(n)>0
 \tag{GP.2}
\]
for an integer $r\ge1$.  Here $r$ is the number of geometric
ramification points of a resulting separable degree-two map
$C_{\overline L}\to E$ with $g(E)=1$.
For genus two, $r=1$, and the entire nontrivial new factor of
$D/C$ in this exceptional case is ordinary.
The geometric rank-zero hypothesis therefore eliminates this
exception.  A rational torsion condition alone is not being
replaced by a geometric one.

\begin{proof}
\emph{1. The projector reduction.}
Extend the constants so that the proposed distinct colliding
points $P,Q$, the cover, its cyclic deck group and a chosen lift
$P'\in D(L)$ are all defined over one finite field $L$.
Geometric $2$-rank zero implies
$P_n(L)[2^\infty]=0$ over every such extension.
The covering formulation of a collision gives a lift $Q'$ over
$Q$ with
\[
                       \delta=[Q'-P']\in J_D(L)[2^\infty].
\]
Hence $A\delta=0$.

We recall the complete geometric reduction from A15,
Steps~2--4; those steps do not use the rank or genus-two
hypothesis on the base except for the lower bound on $g(D)$.
Put $s=\omega(n)$ and $d=2^s$.
The $2^s$ subset products of the elements $\sigma^{n/\ell}$
are distinct, because they lie in independent prime-primary
factors of the cyclic group.  The two deck orbits of $P',Q'$
are free and disjoint.  Thus the expansion of the principal
divisor $A(Q'-P')$ has precisely $d$ distinct simple zeros and
$d$ distinct simple poles.  The resulting function $f$ has
separable degree $d$.

Over $\overline L$, put
\[
 K=\overline L(D),\qquad
 M=\overline L(\gamma f:\gamma\in G),\qquad e=[K:M].
\]
The Castelnuovo--Severi induction in A15 gives
\[
                 e\mid d,\qquad g(M)\le(d/e-1)^2.
 \tag{GP.3}
\]
The exact published input is Poonen,
\emph{Gonality of modular curves in characteristic $p$},
Proposition~2.1, Lemma~2.2 and Corollary~2.3,
author p.~3 \cite{Poonen2007Gonality}.  The induction adjoins functions of degree $d$;
if the previous index is $d/u$ and the adjoining degree is $a$,
the new index is $d/(au)$ and its genus is at most
$a(u-1)^2+(a-1)(au-1)\le(au-1)^2$.
For every odd $n>1$ one has $n>(d-1)^2$:
for $s=1,2$ use $n\ge3,15$, and for $s\ge3$ use
$n\ge3\cdot5^{s-1}>4^s$.
Since $g(D)=n(g-1)+1\ge n+1$, (GP.3) excludes $e=1$.
Thus $e>1$ is a power of two.

Let $Z$ be the smooth projective curve with function field $M$.
The induced $G$-action on $Z$ is faithful: a kernel would have
order dividing both $e$ and $n$.
It is free: a stabilizer containing an element of odd prime
order $\ell$ would partition a fibre divisor of $D\to Z$
of degree $e$ into orbits of length $\ell$, with constant
multiplicities on each orbit, forcing $\ell\mid e$.
Consequently
\[
                 g(Z)-1=n\bigl(g(Z/G)-1\bigr).
\]
Together with $g(Z)\le(d/e-1)^2<n$, this forces
$g(Z)=g(Z/G)=1$.

Write $f=h\circ v$ for $v:D\to Z$.
A fibre above a zero of $h$ is reduced, since the zero divisor
of $f$ is reduced.  It has $e$ points, all in the two original
deck orbits.  Each of these orbits maps injectively to $Z$,
because the action on $Z$ is free.  Hence $e\le2$, so $e=2$.
The extension $K/M$ is separable and quadratic, with involution
$\tau$.  It is separable because it is an intermediate extension
of the separable extension $K/\overline L(f)$.
Since $G$ preserves $M$, it normalizes $\langle\tau\rangle$;
the unique nontrivial element is fixed by conjugation, so
$\tau$ commutes with $G$.

\emph{2. The exceptional diagram is an actual pullback.}
Set
\[
 F=K^G=\overline L(C),\qquad
 E_0=M^G=K^{G\times\langle\tau\rangle}.
\]
Let $E$ be the smooth projective curve with function field $E_0$.
We already know that $g(E)=1$.
The group $G\times\langle\tau\rangle$ acts faithfully on $K$,
and has order $2n$, so
$K/E_0$ is Galois of that degree.  The fields fixed by
$G$ and by $\langle\tau\rangle$ have intersection $E_0$ and
compositum $K$.  Equivalently,
\[
       [F:E_0]=2,\qquad [M:E_0]=n,\qquad FM=K.
 \tag{GP.4}
\]
In particular $C\to E$ is separable of degree two and
$Z\to E$ is the cyclic \'etale cover of degree $n$ already
obtained from the free $G$-action.
The fibre product $C\times_E Z$ is connected by (GP.4).
It is smooth, since its projection to $C$ is the base change
of an \'etale morphism.  Its function field is $K$, so it is
the curve $D$ itself.

For each divisor $m\mid n$, let $H_m$ be the subgroup of $G$
of order $n/m$, and set
\[
       Z_m=Z/H_m,\qquad D_m=D/H_m.
\]
The same field calculation gives
\[
        D_m=C\times_E Z_m,\qquad
        \deg(Z_m/E)=m,\qquad \deg(D_m/C)=m.
 \tag{GP.5}
\]
Every $Z_m\to E$ is \'etale between genus-one curves.
After choosing compatible origins over $\overline L$ it is
an isogeny of odd degree.  Such an isogeny induces an
isomorphism on the geometric $2$-primary Tate modules:
its kernel has odd order, and multiplication by its degree
is invertible on every $2$-primary group.  Therefore
\[
                              f(Z_m)=f(E).
 \tag{GP.6}
\]

\emph{3. Ramification and the exact ranks of all intermediate covers.}
Let $r$ be the number of geometric ramification points of
$C\to E$.  Riemann--Hurwitz gives different of degree
$2g-2>0$, so $r\ge1$.  This is a quadratic Galois extension
in characteristic two.  Its base change by the \'etale map
$Z_m\to E$ has exactly $mr$ ramification points: above each
branch point are $m$ distinct points of $Z_m$, and \'etale
base change preserves the local ramification data.
The Deuring--Shafarevich formula for this degree-two cover,
together with (GP.6), gives
\[
             f(D_m)-1=2\bigl(f(Z_m)-1\bigr)+mr,
 \qquad
             f(D_m)=2f(E)-1+mr.
 \tag{GP.7}
\]
The formula counts the geometric ramification points, since
each has ramification index two.  It does not replace that
count by the degree of the different.
For the exact primary locator of the formula, see
Anbar--Beelen--Nguyen,
\emph{A new tower with good $p$-rank meeting Zink's bound},
Theorem~1, equation~(1), author p.~3 \cite{AnbarBeelenNguyenP}.
Its hypotheses are a Galois $p$-group extension of function
fields over an algebraically closed field, as here.

\emph{4. Conductor decomposition forces positive primitive rank.}
Consider
$V=T_2(J_D)\otimes_{\mathbf Z_2}\mathbf Q_2$,
whose dimension is the geometric $2$-rank of $D$.
The odd group $G$ acts semisimply on $V$.
Pullback and norm for $D\to D_m$ have composition
$[n/m]$ in one direction and the sum of the subgroup
$H_m$ in the other.  Since $n/m$ is odd, they identify
$T_2(J_{D_m})\otimes\mathbf Q_2$ with $V^{H_m}$.

After extending the coefficient field to split $G$, decompose
$V$ according to the exact orders of its characters.
Let $w_j$ be the total dimension of the characters of exact
order $j$, for $j\mid n$.  The invariants just identified give
\[
                           f(D_m)=\sum_{j\mid m}w_j.
 \tag{GP.8}
\]
On character $i$ modulo $n$, the endomorphism $A$ in (GP.1)
acts by
$\prod_{\ell\mid n}(1-\zeta_\ell^i)$.
This is nonzero exactly when $\gcd(i,n)=1$.
Thus its image has the primitive character part as its
$2$-adic Tate module up to isogeny, and
$f(P_n)=w_n$.
M\"obius inversion of (GP.8), using (GP.7) and $n>1$, yields
\[
\begin{aligned}
 f(P_n)
   &=\sum_{m\mid n}\mu(n/m)\bigl(2f(E)-1+mr\bigr)\\
   &=r\sum_{m\mid n}\mu(n/m)m
     =r\varphi(n)>0.
\end{aligned}
 \tag{GP.9}
\]
This contradicts the geometric rank-zero hypothesis.
Hence the proposed collision cannot occur.
\end{proof}

\paragraph{Genus-two consequence and the rational distinction.}
For $g=2$, the different of a separable quadratic map to a
genus-one curve has degree two.  Each ramified point in
characteristic two has different exponent at least two.
Thus $r=1$, with a single ramified point of different exponent
two.  Formula (GP.7) gives
$f(D)-f(C)=n-1$, equal to $g(D)-g(C)$.
The pullback/norm decomposition of $J_D$ into $J_C$ and its
new factor, up to an odd-degree isogeny, therefore shows that
the whole new factor is ordinary.  In particular its primitive
factor is ordinary of dimension $\varphi(n)$.

Consequently, even the rational condition
$P_n(L)[2^\infty]=0$ suffices on a genus-two base if one also
knows that $P_n$ is \emph{not ordinary}: a collision would force
it to be ordinary.  The argument in A18 alone retains this qualification for the
rational-only condition. Finding A46 below removes it altogether
by comparing the integral coefficients in a reduced quadratic fibre,
and applies in every positive characteristic for odd cyclic degree
prime to that characteristic.  In all genera, (GP.9)
describes the positive primitive rank forced by the
exceptional case.

The theta calculation of A15 now gives a smaller geometric
gate on \emph{every} genus-two curve in characteristic two,
of any $2$-rank: an exact odd-order theta point with
$\langle2,-1\rangle=(\mathbf Z/n)^\times$ is sufficient,
and so is the exact same-subgroup orbit-coverage condition
there.  The arithmetic family $n=3q$ of A17 is therefore
available as a family of candidate orders on ordinary
bielliptic bases as well.
This does not assert incidence for every curve.  The
one-dimensional eigenspace argument establishing that
theta criterion remains specific to genus two; A18's
all-genus assertion assumes a primitive factor of geometric
rank zero directly.

\paragraph{Subsequent strengthening.}
Finding A46 in Section~\ref{sec:primitive-integral-sign} removes the
rational primitive-factor elliptic qualification by a separate integral
sign argument. The earlier implications and conductor-rank calculations
remain valid, but no extra ordinarity or elliptic-quotient assumption
is needed in the final rational primitive-factor criterion.


\subsection{A finite exact search for each prescribed odd prime support}
\label{sec:finite-support-theta}

\paragraph{Finding A19 (an explicit field for all primary exponents).}
Let $C/\mathbf F_Q$ be a smooth projective geometrically connected
genus-two curve in characteristic two. Write $X=C^{(1)}$ for its
one-fold relative Frobenius twist, $J=\operatorname{Pic}^0(X)$,
and $\Theta\subset J$ for its Raynaud theta support. Thus
\[
 \Theta=i_B(X),\qquad
 i_B(R)=\mathcal O_X(R)\otimes B^{-1},
                                                               \tag{FS.1}
\]
where $B$ is the locally exact differential line bundle of degree
one on $X$. This is a genus-two Abel curve translated over
$\mathbf F_Q$; a rational effective divisor for $B$ is not required.
Let $S$ be a finite set of odd primes, and write $J(\overline{\mathbf
F}_Q)\{S\}$ for the points whose orders have every prime divisor in
$S$, with arbitrary exponents. Let $\Phi$ be $Q$-power Frobenius on
$J$, and choose $d\ge1$ such that
\[
            \Phi^d=1\quad\hbox{on }J[\ell]
                         \qquad(\ell\in S).                  \tag{FS.2}
\]
Then
\[
 \Theta(\overline{\mathbf F}_Q)\cap J(\overline{\mathbf F}_Q)\{S\}
   =\Theta(\mathbf F_{Q^d})\cap J(\mathbf F_{Q^d})\{S\}.
                                                               \tag{FS.3}
\]
In particular the left side is finite and is computable by a
terminating finite-field search. One may take
\[
 \begin{split}
 d&=\operatorname{lcm}_{\ell\in S}
                  \operatorname{ord}(\Phi\mid J[\ell]),
       \quad\hbox{or the explicit larger integer}\\
 d&=\operatorname{lcm}_{\ell\in S}
                    |\operatorname{GL}_4(\mathbf F_\ell)|,
 \qquad |\operatorname{GL}_4(\mathbf F_\ell)|
                   =\prod_{j=0}^3(\ell^4-\ell^j).
 \end{split}                                                   \tag{FS.4}
\]
For $S=\varnothing$ take $d=1$ and the torsion group to be $\{0\}$.
No GRH or conjecture about prime distribution is a premise.

\paragraph{Literature scope.}
This is a specialized effective finite-support torsion statement
of Boxall type, proved in full below. The bibliographic locator
for Boxall's general work is
\emph{Sous-vari\'et\'es alg\'ebriques de vari\'et\'es
semi-ab\'eliennes sur un corps fini}, in \emph{Number Theory
(Paris, 1992--1993)}, LMS Lecture Note Ser.~215, Cambridge, 1995,
69--80. The primary paper of Scanlon--Voloch discusses related torsion
finiteness results and lists this reference
\cite[Introduction and reference \textnormal{[Box]}]{ScanlonVoloch1999}.
The original Boxall chapter's theorem text was not accessible in
this check; no theorem number, exact bound, or originality claim
is inferred from its bibliographic record. The proof and the
specific field (FS.4) do not depend on that unavailable text.

\begin{proof}
\emph{1. Two points is the exact geometric intersection bound.}
For every nonzero $T\in J(\overline{\mathbf F}_Q)$,
\[
                  \#\bigl(\Theta\cap(\Theta+T)\bigr)\le2
                                                               \tag{FS.5}
\]
for geometric points. Here is a proof using only genus two.
Suppose three distinct points of this intersection correspond to
pairs $(R_i,S_i)$ on $X$ with
$[R_i-S_i]=T$, for $i=1,2,3$. The $R_i$ are distinct, and so
are the $S_i$, because the Abel map is an embedding.
For $i=2,3$ we have
$R_1+S_i\sim R_i+S_1$. These two effective degree-two divisors
are distinct: equality, with $R_1\ne R_i$, would force
$R_1=S_1$ and hence $T=0$.
Thus the corresponding degree-two linear system moves.
Riemann--Roch in genus two forces its divisor class to be
canonical: for a degree-two divisor $D$, $h^0(D)\ge2$ implies
$h^0(K_X-D)\ge1$, and a degree-zero line bundle with a section
is trivial. Therefore both $S_2$ and $S_3$ are the unique
effective degree-one divisor in the class $K_X-R_1$.
This contradicts $S_2\ne S_3$ and proves (FS.5).

\emph{2. An elementary powering lemma.}
Let $\ell$ be odd, let $V=(\mathbf Q_\ell/\mathbf Z_\ell)^4$,
and let $U\in\operatorname{GL}_4(\mathbf Z_\ell)$ satisfy
$U\equiv1\pmod\ell$. If $v\in V$ is a finite-order point and
$w=(U-1)v$ has exact order $\ell^s$, $s\ge1$, then
\[
        (U^{\ell^{s-1}}-1)v=\ell^{s-1}w\ne0,
                                                               \tag{FS.6}
\]
and this point has exact order $\ell$.
To check the equality put $N=\ell^{s-1}$ and $M=U-1$.
The binomial expansion is
\[
 (U^N-1)v=Nw+\sum_{j=2}^{N}\binom Nj M^{j-1}w.
\]
For $j\ge2$,
\[
 v_\ell\!\left(\binom Nj\right)\ge s-1-v_\ell(j),
 \qquad j-2-v_\ell(j)\ge0.
\]
The first follows from
$\binom Nj=(N/j)\binom{N-1}{j-1}$, and the second uses that
$\ell$ is odd. Since $M\in\ell M_4(\mathbf Z_\ell)$, every
summand after $Nw$ is divisible by $\ell^s$ as an operator on
$w$, and is zero. For $s=1$ the sum is empty. This proves
(FS.6). In addition $U^{\ell^{s-1}}$ fixes every point of
$V[\ell]$, again because $U\equiv1\pmod\ell$.

\emph{3. Isolate one prime component of a nonrational point.}
Let $z\in\Theta(\overline{\mathbf F}_Q)\cap J\{S\}$ and
write $z=\sum_{\ell\in S}z_\ell$ for its primary decomposition.
Put $A=\Phi^d$. Its action on every $T_\ell J$ is congruent
to the identity modulo $\ell$. For a component of order
$\ell^a$, the action of $A$ belongs to the kernel of
\[
 \operatorname{GL}_4(\mathbf Z/\ell^a)
                         \longrightarrow\operatorname{GL}_4(\mathbf F_\ell),
\]
which is an $\ell$-group. Hence the orbit length of $z_\ell$
under $A$ is a power of $\ell$.

Suppose $Az\ne z$ and choose $\ell$ for which
$Az_\ell\ne z_\ell$. For every other prime $t\in S$, choose
a power of $t$ equal to the orbit length of $z_t$ under $A$.
Let $c$ be the product of these orbit lengths. Then
$\gcd(c,\ell)=1$, and $U=A^c$ fixes every other component,
but still moves $z_\ell$: raising a permutation of
$\ell$-power order to the prime-to-$\ell$ power $c$ cannot
turn a nontrivial orbit into a fixed point.

Use the identification
$J[\ell^\infty](\overline{\mathbf F}_Q)
 \simeq T_\ell J\otimes(\mathbf Q_\ell/\mathbf Z_\ell)$.
Let $s\ge1$ be the order exponent of $(U-1)z_\ell$, and put
$\sigma=U^{\ell^{s-1}}$. By (FS.6),
\[
                      \sigma z=z+T,
 \qquad T\in J[\ell]\setminus\{0\},\qquad \sigma T=T.
                                                               \tag{FS.7}
\]
The other primary components remain fixed. The map $\sigma$
on geometric points is a positive integral power of $Q$-power
Frobenius. As $\Theta$ is defined over $\mathbf F_Q$, all
\[
                \sigma^jz=z+jT\qquad(0\le j<\ell)
\]
belong to $\Theta$. They are $\ell$ distinct points. Each also
belongs to $\Theta+T$, by using the preceding point in the
cyclic translation orbit. This contradicts (FS.5), since
$\ell\ge3$. Consequently $Az=z$, which is exactly
$z\in J(\mathbf F_{Q^d})$. The reverse inclusion in (FS.3)
is immediate.

\emph{4. An explicit order bound and terminating computation.}
Let
\[
 h_d=\#J(\mathbf F_{Q^d}),\qquad
 N_S=\prod_{\ell\in S}\ell^{v_\ell(h_d)}.
                                                               \tag{FS.8}
\]
Every point in (FS.3) has order dividing $N_S$, and
\[
                 N_S\le h_d\le(Q^{d/2}+1)^4.                 \tag{FS.9}
\]
The last inequality is the genus-two Jacobian Weil bound
\cite[Theorem~19.1(c)]{MilneAV}.
If $P_J(T)=\det(T-\Phi)$ is its integer Weil polynomial, then
\[
 h_d=\prod_{j=1}^4(1-\alpha_j^d)
       =\operatorname{Res}_T(P_J(T),T^d-1),                 \tag{FS.10}
\]
where the $\alpha_j$ are its roots; the sign is positive
because there are four roots
\cite[Theorem~19.1(b)]{MilneAV}.
Thus $h_d$, its valuations at the finitely many primes in
$S$, and $N_S$ are exact finite integer computations.
Using the explicit second choice in (FS.4) does not require
computing the torsion representation first.

Enumerate $X(\mathbf F_{Q^d})$, form the line-bundle classes
$i_B(R)$ in the finite Jacobian, and keep precisely those killed
by $N_S$. This enumerates the entire geometric intersection
in (FS.3). Their exact orders, and any desired doubling-orbit
coverage within their cyclic subgroups, can then be tested by
finite group operations. A smaller bound may be obtained by
replacing $h_d$ with the exponent of its $S$-primary subgroup,
but this is unnecessary for termination. No assertion of
practical running time at the potentially large degree $d$ is made.
\end{proof}

\paragraph{Consequence for the remaining composite-order search.}
For each fixed prime $q$ in A17, taking $S=\{3,q\}$ now makes
the search over \emph{all} orders $3^a q^b$ finite, not only
the search at order $3q$. The actual order must still satisfy
the unit-group test or the same-subgroup orbit-coverage condition.
The primes $q$ range over an unbounded set, so this result does
not prove that one of the finite searches succeeds for every
curve.

There is also an exact reformulation of the squarefree gate.
For $q\ne3$, the curve $[3](\Theta)$ contains a nonzero point
of order $q$ if and only if $\Theta$ contains a point of order
$q$ or $3q$. Indeed a lift $z\in\Theta$ of such a point
satisfies $3qz=0$ but $3z\ne0$, so has one of those two orders;
the converse follows by multiplying by three.
For $q$ in A17 either order meets the genus-two sufficient
condition of A18. Applying Poonen's multiples theorem to the
generating curve $[3](\Theta)$ still permits a varying additional
integer multiplier. It does not put a prime-order point on that
curve itself. Thus neither this reformulation nor (FS.3)
replaces the remaining incidence assertion by prime density.


\subsection{A rank-one curve with no nonzero three-primary theta point}
\label{sec:no-three-primary-theta}

\paragraph{Finding A20 (all powers of three fail on the degree-thirteen example).}
For the genus-two rank-one curve of A09,
\[
 C/\mathbf F_2:\quad Y^2+X^2Y=X^5+X,
\]
its Raynaud theta divisor satisfies
\[
       \Theta_C(\overline{\mathbf F}_2)
          \cap J_C[3^\infty](\overline{\mathbf F}_2)=\{0\}.
 \tag{TP.1}
\]
Thus allowing arbitrary powers of the single prime $3$ does not
repair the universal positive-rank incidence assertion. This does
not obstruct point separation on this curve: A09 gives its exact
order-$13$ theta point and separating cover.

\begin{proof}
The curve and its Frobenius twist have the same equation over
$\mathbf F_2$. Finding A09 gives its Frobenius characteristic
polynomial on every odd-adic Tate module:
\[
                         P(T)=T^4+T^3+2T+4.
 \tag{TP.2}
\]
Let $\overline P$ be its reduction modulo $3$. Exact polynomial
division in $\mathbf F_3[T]$ gives
\[
                       T^{20}-1=\overline P(T)Q(T),
 \tag{TP.3}
\]
where
\[
\begin{aligned}
 Q(T)={}&T^{16}-T^{15}+T^{14}+T^{12}+T^{11}+T^{10}\\
        &-T^6+T^5-T^4-T^2-T-1.
\end{aligned}
\]
Cayley--Hamilton applied to the action of the $2$-power Frobenius
$\Phi$ on $J_C[3]$ therefore proves $\Phi^{20}=1$ there.
All $3^4$ geometric points of $J_C[3]$ are consequently rational
over $\mathbf F_{2^{20}}$.

The determinant formula for the rational points of a Jacobian
gives
\[
\begin{aligned}
 \#J_C(\mathbf F_{2^{20}})
   &=\det(1-\Phi^{20})
     =\operatorname{Res}(P(T),T^{20}-1)\\
   &=1099812538656
     =3^4\cdot13577932576,
\end{aligned}
 \tag{TP.4}
\]
and $13577932576\equiv1\pmod3$.
The resultant has the displayed sign because $P$ has even degree.
One exact computation uses the integral companion matrix
\[
 C_P=\begin{pmatrix}
 0&0&0&-4\\
 1&0&0&-2\\
 0&1&0&0\\
 0&0&1&-1
 \end{pmatrix}
\]
and evaluates $\det(I-C_P^{20})$; an independent polynomial
resultant calculation gives the same integer.
The determinant formula and Frobenius convention are those of
\cite[Theorem~19.1(b), pp.~40--41]{MilneAV}.

The rational $3$-primary subgroup in (TP.4) has exactly $3^4$
points and already contains all of $J_C[3]$, so it equals
$J_C[3]$ and is killed by $3$.
Apply A19 with the prime set $S=\{3\}$ and the Frobenius exponent
$d=20$ established in (TP.3). It puts every geometric
$3$-primary theta point in $J_C(\mathbf F_{2^{20}})$.
Every nonzero such point would thus have exact order $3$.
Finding A10 excludes order-three theta points on every nonordinary
genus-two curve, and A09 gives $f_C=1$. This is a contradiction.
Finally the origin is on theta, since in A09 the locally exact
differential line bundle is $\mathcal O_C(P)$ for the point at
infinity. This proves equality in (TP.1).
\end{proof}

The standard-library certificate \texttt{check\_three\_primary.py}
performs the companion-matrix and modular polynomial computations.
The independently written SymPy certificate
\texttt{verify\_three\_primary.py} uses a polynomial resultant
and imports no producer code. These verify the finite arithmetic
in (TP.3)--(TP.4); the geometric exclusion also uses the separately
proved A19 and A10.


\subsection{Generic avoidance of every odd torsion order}
\label{sec:generic-all-odd}

\paragraph{Finding A21 (both positive-rank genus-two strata).}
Let $k$ be an algebraically closed field of characteristic two,
let $f\in\{1,2\}$, and let $Z$ be an irreducible component of
$\mathcal M_2^f$.  Let $C$ denote its geometric generic curve,
put $X=C^{(1)}$, and let $B$ be the locally exact differential
line bundle on $X$.  Its Raynaud theta support
\[
 \Theta=\{L\in\operatorname{Pic}^0(X):
                       H^0(X,B\otimes L)\ne0\}
\]
contains no nonzero point of odd order:
\[
 \Theta\cap
 \bigl(\operatorname{Pic}^0(X)[n]\setminus\{0\}\bigr)
             =\varnothing\qquad(n\ge1\text{ odd}).             \tag{GO.1}
\]
Consequently, for every odd $n>1$, every connected cyclic
degree-$n$ etale cover of $C$ has ordinary whole new factor.
This includes prime powers and composite degrees on the
$2$-rank-one stratum.  All assertions are unconditional.

The statement about one geometric generic curve does not
assert that a single nonempty open subset avoids all odd
orders simultaneously.  For any prescribed finite list of
odd integers the avoidance conditions do hold on a common
nonempty open, and the finite-field specialization argument
of A08 then applies for either $f=1$ or $f=2$.

\paragraph{Exact published inputs.}
Achter--Pries, \emph{Monodromy of the $p$-rank strata of the
moduli space of curves}, IMRN 2008, Art.~rnn053,
Theorem~4.5 (author version p.~14), gives full
$\operatorname{Sp}_{2g}(\mathbf F_\ell)$ and
$\operatorname{Sp}_{2g}(\mathbf Z_\ell)$ monodromy for
$\ell\ne p$, including $g=2$ and $f=1,2$
\cite{AchterPries2008}.
For the generic-point interpretation used here, the proof
of Ozman--Pries, Proposition~4.3, printed p.~465, explicitly
states that the full symplectic level-$\ell$ torsor $Z[\ell]$
is irreducible \cite[Proposition~4.3, proof]{OzmanPries2019}.
Thus this input is stronger than mere
connectedness of a cover over a possibly nonnormal base.
The tricanonical framing construction in Achter--Pries,
Section~5.1 and Lemma~5.2, pp.~17--18, passes to a scheme
without reducing monodromy.  The lifting assertion for a
closed subgroup of $\operatorname{Sp}_4(\mathbf Z_\ell)$
surjecting modulo $\ell$ is also
Landesman--Swaminathan--Tao--Xu,
\emph{Lifting subgroups of symplectic groups over
$\mathbf Z/\ell\mathbf Z$}, Research in Number Theory \textbf{3}
(2017), Art.~14, Theorem~1; their Proposition~1(b) says that
$\operatorname{Sp}_4(\mathbf Z/\ell^a)$ is perfect for every
odd $\ell$ and $a\ge1$ \cite{LSTX2017}.  Finally, Hall,
\emph{Classical Groups and Geometry}, version 6 May 2015,
Theorem~12.26, printed p.~173, gives the center and simplicity
of $\operatorname{PSp}_4(\mathbf F_\ell)$ for all odd $\ell$,
including $\ell=3$ \cite{Hall2015}.

\begin{proof}
\emph{1. The monodromy statement at the geometric generic
point.}
Use tricanonical frames to replace $Z$ by a scheme carrying
its universal curve.  The frame morphism has geometrically
integral $\operatorname{GL}_5$ fibres.  In particular its
pullback to the irreducible full level-$\ell$ torsor
$Z[\ell]$ is irreducible.  Restrict to a nonempty smooth
affine open $T$ of the resulting integral scheme.  The
level-$\ell$ torsor remains irreducible on $T$, so its
geometric generic Galois image is the full
$\operatorname{Sp}_4(\mathbf F_\ell)$.
The closed $\ell$-adic image over this same generic field
is consequently $\operatorname{Sp}_4(\mathbf Z_\ell)$ by
the lifting theorem just cited.  This reasoning works for
every odd $\ell$ on the same generic curve.  It does not
assume that deleting a singular locus from an arbitrary
irreducible base preserves its fundamental group.

These statements apply equally to $J=\operatorname{Pic}^0(X)$.
Indeed relative Frobenius induces an isogeny between the
Jacobians of $C$ and $X$ of degree a power of two.  It is an
isomorphism on every odd-order Tate module, and is a
similitude of the corresponding alternating pairings.
Conjugating a full symplectic group by a similitude leaves
that symplectic group unchanged.

\emph{2. Independence at distinct odd primes.}
Write $n=\prod_{\ell\mid n}\ell^{a_\ell}$ and
$G_\ell=\operatorname{Sp}_4(\mathbf Z/\ell^{a_\ell})$.
Each $G_\ell$ is perfect.  Its reduction kernel modulo
$\ell$ is an $\ell$-group, as is seen by filtering the
congruence kernel by successive powers of $\ell$.
If $G_\ell$ has a nontrivial simple quotient, that quotient
is nonabelian, since $G_\ell$ is perfect.  The image of the
congruence kernel is a normal $\ell$-subgroup of this simple
quotient and must be trivial.  The quotient therefore
factors through $\operatorname{Sp}_4(\mathbf F_\ell)$.
Its center has trivial image in a nonabelian simple group,
and Hall's theorem now implies that the simple quotient is
exactly $\operatorname{PSp}_4(\mathbf F_\ell)$.
The orders
\[
 \bigl|\operatorname{PSp}_4(\mathbf F_\ell)\bigr|
       =\frac{\ell^4(\ell^2-1)(\ell^4-1)}2                  \tag{GO.2}
\]
strictly increase with the positive integer $\ell$.
Thus factors belonging to distinct odd primes have no
common nontrivial simple quotient.

For clarity, this implies the required product assertion
without an independence assumption.  If a subgroup
$H\subset A\times B$ surjects to $A$ and $B$, Goursat's
argument identifies it with the fibre product over
isomorphic quotients $A/N_A\simeq B/N_B$.  If $H$ is
proper, this common quotient is nontrivial and has a
nontrivial simple quotient.  This would be a common simple
quotient of $A$ and $B$.  Induct on the number of prime
factors.  A nonabelian simple quotient of a product of
perfect groups factors through one factor: the images of
the factors are commuting normal subgroups, and at most
one can be the whole nonabelian simple group.  There are
no abelian simple quotients.  The obstruction just
described is therefore absent at each induction step.

The generic Galois image on $J[n]$ surjects to every
$G_\ell$ by Step~1 and is contained in their product.
Consequently it is
\[
            \prod_{\ell\mid n}G_\ell
                  =\operatorname{Sp}_4(\mathbf Z/n).         \tag{GO.3}
\]
The equality on the right uses the Chinese remainder
decomposition of matrices and their defining symplectic
equations.

\emph{3. Transitivity and the genus-two intersection bound.}
Over the local ring $R=\mathbf Z/\ell^a$, a vector of exact
order $\ell^a$ is primitive.  If $v$ is primitive, the
perfect alternating form admits $w$ with
$\langle v,w\rangle=1$.  The span of $v,w$ is a
nondegenerate free direct summand; its orthogonal
complement is a free rank-two symplectic module.
Completing $v,w$ to a symplectic basis proves that
$\operatorname{Sp}_4(R)$ is transitive on primitive
vectors.  Applying this in each Chinese remainder factor
proves that the group in (GO.3) is transitive on the
points of exact order $n$.

Because $B$ has degree one, $\Theta$ is the translated
Abel curve
\[
                 R\longmapsto\mathcal O_X(R)\otimes B^{-1}.
                                                               \tag{GO.4}
\]
For every nonzero $T\in J$, it satisfies
\[
                   \#\bigl(\Theta\cap(\Theta+T)\bigr)\le2.
                                                               \tag{GO.5}
\]
This is the bound proved in A19.  A short proof is included
to specify its exact scope.  Three distinct intersection
points give pairs $(R_i,S_i)$ with $[R_i-S_i]=T$.
The $R_i$ and the $S_i$ are separately distinct.  For
$i=2,3$, the distinct effective divisors
$R_1+S_i\sim R_i+S_1$ move in a degree-two linear system.
Riemann--Roch in genus two makes this class canonical.
Both $S_2$ and $S_3$ would then be the unique effective
degree-one divisor in the class $K_X-R_1$, a contradiction.
If the two displayed degree-two divisors were equal as
divisors, their equality and $R_1\ne R_i$ would instead
force $T=0$.  This also is excluded.

Suppose now that $\Theta$ contains a point of exact odd
order $n>1$.  The line bundle $B$, and hence its theta
support, is defined over the generic field of the family.
Galois invariance and (GO.3) force $\Theta$ to contain
every point of exact order $n$.  Choose coordinates
$J[n]\simeq(\mathbf Z/n)^4$, let $e_1,e_2$ be two
distinct coordinate vectors, and choose an odd prime
$\ell\mid n$.  Set
\[
                   P=e_1,\qquad T=(n/\ell)e_2.
\]
The point $T$ has exact order $\ell$ and all
$P+jT$, $0\le j<\ell$, have exact order $n$, since their
first coordinate is $1$.  They therefore lie on $\Theta$.
Each also lies on $\Theta+T$, using the preceding point
in this cyclic list.  There are $\ell\ge3$ such distinct
points, contradicting (GO.5).  This excludes every exact
odd order and proves (GO.1).

\emph{4. The whole new factor of an odd cyclic cover.}
Work over the algebraic closure of the generic field.
Let $\pi:D\to C$ be a connected cyclic etale cover of
odd degree $n$.  Choose a line bundle $M$ of exact order
$n$ defining the cover, so that
\[
 \pi_*\mathcal O_D=\bigoplus_{i=0}^{n-1}M^i,\qquad
 H^1(D,\mathcal O_D)=\bigoplus_{i=0}^{n-1}H^1(C,M^i).
                                                               \tag{GO.6}
\]
For $i\ne0$, the latter summand has dimension one.
Put $L_i=(M^i)^{(1)}$ on $X=C^{(1)}$.  It is nontrivial
of odd order and $F_{C}^*L_i=(M^i)^2$.  Tensoring
$0\to\mathcal O_X\to F_{C*}\mathcal O_C\to B\to0$
by $L_i$ and using the projection formula gives
\[
 0\longrightarrow L_i\longrightarrow F_{C*}(M^i)^2
                 \longrightarrow B\otimes L_i
                 \longrightarrow0.                          \tag{GO.7}
\]
By (GO.1), $H^0(X,B\otimes L_i)=0$.  Thus
$H^1(X,L_i)\to H^1(C,M^{2i})$ is injective, hence
bijective between its one-dimensional source and target.
Under Frobenius-semilinear identification this is precisely
the arrow from character $i$ to character $2i$ in (GO.6).
Multiplication by two permutes all nonzero residues
modulo $n$, so Frobenius is bijective on the entire
nontrivial-character sum.

Its stable rank is therefore $n-1$; the trivial-character
sum has stable rank $f$.  Consequently
$f(D)=f+n-1$.  The norm--pullback isogeny decomposition
$J_D\sim J_C\times\operatorname{Prym}(D/C)$ and
Riemann--Hurwitz give
\[
 f\bigl(\operatorname{Prym}(D/C)\bigr)=n-1
       =\dim\operatorname{Prym}(D/C),
\]
which proves the final assertion.
\end{proof}

\paragraph{Scope relative to the earlier results.}
For prime $n$, the conclusion agrees with Ozman--Pries,
Theorem~1.1; for the ordinary stratum, Yang's
new-ordinarity theorem already supplies the arbitrary
odd-degree conclusion used in A08.
The argument here is a direct genus-two consequence of
published full monodromy and the intersection bound.  It
does not require an extension of the compact-type
degeneration argument from prime to composite degrees.
The full joint monodromy in (GO.3) also supplies the finite
level needed in A22 below.


\subsection{No fixed finite odd prime support works uniformly}
\label{sec:no-fixed-prime-support}

\paragraph{Finding A22 (all exponents and all products).}
Let $S$ be any finite set of odd primes and let
$f\in\{1,2\}$.  There is a finite field $\mathbf F_Q$ of
characteristic two and a smooth projective geometrically
connected genus-two curve $C/\mathbf F_Q$ of $2$-rank $f$
such that, for $X=C^{(1)}$, $J=\operatorname{Pic}^0(X)$ and
its Raynaud theta support $\Theta$,
\[
       \Theta(\overline{\mathbf F}_Q)\cap
       \bigl(J(\overline{\mathbf F}_Q)\{S\}\setminus\{0\}\bigr)
                         =\varnothing.                       \tag{SA.1}
\]
Here $J\{S\}$ allows arbitrary powers of the primes in
$S$ and arbitrary products of those powers.
The curve can be chosen in any prescribed nonempty open
of any geometric irreducible component of
$\mathcal M_2^f$.  In particular these examples are
Zariski dense in each such component.

This strengthens finite-list avoidance: no fixed finite
prime support, even with all exponents and mixed orders
allowed, supplies a universal pure odd-torsion theta
witness on either positive-rank stratum.  The assertion
is unconditional.  The field and curve may depend on $S$;
no uniform degree bound for that field is asserted.

\begin{proof}
The case $S=\varnothing$ is immediate, so assume that $S$
is nonempty and put $N=\prod_{\ell\in S}\ell$.

\emph{1. A family with a finite avoidance condition.}
Start with a component of $\mathcal M_2^f$ over
$\overline{\mathbf F}_2$.  Use the tricanonical frame scheme
and choose a nonempty smooth affine open $T$ as in A21.
There is a universal smooth proper genus-two curve
$\mathcal C\to T$, every fibre has $2$-rank $f$, and the
geometric generic monodromy on every odd level is full.
One can also restrict to the inverse image of any
prescribed nonempty moduli open.

Write $\mathcal X=\mathcal C^{(1)}$ for relative Frobenius
twist over $T$ and
$\mathcal J=\operatorname{Pic}^0(\mathcal X/T)$.
Relative Frobenius is finite locally free of degree two.
In the sequence
\[
 0\longrightarrow\mathcal O_{\mathcal X}
  \longrightarrow F_*\mathcal O_{\mathcal C}
  \longrightarrow\mathcal B\longrightarrow0,                 \tag{SA.2}
\]
the quotient $\mathcal B$ is locally free of rank one.
This can be checked in a relative etale coordinate, where
$1,t$ is the usual basis for relative Frobenius.
It has degree one on every fibre.  Consequently
\[
 \mathcal X\longrightarrow\mathcal J,\qquad
 R\longmapsto\mathcal O_{\mathcal X_t}(R)
                         \otimes\mathcal B_t^{-1}             \tag{SA.3}
\]
is a relative translated Abel map.  Its construction uses
the relative diagonal and the Picard functor; a rational
base point on the fibres is unnecessary.  Its image is
closed, because $\mathcal X$ is proper over $T$ and
$\mathcal J$ is separated over $T$, and its geometric
fibre image is precisely the Raynaud theta support.

The scheme
$\mathcal J[N]^\times=\mathcal J[N]\setminus\{0\}$
is finite etale over $T$: the zero section is open and
closed in this finite etale scheme.  The inverse image
of the closed image of (SA.3) in
$\mathcal J[N]^\times$ is closed, and its projection to
$T$ is closed because the projection is finite.
It misses the generic point by A21.  Remove this proper
closed subset and replace what remains by a nonempty
affine open $U$.  Then every geometric fibre over $U$
satisfies
\[
                  \Theta_t\cap
                  (\mathcal J_t[N]\setminus\{0\})
                           =\varnothing.                     \tag{SA.4}
\]
This is a finite incidence condition, so no passage from
infinitely many open subsets to one open has occurred.

All data descend to a sufficiently large finite field
$k_0=\mathbf F_{q_0}$.  Enlarge that constant field so
that $U$ is geometrically integral and
$q_0\equiv1\pmod{N^2}$.  The latter is possible because
$N$ is odd.  The family, its fixed-rank condition,
and (SA.4) are preserved by this constant extension.

\emph{2. A connected symplectic level torsor.}
Choose a primitive $N^2$th root of unity in $k_0$ to
identify the values of the Weil pairing with the standard
constant pairing.  Let
\[
 V=\operatorname{Isom}_{\mathrm{sympl}}
       \bigl((\mathbf Z/N^2)^4_U,\mathcal J[N^2]\bigr)
                      \longrightarrow U.                    \tag{SA.5}
\]
This is a finite etale torsor for the constant group
\[
       G=\operatorname{Sp}_4(\mathbf Z/N^2)
        =\prod_{\ell\in S}
                       \operatorname{Sp}_4(\mathbf Z/\ell^2).
\]
Its geometric generic monodromy is all of $G$ by A21.
The restriction of $V_{\overline{k}_0}$ to the generic
point of $U_{\overline{k}_0}$ is therefore connected.
Since $U$ is smooth and
geometrically integral, $V$ is geometrically integral.
Equivalently, its geometric connectedness follows from
the full monodromy, and smoothness makes distinct
irreducible components disjoint.
The total arithmetic image is also $G$: it contains
the full geometric image, and it is contained in $G$
because $\mu_{N^2}$ is constant over $k_0$.

\emph{3. Realizing a precise Frobenius class by a twist.}
For $\ell\in S$ put $u_\ell=1+\ell$ modulo $\ell^2$ and
choose
\[
 \gamma_\ell=
 \operatorname{diag}(u_\ell,u_\ell,
                         u_\ell^{-1},u_\ell^{-1})
     \in\operatorname{Sp}_4(\mathbf Z/\ell^2).
                                                               \tag{SA.6}
\]
Use the standard symplectic basis ordered as
$e_1,e_2,f_1,f_2$.  Then
\[
 \gamma_\ell\equiv I\pmod\ell,\qquad
 (\gamma_\ell-I)/\ell
                   \equiv\operatorname{diag}(1,1,-1,-1)
                             \pmod\ell.                       \tag{SA.7}
\]
The second matrix is invertible.
Let $\gamma=(\gamma_\ell)_{\ell\in S}\in G$.
Each $\gamma_\ell$ has order $\ell$, so $\gamma$ has
order $N$.

Here is a complete existence argument for a specialization
with this Frobenius class.  It is the finite-field
twisting proof of Chebotarev, with only qualitative
Lang--Weil needed.  Let $\mathcal F$ denote arithmetic
$q_0$-power Frobenius acting on the geometric points of
$V$.  For a symplectic frame $v$ define a new descent
action of the arithmetic generator by
\[
                       v\longmapsto
                               \mathcal F(v)\gamma^{-1}.
                                                               \tag{SA.8}
\]
Right multiplication by $\gamma^{-1}$ is a deck
automorphism defined over $k_0$ and commutes with
$\mathcal F$.  Since $\gamma^N=1$, this descent action
agrees with the original action on
$\operatorname{Gal}(\overline{k}_0/\mathbf F_{q_0^N})$.
It therefore defines a finite Galois descent datum on
$V_{\mathbf F_{q_0^N}}$.  The datum is effective because
$V$ is affine; the exact general reference is the
Stacks Project, Lemma~35.37.1, Tag~\texttt{0245}
\cite{Stacks0245}.
Write $V^\gamma$ for the descended twist.
The projection to $U$ descends because deck
automorphisms act trivially there.

The twist $V^\gamma$ is geometrically integral, being
geometrically isomorphic to $V$.  Lang--Weil gives
$V^\gamma(\mathbf F_{q_0^m})\ne\varnothing$ for every
sufficiently large $m$.  One precise reference for this
form, including affine varieties, is Poonen,
\emph{Rational points on varieties}, Theorem~7.7.1(ii)--(iii),
printed p.~222 \cite{Poonen2017RPV}.  Choose such an $m$
with $m\equiv1\pmod N$ and a rational point of
$V^\gamma$ over $\mathbf F_{q_0^m}$.  Under the geometric
identification with $V$, it gives a frame $v$ over a
point $t\in U(\mathbf F_{q_0^m})$ such that
\[
                \mathcal F^m(v)\gamma^{-m}=v,
       \qquad\text{hence}\qquad
                \mathcal F^m(v)=v\gamma.                      \tag{SA.9}
\]
Thus the $Q$-power Frobenius of the fibre, where
$Q=q_0^m$, has matrix $\gamma$ on its $N^2$-torsion in
the frame $v$.  Changing the frame conjugates this
matrix, which leaves the properties in (SA.7) unchanged.
This argument does not require $t$ to have minimal
residue-field degree $m$; the fibre is considered over
the stated field $\mathbf F_Q$.

\emph{4. All rational $S$-primary torsion has squarefree
exponent.}
Let $C/\mathbf F_Q$ be this fibre and
$\Phi$ its $Q$-power Frobenius on $J$.
For each $\ell\in S$, lift the chosen torsion frame to
a basis of $T_\ell J$.  Equations (SA.7)--(SA.9) imply
\[
             \Phi-I=\ell A_\ell,\qquad
                         A_\ell\in
                         \operatorname{GL}_4(\mathbf Z_\ell).
                                                               \tag{SA.10}
\]
On
$J[\ell^\infty](\overline{\mathbf F}_Q)
 \simeq(\mathbf Q_\ell/\mathbf Z_\ell)^4$,
the matrix $A_\ell$ is an automorphism.  Consequently
\[
 \begin{split}
 J(\mathbf F_Q)[\ell^\infty]
     &=\ker\bigl(\Phi-I\mid
                  J[\ell^\infty](\overline{\mathbf F}_Q)\bigr)\\
     &=J[\ell](\overline{\mathbf F}_Q).
 \end{split}                                                   \tag{SA.11}
\]
Indeed $\ell A_\ell v=0$ if and only if
$\ell v=0$.  In particular all the $\ell$-torsion
points are rational, and there are no rational points
of order $\ell^2$ or higher.  Taking the direct sum
over $S$ gives
\[
                  J(\mathbf F_Q)\{S\}
                          =J[N](\overline{\mathbf F}_Q).
                                                               \tag{SA.12}
\]
This equation includes mixed orders dividing $N$.

\emph{5. Excluding arbitrary exponents geometrically.}
Since $\Phi$ is the identity on $J[\ell]$ for every
$\ell\in S$, Finding~A19 applies with $d=1$.
It puts every geometric $S$-primary theta point in
$J(\mathbf F_Q)\{S\}$.  By (SA.12) this is $J[N]$,
and (SA.4) excludes every nonzero point of that group
from theta.  This proves (SA.1).

The construction started inside an arbitrary prescribed
nonempty open of a component.  Repeating it for each
such open proves the asserted Zariski density.
\end{proof}

\paragraph{What the result leaves open.}
A22 is a nonexistence result for a \emph{fixed finite}
choice of prime factors across all curves.  It does not
exclude a witness whose prime factors depend on the
curve, and it does not contradict A19's terminating
search for a specified finite support on one specified
curve.  It also does not establish the variable-prime
theta-incidence condition in A18.  An Artin-prime
density theorem continues to supply arithmetic
candidates, while an actual theta incidence at an
appropriate order remains a separate geometric
condition.


\subsection{An ordinary obstruction to the large-prime quotient test}
\label{sec:b3-ordinary-obstruction}

\paragraph{Finding A23 (the large-prime quotient test can fail over every extension).}
Let $C/\mathbf F_4$ be the smooth projective model of
\[
                   y^2+y=x+x^{-1}+(x-1)^{-1}.
 \tag{B3O.1}
\]
It is an ordinary genus-two curve.  Put
$Q=4^m$, $N_m=\#C(\mathbf F_Q)$ and
$J_m=\#J_C(\mathbf F_Q)$ for $m\ge1$.
For every prime $\ell\mid J_m$ one has
\[
                              \ell N_m<2J_m.
 \tag{B3O.2}
\]
Consequently the sufficient conditions in recovered A1
B.3, formula (B.CL.2), can never hold for a prime quotient
of this actual Jacobian over any constant extension of
$\mathbf F_4$, even if the primitive-root requirement is omitted.
Nevertheless the Raynaud theta support of this same curve has
nonzero points of exact order $4159$ over $\mathbf F_{4^7}$,
and $\langle2,-1\rangle=(\mathbf Z/4159)^\times$.
Thus (B3O.2) is a failure of a sufficient numerical test, not
an absence of geometric theta witnesses.

\begin{proof}
The equation is defined over $\mathbf F_2$.  Its three simple
poles at $0,1,\infty$ give genus two and $2$-rank two by the
Artin--Schreier formulas \cite[Lemma~2.6]{PriesZhu2012}.
There are three rational points over each of $\mathbf F_2$
and $\mathbf F_4$: the three ramified points account for all
$\mathbf F_2$-points; if $u^2+u+1=0$, the remaining
$\mathbf F_4$-values of $x$ are $u,u^2$, and the right side
of (B3O.1) has respective values $u^2,u$, both of absolute
trace one.  These have no Artin--Schreier preimages.
These are also the exact inputs (LW.12) of A13.

Let $\alpha_1,\ldots,\alpha_4$ be the Frobenius eigenvalues
over $\mathbf F_2$.  The two point counts imply
$\sum_i\alpha_i=0$ and $\sum_i\alpha_i^2=2$.
The genus-two functional equation therefore yields the
characteristic polynomial
\[
                         P_2(T)=T^4-T^2+4,
 \tag{B3O.3}
\]
as in (LW.13).  If $\beta_1,\beta_2$ are the roots of
$T^2-T+4$, the four squares $\alpha_i^2$ are
$\beta_1,\beta_1,\beta_2,\beta_2$.
Thus the Frobenius polynomial over $\mathbf F_4$ is
$(T^2-T+4)^2$; no isogeny classification is needed here.
Define
\[
 t_m=\beta_1^m+\beta_2^m,\qquad e_m=Q+1-t_m.
 \tag{B3O.4}
\]
The recurrence $t_0=2$, $t_1=1$ and
$t_m=t_{m-1}-4t_{m-2}$ shows that $t_m,e_m$ are integers.
The curve and Jacobian Frobenius point-count formulas give
\[
 \begin{aligned}
 N_m&=Q+1-2t_m=2e_m-(Q+1),\\
 J_m&=(1-\beta_1^m)^2(1-\beta_2^m)^2=e_m^2.
 \end{aligned}
 \tag{B3O.5}
\]
For the Jacobian formula see
\cite[Theorem~19.1(b)]{MilneAV}.
Since $|\beta_i|=2$, the Weil estimate gives
$e_m\ge(\sqrt Q-1)^2>0$.

If a prime $\ell$ divides $J_m=e_m^2$, then
$\ell\mid e_m$ and $\ell\le e_m$.
Also $N_m\ge0$, since it is a point count.  Hence
\[
 \ell N_m\le e_mN_m
       =2e_m^2-(Q+1)e_m
       <2e_m^2=2J_m,
 \tag{B3O.6}
\]
which proves (B3O.2) for every $m$ and every such prime.

The recovered source \cite[B.3, formula~(B.CL.2)]{A1run}
requires, for a surjection
$J_C(\mathbf F_Q)\twoheadrightarrow\mathbf Z/\ell$
with kernel $K$, both $N_m\ge2|K|$ and $\ell>N_m$,
in addition to $2$ being primitive modulo $\ell$.
Since $|K|=J_m/\ell$, its first inequality is exactly
$\ell N_m\ge2J_m$, contradicted by (B3O.2).
No primitive-divisor theorem or large-prime-factor theorem
can universally force these sufficient inequalities:
they are false for this ordinary curve over every extension
of its stated base field.

Finally A13 proves that the same model over $\mathbf F_2$
has at least four Raynaud-theta points of exact order
$4159$ over $\mathbf F_{128}=\mathbf F_{2^7}$, and that
$\langle2,-1\rangle=(\mathbf Z/4159)^\times$.
The inclusion
$\mathbf F_{2^7}\subset\mathbf F_{2^{14}}=\mathbf F_{4^7}$
preserves their order and theta incidence.  They therefore
remain valid geometric witnesses on the curve (B3O.1)
over its specified base $\mathbf F_4$.
The primitive-root condition in B.3 is stronger than this
signed-doubling condition: A13 has
$\operatorname{ord}_{4159}(2)=2079$, not $4158$.
The two arguments are consistent because B.3 gives only
one sufficient test and asserts no converse.
\end{proof}

\paragraph{Exact source locator for A23.}
The recovered B.3 is the complete claim and proof titled
\emph{The corrected pair estimate and its precise limitation}
in \path{ANABELIAN_A1_paper.tex}, lines 2488--2546 of the
retained checkpoint-03 text.  A13 supplies formulas
(LW.12)--(LW.17) and its complete theta-incidence proof.
The all-extension inequality (B3O.6) is symbolic; the
accompanying finite recurrence checks are arithmetic
cross-checks, not the proof of its universal quantifier.


\subsection{A Frobenius-orbit obstruction at inert prime powers}
\label{sec:frobenius-sidon}

\paragraph{Finding A24 (a necessary congruence for the A13 curve).}
Let $C/\mathbf F_2$ be the ordinary genus-two curve of A13,
\[
             y^2+y=x+x^{-1}+(x-1)^{-1},
\]
and let $X=C^{(1)}$, $J=\operatorname{Pic}^0(X)$ and $\Theta$
be its Raynaud theta support.  Let $n=\ell^r$, where $\ell$
is an odd prime and $r\ge1$, and assume the actual finite-level
condition
\[
                   \langle2,-1\rangle=(\mathbf Z/n)^\times.
                                                        \tag{SO.1}
\]
If $\Theta$ contains a point of exact order $n$, then $\ell$
is not inert in $\mathbf Q(\sqrt{-15})$.  Equivalently,
\[
        \ell\in\{3,5\}\quad\hbox{or}\quad
                    \ell\equiv1,2,4,8\pmod {15}.          \tag{SO.2}
\]
This is an unconditional geometric restriction.  In particular,
allowing higher powers satisfying (SO.1) does not restore an
inert underlying prime.  No assertion that every prime lifts to
all powers, and no general non-Wieferich hypothesis, is used.

The surviving arithmetic candidates are not exhausted by this
restriction.  Under premise M, the set of split primes satisfying
$\langle2,-1\rangle=(\mathbf Z/\ell)^\times$ has natural density
\[
                         \frac{72}{95}\mathcal A>0,        \tag{SO.3}
\]
where $\mathcal A$ is Artin's constant.  The smaller set of split
primes for which $2$ itself is primitive has density
$48\mathcal A/95$.  These conditional density statements supply
arithmetic candidates and do not assert theta incidence.

\begin{proof}[Proof of the unconditional restriction]
\emph{1. The finite-subgroup bound.}
For every nonzero $T\in J(\overline{\mathbf F}_2)$, A19 proves
\[
                   \#\bigl(\Theta\cap(\Theta+T)\bigr)\le2.
                                                        \tag{SO.4}
\]
Here and below these are numbers of geometric points.  If $H$ is
any finite subgroup of $J(\overline{\mathbf F}_2)$ and
$A\subseteq\Theta\cap H$, count the ordered pairs of distinct
points of $A$ by their nonzero difference.  Each difference has
at most two representations by (SO.4).  Thus
\[
                         |A|(|A|-1)\le2(|H|-1).           \tag{SO.5}
\]
In particular, for a prime-order point $z\in\Theta$ and its
Frobenius-generated $\mathbf F_\ell$-space $W$ of dimension $d$,
its Frobenius orbit $O$ satisfies
$|O|(|O|-1)\le2(\ell^d-1)$.  The same counting proof applies
to finite modules of prime-power exponent; it is not restricted
to vector spaces.

We will also use symmetry.  The characteristic-two locally exact
differential line bundle satisfies $B^{\otimes2}\simeq\omega_X$.
Riemann--Roch in degree one therefore gives
\[
 h^0(X,B\otimes L)=h^0(X,B\otimes L^{-1})
                        \qquad(L\in J),
\]
so $\Theta=-\Theta$.  This uses the same Raynaud line bundle as
in A03 and A11, not an arbitrary translate of the Abel curve.

\emph{2. The exact quadratic Frobenius relation.}
The equation and its Frobenius twist have the same model over
$\mathbf F_2$.  Let $\psi$ be $2$-power Frobenius and put
$\Phi=\psi^2$, the $4$-power Frobenius.  The checked A13 point
counts give the characteristic polynomial
\[
                    P_2(T)=T^4-T^2+4.
\]
Cayley--Hamilton on each odd Tate module gives the integral
identity
\[
                      \Phi^2-\Phi+4=0.                    \tag{SO.6}
\]
Thus the quadratic annihilator is established directly; it is
not inferred from the squared characteristic polynomial over
$\mathbf F_4$ without a semisimplicity argument.

Suppose henceforth that $\ell$ is inert.  Then
$Q(T)=T^2-T+4$ is irreducible modulo $\ell$, since its
discriminant is $-15$.  In particular $\ell\notin\{3,5\}$.
Let $\beta$ denote the class of $T$ in
$\mathbf F_{\ell^2}=\mathbf F_\ell[T]/(Q)$.  Its norm is
\[
                         \beta^{\ell+1}=4.                \tag{SO.7}
\]
On every nonzero vector of $J[\ell]$, the $\mathbf F_\ell[\Phi]$
span is a one-dimensional $\mathbf F_{\ell^2}$-space, hence a
two-dimensional $\mathbf F_\ell$-space, and the $\Phi$-orbit has
length
$s_1=\operatorname{ord}_{\mathbf F_{\ell^2}^{\times}}(\beta)$.

\emph{3. A uniform signed-orbit bound at the residue level.}
Write $h=(\ell-1)/2$.  Reduction of (SO.1) modulo $\ell$ gives
the signed-unit condition there.  In a cyclic group of even
order, adjoining $-1$ can enlarge a subgroup by at most two.
Consequently the order of $2$ is either $\ell-1$, or $h$ with
$h$ odd.  In either case
\[
                       \operatorname{ord}_\ell(4)=h.       \tag{SO.8}
\]
Let $H_1=\langle\beta,-1\rangle\subset\mathbf F_{\ell^2}^\times$.
Its norm image is $\langle4\rangle$, of order $h$, while the
norm kernel contains $-1$ and therefore has even order.
If this kernel had order two, then $|H_1|=2h=\ell-1$.
The unique subgroup of that order in the cyclic group
$\mathbf F_{\ell^2}^\times$ is $\mathbf F_\ell^\times$.
This would put $\beta$ in $\mathbf F_\ell$, contradicting
irreducibility.  The kernel consequently has order at least
four, and
\[
                          |H_1|\ge4h=2(\ell-1).            \tag{SO.9}
\]
This argument includes every inert eligible prime; it does not
require a large-prime cutoff or a list of exceptional cases.

\emph{4. Actual prime-power levels.}
Put $\mathcal O=\mathbf Z_\ell[T]/(Q)$, the ring of integers
of the unramified quadratic extension of $\mathbf Q_\ell$.
Identity (SO.6) makes $T_\ell J$ an $\mathcal O$-module.
The quadratic polynomial is irreducible over $\mathbf Q_\ell$,
so a nonzero element of $\mathcal O$ acts invertibly on
$T_\ell J\otimes\mathbf Q_\ell$.  The Tate module is therefore
torsion free over the discrete valuation ring $\mathcal O$,
and is free of rank two over it.  For
$R=\mathcal O/\ell^r\mathcal O$ this gives
\[
                              J[\ell^r]\simeq R^2.
\]
Let $z$ have exact additive order $\ell^r$.  In these coordinates
at least one component of $z$ is a unit of the local principal
ring $R$; otherwise its order would be at most $\ell^{r-1}$.
Thus $a\mapsto az$ is an injection $R\hookrightarrow R^2$.
The Frobenius-generated subgroup $H=Rz$ consequently has
exactly $\ell^{2r}$ points, and the $\Phi$-orbit of $z$ has
length $s_r=\operatorname{ord}_{R^\times}(\beta)$.

The reduction kernel $R^\times\to\mathbf F_{\ell^2}^\times$
is an $\ell$-group of exponent dividing $\ell^{r-1}$.
The latter assertion follows by repeated binomial expansion of
$1+\ell a$, with $\ell$ odd.  Therefore
\[
                           s_r=s_1\ell^b,
                         \qquad0\le b\le r-1.              \tag{SO.12}
\]
The norm $R^\times\to(\mathbf Z/\ell^r)^\times$ sends
$\beta$ to $4$.  The cyclic-unit argument used in (SO.8),
now applied to the actual condition (SO.1), gives
\[
                   \operatorname{ord}_{\ell^r}(4)
                         =h\ell^{r-1}.                    \tag{SO.13}
\]
This norm order divides $s_r$.  As $s_1$ is prime to $\ell$,
(SO.12)--(SO.13) force $b=r-1$.

Put $H_r=\langle\beta,-1\rangle\subset R^\times$.
If $s_1$ is odd, then $s_r=s_1\ell^{r-1}$ is odd, so adjoining
$-1$ doubles both orders.  If $s_1$ is even, then $-1$ already
belongs to both cyclic groups: the unique element of order two
in the units of this odd-residue-characteristic local ring is
$-1$.  To see uniqueness, $a^2=1$ implies
$(a-1)(a+1)=0$, and one of the factors is a unit because
$2$ is a unit; hence $a=1$ or $a=-1$.
It follows in both cases, with $k=\ell^{r-1}$, that
\[
                      |H_r|=k|H_1|\ge2(\ell-1)k.           \tag{SO.14}
\]
If $z\in\Theta$, Frobenius invariance and symmetry put its
entire signed orbit $H_rz$ in $\Theta\cap Rz$.  The injection
$R\hookrightarrow Rz$ makes its size exactly $|H_r|$.
Every inert prime here is at least seven, since $3,5$ ramify.
For $\ell\ge5$,
\[
\begin{split}
 |H_r|(|H_r|-1)
  &\ge k^2\,2(\ell-1)(2\ell-3)\\
  &>2\ell^2 k^2
   >2(\ell^{2r}-1).
\end{split}                                                \tag{SO.15}
\]
The middle inequality is
$2(\ell-1)(2\ell-3)-2\ell^2=2\ell^2-10\ell+6>0$.
The first follows from $k\ge1$ and
$(ak)(ak-1)\ge k^2a(a-1)$ for $a\ge1$.
This contradicts (SO.5), since $|Rz|=\ell^{2r}$.
This proves the stated
prime-power exclusion.  Quadratic reciprocity, or the two
quadratic characters of conductors $3$ and $5$, gives the split
classes $1,2,4,8$ modulo $15$, proving (SO.2).
\end{proof}

\paragraph{A finite consistency check, not a proof premise.}
Exact orders of $\beta$ in the three smallest inert residue fields
are
\[
\begin{array}{c|r|r|r}
 \ell&\operatorname{ord}(\beta)&s_1(s_1-1)&2(\ell^2-1)\\\hline
 7&24&552&96\\
 11&60&3540&240\\
 13&84&6972&336
\end{array}
\]
The accompanying \texttt{check\_A24\_orbits.py} verifies these
orders by quadratic-ring multiplication and an independent
$2$-by-$2$ multiplication-matrix implementation.  The uniform
norm-kernel proof above does not depend on this finite check.

\begin{proof}[Proof of the conditional residual densities]
Apply Lenstra's fixed finite Frobenius-condition theorem to
$W=\langle2,-1\rangle\subset\mathbf Q^\times$ and the
splitting condition in $F=\mathbf Q(\sqrt{-15})$.
The sufficient analytic family is
\[
 \zeta_{L_n}(s),\qquad
 L_n=\mathbf Q\bigl(\zeta_n,2^{1/n},(-1)^{1/n}\bigr),
                   \qquad n\ge1\text{ squarefree}.          \tag{SO.16}
\]
These Dedekind zeta functions are contained in premise M.
Lenstra's Section~2 permits an infinite finitely generated
subgroup with torsion, so it applies to this $W$ directly.
Theorem~3.1 and its Section~3.2 argument allow the fixed field
$F$ without requiring a separately enlarged GRH family
\cite[\S\S2--3, Theorem~3.1]{Lenstra1977}.

For odd squarefree $m$, write
$K_m=\mathbf Q(\zeta_m,2^{1/m})$.  As proved in A17,
\[
 [K_m:\mathbf Q]=m\varphi(m),\qquad
 K_m\cap\mathbf Q(\zeta_8)=\mathbf Q,
\]
and the maximal abelian subfield of $K_m$ is
$\mathbf Q(\zeta_m)$.  For completeness, Eisenstein at $2$
gives the degree, while the affine Galois group
$C_m\rtimes(\mathbf Z/m)^\times$ has commutator $C_m$:
the cyclotomic element $2$ conjugates a radical generator
$\tau$ to $\tau^2$.  The assertion about the abelian subfield
and the coprime-conductor intersection follow.

Taking odd roots of $-1$ adds nothing.  Taking a $2m$th root
of $-1$ adds $i$, and a $2m$th root of $2$ is generated by
$\sqrt2$ and $2^{1/m}$, using Bezout for $2,m$.  Thus
\[
 L_m=K_m,\qquad L_{2m}=\mathbf Q(\zeta_8)K_m,\qquad
 [L_m:\mathbf Q]=m\varphi(m),\quad
 [L_{2m}:\mathbf Q]=4m\varphi(m).                           \tag{SO.17}
\]
The maximal abelian subfields are respectively
$\mathbf Q(\zeta_m)$ and $\mathbf Q(\zeta_{8m})$.
The quadratic field $F$ has discriminant $-15$, is contained
in $\mathbf Q(\zeta_{15})$, and ramifies at both $3$ and $5$.
It follows that, for either $n=m$ or $n=2m$,
\[
 F\cap L_n=
 \begin{cases}
 F,&15\mid m,\\
 \mathbf Q,&15\nmid m.
 \end{cases}                                               \tag{SO.18}
\]
Indeed the ramification at a missing one of $3,5$ excludes
containment, while $15\mid m$ gives the containment directly.

For each finite collection of prime obstructions, ordinary
Chebotarev and inclusion-exclusion give the corresponding
finite sum of $\mu(n)/[FL_n:\mathbf Q]$.  Lenstra's theorem,
under (SO.16), identifies its limit with natural density.
The absolutely convergent sums can now be evaluated.  First
the total eligible density, without the splitting condition, is
\[
 D=\sum_{n\ {\rm squarefree}}\frac{\mu(n)}{[L_n:\mathbf Q]}
   =\frac34\prod_{q\ {\rm odd\ prime}}
                     \left(1-\frac1{q(q-1)}\right)
   =\frac32\mathcal A.
\]
Forcing $15\mid n$ changes the Euler factors at $3,5$ by
\[
 \frac{-1/6}{1-1/6}\cdot\frac{-1/20}{1-1/20}
                            =\frac1{95}.
\]
By (SO.18), the split density is therefore
\[
 \sum_n\frac{\mu(n)}{[FL_n:\mathbf Q]}
       =\frac12D+\frac12\frac{D}{95}
       =\frac{72}{95}\mathcal A.
\]
For $W=\langle2\rangle$ the even field is instead
$K_{2m}=K_m(\sqrt2)$ of degree $2m\varphi(m)$; the same
intersection argument applies.  The total density is then
$\mathcal A$, and its split part is $48\mathcal A/95$.
A sufficient GRH family for this smaller-set assertion is
$\zeta_{K_n}$ for squarefree $n$, as in A17.  Premise M
contains both stated sufficient families.
\end{proof}

\paragraph{Why the remaining candidates are not eliminated.}
The actual A13 witness $4159$ is split: $4159\equiv4\pmod{15}$.
Thus this same curve already prevents interpreting the
restriction as the absence of all eligible witnesses.
More generally, at a split prime away from the discriminant,
the quadratic Frobenius algebra is
$\mathbf F_\ell\times\mathbf F_\ell$.  For a point with
both eigencomponents nonzero, its generated space has dimension
two but its Frobenius orbit has at most $\ell-1$ points.
Such an orbit always satisfies the numerical inequality
$(\ell-1)(\ell-2)\le2(\ell^2-1)$, so this counting test
alone gives no contradiction.  An eigenline point can have a
shorter orbit; the A13 points over $\mathbf F_{128}$ have
$4$-power Frobenius orbit length dividing $7$.
The signed-unit condition controls the norm order; it does not
make theta invariant under multiplication by $2$.
Consequently A24 gives a necessary curve-dependent congruence
condition for prime and prime-power searches, while the
variable-prime theta-incidence gate of A18 remains open.


\subsection{A pairing obstruction to enlarging the rational torsion subgroup}
\label{sec:lang-pairing-rank-bound}

\paragraph{Finding A25 (a prime-level Lang test forces torsion rank at most two).}
Let $C/\mathbf F_Q$ be a genus-two curve in characteristic two,
let $J$ be its Jacobian, and put $G=J(\mathbf F_Q)$.
Let $r$ be an odd prime and
$\kappa=\dim_{\mathbf F_r}G[r]$.
Suppose either the nonzero-point numerical test (LW.2), or
the weaker ordinary-case positivity test following (LW.1),
holds for the subgroup $H$ under consideration.
Then the following implications hold:
\[
\begin{array}{c|c|c}
\text{hypothesis on }2\pmod r&H&\text{range forcing }\kappa\le2\\
\hline
\operatorname{ord}_r(2)=r-1&G[r]&r\ge11\\
\operatorname{ord}_r(2)=r-1&G[3r]&r\ge17\\
\langle2,-1\rangle=(\mathbf Z/r)^\times&G[r]&r\ge29\\
\langle2,-1\rangle=(\mathbf Z/r)^\times&G[3r]&r\ge41.
\end{array}
\tag{LP.1}
\]
In each range, writing $s=\sqrt Q$ and $J_Q=|G|$, the full
test (LW.2) therefore requires respectively
\[
 r>\left(J_Q\frac{2s+1}{(s-1)^2}\right)^{1/2}
 \quad\text{or}\quad
 r>\frac19\left(J_Q\frac{2s+1}{(s-1)^2}\right)^{1/2}.
\tag{LP.2}
\]
These thresholds have asymptotic sizes
$\sqrt2 Q^{3/4}$ and $(\sqrt2/9)Q^{3/4}$.
The corresponding ordinary thresholds replace $2s+1$ by $2s$.
These are necessary conditions for the specified sufficient
numerical tests, not for the existence of theta points.

\begin{proof}
The canonical principal polarization of the Jacobian gives a
nondegenerate alternating pairing
\[
                e_r:J[r]\times J[r]\longrightarrow\mu_r
\tag{LP.3}
\]
on its four-dimensional geometric $r$-torsion space, equivariant
for Frobenius.  For the pairing and its equivariance see
\cite[\S16, p.~29, and Lemma~16.2(e), p.~30]{MilneAV};
the principal polarization identifies the dual torsion space
with the original one.
If $\kappa\ge3$, the subspace $G[r]$ is not totally isotropic:
a totally isotropic subspace $W$ satisfies
$W\subseteq W^\perp$ and hence
$\dim W\le4-\dim W$.
Choose $P,T\in G[r]$ with $e_r(P,T)\ne1$.
Since $r$ is prime, this pairing value has exact order $r$.
Both points are fixed by Frobenius, so equivariance gives
\[
                         e_r(P,T)^Q=e_r(P,T).
\]
It follows that $r\mid Q-1$.
Writing $Q=2^v$ and $d=\operatorname{ord}_r(2)$, we obtain
\[
                         d\mid v,\qquad Q\ge2^d.
\tag{LP.4}
\]

The weaker of the two Lang tests in the statement already
implies
\[
 |H|>\frac{2sJ_Q}{(s-1)^2}
       \ge 2s(s-1)^2,
\tag{LP.5}
\]
where the last inequality uses $J_Q\ge(s-1)^4$.
For $Q\ge16$, the last expression is at least
$(9/8)Q^{3/2}$, because
$2(1-1/s)^2\ge2(3/4)^2=9/8$.
In particular, either successful test requires
$|H|>Q^{3/2}$ in this range.

If $H=G[r]$, then $|H|=r^\kappa\le r^4$.
If $H=G[3r]$, the ranges in (LP.1) exclude $r=3$, and
primary decomposition gives
$|H|=|G[3]|\,|G[r]|\le81r^4$.
If $\kappa\ge3$, (LP.4) now forces failure of (LP.5)
whenever
\[
                         c r^4<2^{3d/2},
\tag{LP.6}
\]
where $c=1$ or $81$ respectively.

For a primitive root, $d=r-1$.
The two starting inequalities are
\[
 11^4=14641<2^{15},\qquad
 81\cdot17^4=6765201<2^{24}.
\]
They persist on increasing the odd integer $r$ by two:
the left side grows by a factor at most
$(13/11)^4<2$, whereas the right side grows by $8$.
Thus (LP.6) holds in the first two ranges of (LP.1).

Under the signed-doubling hypothesis,
$r-1=|\langle2,-1\rangle|\le2d$, so
$d\ge(r-1)/2$.
For the other two ranges it suffices to check
\[
 29^4=707281<2^{21},\qquad
 81\cdot41^4=228886641<2^{30}.
\]
On increasing $r$ by two, the lower bound
$2^{3(r-1)/4}$ on the right of (LP.6) grows by
$2^{3/2}>2$, whereas the left side again grows by a
factor less than $2$.
All four ranges also imply $Q\ge16$ through (LP.4).
Consequently $\kappa\ge3$ contradicts either successful
Lang test, proving (LP.1).

Finally, solving (LW.2) for the subgroup size gives
$|H|>J_Q(2s+1)/(s-1)^2$.
With $\kappa\le2$, one has $|G[r]|\le r^2$ and
$|G[3r]|\le81r^2$, which proves (LP.2).
The asymptotics follow from $J_Q=Q^2+O(Q^{3/2})$.
\end{proof}

\paragraph{Implication for the proposed arithmetic route.}
In particular, first passing to an extension where all of
$J[r]$ is rational cannot force the prime-level Lang test
for an Artin prime $r\ge11$: it creates $\kappa=4$, the
excluded situation in (LP.1).
The factor $|G[3]|\le81$ can improve an individual index,
but the analogous strategy with $G[3r]$ fails for Artin
primes $r\ge17$.
This argument concerns these rational subgroup estimates.
It does not remove geometric theta witnesses, and does not
assert the same rank bound for truncated prime-power
subgroups $G[r^b]$ without an additional argument.


\subsection{A parity improvement for prescribed-point rational covers}
\label{sec:rank-one-augmentation}

\paragraph{Finding A26 (the rank-one augmentation criterion).}
Let $C/\mathbf F_Q$ be a smooth proper geometrically connected
genus-two curve in characteristic two, let $P\in C(\mathbf F_Q)$,
and suppose $N=\#C(\mathbf F_Q)$ is even.
Let $r$ be an odd prime with $\operatorname{ord}_r(2)=r-1$, and set
\[
 G=J_C(\mathbf F_Q),\qquad V=G/rG.
\tag{RA.1}
\]
If $|V|>N$, some surjection $G\twoheadrightarrow\mathbf Z/r$
gives a constant, geometrically connected, $P$-split cyclic
\'etale Lang cover whose relative class number is odd.
If $r\mid|G|$ and $r>N$, every such surjection has this property.
In both statements the prescribed point $P$ is arbitrary.
No lower bound on $N$ in terms of the kernel size is required.

Every genus-two curve of $2$-rank one satisfies the stated
parity hypothesis over its finite field of definition and
over every finite extension.  Thus both assertions apply
to every such curve whenever their remaining numerical
hypotheses hold.

\begin{proof}
We first prove the parity assertion.  The hyperelliptic
degree-two map in characteristic two is an Artin--Schreier
cover.  If its reduced geometric poles have odd orders
$d_1,\ldots,d_t$, its genus and $2$-rank are
\[
       g=\frac{\sum_{i=1}^t(d_i+1)-2}{2},
       \qquad f_C=t-1;
\]
see \cite[Lemma~2.6, p.~712]{PriesZhu2012}.
For $g=2$ and $f_C=1$, there are two poles, with
orders $1$ and $3$.  Frobenius preserves the ramification
orders, so it fixes both branch points individually.
Each has one geometric point above it, which is therefore
rational.  These are exactly the two fixed points of the
hyperelliptic involution.  Every other rational point is
paired with a distinct rational point by that involution.
Consequently $N$ is even.  The same argument applies over
every extension; in particular the two branch points remain
rational.

Now assume only the evenness of $N$.
In the odd-order group algebra $\mathbf F_2[V]$, put
\[
 \Theta_V=[0]+\sum_{R\in C(\mathbf F_Q)}[[R-P]\bmod rG].
\tag{RA.2}
\]
Its augmentation is $N+1=1$ in $\mathbf F_2$.
Recovered A1 B.2 gives the following exact criterion:
for a surjection $a:G\to\mathbf Z/r$, the associated
$P$-split Lang cover has odd relative class number if and
only if every nontrivial character of $\mathbf Z/r$ has
nonzero evaluation on $a(\Theta_V)$
\cite[B.2]{A1run}.
The connected Lang-pullback argument of A12, applied
to the Abel map based at $P$, proves geometric connectedness.
The zero point above $[P-P]=0$ gives the rational point
above the prescribed $P$, and the constant deck group
makes that fibre split.

Suppose all surjections to $\mathbf Z/r$ fail the criterion.
The irreducibility of $\Phi_r$ over $\mathbf F_2$, which
follows from $\operatorname{ord}_r(2)=r-1$, means that
one vanishing nontrivial character evaluation for such a
quotient forces all its nontrivial character evaluations
to vanish.  Every nontrivial character of $V$ factors
through one of these quotients.  Thus all nontrivial
Fourier coefficients of $\Theta_V$ vanish.
Fourier inversion is valid because $|V|$ is odd, and gives
\[
        \Theta_V=c\sum_{b\in V}[b],\qquad
        c=\frac{N+1}{|V|}=1\quad\text{in }\mathbf F_2.
\tag{RA.3}
\]
But $|V|>N$ leaves at least one empty Abel-image coset.
The zero coset is occupied by $P$, so an empty coset is
nonzero.  Its coefficient in (RA.2) is zero, contradicting
(RA.3).  Therefore some quotient succeeds.

For the assertion about every surjection, fix one
$a:G\to\mathbf Z/r$ and suppose it fails.
Apply the same argument to
$\Theta_a=[0]+\sum_R[a([R-P])]$ in
$\mathbf F_2[\mathbf Z/r]$.
Irreducibility and the augmentation force
$\Theta_a=\sum_{b\in\mathbf Z/r}[b]$.
Since $r>N$, an empty nonzero coset again contradicts
that equality.  As $a$ was arbitrary, every quotient
succeeds.  A surjection exists when $r\mid|G|$ by the
structure theorem for finite abelian groups.
\end{proof}

\paragraph{Relation to the recovered row statements.}
The Fourier identity (RA.3) is already the displayed
identity in recovered successor A.35, and the empty-coset
argument also occurs in A.10 \cite[A.10 and A.35]{ASrun}.
Those statements retain an additional pair-count
inequality; the present even-$N$ specialization obtains
an immediate contradiction from augmentation and does
not require that inequality.
Recovered A1 B.1 then gives the prescribed-point
separation implication for the resulting odd cyclic cover.
This is a rational relative-class-number condition over
the displayed field, not geometric primitive $2$-rank
zero; it need not survive an arbitrary constant extension.
The existence of a prime satisfying (RA.1) and $|V|>N$
is an additional requirement, not a consequence of
Hooley's positive density statement alone.


\subsection{A rank-one curve that defeats every augmentation size test}
\label{sec:rank-one-augmentation-obstruction}

\paragraph{Finding A27 (the augmentation size test can fail over every extension).}
Let $u\in\mathbf F_4$ satisfy $u^2+u+1=0$, and let
$C/\mathbf F_4$ be the smooth projective model of
\[
                    y^2+y=ux^3+x+u/x.
\tag{RO.1}
\]
This is a genus-two curve of $2$-rank one.
For every $m\ge1$, put $Q=4^m$, $G_m=J_C(\mathbf F_Q)$
and $N_m=\#C(\mathbf F_Q)$.
For every odd prime $r$ with $\operatorname{ord}_r(2)=r-1$,
one has
\[
                         |G_m/rG_m|<N_m.
\tag{RO.2}
\]
In particular, no prime and constant extension of this
curve can satisfy either sufficient numerical criterion
of A26.  This statement does not exclude rational
class-number witnesses that fail those sufficient
inequalities, or geometric theta witnesses.

\begin{proof}
\emph{1. Exact Frobenius data.}
The two Artin--Schreier poles of (RO.1) have odd orders
$3$ and $1$, so the genus and $2$-rank formulas give
$g=2$ and $f_C=1$ \cite[Lemma~2.6]{PriesZhu2012}.
The two ramified points are rational.
Direct finite-field enumeration gives
\[
                 N_1=6,\qquad N_2=28.
\tag{RO.3}
\]
The accompanying independent checker enumerates the
solutions in $\mathbf F_4$ and $\mathbf F_{16}$ directly,
including the two ramified points.  It embeds $u$ as
$z^2+z$ in $\mathbf F_{16}=\mathbf F_2[z]/(z^4+z+1)$.
The trace-zero method in the retained 72-model search
gives the same counts by a separate enumeration.

If $P_4(T)$ is the Frobenius characteristic polynomial,
(RO.3) gives $a_1=N_1-5=1$ and
$a_2=(a_1^2-(17-N_2))/2=6$.
The genus-two functional equation therefore gives
\[
 \begin{aligned}
 P_4(T)&=T^4+T^3+6T^2+4T+16\\
       &=(T^2-T+4)(T^2+2T+4)=A(T)B(T).
 \end{aligned}
\tag{RO.4}
\]
Only this exact polynomial identity is used below;
no elliptic isogeny classification is needed.

Let $t_m$ and $w_m$ be the sums of the $m$th powers of
the roots of $A$ and $B$, respectively.  They are integers,
with recurrences
\[
\begin{gathered}
 t_0=2,\quad t_1=1,\quad t_m=t_{m-1}-4t_{m-2},\\
 w_0=2,\quad w_1=-2,\quad w_m=-2w_{m-1}-4w_{m-2}.
\end{gathered}
\tag{RO.5}
\]
Put $e_m=Q+1-t_m$ and $h_m=Q+1-w_m$.
The Frobenius formulas \cite[Theorem~19.1(b)]{MilneAV} yield
\[
             |G_m|=e_mh_m,\qquad N_m=Q+1-t_m-w_m.
\tag{RO.6}
\]
All roots of $A$ and $B$ have complex absolute value $2$,
so
\[
 e_m\le Q+1+2\sqrt Q,
 \qquad N_m\ge L_m:=4^m+1-4\cdot2^m.
\tag{RO.7}
\]
Moreover $t_m\equiv1\pmod4$ for $m\ge1$ by (RO.5),
and hence
\[
                               4\mid e_m.
\tag{RO.8}
\]

\emph{2. The rational torsion rank for an eligible prime.}
Write $\kappa_m(r)=\dim_{\mathbf F_r}G_m[r]$.
Finite abelian-group structure gives
$|G_m/rG_m|=r^{\kappa_m(r)}$, and
$0\le\kappa_m(r)\le4$.
Let $F$ be the $4$-power Frobenius acting on $J_C[r]$.
Then $G_m[r]=\ker(F^m-1)$.

For $r\ge5$, the polynomials $A,B$ are coprime modulo $r$,
because their resultant is $36$: their difference is $3T$
and $A(0)=4$.
The characteristic polynomial (RO.4), Cayley--Hamilton and
Bezout therefore split the four-dimensional torsion space
as two $F$-stable, two-dimensional spaces
$W_A\oplus W_B$ on which $A(F)=0$ and $B(F)=0$,
respectively.  This uses no semisimplicity assumption.
Since
\[
                         (T-2)B(T)=T^3-8,
\]
one has $F^3=8$ on $W_B$.
If $F^m$ fixes a nonzero vector in $W_B$, it follows that
$2^{3m}\equiv1\pmod r$.
For an eligible prime this means
\[
                              r-1\mid3m.
\tag{RO.9}
\]
If $\kappa_m(r)\ge2$ and $W_B$ contributes no fixed
vector, all of $W_A$ is fixed by $F^m$.
Taking its determinant gives $4^m\equiv1\pmod r$,
so $r-1\mid2m$.
Consequently
\[
 r\ge5,\quad\kappa_m(r)\ge2
                  \quad\Longrightarrow\quad r\le3m+1.
\tag{RO.10}
\]

\emph{3. Every extension degree $m\ge10$.}
At $m=10$,
\[
       (3m+1)^4=923521<1044481=L_{10}.
\]
Also
$L_{m+1}-4L_m=8\cdot2^m-3>0$, while for $m\ge10$
\[
 \frac{(3(m+1)+1)^4}{(3m+1)^4}
          \le(34/31)^4<2.
\]
Induction gives
\[
                            (3m+1)^4<L_m\le N_m
                            \qquad(m\ge10).
\tag{RO.11}
\]
For $r=3$, $|G_m/3G_m|\le81<(3m+1)^4<N_m$.
For $r\ge5$ and $\kappa_m(r)\ge2$, (RO.10)--(RO.11)
prove (RO.2).  If $\kappa_m(r)=0$, then
$|G_m/rG_m|=1<N_m$.

It remains to consider $\kappa_m(r)=1$.
If $r\mid e_m$, (RO.8) and the fact that $r$ is odd give
\[
 r\le e_m/4\le\frac{Q+1+2\sqrt Q}{4}
                    <Q+1-4\sqrt Q\le N_m.
\tag{RO.12}
\]
The strict middle inequality holds for $\sqrt Q\ge8$:
after multiplying by four it is
$3((\sqrt Q)^2-6\sqrt Q+1)>0$.
If $r\nmid e_m$, then $r\mid h_m$ by (RO.6), and
$F^m$ has a fixed vector on $W_B$.
Equation (RO.9) and (RO.11) again give $r<N_m$.
Since $|G_m/rG_m|=r$ in this case, (RO.2) follows.

\emph{4. The nine remaining extension degrees.}
The exact recurrences (RO.5) and trial division give the
following complete table.  The middle column lists
\emph{all} eligible odd primes dividing $|G_m|$.
The last column is the maximum of
$r^{\min\{4,v_r(|G_m|)\}}$ over those primes, or $1$
when there are none.  It is an upper bound on
$|G_m/rG_m|$ for every eligible prime, including those
not dividing $|G_m|$.
\[
\begin{array}{c|r|l|r}
m&N_m&\text{eligible prime divisors}&\text{upper bound}\\
\hline
1&6&\varnothing&1\\
2&28&3&9\\
3&60&19&19\\
4&256&3,5,13&13\\
5&996&\varnothing&1\\
6&3976&3,19&81\\
7&16764&\varnothing&1\\
8&66016&3,5,13&13\\
9&260340&19,181&361.
\end{array}
\tag{RO.13}
\]
Every entry in the last column is strictly smaller than
$N_m$.  The independent checker retains the full
factorizations, prime certificates by trial division,
and multiplicative orders of $2$, as well as the
point enumerations and both induction anchors.
Thus (RO.2) holds for the remaining degrees and hence
for all $m\ge1$.
\end{proof}

\paragraph{Scope of the obstruction.}
The result refutes universal success of the numerical
augmentation criterion, even on $2$-rank-one curves.
The conclusion of A26 itself is a sufficient conditional
statement and remains valid.
No class-number calculation for every actual cyclic
cover is asserted here, and the numerical failure does
not imply absence of a rational witness or of a geometric
primitive-zero cover.


\subsection{The recovered canonical-class character criterion}
\label{sec:recovered-canonical-character}

\paragraph{Cited and rechecked result A28 (recovered successor A.26).}
Let $C/k$ be a smooth proper geometrically connected genus-two
curve over a finite field of characteristic $p$, let $P\in C(k)$,
and put
\[
             G=J_C(k),\qquad\kappa_P=[K_C-2P]\in G.
\tag{RC.1}
\]
Suppose $r\ne p$ is an odd prime and
$-1\in\langle p\rangle\subset(\mathbf Z/r)^\times$.
If $\kappa_P\notin rG$, there is a constant, geometrically
connected, $P$-split cyclic \'etale cover of degree $r$
over $k$ whose relative class number is prime to $p$.
More precisely, every surjection $a:G\to\mathbf Z/r$
with $a(\kappa_P)\ne0$ gives such a witness.
No size condition on $r$, $G/rG$ or $C(k)$ is required.

This is recovered \texttt{ANABELIAN\_SUCCESSOR}, final row
A.26, \emph{canonical-class sufficient condition}, not a
new prime-level theorem.  Its complete proof is retained
in the source, with the prescribed-point and character
hypotheses explicit \cite[A.26]{ASrun}.

The complete retained source proof was read in this run, including
its middle-coefficient Riemann--Roch symmetry, Frobenius
permutation of reduced characters, and faithful-character
contradiction.  The Lang-pullback connectedness and the
prescribed split fibre are the construction checked in A12.
The existing prime theorem is cited here rather than reproved.


\paragraph{The rank-one and Weierstrass interpretation.}
In characteristic two, if $W\in C(k)$ is a rational
Weierstrass point, then $K_C\sim2W$ and
\[
                           \kappa_P=2[W-P].
\tag{RC.2}
\]
Multiplication by two is invertible on $G/rG$ for odd $r$.
Thus the condition in A28 is exactly
$[W-P]\notin rG$.
For a prime $r$, inversion is a power of two if and
only if $\operatorname{ord}_r(2)$ is even.
For an arbitrary odd composite modulus, even order
alone is insufficient: $2$ has order four modulo $15$,
but its powers never equal $-1$ modulo $15$.
This distinction is retained in using higher conductors.

Any two rational Weierstrass points $W,W'$ satisfy
$2[W-W']=0$.  Every odd-order quotient kills their
difference, so the condition is independent of the
choice of $W$.
Rank-one genus-two curves have the two rational
ramification points proved in A26, so this interpretation
always applies to that stratum.
When $P$ is one of those points, $\kappa_P=0$ and
this criterion supplies no witness.
More generally, failure of its detecting-class
hypothesis does not assert that all actual covers fail.
Hooley supplies no missing prime factor of the fixed
finite-order class $\kappa_P$; extending constants
preserves that class's order.


% A29: canonical-class detection on a cyclic prime-power packet.
% Authored in work/; never edits final/.
\subsection{A prime-power canonical-class criterion at a prescribed point}
\label{sec:prime-power-canonical}

\paragraph{Finding A29 (detecting a nonzero primary canonical class).}
Let $C/k$ be a smooth projective geometrically connected genus-two
curve over a finite field $k=\mathbf F_q$ of characteristic $p$, let
$P\in C(k)$, and put
\[
 G=J_C(k),\qquad \kappa_P=[K_C-2P]\in G.
 \tag{PC.1}
\]
Suppose that an odd prime $r\ne p$ satisfies
\[
 \operatorname{ord}_r(p)\ \text{is even},\qquad
 (\kappa_P)_r\ne0,
 \tag{PC.2}
\]
where the subscript denotes the $r$-primary component in the finite
abelian group $G$. Then there are an integer $a\ge1$ and a constant,
geometrically connected cyclic etale $P$-split cover
$D\to C$ of degree $n=r^a$ over $k$ itself such that, for its unique
intermediate quotient $E\to C$ of degree $r^{a-1}$,
\[
 \begin{aligned}
 \mathcal P&=\ker(\operatorname{Nm}:J_D\longrightarrow J_E)^0,\\
 \#\mathcal P(k)&=\frac{\#J_D(k)}{\#J_E(k)},
 \qquad p\nmid\#\mathcal P(k).
 \end{aligned}
 \tag{PC.3}
\]
Consequently the covering formulation of A1 separates the point
section of $P$ from the point section of every distinct $Q\in C(k)$.
The statement only needs the rational $p$-primary group of this
primitive factor to vanish; its geometric $p$-rank need not be zero.
It does not assert that $\#J_D(k)/\#J_C(k)$ is prime to $p$ when
$a>1$.

\paragraph{Recovered inputs and the refinement.}
The prime-degree canonical-class criterion is already proved in
recovered successor \emph{final row A.26}, in every characteristic:
if $-1\in\langle p\rangle\subset(\mathbf Z/r)^\times$ and
$\kappa_P\notin rG$, a prime-degree rational witness exists
\cite[A.26]{ASrun}. Its proof uses the middle-coefficient
Riemann--Roch identity used below. Recovered A1 B.13--B.14 already
identify the primitive factor in a cyclic prime-power tower as the
appropriate rational norm obstruction \cite[B.13--B.14]{A1run}.
The refinement here combines these inputs to replace
$\kappa_P\notin rG$ by the strictly weaker finite-group condition
$(\kappa_P)_r\ne0$. We prove the whole combination, including
its pointed-cover consequence. This attribution does not make a
claim of bibliographic priority beyond the recovered source audit.

\begin{proof}
\emph{1. A cyclic quotient detecting the primary component.}
Decompose the $r$-primary subgroup as
$G_r\simeq\bigoplus_i\mathbf Z/r^{e_i}$.
Since $(\kappa_P)_r\ne0$, at least one of its coordinates is
nonzero. Projection to that coordinate, preceded by the canonical
primary projection of $G$, gives a surjection
\[
 b:G\twoheadrightarrow A=\mathbf Z/r^a,
 \qquad b(\kappa_P)\ne0,
 \tag{PC.4}
\]
with $a=e_i$. Conversely, any cyclic $r$-power quotient detecting
$\kappa_P$ implies $(\kappa_P)_r\ne0$.
In particular (PC.4) does not require $\kappa_P$ to survive in
$G/rG$: an element can be nonzero in a cyclic $r^a$-factor while
being divisible by $r$ in that factor.

Let $d_a=\operatorname{ord}_{r^a}(p)$. Reduction modulo $r$
shows that $\operatorname{ord}_r(p)$ divides $d_a$, so $d_a$ is
even. The group $(\mathbf Z/r^a)^\times$ is cyclic and its
unique element of order two is $-1$. Hence
\[
                  p^{d_a/2}\equiv-1\pmod{r^a}.
 \tag{PC.5}
\]
This uses the actual order at the chosen power. No assertion of
maximal order growth, or non-Wieferich hypothesis, is needed.

\emph{2. The constant pointed Lang cover.}
Write $\Phi$ for $q$-power Frobenius on $J=J_C$ and
$H=\ker b$. The map $\Phi-1$ is an etale isogeny with constant
kernel $J(k)=G$. It factors as
\[
 J\longrightarrow J/H\xrightarrow{\psi}J,
 \qquad \ker\psi=G/H=A.
\]
Pull $\psi$ back by $i_P:C\to J$, $R\mapsto[R-P]$, to obtain
$D=C\times_{J}(J/H)$. This is an etale cover of degree $n$ with
constant deck group $A$. Its fibre over $P$ is the constant kernel
of $\psi$, so it has a chosen rational point $P'$ above $P$ and
the whole fibre is rational.

For completeness, this pullback is geometrically connected.
Otherwise its geometric monodromy would be a proper subgroup
$A_0\subset A$. Quotienting by $A_0$ gives a nontrivial intermediate
etale isogeny $\psi':J'\to J$ whose pullback to $C$ has trivial
geometric monodromy. Over an algebraic closure, $i_P$ therefore
lifts to $v:C\to J'$. Subtract $v(P)$, which lies in $\ker\psi'$.
The Jacobian universal property gives a homomorphism $u:J\to J'$
with $u i_P=v-v(P)$ and $\psi'u=\operatorname{id}_J$.
The homomorphism $u$ is an isogeny, since its composition with
$\psi'$ is an isomorphism and both varieties have dimension two.
Multiplicativity of degree now contradicts $\deg\psi'>1$.
This is the connected Lang-pullback argument of A12.

\emph{3. The genus-two character polynomial and its symmetry.}
Let $\chi:A\to\mu_n$ be a nontrivial characteristic-zero
character. The unramified character Euler product for this cover is
\[
 L_P(\chi,T)
 =\sum_{d\ge0}\left(
   \sum_{Z\in\operatorname{Div}^{+}_d(C)(k)}
      \chi\bigl(b([Z-dP])\bigr)\right)T^d.
\]
The convention for Frobenius translation is the Lang convention
of recovered A1 B.2; reversing that convention replaces every
character by its inverse and leaves the argument unchanged.
In genus two its exact polynomial is
\[
 \begin{aligned}
 L_P(\chi,T)&=1+c_1(\chi)T
                  +q\chi\bigl(b(\kappa_P)\bigr)T^2,\\
 c_1(\chi)&=\sum_{R\in C(k)}\chi\bigl(b([R-P])\bigr).
 \end{aligned}
 \tag{PC.6}
\]
Here is a direct verification of the truncation. For $d>2$,
Riemann--Roch makes the number of effective divisors in each
degree-$d$ line-bundle class independent of the class. Summing
the nontrivial character over $\operatorname{Pic}^d(C)(k)$ gives
zero. In degree two, every line-bundle class has one effective
divisor except the canonical class, which has $q+1$. The excess
is $q\chi(b(\kappa_P))$. The degree-zero coefficient is one,
and degree one gives the displayed point sum. The rational point
$P$ identifies every degree component with $J(k)$ in these sums.

Riemann--Roch in degree one also gives
\[
 c_1(\chi)=\chi\bigl(b(\kappa_P)\bigr)c_1(\chi^{-1}).
 \tag{PC.7}
\]
Indeed $R\mapsto\iota R$, where $R+\iota R\sim K_C$, is a
bijection on $C(k)$. This follows directly because $K_C-R$ has
a unique effective divisor of degree one, defined over $k$.
The identity
$[\iota R-P]=\kappa_P-[R-P]$ proves (PC.7).
Thus no rational Weierstrass point is required.

\emph{4. Every faithful character is a unit at $p$.}
Reduce the cyclotomic integers modulo a fixed prime above $p$,
and denote the resulting character by $\bar\chi$. Since
$p\nmid n$, reduction preserves its order and its faithfulness.
Put $S_\chi=\overline{c_1(\chi)}$. If $j=d_a/2$, (PC.5)
and the integer coefficients of the point sum give
\[
 S_{\chi^{-1}}=S_\chi^{p^j},\qquad
 S_\chi=\bar\chi\bigl(b(\kappa_P)\bigr)S_\chi^{p^j}.
 \tag{PC.8}
\]
Suppose that $L_P(\chi,1)$ vanishes modulo that prime.
Equation (PC.6) gives $S_\chi=-1$, since $q=0$ in the residue
field. Its $p^j$th power is also $-1$, so (PC.8) forces
$\bar\chi(b(\kappa_P))=1$.
For a faithful $\chi$ this contradicts (PC.4). Therefore
\[
 L_P(\chi,1)\ \text{is a unit at }p
       \quad\text{for every character of exact order }r^a.
 \tag{PC.9}
\]
Nonfaithful characters may have zero reduction. The argument
deliberately makes no assertion about those factors.

\emph{5. The top-step Prym has prime-to-$p$ rational order.}
Choose a generator $\sigma$ of $A$ and set
$\tau=\sigma^{n/r}$, of order $r$; let $E=D/\langle\tau\rangle$.
Pullback and norm for $D\to E$ have composition $[r]$ on $J_E$.
They yield a $k$-isogeny
$J_E\times\mathcal P\to J_D$, so finite-field point counts give
the first equality in (PC.3).

For the exact character factorization, the Euler products of
$D/C$ and $E/C$ are products of the corresponding character
Euler factors. The characters for $E/C$ are exactly those of $A$
trivial on $\langle\tau\rangle$. Since $n$ is a prime power,
the remaining characters are exactly the faithful ones. Hence
\[
 \frac{\#J_D(k)}{\#J_E(k)}
       =\prod_{\substack{\chi\in\widehat A\\
                          \operatorname{ord}(\chi)=r^a}}
                       L_P(\chi,1).
 \tag{PC.10}
\]
This also follows by cancelling the imprimitive character factors
in the zeta-function ratios of recovered A1 B.2. The right side
is an integer, by its identification with $\#\mathcal P(k)$.
Each factor is a unit at the chosen prime above $p$ by (PC.9),
so the integer is prime to $p$. Thus
$\mathcal P(k)[p^\infty]=0$. The dimension check is
$\dim\mathcal P=(n-n/r)(g(C)-1)=\varphi(n)$.

\emph{6. The prime-power deck argument separates the point.}
Suppose that a distinct $Q\in C(k)$ collided with $P$ in the
covering formulation of A1. Applied to the pointed cover just
constructed, that formulation supplies $Q'\in D(k)$ above $Q$
with
\[
                  \delta=[Q'-P']\in J_D(k)[p^\infty].
 \tag{PC.11}
\]
This is the same collision-to-cover implication used in recovered
A1 B.1 and in A11; no converse to that implication is used.
The point $(\tau-1)\delta$ is rational, is $p$-primary, and lies in
\[
             \operatorname{im}(\tau-1)=\mathcal P.
 \tag{PC.12}
\]
To verify the equality, the image is connected and lies in the
norm kernel. Conversely, on $\mathcal P$ the endomorphism
$[r]-\sum_{i=0}^{r-1}\tau^i$ is $[r]$ and is divisible by
$\tau-1$ as a polynomial in $\tau$. Multiplication by $r$ is
surjective on $\mathcal P$, proving the reverse inclusion.
Equation (PC.3) therefore gives $(\tau-1)\delta=0$.

It follows that
$\tau Q'+P'-Q'-\tau P'$ is principal. Its four points are
distinct: the deck action is free, and the fibres over $P$ and
$Q$ are disjoint. Its positive and negative parts each have
degree two with simple support. The resulting rational function
gives a separable degree-two map $D_{\bar k}\to\mathbf P^1$;
the simple poles rule out a $p$th power if $p=2$, and the degree
itself excludes inseparability if $p\ne2$.

Since $g(D)=n+1\ge4$, its hyperelliptic pencil is unique.
The order-$r$ automorphism $\tau$ descends to an automorphism
of that projective line of order $r$: the kernel of the descent
action has order at most two. Over $\bar k$ this odd-order
automorphism has a fixed point. Its fibre contains at most two
geometric points, so an order-$r$ cyclic group acts trivially on
that fibre. This gives a fixed point of $\tau$ on $D$, contrary
to its free deck action. The assumed collision is impossible.
\end{proof}

\paragraph{The characteristic-two Weierstrass form and its exact residue.}
If $p=2$ and $W\in C(k)$ is a Weierstrass point, then
$K_C\sim2W$, so $\kappa_P=2[W-P]$. Multiplication by two is
an automorphism on every odd-primary subgroup. Consequently A29
applies precisely when $[W-P]$ has a nonzero $r$-primary component
for some odd prime $r$ with even $\operatorname{ord}_r(2)$.
Any two rational Weierstrass points differ by a class killed by
two, so this condition is independent of their choice.
It supplies a finite-field rational primitive-factor witness with
no size inequality and no analytic hypothesis.

This detector leaves unchanged those points for which the odd part
of $[W-P]$ is supported entirely on primes with odd
$\operatorname{ord}_r(2)$, including the case where $[W-P]$ is
purely $2$-primary. In particular it says nothing at a Weierstrass
point $P$, where $\kappa_P=0$. A constant extension cannot remove
this obstruction to the detector: the map on rational points of
the Jacobian is injective, so the finite order and primary support
of the fixed class $\kappa_P$ do not change. This is a limitation
of the sufficient construction, not a nonexistence theorem for
other cover witnesses. The A09 example already has non-Weierstrass
rational points whose differences from a rational Weierstrass
point are purely $2$-primary, because its rational Jacobian has
order eight.

\paragraph{Strictness on an actual curve.}
The difference from the prime-level canonical detector occurs on
the explicit curve of A09, without an abstract group realization
assumption. Write its point at infinity as $W$, and retain the
point $R=(u,v)\in C(\mathbf F_{16})$ whose class
$L=[R-W]$ has exact order $13$. For the prescribed point $R$,
\[
                   \kappa_R=-2L
 \tag{PC.13}
\]
also has exact order $13$. The finite etale surjective isogeny
$[13]:J_C\to J_C$ has a nonempty fibre over $\kappa_R$.
A point of this finite fibre is rational over some finite
extension $k'/\mathbf F_{16}$; write it as $U$, so that
$13U=\kappa_R$. Thus
$\kappa_R\in13J_C(k')$ while its $13$-primary component remains
nonzero. For any odd prime $r\ne13$, choose an integer $s$ with
$rs\equiv1\pmod{13}$. Then
$\kappa_R=r(s\kappa_R)$ in $J_C(k')$. Consequently
\[
                  \kappa_R\in rJ_C(k')
                    \quad\text{for every odd prime }r.
 \tag{PC.14}
\]
Every prime-level canonical detector in A28 therefore fails over
this field. Nevertheless $\operatorname{ord}_{13}(2)=12$, so
A29 constructs a cyclic $13^a$ pointed witness with the top-step
rational Prym of odd order. Any quotient detecting $\kappa_R$
must have $a\ge2$, because its image in every quotient of order
$13$ is zero. This proves strictness of the canonical detection
criteria on an actual curve. It does not assert that all actual
prime-degree covers fail: the independently constructed geometric
order-$13$ witness of A09 remains available. No explicit degree
for $k'$ is needed or claimed.

Recovered successor A.44 concerns a different pointed-involution
constraint at a Weierstrass base point and its primitive Frobenius
multipliers \cite[A.44]{ASrun}. It is compatible with the present
residual case. In general composite degree an even multiplicative
order alone does not force $-1\in\langle p\rangle$; the use of a
single odd prime power in (PC.5) is essential to the stated
hypothesis and to the simple order-$r$ deck endgame.


\subsection{The unramified double-cover repair and a genus-three obstruction}
\label{sec:genus-three-pointing-counterexample}

\paragraph{Finding A30 (the genus-two character truncation is essential).}
Over $k=\mathbf F_4$, let $D$ be the smooth projective model of
\[
                   z^2+z=\frac{1}{(t^2+t)^3}.
\tag{DC.1}
\]
It has genus three and $2$-rank one.  Let $P$ be the point
above $t=\infty$ with $z=0$, and let $W$ be the unique point
above $t=0$.  Both are rational; $W$ is Weierstrass and $P$
is not.  There is a surjection
\[
            b:J_D(k)\longrightarrow\mathbf Z/3,
                 \qquad b([W-P])\ne0,
\tag{DC.2}
\]
whose associated constant, geometrically connected, $P$-split
cyclic \'etale cover $D'\to D$ has \emph{even} relative class
number.  Since $\operatorname{ord}_3(2)=2$, the genus-two
conclusion of A28 that every such detecting quotient gives
an odd relative class number does not extend to genus three.
This does not rule out witnesses from other quotients or
higher powers for the same pointed genus-three curve.
The canonical class is also
detected: $K_D\sim4W$, so $b([K_D-4P])\ne0$.

More precisely, if $H_Y(T)$ denotes the numerator of the zeta
function of a curve $Y/k$, and
$B=(\ker(J_{D'}\to J_D))^0$, then
\[
 \overline{H_D}(T)=1-T,\qquad
 \overline{H_{D'}}(T)=(1-T)^2(1-T^3),\qquad
 \overline{H_B}(T)=(1-T)(1-T^3)
                         \quad\text{in }\mathbf F_2[T].
\tag{DC.3}
\]
The cover has degree three and $g(D')=7$.

\begin{proof}
\emph{1. The curve and its elliptic quotient.}
The rational function in (DC.1) has poles of order three
at $t=0,1$, and no others.  The Artin--Schreier genus and
$2$-rank formulas give $g(D)=3$ and $f_D=1$
\cite[Lemma~2.6, p.~712]{PriesZhu2012}.
Its hyperelliptic involution is $z\mapsto z+1$.
The two points above infinity have $z=0,1$ and are
unramified; the points above $0,1$ are its two ramification
points.  This proves the assertions about $P$ and $W$.
Since the hyperelliptic degree-two divisor class is $2W$,
the canonical class in genus three is $4W$.

Put
\[
 E/k:\ y^2+y=x^3,\qquad O=\infty,\qquad Q=(0,0).
\]
There is a separable degree-two map
\[
 f:D\longrightarrow E,\qquad
                 x=(t^2+t)^{-1},\quad y=z.
\tag{DC.4}
\]
Indeed its function-field extension is
$t^2+t=1/x$.  The two simple poles of $1/x$ on $E$ are
$Q=(0,0)$ and $-Q=(0,1)$, so these and only these are its
branch points.  Furthermore $f(P)=Q$ and $f(W)=O$.

There are nine points on $E(k)$: each of the four values
of $x$ has two values of $y$, and there is $O$.
Every affine point $(a,c)$ has order dividing three.
Its tangent line is $y=c+a^2(x-a)$; substituting
$x=a+v$ and using $a^4=a$ leaves $v^3=0$.
Thus the tangent meets the cubic three times at that point.
Consequently
\[
              E(k)\simeq(\mathbf Z/3)^2,
              \qquad \operatorname{ord}(Q)=3.
\tag{DC.5}
\]
The point count also gives
$H_E(T)=1+4T+4T^2$, in particular $\overline{H_E}=1$.

\emph{2. A translated Lang isogeny with the required fibres.}
Choose $a:E(k)\twoheadrightarrow\mathbf Z/3$ with $a(Q)=1$,
and write $H=\ker a$.
Let $F$ be $4$-power Frobenius on $E$.  Factoring the Lang
isogeny through $H$ gives
\[
 E\longrightarrow Y:=E/H\xrightarrow{\psi}E,
             \qquad \psi\pi_H=F-1.
\tag{DC.6}
\]
The map $\psi$ is an \'etale isogeny of degree three with
constant kernel $E(k)/H$.  The usual Lang-fibre computation
of A12 shows that a fibre $\psi^{-1}(R)$ is rational exactly
when $R\in H$; if rational, all three points are rational.
Use the translated cover
\[
                  \eta:Y\longrightarrow E,
                           V\longmapsto\psi(V)+Q.
\tag{DC.7}
\]
Its fibre above $Q$ has three rational points.  Its fibre
above $-Q$ has Frobenius translation
$a(-Q-Q)=a(Q)=1$, because $3Q=0$.
It is therefore one closed point of degree three.
The fibre above $O$ likewise has nontrivial translation
$a(-Q)=-1$.

Define $D'=D\times_{E,\eta}Y$.  This is a degree-three
constant \'etale cover of $D$, and its fibre above $P$ is
split.  It is geometrically connected: over an algebraic
closure the function-field extensions of $E$ defining $D$
and $Y$ have coprime degrees two and three, and the latter
is Galois.  They are linearly disjoint.
Unramified Riemann--Hurwitz gives $g(D')=3(3-1)+1=7$.

The rational quotient defining this pointed cover is
$b=a\circ f_*:J_D(k)\to\mathbf Z/3$.
Indeed $f_*[R-P]=[f(R)-Q]$, which is exactly the Lang
translation used in (DC.7).  In particular
\[
                 b([W-P])=a(O-Q)=-1\ne0.
\tag{DC.8}
\]
This also proves that $b$ is surjective.

\emph{3. The exact zeta calculation modulo two.}
For a separable degree-two cover $X\to Z$ of curves in
characteristic two, let $\mathcal R$ be its set of closed
branch points on $Z$.  Comparing Euler factors gives
\[
 Z_X(T)\equiv Z_Z(T^2)
                   \prod_{R\in\mathcal R}(1-T^{\deg R})
                                \pmod2.
\tag{DC.9}
\]
At an unramified closed point of degree $d$, either there
are two degree-$d$ points above it or one degree-$2d$ point.
Both local factors reduce to $(1-T^{2d})^{-1}$.
At a ramified point the local factor is $(1-T^d)^{-1}$;
multiplication by $1-T^d$ changes the corresponding
$Z_Z(T^2)$ factor to this one.  This proves (DC.9) as an
identity of formal power series.

Since the constant field has even cardinality, passing to
zeta numerators in (DC.9) gives the polynomial identity
\[
 \overline{H_X}(T)=\overline{H_Z}(T^2)
             \frac{\prod_{R\in\mathcal R}(1-T^{\deg R})}{1-T}.
\tag{DC.10}
\]
For $D\to E$, the two branch points are rational and
$\overline{H_E}=1$, giving $\overline{H_D}=1-T$.
The elliptic curve $Y$ is isogenous to $E$, so
$H_Y=H_E$ and $\overline{H_Y}=1$.
For $D'\to Y$, the branch divisor is the \'etale pullback
of $Q+(-Q)$.  Its closed-point degrees are $1,1,1,3$ by
Step~2.  Formula (DC.10) therefore gives
\[
       \overline{H_{D'}}(T)
              =\frac{(1-T)^3(1-T^3)}{1-T}
              =(1-T)^2(1-T^3).
\]
Pullback and norm for the degree-three cover show that
$J_D\times B$ is isogenous to $J_{D'}$.  Dividing their
Frobenius polynomials proves the last identity in (DC.3).
Finally
$\#B(k)=H_B(1)\equiv0\pmod2$.
The point-count and isogeny identities are the usual
Frobenius identities, as in
\cite[Theorem~19.1(b) and its Frobenius proof, pp.~40--41]{MilneAV}.

\emph{4. The individual reduced character polynomials.}
The failure is visible before taking their product.
Let $\chi:\mathbf Z/3\to\overline{\mathbf F}_2^\times$
send $1$ to $\zeta$, a primitive cube root.
The nontrivial unramified character sheaf on $E$ defined
by (DC.7) has $L$-polynomial $1$.
For example, Riemann--Roch on the elliptic curve makes
the number of effective divisors in a rational divisor
class of each positive degree independent of the class,
so every positive-degree nontrivial character coefficient
vanishes.

The character Euler-factor argument of Step~3 applies
with a character value $\xi$ at a closed point of $E$:
both unramified splitting types give
$(1-\xi^2T^{2d})^{-1}$ modulo two, while a branch point
contributes the correction $1-\xi T^d$.
The character values at $Q,-Q$ are $1,\zeta$ respectively.
Thus the degree-four unramified character polynomial on
the genus-three curve $D$ has reduction
\[
       \overline{L_{\chi,D,P}}(T)=(1-T)(1-\zeta T),
       \qquad
       \overline{L_{\chi^{-1},D,P}}(T)
                                      =(1-T)(1-\zeta^2T).
\tag{DC.11}
\]
Both values at $T=1$ vanish.  The genus-two reduction is
linear, whereas (DC.11) has a genuine quadratic term.
This is precisely why its inversion symmetry does not
give the genus-two contradiction.
\end{proof}

\paragraph{What happens to the proposed unramified double-cover repair.}
The underlying geometric observation is correct.
For a rank-one genus-two Artin--Schreier curve written,
after choosing the two rational branch points as $0,\infty$,
as
\[
 C:\ y^2+y=f_1(x)+f_3(x),
\]
where $f_1$ has just the simple pole at $0$ and $f_3$
just the triple pole at $\infty$, form the connected
biquadratic fibre product
\[
 z_1^2+z_1=f_1(x),\qquad z_2^2+z_2=f_3(x),
                         \qquad y=z_1+z_2.
\tag{DC.12}
\]
Over an algebraic closure the two Artin--Schreier classes
are independent, since their ramification supports differ.
The inertia groups over $0,\infty$ are respectively the
$z_1$- and $z_2$-translation groups.
The diagonal translation is in neither inertia group,
so its quotient map $\widetilde C\to C$ is \'etale of
degree two.  It is the unique nontrivial geometric etale
double cover: $f_C=1$ means that
$H^1_{\mathrm{et}}(C_{\bar k},\mathbf F_2)$ has dimension one.
Its source has genus three by Riemann--Hurwitz.
The quotient by $z_2$-translation is
the genus-zero $z_1$-curve.  Hence $z_2$-translation is
the hyperelliptic involution of $\widetilde C$.

At the simple pole its two lifts are interchanged by
this involution, so neither is Weierstrass.  At the
triple pole both lifts are fixed and are Weierstrass.
Equivalently, the hyperelliptic presentation on the
genus-zero $z_1$-curve has two poles of order three,
so the same formula gives $2$-rank one.
Finite constant extensions, or the appropriate constant
twist, can make a prescribed fibre rational, but do not
change these geometric conclusions.

This construction nevertheless does not repair the A1
argument.  First, its nontrivial geometric degree-two
quotient does not factor through the geometric
prime-to-two fundamental group, whose finite quotients
have odd order.  A collision in the quotient used by A1
therefore does not supply lifts to this cover.
Second, if $P_s,P_t$ are the simple- and triple-pole
points on $C$, their difference has exact order two:
$\operatorname{div}(x)=2P_s-2P_t$, and a nonzero
degree-one principal divisor would make $C$ rational.
For lifts $\widetilde P_s,\widetilde P_t$ one has
\[
 \pi_*[\widetilde P_t-\widetilde P_s]=[P_t-P_s].
\tag{DC.13}
\]
Every odd cyclic character pulled back from $C$ thus
kills this difference.  A new odd character detecting
it upstairs cannot supply the A28 detector by descent
from such a character on $C$.

The explicit curve (DC.1) is exactly this construction
for $C:y^2+y=x^3+x^{-1}$, with $z_1=t$ and $z_2=z$.
Its example (DC.11) gives an additional obstruction:
even as an argument about the upstairs curve alone,
the proposed unrestricted higher-genus symmetry rule
is false.


\subsection{An exact theta witness on the numerical-obstruction curve}
\label{sec:a27-theta-seven}

\paragraph{Finding A31 (a degree-seven geometric witness on A27).}
Let $k=\mathbf F_4=\mathbf F_2(u)$, with $u^2+u+1=0$, and let
\[
 C/k:\qquad y^2+y=u t^3+t+u/t
 \tag{TS.1}
\]
be the smooth proper rank-one genus-two curve of A27.  Let $W$ be
its point over $t=0$ and let $R=(u,u^2)$.  Then
$M=\mathcal O_C(R-W)$ has exact order seven.  On the Frobenius
twist $C^{(1)}$, its twist $M^{(1)}$ lies on the Raynaud theta
divisor.  Consequently A11 gives a geometrically connected cyclic
degree-seven cover with geometric primitive $2$-rank zero.  After
a finite constant extension it gives a constant split witness at
any prescribed point of $C$.

For the prescribed point $R$, the canonical class
$\kappa_R=[K_C-2R]$ has exact order seven.  Thus A29's detector
does not apply to this class, even after extending constants,
because $\operatorname{ord}_7(2)=3$ is odd.  This is a concrete
instance of the remaining detector case that the geometric method
does handle.  It is not a universal incidence theorem.  It is also
compatible with every numerical failure proved in A27.

\begin{proof}
\emph{1. An explicit isomorphic model.}
The coordinate change
\[
 \begin{aligned}
 X&=u/t, &
 Y&=u^2t^{-2}(y+u)+u^2,\\
 t&=u/X, &
 y&=u+(Y+u^2)/X^2
 \end{aligned}
 \tag{TS.2}
\]
and its displayed inverse identify the function field of (TS.1)
with that of
\[
 \begin{aligned}
 C':\quad Y^2+X^2Y&=f(X),\\
 f(X)&=X^5+X^4+uX^3+u^2X^2+uX+u.
 \end{aligned}
 \tag{TS.3}
\]
Here is a coefficient check, retaining all constant terms.
Put $\alpha=u^2t^{-2}$ and $\beta=t^{-2}+u^2$, so
$Y=\alpha y+\beta$.  Since $X^2\alpha=\alpha^2$,
the equation on the left becomes
\[
 \begin{aligned}
 Y^2+X^2Y
 &=\alpha^2(ut^3+t+u/t)+\beta^2+X^2\beta\\
 &=u^2t^{-5}+ut^{-4}+ut^{-3}
       +ut^{-2}+u^2t^{-1}+u\\
 &=f(u/t).
 \end{aligned}
 \tag{TS.4}
\]
Both are models of their smooth proper function-field curves, so
this birational identification extends to an isomorphism.
The two odd Artin--Schreier pole orders in (TS.1) are three and
one; thus the genus and $2$-rank are those already checked in A27.
The point $W$ becomes the unique point at infinity in (TS.3),
with $\operatorname{ord}_W(X)=-2$ and
$\operatorname{ord}_W(Y)=-5$.  The point $R$ becomes $(1,u)$.
The affine model is smooth: the only possible singular
$X$-coordinate is zero, and $f'(0)=u\ne0$ there.

\emph{2. A principal divisor of degree seven.}
Set
\[
                F=(X+u)Y+X.
 \tag{TS.5}
\]
The hyperelliptic involution sends $Y$ to $Y+X^2$.
Taking its function-field norm gives the exact polynomial identity
\[
 \begin{aligned}
 F\,\iota(F)
 &=X^2+X^3(X+u)+(X+u)^2f(X)\\
 &=(X+1)^7.
 \end{aligned}
 \tag{TS.6}
\]
Indeed expansion yields coefficient one in every degree from
zero through seven.  At the two points above $X=1$,
\[
       F(1,u)=0,\qquad F(1,u^2)=u^2\ne0.
 \tag{TS.7}
\]
They are distinct unramified points since $X^2=1$ there.
Thus (TS.6) assigns all seven zeros of $F$ to $R$, with no
zero at its conjugate.  The function is regular on the affine
curve and has a pole of exact order seven at $W$:
the term $XY$ has pole order seven, whereas $uY$ and $X$
have pole orders five and two.  Consequently
\[
                       \operatorname{div}(F)=7(R-W).
 \tag{TS.8}
\]
The class $[R-W]$ is nonzero.  Otherwise a function with divisor
$R-W$ would give a degree-one map from a genus-two curve to
$\mathbf P^1$, which is impossible.  Since seven is prime,
(TS.8) proves that its exact order is seven.

\emph{3. The theta bundle and the Frobenius twist.}
The differential
\[
                    \omega=dX/X^2=d(1/X)
 \tag{TS.9}
\]
is exact and has divisor $2W$.  At a finite point with $X\ne0$,
$X$ is an etale local coordinate, so $\omega$ is nonvanishing.
At $X=0$, the differentiated equation gives
$\omega=dY/f'(X)$; since $f'(0)=u\ne0$, this is again a
nonvanishing regular differential.  Its divisor is therefore
supported at $W$ and has degree $2g-2=2$.

Let $B_C$ denote the line of locally exact differentials on
$C^{(1)}$.  The section supplied by (TS.9) has divisor
$W^{(1)}$ on that twist: a local zero of order two in an exact
differential corresponds to a zero of order one in the
$\mathcal O_{C^{(1)}}$-line.  Hence
\[
 B_C\simeq\mathcal O_{C^{(1)}}(W^{(1)}),\qquad
 B_C\otimes M^{(1)}
       \simeq\mathcal O_{C^{(1)}}(R^{(1)}).
 \tag{TS.10}
\]
The second line bundle has a nonzero section, proving the asserted
theta incidence.  This keeps the twist explicit over $\mathbf F_4$:
the coefficients of (TS.3) are squared, and
$R^{(1)}=(1,u^2)$ in the resulting coordinates.
Twisting (TS.8) preserves the exact order seven.

\emph{4. The cover and the exact limitation of the detectors.}
The powers of two modulo seven are $1,2,4$.  Adjoining their
negatives gives every nonzero residue, so
\[
                  \langle2,-1\rangle=(\mathbf Z/7)^\times.
 \tag{TS.11}
\]
Finding A11 applies to (TS.10) and gives a cyclic degree-seven
cover with primitive geometric $2$-rank zero.  Since the degree
is prime, this is the whole new factor.  Descending the cover and
its deck transformations and rationalizing a point above any
prescribed base point require only a finite constant extension.
The geometric rank-zero property survives these operations; the
relative rational class number is then odd, and A11's pointed
deck argument separates any putative collision.

Finally (TS.9) also gives $K_C\sim2W$, so
$\kappa_R=-2[R-W]$ has exact order seven.  Its only nonzero odd
primary part is the one for seven, whose order of two is odd.
Injectivity of rational Jacobian points under constant extension
preserves that order and prime support.  Thus this prescribed
point remains outside A29's sufficient hypothesis.  In fact A27
gives $\#J_C(\mathbf F_4)=28$ and $\#C(\mathbf F_4)=6$, so for
every one of those six prescribed points the odd primary support
of $\kappa_P$ is contained in $\{7\}$.  A29's detector therefore
fails for all six points, and continues to fail for these same
points after every constant extension, whereas the geometric
cover just constructed becomes split at each prescribed point
after a finite extension.  A27 refuted
universal success of its stated size inequalities at primes with
two primitive; seven does not satisfy that primitive-root
hypothesis.  Neither A27 nor the failure of A29's detector
asserted absence of the present geometric witness.
\end{proof}

The model (TS.3), norm (TS.6) and isomorphism (TS.2) were found
in this run by the root investigator.  The compact certificate
\texttt{verify\_A31\_norm7.py} checks their coefficients directly
over $\mathbf F_4$, including the inverse substitution, both point
values, smoothness at $X=0$ and the signed unit orbit.  It uses no
Jacobian arithmetic package.  The divisor, differential, exact-order
and covering arguments are the written proof above, not an
inference from a numerical search alone.


\subsection{An ordinary new factor with an odd relative class number}
\label{sec:ordinary-rational-witness}

\paragraph{Finding A32 (the rational witness can be geometrically ordinary).}
Let $C/\mathbf F_4$ be the curve of A27 and A31, let $W$ be its
point over $t=0$, and put $K=\mathbf F_{256}$.  There is a constant,
geometrically connected, $W$-split cyclic etale cover $D\to C_K$
of degree thirteen whose Prym
\[
             \mathcal P=\ker(J_D\longrightarrow J_C)^0
 \tag{OR.1}
\]
is ordinary of dimension twelve, while $\#\mathcal P(K)$ is odd.
Writing $K_m=\mathbf F_{256^m}$, the exact extension rule is
\[
 \frac{\#J_D(K_m)}{\#J_C(K_m)}\ \text{is odd}
             \quad\Longleftrightarrow\quad 7\nmid m.
 \tag{OR.2}
\]
All $2^{12}$ geometric etale $2$-torsion points of $\mathcal P$
are rational over $K_7$; seven is the least extension degree
with any nonzero rational $2$-torsion in that factor.
The reduced Frobenius numerator is
\[
                H_{\mathcal P}(T)\bmod2=1+T^4+T^{12}.
 \tag{OR.3}
\]
Thus this rational class-number witness uses a geometrically
ordinary new factor.  It is a different mechanism from A31's
geometric rank-zero witness, and it applies at a Weierstrass
point where the canonical class detector of A29 is zero.

\begin{proof}
\emph{1. The exact cyclic quotient and its point-image table.}
For the calculation use the isomorphic model
\[
 C_0:\quad Y^2+X^2Y=uX^5+X^3+uX,
 \qquad X=1/t,\quad Y=y/t^2.
 \tag{OR.4}
\]
Its point at infinity is $W$.  Realize
\[
 K=\mathbf F_2[z]/(z^8+z^4+z^3+z+1),\qquad
 u=z^7+z^5+z^4+z^3+z^2.
 \tag{OR.5}
\]
Then $u^2+u+1=0$.  For the finite table only, an integer
$\sum_i a_i2^i$ encodes the field element $\sum_i a_iz^i$.
Thus $u$ is encoded by $188$.

Finding A27 gives
\[
        G=J_C(K),\qquad |G|=65520=13\cdot5040,
        \qquad |C(K)|=256.
 \tag{OR.6}
\]
The thirteen-primary subgroup is therefore cyclic of order
thirteen.  In Mumford coordinates for (OR.4), the element
\[
        \Delta=(X^2+80X+237,\ 106X+207)
 \tag{OR.7}
\]
has exact order thirteen; the coefficients in this display use
the field encoding just specified.  Define the surjection
$a:G\to\mathbf Z/13$ by
\[
                      [5040]Q=a(Q)\Delta.
 \tag{OR.8}
\]
Multiplication by $5040$ kills every other primary factor and
is an automorphism on the thirteen-primary subgroup.  Exact
enumeration gives the following complete distribution:
\[
\begin{array}{c|rrrrrrrrrrrrr}
 j&0&1&2&3&4&5&6&7&8&9&10&11&12\\ \hline
 \#\{R:a([R-W])=j\}
 &44&19&18&18&16&19&16&16&19&16&18&18&19.
\end{array}
 \tag{OR.9}
\]
The entries add to $256$, including the point at infinity.

Here the finite computation has two separate implementations.
The first uses the extended-GCD composition formula for Mumford
divisors.  The second imports none of its arithmetic and obtains
composition by binary linear algebra as follows.  Write a reduced
class as $(U,V)$, where $U$ is monic, $\deg U\le2$,
$\deg V<\deg U$, and
\[
                         U\mid V^2+X^2V+f,
                \qquad f=uX^5+X^3+uX.
 \tag{OR.10}
\]
To add two such classes, cancel their common opposite support
using $\gcd(U_1,U_2,V_1+V_2+X^2)$.  This removes pairs whose
sum is the divisor of their $X$-polynomial, hence does not change
the resulting Jacobian class.  More explicitly, a nonramified
pair $Q+\iota Q-2W$ is the divisor of $X-X(Q)$, with the
same cancellation repeated for each common multiplicity.  The
only affine branch point is $(0,0)$; it occurs at most once in
a reduced pair, since $f'(0)=u\ne0$ prevents $X^2$ from
dividing $V^2+X^2V+f$.  A double occurrence therefore cancels
by $\operatorname{div}(X)=2(0,0)-2W$.
For the remaining polynomials,
solve for $V$ of degree below $\deg(U_1U_2)$ subject to
\[
 \begin{aligned}
 V&\equiv V_i\pmod{U_i}\quad(i=1,2),\\
 V^2+X^2V+f&\equiv0\pmod{U_1U_2}.
 \end{aligned}
 \tag{OR.11}
\]
This is a linear system over $\mathbf F_2$ in the binary
coefficients of $V$.  Common remaining support is on the same
unramified branch, where the derivative in $Y$ is nonzero;
the required higher-multiplicity root interpolation is therefore
unique, by recursively solving the local root equation using
that nonzero derivative.  This applies also to a doubled
degree-two divisor; it is not restricted to distinct support.
At the ramified point any opposite pair was already cancelled,
and a remaining simple occurrence fixes only its value.
Consequently the resulting pair represents the sum divisor
with its correct multiplicities.

If its degree is above two, replace it by
\[
     U'=(V^2+X^2V+f)/U\quad\text{made monic},\qquad
     V'=V+X^2\pmod{U'}.
 \tag{OR.12}
\]
Here the divisor of $Y+V$ has finite part $D_U+E$ and pole
order $\max\{5,2\deg V\}$ at $W$.  The two possible pole
orders are odd and even, so there is no leading cancellation.
Its norm is $U$ times the polynomial for $E$.  Thus
$[D_U-(\deg U)W]=-[E-(\deg E)W]$, and replacing $E$ by its
hyperelliptic conjugate gives precisely the second formula in
(OR.12).  The leading coefficient $u$ of $f$ is retained in
the division; making $U'$ monic multiplies only its defining
polynomial by a nonzero constant and changes no divisor.
The new degree is at most two.  A reduced class is unique:
pad its effective representative to degree two with $W$ and
apply genus-two Riemann--Roch.  Two distinct representatives
could occur only in the canonical degree-two pencil; its
conjugate pairs have already been cancelled to the identity.
The second implementation verifies (OR.10) after every step,
uses a different scalar-multiplication chain, enumerates all
affine solutions of (OR.4) plus infinity, and independently
reproduces (OR.7)--(OR.9).  Its companion point file retains all
$256$ coordinates and quotient images.  Thus the finite input
to the character calculation is fully specified and reproducible.

\emph{2. The pointed cover and the reduced character values.}
Take the connected constant Lang quotient cover associated to
(OR.8), based at $W$, as in A12.  Its fibre over $W$ is split.
The genus is fourteen, so the dimension of $\mathcal P$ is twelve.
Since $K_C\sim2W$, the exact genus-two polynomial of recovered
A1 B.2 is
\[
 L_W(\chi,T)=1+c_1(\chi)T+256T^2,
 \qquad
 c_1(\chi)=\sum_{R\in C(K)}\chi(a([R-W]))
 \tag{OR.13}
\]
for every nontrivial character of $\mathbf Z/13$
\cite[B.2]{A1run}.  Put $\zeta$ for a primitive thirteenth root
in characteristic two.  Formula (OR.9) gives
\[
 \begin{aligned}
 S(Z)&=Z+Z^5+Z^8+Z^{12},\\
 \overline{L_W(\chi_i,T)}&=1+S(\zeta^i)T
                         \qquad(1\le i\le12).
 \end{aligned}
 \tag{OR.14}
\]
In particular the B.2 group-ring test is $1+S(Z)$.
Two has order twelve modulo thirteen, so
$\Phi_{13}(Z)=1+Z+\cdots+Z^{12}$ is irreducible over
$\mathbf F_2$.  Neither $S$ nor $1+S$ is that polynomial or
zero.  Therefore
\[
                        S(\zeta^i)\notin\{0,1\}
                           \quad(1\le i\le12).
 \tag{OR.15}
\]
The simultaneous test of B.2 now gives an odd relative class
number.  The norm-pullback isogeny identifies that integer with
$\#\mathcal P(K)$, as in B.1.

\emph{3. Exact character scalars and ordinary geometric rank.}
Set $\beta=S(\zeta)$.  Multiplication of the four exponents
$\{1,5,8,12\}$ by eight permutes them modulo thirteen, so
$\beta^8=\beta$.  Thus $\beta\in\mathbf F_8$, and (OR.15)
shows that it has exact multiplicative order seven.
The three cosets obtained by multiplying those exponents by
$1,2,4$ partition $(\mathbf Z/13)^\times$.  It follows that
\[
              \beta+\beta^2+\beta^4
                        =\sum_{i=1}^{12}\zeta^i=1.
 \tag{OR.16}
\]
Its minimal polynomial is therefore the trace-one irreducible
cubic $Z^3+Z^2+1$.  Consequently
\[
 \prod_{j=0}^{2}(1+\beta^{2^j}T)=1+T+T^3,
 \qquad
 \prod_{i=1}^{12}(1+S(\zeta^i)T)=(1+T+T^3)^4.
 \tag{OR.17}
\]
The twelve character scalars comprise the three conjugates,
each repeated four times.  Character factorization in B.2
now proves (OR.3) directly.

For clarity, the cohomological interpretation is a cited result,
not a new comparison theorem.  Recovered successor A.20 identifies
the reduced character polynomial with the determinant of the
linear extension of $256$-power coherent Frobenius on the
character summand of $H^1(D,\mathcal O_D)$
\cite[A.20]{ASrun}.  Its etale determinant step is Deligne,
SGA~$4\frac12$, \emph{Fonctions $L$ modulo $\ell^n$ et modulo
$p$}, Theorem~2.2(b) and formula~(2.4.1), with the hypotheses
of \S2.1 \cite{DeligneModP}.  Here the coefficient ring is a
finite field of characteristic two containing the character
values, hence is noetherian and reduced; the rank-one lisse
character sheaf is constructible and has finite Tor dimension
on the smooth proper curve $C_K$.  Each nontrivial summand has dimension one
because the base genus is two.  All its displayed scalars are
nonzero, so the full twelve-dimensional new part is stable under
Frobenius.  The Artin--Schreier stable-part comparison therefore
gives geometric $2$-rank twelve; equivalently, the new factor is
ordinary.  The exact primary comparison locator is Katz,
SGA~7~II, Expos\'e~XXII, Proposition~2.2.5, with the
characteristic-polynomial convention in Proposition~2.2.4
\cite{KatzCongruence}.

\emph{4. Constant extensions and rational torsion.}
After extension to $K_m$, the linear Frobenius on those
one-dimensional summands has scalars $S(\zeta^i)^m$.
Every one of the original scalars has exact order seven.
Thus the relative class-number reduction is
\[
        \prod_{i=1}^{12}\bigl(1+S(\zeta^i)^m\bigr),
 \tag{OR.18}
\]
which is nonzero exactly when $7\nmid m$.  This proves (OR.2).
The deck group is constant, so Frobenius preserves the twelve
character lines; there is no extra Jordan block hidden by their
repeated characteristic-polynomial roots.  At $m=7$ its action
on the entire new stable part is the identity.  The functorial
stable-part comparison, or its dual on etale Jacobian torsion,
then makes all $2^{12}$ geometric etale $2$-torsion points of
$\mathcal P$ rational.  At smaller positive degrees none is
rational, by (OR.18).  This proves the claimed exact first
extension degree.
\end{proof}

At this same level the numerical premises are very far from
necessary: $|G/13G|=13<256=|C(K)|$, and B.3's pair bound would
require $256\ge2\cdot5040$.  Finding A32 supplies an actual
primitive-two prime witness despite both failures.  Its ordinary
new factor also shows why rational class-number units must not
be equated with geometric rank zero.  No absence of other
geometric theta points of order thirteen is asserted.


\subsection{Two reflection maps reduce genus two to the ordinary case}
\label{sec:quintic-ordinary-reduction}

\paragraph{Finding A33 (a quintic transfer to an ordinary genus-two curve).}
For a curve over a finite field of characteristic two, let $\Pi_C$ have
the meaning of A1: the geometric kernel is replaced by its maximal
prime-to-two quotient, and the arithmetic quotient remains the whole
absolute Galois group.  Consider the following restricted assertion:
\[
 \begin{gathered}
 \text{For every finite field $k$ of characteristic two and every ordinary}\\
 \text{smooth proper geometrically connected genus-two curve $E/k$,}\\
 E(k)\longrightarrow\operatorname{Sec}(\Pi_E\to G_k)/\Delta_E
 \text{ is injective.}
 \end{gathered}
 \tag{DR.1}
\]
If (DR.1) holds, the same point-section assertion holds for every
smooth proper geometrically connected genus-two curve over every finite
field of characteristic two, without a restriction on its $2$-rank.
In fact the hypothesis can be restricted further to ordinary $E/k$
equipped with an embedding
$\mathcal O_5=\mathbf Z[(1+\sqrt5)/2]\hookrightarrow
\operatorname{End}_k(J_E)$ whose image is fixed by the Rosati involution
of the canonical principal polarization.  Thus only this real
multiplication case of the ordinary assertion is needed.

More precisely, let $C/\overline k$ have genus two and let
$\pi:D\to C$ be any connected cyclic etale cover of degree five, with
deck generator $\sigma$.  A lift $\tau$ of the hyperelliptic involution
can be chosen, and it satisfies
\[
 \tau^2=1,\qquad \tau\sigma\tau^{-1}=\sigma^{-1}.
 \tag{DR.2}
\]
The quotient $v:D\to E=D/\langle\tau\rangle$ has genus two.  If
$\mathcal B=\ker(\operatorname{Nm}_\pi:J_D\to J_C)^0$, then
\[
 \dim\mathcal B=4,\qquad
 \mathcal B\sim J_E\times J_E,\qquad
 f_{\mathcal B}\in\{0,4\},\qquad f_E\in\{0,2\}.
 \tag{DR.3}
\]
Here $\sim$ is isogeny over the algebraic closure and $f$ denotes
geometric $2$-rank.  After a finite constant extension all the maps and
deck transformations can be defined and a chosen base-point fibre can
be split.  If $f_{\mathcal B}=0$, this cover itself supplies the
geometric primitive-zero witness.  If $f_{\mathcal B}=4$, the two maps
\[
                         v,\quad v\circ\sigma:D\longrightarrow E
 \tag{DR.4}
\]
transfer any one actual point-section collision on $C$ to two collisions
on the same ordinary genus-two curve $E$.  Their fibres together
separate geometric points of $D$.

\paragraph{Source attribution and scope.}
The reflection-lift, completed-local-ring, and Prym-image arguments
already occur in recovered A1 B.40 for cubic covers in genera two and
three \cite[B.40]{A1run}.  Recovered successor A.37 studies a pointed
cubic reflection quotient and explains why a reflection twist does not
remove a two-primary obstruction on its elliptic quotient
\cite[A.37]{ASrun}.  The complete proof of successor A.4 explicitly
lifts equal point sections to equal point sections on the same chosen
odd cover \cite[A.4]{ASrun}.  We use that machinery explicitly here.  The
degree-five choice retains genus two, and the two quotient maps below
transport one common lifted section collision.  This is a reduction of
the older route to its ordinary genus-two case; it supplies neither a
universal theta-incidence assertion nor a proof of (DR.1) by Hooley.
No analytic hypothesis is used in the reduction.

\begin{proof}
\emph{1. The reflection quotient and its genus.}
Work first over an algebraic closure.  Connected cyclic etale covers of
degree five exist because $J_C[5]\simeq(\mathbf Z/5)^4$ geometrically.
One may use an exact-order-five line bundle and a trivialization of its
fifth power.  For any such cover, the hyperelliptic involution $\iota$
acts on $J_C$ by $[-1]$, and hence takes its defining torsion line
bundle to the inverse.  It therefore lifts to an automorphism $\tau$
of $D$ that inverts the deck group.  Its square is some $\sigma^a$.
Conjugating that square by $\tau$ both leaves it unchanged and replaces
it by $\sigma^{-a}$.  Thus $2a=0$ in $\mathbf Z/5$, so $a=0$ and
(DR.2) follows.

The hyperelliptic map $C\to\mathbf P^1$ is a separable double cover in
characteristic two.  Above each of its geometric branch points,
$\iota$ has one fixed point $W$ on $C$.  The five points of
$\pi^{-1}(W)$ form a deck torsor, on which $\tau$ acts as
$j\mapsto b-j$.  This has exactly one fixed point, because two is
invertible modulo five.  Every fixed point of $\tau$ occurs this way,
since its image is fixed by $\iota$.

At each such fixed point $T$, the etale map $\pi$ identifies the
completed local rings at $T$ and $W$.  The actions of $\tau$ and
$\iota$ agree under that identification.  Consequently the local
different exponent of $D\to E$ at $T$ is exactly the different
exponent of $C\to\mathbf P^1$ at $W$.  There is no ramification away
from these fixed points.  The total different of the latter
hyperelliptic map is six by Riemann--Hurwitz.  Since $g(D)=6$ by the
unramified degree-five formula, Riemann--Hurwitz for $D\to E$ gives
\[
                 10=2g(D)-2=2(2g(E)-2)+6,
                 \qquad g(E)=2.
 \tag{DR.5}
\]
This argument works for all three possible $2$-ranks of $C$; it does
not assume the rank-one branch pattern.

\emph{2. The Prym isogeny and the rank dichotomy.}
Let $A=v^*J_E\subset J_D$ and let $A'=\sigma A$.  Pullback has
finite kernel because its composition with the corresponding norm is
$[2]$, so both images have dimension two.  The norm of $A$ to $J_C$
is fixed by the induced hyperelliptic involution, which acts by $[-1]$.
The homomorphism $\operatorname{Nm}_\pi\circ v^*$ is therefore killed
by two.  A homomorphism of abelian varieties killed by a nonzero integer
is zero, so $A\subset\mathcal B$; the same is true of $A'$.

Every point of $A$ is fixed by $\tau$, and every point of $A'$ is fixed
by $\sigma\tau\sigma^{-1}$.  These two reflections generate the
degree-five rotation group.  Thus a point in their intersection is
fixed by $\sigma$.  On the other hand,
\[
 \pi^*\operatorname{Nm}_\pi=1+\sigma+\sigma^2+\sigma^3+\sigma^4.
 \tag{DR.6}
\]
Its norm is zero, so (DR.6) shows that the intersection is killed by
five.  This also proves scheme-theoretic finiteness: the intersection
is contained in $J_D[5]$.  The sum of $A$ and $A'$ therefore has
dimension four and equals $\mathcal B$.  Thus the map
$(x,y)\mapsto v^*x+\sigma v^*y$ from $J_E\times J_E$ to
$\mathcal B$ is an isogeny, as asserted in (DR.3).

There is a short group-algebra proof of the rank restriction.  The
geometric group $\mathcal B[2](\overline k)$ is an
$\mathbf F_2[\mathbf Z/5]$-module of dimension $f_{\mathcal B}\le4$.
It has no nonzero invariant vector: such a vector $x$ would satisfy
$0=(1+\sigma+\cdots+\sigma^4)x=5x=x$ by (DR.6).
The group algebra is semisimple, and its only nontrivial simple module
has dimension four because
\[
             \operatorname{ord}_5(2)=4,
             \qquad \Phi_5(T)=T^4+T^3+T^2+T+1
             \text{ is irreducible over }\mathbf F_2.
 \tag{DR.7}
\]
It follows that $f_{\mathcal B}$ is zero or four.  Isogeny invariance
and additivity of $2$-rank give $f_{\mathcal B}=2f_E$, proving all of
(DR.3).  In particular a nonzero new rank forces $E$ to be ordinary.

The quotient Jacobian also has the asserted integral real
multiplication.  Define the integral correspondence endomorphism
\[
                   T=v_*\sigma v^*\in\operatorname{End}(J_E).
 \tag{DR.7a}
\]
The pullback--norm identity $v^*v_*=1+\tau$ and the relation
$\tau\sigma=\sigma^{-1}\tau$ give
\[
            v^*T=(\sigma+\sigma^{-1})v^*.
 \tag{DR.7b}
\]
On $\mathcal B$, the deck sum in (DR.6) vanishes.  Expanding
$(\sigma+\sigma^{-1})^2+(\sigma+\sigma^{-1})-1$ gives exactly
$1+\sigma+\sigma^2+\sigma^3+\sigma^4$.  Therefore
$v^*(T^2+T-1)=0$.  Since $v^*$ has finite kernel, the homomorphism
$T^2+T-1$ has finite image and hence is zero.  The polynomial
$X^2+X-1$ is irreducible over $\mathbf Q$.  The induced unital map
from its quadratic integer ring to $\operatorname{End}(J_E)$ is
injective: any nonzero ideal in that ring contains a nonzero integer,
whereas multiplication by a nonzero integer on $J_E$ is not zero.
This is the desired copy of $\mathcal O_5$.

Under the canonical principal polarizations, norm and pullback are
adjoint and the adjoint of a curve automorphism is its inverse.
Consequently
\[
 T^{\dagger}=v_*\sigma^{-1}v^*
            =v_*\tau\sigma\tau v^*=v_*\sigma v^*=T.
 \tag{DR.7c}
\]
Here $v_*\tau=v_*$ and $\tau v^*=v^*$ were used.  The entire
quadratic ring is therefore fixed by Rosati.  This correspondence is
defined over any field over which $v$ and $\sigma$ are defined.

\emph{3. One common lift of a section collision.}
Let $P,Q\in C(k)$ give conjugate sections in the precise quotient
$\Pi_C$ used by A1.  Choose one geometric degree-five cover and the
reflection data above.  Since all coefficients and geometric points
are algebraic over a finite field, a finite constant extension $k'/k$
makes the cover, its deck transformations, and $\tau$ defined and makes
some chosen $P'\in D(k')$ over $P$ rational.  The constant degree-five
deck group then makes the entire fibre above $P$ split.  The assumed
section conjugacy persists on restriction to $G_{k'}$.

The arithmetic group of this pointed cover is an open subgroup of
$\Pi_{C_{k'}}$, and its geometric kernel has the full prime-to-two
topology.  Here is the normal-core justification needed for that
last assertion.  For the normal index-five subgroup defining this
cover, take the kernel of any additional finite odd quotient and
intersect its five conjugates in the original geometric fundamental
group.  The resulting quotient of the index-five subgroup embeds in
a product of finite odd groups.  Its extension by the original cyclic
deck group of order five again has odd order.  Thus that quotient
already factors through the original maximal pro-odd geometric
quotient.  This proves the compatibility with $\Pi_D$ used here and
in A01; the arithmetic Galois quotient has not been replaced.

Equivalently, represent the assumed geometric section conjugacy by a
$G_{k'}$-fixed prime-to-two path from $P$ to $Q$.  Its unique lift with
initial point $P'$ has a rational endpoint $Q'\in D(k')$ over $Q$ and
is a $G_{k'}$-fixed path between $P'$ and $Q'$.  Hence
\[
             [s_{P'}]=[s_{Q'}]
             \quad\text{in }\operatorname{Sec}(\Pi_D\to G_{k'})/\Delta_D.
 \tag{DR.8}
\]
This is the full-section lift checked in recovered successor A.4
\cite[A.4]{ASrun}.  There is one fixed pair of lifts, chosen before
either quotient map is applied.  It is not an inference from two
separately selected abelian Kummer lifts.

\emph{4. The zero-rank alternative already separates the points.}
Suppose $P\ne Q$ and $f_{\mathcal B}=0$.  Abelian Kummer theory applied
to (DR.8), as checked in A01, gives
\[
                     \delta=[Q'-P']\in J_D(k')[2^\infty].
 \tag{DR.9}
\]
The connected image of $\sigma-1$ is $\mathcal B$.  It lies there by
the norm identity, and conversely $5-(1+T+\cdots+T^4)$ is divisible
by $T-1$; on $\mathcal B$, substitution $T=\sigma$ expresses $[5]$
as a multiple of $\sigma-1$.  Surjectivity of $[5]$ on
$\mathcal B$ proves the reverse inclusion.  Since this abelian variety
has no nonzero geometric two-primary point, (DR.9) gives
$(\sigma-1)\delta=0$.

It follows that $\sigma Q'+P'-Q'-\sigma P'$ is principal.  Its four
points are distinct: the two base fibres are disjoint and the deck
action is free.  A function with this divisor has degree two and
simple poles, so it is separable in characteristic two.  Thus $D$ is
hyperelliptic.  Its genus is six, so its hyperelliptic pencil is unique.
The order-five deck automorphism descends to an order-five
automorphism of that projective line, because the kernel of the
descent action has order at most two.  Over the algebraic closure it
has a fixed point.  Its fibre has at most two geometric points, and a
group of order five acts trivially on such a set.  This gives a fixed
point of $\sigma$ on $D$, contrary to etaleness of the cyclic cover.
This is the primitive-zero deck argument of A11/A18, with its complete
degree-five specialization included here.

\emph{5. The ordinary alternative and the two maps.}
It remains to consider $f_{\mathcal B}=4$.  By Step 2, $E/k'$ is an
ordinary genus-two curve with the specified Rosati-invariant
$\mathcal O_5$ action defined over $k'$.  Every morphism of curves induces a map on
the arithmetic fundamental groups used here.  Indeed a pullback of
an etale cover is etale, regardless of whether the morphism itself is
ramified, and the pullback of a geometric odd quotient still has odd
monodromy.  Thus both maps in (DR.4) take (DR.8) to equal point-section
classes on $E$.

Assume (DR.1), or only its stated real multiplication restriction.
Applying it first to $v$ and then to $v\circ\sigma$
gives
\[
                 v(P')=v(Q'),\qquad
                 v(\sigma P')=v(\sigma Q').
 \tag{DR.10}
\]
The geometric fibres of $v$ are exactly the $\tau$-orbits.  If
$Q'\ne P'$, the first equality gives $Q'=\tau P'$.  The second then
gives $Q'=\sigma^{-1}\tau\sigma P'=\sigma^{-2}\tau P'$.
Consequently $\sigma^2$ fixes $\tau P'$, impossible for the free
degree-five deck action.  Therefore $P'=Q'$, and hence $P=Q$.
This proves the asserted reduction.  The ordinary assertion is also
plainly a special case of the unrestricted genus-two assertion, so
these two universal finite-field assertions are equivalent.
\end{proof}

\paragraph{A finite canonical-class sufficient condition after transfer.}
The ordinary universal assertion in (DR.1) can be replaced, for a
particular prescribed $P$ and a particular cover in the ordinary
alternative, by the two individual hypotheses of A29 at
\[
                 R_0=v(P'),\qquad R_1=v(\sigma P').
 \tag{DR.11}
\]
Namely, for each $i=0,1$ suppose the class
$\kappa_{R_i}=[K_E-2R_i]\in J_E(k')$ has a nonzero $r_i$-primary
component for some odd prime $r_i$ with even
$\operatorname{ord}_{r_i}(2)$.  The primes may differ.  A29 separates
the section of each $R_i$ from every other rational point of $E$.
Applying that result to the two pushed collisions proves (DR.10), so
the same proof separates $P$ on $C$.

If $P$ is Weierstrass, choose the reflection that fixes $P'$.  Its
quotient fibre over $R_0$ has the single geometric point $P'$, so the
one hypothesis at $R_0$ suffices; no second detector is then needed.
Such a reflection exists because the five reflections act on the
five-point fibre as all maps $j\mapsto b-j$.  This gives a legitimate
odd-cover transfer even when $\kappa_P=0$ on the original curve.
An A29 cover of $E$ pulls back to an etale odd cover of $D$; after
taking the pointed connected component and an odd normal closure if
needed, its composition with $D\to C$ remains geometrically of odd
degree.  The ramified double quotient is used only as a morphism
through which to pull back such covers, not as an allowed etale
double cover in the original prime-to-two tower.

The displayed canonical-class hypotheses are additional finite-group
conditions.  The reflection construction does not prove them.  Nor
does it force a theta point on $C$ of order five or of any other
specified order.  Thus the reduction avoids the genus-three
truncation obstruction of A30, but leaves an ordinary genus-two
existence question in the older route.


\paragraph{The ordinary real-multiplication restriction is proper.}
There are ordinary genus-two curves over finite fields of characteristic
two whose geometric Jacobians admit no Rosati-invariant action of
$\mathcal O_5$, even after extending constants. An explicit example is
the smooth proper curve over $\mathbf F_4=\mathbf F_2(u)$ given by
\[
 C_0:\qquad y^2+y=x+\frac{u}{x}+\frac{u^2}{x+1},
                         \qquad u^2+u+1=0.
 \tag{DRP.1}
\]
Its Jacobian is absolutely simple and has geometric rational
endomorphism algebra a quartic field whose Rosati-fixed subfield is
$\mathbf Q(\sqrt2)$. Thus the additional restriction in A33 is an
actual proper class of ordinary curves, not merely a formally added
condition. This example makes no claim against point-section
injectivity on any curve.

\begin{proof}
The three poles $0,1,\infty$ in (DRP.1) are simple. The
Artin--Schreier genus and rank formula gives genus two and $2$-rank
two \cite[Lemma~2.6]{PriesZhu2012}. Each pole has one rational point
above it. For $x=u,u^2$, the right side is $u$, whose trace to
$\mathbf F_2$ is one. Therefore $\#C_0(\mathbf F_4)=3$.

For the second count write
$\mathbf F_{16}=\mathbf F_2[z]/(z^4+z+1)$ and $u=z^2+z$.
Encode a field element by its four binary coefficients in increasing
powers of $z$, so $u=6$ and $u^2=7$. Among the fourteen values
$x\in\mathbf F_{16}\setminus\{0,1\}$, the right side of
(DRP.1) has absolute trace zero exactly at
\[
                         x\in\{6,7,8,11,13,15\}.
 \tag{DRP.2}
\]
Each supplies two rational points. Hence
$\#C_0(\mathbf F_{16})=3+2\cdot6=15$.
The accompanying integer-only certificate retains all fourteen
right-side values and traces, and independently enumerates the
solutions of $y^2+y$ to verify the trace counts.

Let $\pi$ be $4$-power Frobenius on $A=J_{C_0}$. The two point
counts and the genus-two Newton identities give its characteristic
polynomial
\[
                    P(X)=X^4-2X^3+X^2-8X+16.
 \tag{DRP.3}
\]
Indeed its cubic coefficient is $3-5=-2$, and its quadratic
coefficient is $((-2)^2-(17-15))/2=1$.
Modulo three, $P$ is $\Phi_5(X)$, irreducible because
$\operatorname{ord}_5(3)=4$. Thus $P$ is irreducible over
$\mathbf Q$, and $K=\mathbf Q(\pi)$ has degree four.

We next exclude root-of-unity ratios between distinct eigenvalues.
The exact polynomial identity
\[
 \begin{aligned}
 \operatorname{Res}_X(P(X),P(tX))
   &=256(t-1)^4 A_0(t)^2 B_0(t),\\
 A_0(t)&=4t^4+7t^3+10t^2+7t+4,\\
 B_0(t)&=16t^4-8t^3+t^2-8t+16
 \end{aligned}
 \tag{DRP.4}
\]
can be checked as the determinant of $P(tM)$, where $M$ is the
four-by-four companion matrix of $P$. The certificate evaluates that
determinant by its complete permutation formula over $\mathbf Z[t]$.
Both $A_0$ and $B_0$ are primitive and reduce modulo three to
$\Phi_5$, so both are irreducible over $\mathbf Q$. Their monic
normalizations have nonintegral coefficients, so their roots are not
algebraic integers and in particular are not roots of unity.
The roots of the resultant are the quotients $\alpha_j/\alpha_i$
of the four roots of $P$. Since $P$ has distinct roots, the four
factors $t-1$ account exactly for the diagonal quotients. No quotient
between distinct roots is a root of unity. Consequently $\pi^m$
has four distinct eigenvalues for every $m\ge1$.

Here the standard integral Tate-module injectivity, rather than mere
faithfulness over $\mathbf Q$, is relevant. Milne's
Theorem~12.5 states that the geometric endomorphism group is a finite
free $\mathbf Z$-module and that its tensor product with
$\mathbf Z_\ell$ injects into the endomorphisms of $T_\ell A$
for $\ell\ne2$ \cite[Theorem~12.5, pp.~21--22]{MilneAV}.
Choose a finite field over which a basis of this endomorphism group
is defined; every basis element is defined over some finite field,
so one common extension exists. Its endomorphisms commute with some
$\pi^m$. On the four-dimensional rational Tate module the
centralizer of $\pi^m$ has dimension four: after extending scalars
it is the diagonal algebra for four distinct eigenvalues. The stated
injectivity therefore gives
\[
 \dim_{\mathbf Q}\operatorname{End}^0_{\overline{\mathbf F}_4}(A)
       \le4.
\]
Since this algebra already contains the quartic field $K$, it equals
$K$. In particular $A$ is absolutely simple: a nontrivial geometric
isogeny decomposition would give a nontrivial idempotent in this
field, by Poincare reducibility
\cite[Proposition~12.1, p.~19]{MilneAV}.

The Rosati involution for the canonical polarization satisfies
$\pi^\dagger=4/\pi$ \cite[Lemma~19.2, p.~41]{MilneAV}.
This is a nontrivial involution on $K$. Put $s=\pi+4/\pi$.
Dividing (DRP.3) by $\pi^2$ gives
\[
                         s^2-2s-7=0.
 \tag{DRP.5}
\]
Its discriminant is $32$, so $\mathbf Q(s)=\mathbf Q(\sqrt2)$.
This quadratic field is fixed by Rosati and is the entire fixed
field, since $[K:\mathbf Q]=4$ and the involution is nontrivial.
A Rosati-invariant embedding of $\mathcal O_5$ would extend to an
embedding of $\mathbf Q(\sqrt5)$ into $\mathbf Q(\sqrt2)$,
which is impossible. The assertion already concerns all geometric
endomorphisms, so no constant extension changes the conclusion.
\end{proof}


\paragraph{An ordinary Jacobian in the real-multiplication class.}
The ordinary curve already used in A13,
\[
 C/\mathbf F_2:\qquad y^2+y=x+x^{-1}+(x-1)^{-1},
 \tag{RME.1}
\]
has a Rosati-invariant integral action of
$\mathcal O_5=\mathbf Z[(1+\sqrt5)/2]$ on its Jacobian $A$,
defined over $\mathbf F_2$. Thus the ordinary Jacobian class
appearing in the refined assertion of A33 is nonempty.

\begin{proof}
The Frobenius polynomial checked in A13 is
$P_2(X)=X^4-X^2+4$. Let $\pi$ be $2$-power Frobenius and
let $V=\pi^\dagger=2/\pi$ be its Verschiebung, using the
canonical principal polarization
\cite[Lemma~19.2]{MilneAV}. Put $s=\pi+V$. The polynomial
identity for $\pi$ gives
\[
                              s^2=5.
 \tag{RME.2}
\]
It remains to justify integrality of $(s-1)/2$, including the
connected part of the two-torsion group scheme.

The three Weierstrass points over $0,1,\infty$ are rational.
Their three distinct pairwise differences give all nonzero
geometric two-torsion points, since the curve is ordinary of
genus two. Thus $A[2]_{\mathrm{red}}$ is the constant group
scheme $(\mathbf Z/2)^2$. Over the perfect field $\mathbf F_2$,
the reduced subscheme of a finite group scheme is a subgroup;
its map to the etale quotient is an isomorphism. In the
commutative case this splits the group scheme into its connected
and reduced factors. Cartier self-duality of $A[2]$, supplied
by the principal polarization, then identifies its connected
factor with the dual of $(\mathbf Z/2)^2$. Equivalently, the
dual split decomposition contains the connected subgroup
$\mu_2^2$ of rank four, which is already the full connected
rank. Consequently, over the stated field,
\[
                       A[2]\simeq(\mathbf Z/2)^2\times\mu_2^2.
 \tag{RME.3}
\]
Here Cartier duality is at the group-scheme level, not just on
geometric points; see \cite[Theorem~11.1 and the opening of
Section~16, pp.~18 and~29]{MilneAV}.

On the constant etale factor, $\pi$ is the identity, while on
$\mu_2^2$ it is zero, since Frobenius squares its coordinates.
Under Cartier duality and the principal polarization, $V$ has
the opposite two restrictions: zero on the etale factor and
identity on the multiplicative factor. Thus $s$ acts identically
on the entire group scheme $A[2]$.

The endomorphism $s-1$ therefore kills the kernel of $[2]$ as
a finite group scheme. By the quotient property of the isogeny
$[2]:A\to A$, it factors uniquely as $[2]T$ for an integral
endomorphism $T$ defined over $\mathbf F_2$. Substituting
$s=2T+1$ in (RME.2) and using torsion-freeness of the
endomorphism group gives $T^2+T-1=0$. Also $s^\dagger=s$, so
$T^\dagger=T$. As $X^2+X-1$ is irreducible over $\mathbf Q$,
the resulting unital action of its quadratic integer ring is
injective. This is the required integral $\mathcal O_5$ action.
\end{proof}

\paragraph{Its purely two-primary theta candidates.}
For this same curve the Raynaud theta divisor has exactly three
geometric two-primary points, namely the classes
$\mathcal O(W_i)\otimes B^{-1}$ of its three Weierstrass points.
Indeed A24 proves the exact relation
$\Phi^2-\Phi+4=0$ for $\Phi=\pi^2$, the $4$-power Frobenius.
On the geometric etale Tate module, $\Phi$ is invertible, so
this relation gives $\Phi=1$ modulo four. Applying A35 over
$\mathbf F_4$ places every two-primary theta point in
$A(\mathbf F_4)$. The verified count
$\#C(\mathbf F_4)=3$ leaves only those three Weierstrass
points in the Abel image. Their theta classes are distinct
two-torsion points because $2W_i\sim2B\sim K_C$; they are
nonzero since $B$ is non-effective on an ordinary curve.
In particular a non-Weierstrass point of this curve has a
nonzero odd-primary component in its canonical class. This
does not prescribe the prime divisors of that component or
prove that any one of them has even order of two.


\subsection{A bounded rational Prym exponent gives a cyclic-cover criterion}
\label{sec:bounded-prym-exponent}

\paragraph{Finding A34 (the actual-exponent criterion).}
Let $C/k$ be a smooth projective geometrically connected genus-two
curve over a finite field of characteristic two, let $P\in C(k)$,
and let $\pi:D\to C$ be a geometrically connected cyclic etale cover
of odd degree $n>1$, with constant deck group
$G=\langle\sigma\rangle$ and a rational point $P'\in D(k)$ above $P$.
Put
\[
 B=\ker(\operatorname{Nm}:J_D\longrightarrow J_C)^0,
 \qquad M=\exp\bigl(B(k)[2^\infty]\bigr),
 \tag{BE.1}
\]
where the exponent of the trivial group is one. If
\[
                         n>(2M-1)^2,
 \tag{BE.2}
\]
then the covering formulation of A1 separates the point section
of $P$ from every distinct $Q\in C(k)$.
There is no rank-one hypothesis, no exclusion of a separable
elliptic double quotient, and no requirement that $n$ be prime.
The exponent in (BE.1) is measured over the stated field $k$;
it must be measured again after a constant extension.

The conclusion is a sufficient criterion. The final corollary
below realizes $M=2$ and $n=19$: the criterion applies to that
cover even though its relative class number is even. This is
a strict extension of the rational-unit test for a specified
pointed cover over its stated field; it makes no nonexistence
claim about other covers of the same pointed curve.

\paragraph{Recovered and published infrastructure.}
The field generated by deck translates of a small function, its
Castelnuovo--Severi genus bound, and the resulting free elliptic
action already occur in recovered successor A.53
\cite[A.53]{ASrun}. Its stated nonabelian contradiction does not
apply to the cyclic group here. We retain that infrastructure and
give the cyclic four-point and wild-ramification argument in full.
The published input for the field bound is Poonen,
\emph{Gonality of modular curves in characteristic $p$},
Proposition~2.1, Lemma~2.2 and Corollary~2.3, author p.~3
\cite{Poonen2007Gonality}.

\begin{proof}
\emph{1. An exact torsion order gives a separable function.}
Suppose that $P$ and a distinct $Q\in C(k)$ collide in the
covering formulation of A1. Applied to the pointed cover in the
statement, that formulation supplies $Q'\in D(k)$ above $Q$ with
\[
                    \delta=[Q'-P']\in J_D(k)[2^\infty].
\]
Only this direction of the covering implication is used, as in
A01 and recovered A1 B.1 \cite[B.1]{A1run}.

One has $\operatorname{im}(\sigma-1)=B$.
Indeed the image is connected and lies in the norm kernel.
Conversely, $[n]-\sum_{i=0}^{n-1}\sigma^i$ is divisible by
$\sigma-1$ in the integral group ring; on $B$ it is $[n]$,
which is surjective. Thus
$t=(\sigma-1)\delta\in B(k)[2^\infty]$.
Let $m$ be its exact order, with $m=1$ when $t=0$.
Then $m\mid M$ and $m$ is a power of two. There is a function
$f$ on $D_{\bar k}$ with
\[
 \operatorname{div}(f)
       =m\bigl(\sigma Q'+P'-Q'-\sigma P'\bigr),
                  \qquad d:=\deg(f)=2m.
 \tag{BE.3}
\]
The four points are distinct, because the deck action is free
and the two fibres of $\pi$ are disjoint. If $m>1$ and $f$ were
a square in $\bar k(D)$, half the divisor in (BE.3) would be
principal, contradicting the exact order of $t$. For $m=1$,
the simple poles already exclude a square. Over the perfect
field $\bar k$, this proves that $f$ is separable.

\emph{2. The conjugate functions generate a small intermediate field.}
Set
\[
 K=\bar k(D),\qquad F=\bar k(\gamma f:\gamma\in G),
       \qquad e=[K:F],
\]
and let $Z$ be the smooth projective curve with function field $F$.
All generators have degree $d$. The cited Castelnuovo--Severi
consequence gives
\[
                 e\mid d,\qquad g(Z)\le(d/e-1)^2.
 \tag{BE.4}
\]
Precisely, Poonen's Corollary~2.3 applies because $F$ is
generated by degree-$d$ functions; its controlled-field
integer is $d/e$. Every extension here lies inside the
separable extension $K/\bar k(f)$.

Since $D\to C$ is etale, $g(D)=n+1$.
By (BE.2), $n>(d-1)^2$, so (BE.4) excludes $e=1$.
Thus $e>1$ is a power of two. The map $v:D_{\bar k}\to Z$
is separable of degree $e$, and $f=h\circ v$ for a nonconstant
function $h$ on $Z$.

\emph{3. The induced deck action on the intermediate curve is free.}
The field $F$ is $G$-stable. A kernel of the induced action
would have order dividing both $n$ and $e$, so the action on
$Z$ is faithful. It is free as well. A nontrivial stabilizer
would contain a subgroup of odd prime order $\ell$. That
subgroup acts freely on $D$ and preserves the fibre divisor
of $v$ above the fixed point. Its points are partitioned into
orbits of length $\ell$, with multiplicities constant on each
orbit. This would imply $\ell\mid e$, a contradiction.

Riemann--Hurwitz for $Z\to Z/G$ now gives
\[
                       g(Z)-1=n\bigl(g(Z/G)-1\bigr).
\]
The bound $0\le g(Z)<n$ from (BE.4) forces
$g(Z)=g(Z/G)=1$.
For the separable map $v$ its different divisor therefore has
degree
\[
                    \deg\operatorname{Diff}(v)=2g(D)-2=2n.
 \tag{BE.5}
\]

\emph{4. Wild ramification forces $e=2$.}
Choose a zero $z$ of $h$, and write $s=\operatorname{ord}_z(h)>0$.
Every point in the fibre $v^{-1}(z)$ must be one of the two
positive points in (BE.3), and the multiplicity of $f$ at either
one is $m$. Thus the fibre has one or two distinct points,
each with the same ramification index $m/s$, a power of two.
If $e>2$, that index is at least two. It is wild, and its
different exponent is at least its ramification index.

Here the last elementary local fact can be seen directly.
For uniformizers $u$ downstairs and $t$ upstairs over the
algebraically closed constant field, write
$u=a_qt^q+a_{q+1}t^{q+1}+\cdots$, where $q$ is the local
ramification index. If $q$ is even, the derivative of the
leading term vanishes, so
$\operatorname{ord}_t(du/dt)\ge q$; this derivative is nonzero
by separability. Its valuation is the different exponent.

Consequently the different contribution over $z$ is at least
the sum of the ramification indices, namely $e$.
The $n$ translates of $z$ under $G$ are distinct, since the
action on $Z$ is free. Equivariance gives the same lower bound
over each of them. If $e>2$, this would give
$\deg\operatorname{Diff}(v)\ge ne>2n$, contradicting (BE.5).
Hence $e=2$.

\emph{5. The two positive and two negative points exclude $e=2$.}
If the positive points $P'$ and $\sigma Q'$ had different
images in $Z$, the fibre over either image could contain
no other point: every such point would be an additional
zero of $f$. Since $v$ has degree two, both points would
be ramified. Their $G$-orbits are disjoint and each has $n$
points, because they lie above distinct points of $C$.
Every point in these two orbits has different exponent at
least two. This would give different degree at least $4n$,
again contradicting (BE.5).

Thus $P'$ and $\sigma Q'$ are the two distinct unramified
points of a single fibre. Ramification is invariant under
the deck action, so $Q'$ and $\sigma P'$ are unramified too.
Their pole fibres must likewise coincide: an unramified
degree-two fibre cannot be supported on only one pole.
Putting $x=v(P')$ and $y=v(Q')$, equivariance gives
\[
                         x=\sigma y,\qquad y=\sigma x.
 \tag{BE.6}
\]
Hence $\sigma^2x=x$. Since $n>1$ is odd, $\sigma^2$ is a
nonidentity element of $G$, contrary to the free action on
$Z$. This final contradiction proves A34.
\end{proof}

\paragraph{An exact way to measure the exponent at an inert cyclotomic prime.}
There is a useful arithmetic specialization, which does not
produce a cover. Let now $C/\mathbf F_Q$ have genus two in
characteristic $p$, let $P$ be rational, and let a constant
$P$-split cyclic cover have odd prime degree $r\ne p$, with
$\operatorname{ord}_r(p)=r-1$. Write $b:J_C(\mathbf F_Q)\to
\mathbf Z/r$ for its rational quotient and let $\chi$ be a
faithful character. The genus-two Riemann--Roch calculation
of A29 gives
\[
 L_\chi(T)=1+c_\chi T+Q\chi(b(\kappa_P))T^2,
 \quad c_\chi=\sum_{R\in C(\mathbf F_Q)}\chi(b([R-P])),
 \quad \kappa_P=[K_C-2P].
 \tag{BE.7}
\]
Put $E=\mathbb Q(\zeta_r)$ and normalize the valuation at
its unique prime above $p$ by $v(p)=1$. Then the rational
$p$-primary group of the whole Prym has the exact form
\[
 B(\mathbf F_Q)[p^\infty]\simeq
       \mathbb Z_p[\zeta_r]/p^e\mathbb Z_p[\zeta_r],
       \qquad e=v\bigl(L_\chi(1)\bigr),
 \tag{BE.8}
\]
including the zero group when $e=0$. Its exponent is $p^e$.

Here is a proof that avoids assuming a rank-one crystalline
comparison. The primitive-root condition makes
$\mathcal O=\mathbb Z_p[\zeta_r]$ the ring of integers of
the unramified extension of degree $r-1$. It acts on $B$,
because $1+\sigma+\cdots+\sigma^{r-1}$ vanishes on
$B=\operatorname{im}(\sigma-1)$. A finite $\mathcal O$-module
is a direct sum of modules $\mathcal O/p^{a_i}$.
Each nonzero summand contributes $r-1$ dimensions to its
subgroup killed by $p$. But
$B(\mathbf F_Q)[p]\subset B(\overline{\mathbf F}_p)[p]$,
whose dimension is the geometric $p$-rank, at most
$\dim B=r-1$. Thus at most one summand occurs.
Equivalently, the geometric etale Tate module is free over
$\mathcal O$ of rank zero or one; the finite-module argument
already suffices for the exponent assertion.

Finally the character product and the Jacobian norm isogeny
give
\[
 \#B(\mathbf F_Q)=N_{E/\mathbb Q}\bigl(L_\chi(1)\bigr).
\]
Since $p$ is inert and unramified in $E$, its valuation is
$(r-1)v(L_\chi(1))$. Comparing this with the order
$p^{(r-1)e}$ of the unique possible module proves (BE.8).
The same reasoning also proves that the exponent is one
when the rational $p$-primary group is zero.

\paragraph{The Weil valuation bound and its exceptions.}
Every conjugate of $c_\chi$ has absolute value at most
$2\sqrt Q$. Indeed its degree-two character factor is a
factor of the zeta numerator of the smooth etale cover,
whose Frobenius eigenvalues have absolute value $\sqrt Q$;
see \cite[Theorem~1.6]{WeilI}. Hence, putting $M_p=p^e$,
the algebraic-integer norm argument always gives
\[
                         M_p\le(\sqrt Q+1)^2.
 \tag{BE.9}
\]
For characteristic two and $Q\ge8$ it gives the sharper
dichotomy
\[
 \begin{cases}
 c_\chi=-1,\quad b(\kappa_P)=0,
                  \quad L_\chi(1)=Q,\quad M_2=Q;\\[2pt]
 \text{or}\qquad M_2\le2\sqrt Q+1.
 \end{cases}
 \tag{BE.10}
\]
The exceptional first alternative has all faithful
$L$-polynomials equal to $1-T+QT^2$.

For completeness, write $Q=p^a$.
Since $p$ is the unique prime over $p$, divisibility by
$p^e$ in the local ring means divisibility by the rational
integer $p^e$ in the ring of integers of $E$.
If a nonzero algebraic integer has all conjugates of
absolute value less than one, its nonzero integral norm
would have absolute value less than one, which is impossible.
Applying this to $L_\chi(1)/p^e$ proves (BE.9).
If $e\le a$, then $p^e$ divides $1+c_\chi$. Unless
$c_\chi=-1$, the same argument gives $p^e\le2\sqrt Q+1$.
If $e>a$ and $Q>2\sqrt Q+1$, applying it instead to
$(1+c_\chi)/Q$ forces $c_\chi=-1$; then
$L_\chi(1)=Q\chi(b(\kappa_P))$ has valuation $a$, a
contradiction. For $p=2$, $Q\ge8$ ensures the required
strict inequality. When $c_\chi=-1$, all its conjugates,
including $c_{\chi^{-1}}$, equal $-1$. The exact identity
$c_\chi=\chi(b(\kappa_P))c_{\chi^{-1}}$ from A29 forces
$b(\kappa_P)=0$, proving the exceptional alternative.
For $Q=2,4$, (BE.9) only gives the respective power-of-two
bounds $M_2\le4,8$; these small-field exceptions must not
be silently absorbed into (BE.10).

The exceptional alternative is already restricted by the
split fibre. Let $N=\#C(\mathbf F_Q)$ and let $N_i$ be its
Abel-image occupancies in the chosen quotient $\mathbf Z/r$.
The equality $1+\sum_{i=0}^{r-1}N_i\zeta_r^i=0$ implies
\[
 1+N_0=N_1=\cdots=N_{r-1}=\lambda,\qquad
 N+1=\lambda r.
 \tag{BE.11}
\]
Since $P$ occupies the zero coset, $N_0\ge1$ and
$\lambda\ge2$, so $r\le(N+1)/2$. If $N$ is even, then
$\lambda$ is odd and $\lambda\ge3$, giving
$r\le(N+1)/3$.

\paragraph{What the uniform valuation estimate does not improve.}
For $Q\ge8$ outside the exceptional case, replacing the
actual exponent in A34 by the bound in (BE.10) gives the
sufficient but coarse inequality
\[
                         r>(4\sqrt Q+1)^2.
 \tag{BE.12}
\]
This yields no improvement over the earlier B.3 size test.
To check that claim exactly, write $J=\#J_C(\mathbf F_Q)$,
$N=\#C(\mathbf F_Q)\ge1$, and let $a_2$ be the middle
coefficient of its genus-two zeta numerator. Weil gives
$N\le Q+1+4\sqrt Q$ and $a_2\le6Q$, while direct evaluation
of the numerator gives
\[
             J=(Q+1)N+a_2-2Q
                   \le(Q+1)N+4Q\le(5Q+1)N.
 \tag{BE.13}
\]
Condition (BE.12) implies both $r>N$ and $r>10Q+2$.
Consequently $N>2J/r$. If the rational quotient exists,
its kernel has order $J/r$, so both hypotheses
$r>N$ and $N\ge2|\ker b|$ of recovered A1 B.3 hold
\cite[B.3]{A1run}. In the even-$N$ case A26 already uses
the simpler inequality $r>N$.

For $Q=2,4$, the coarse bounds $M_2\le4,8$ would demand
$r>49,225$, respectively, whereas the Weil bound gives
$J<34$ and $J\le81$, respectively; no such rational
quotient can exist. In the exceptional case $M_2=Q$,
(BE.11) and the Weil bound give
$r\le(Q+2+4\sqrt Q)/2<(2Q-1)^2$ for $Q\ge2$, so A34's
inequality cannot hold there either.

Thus the potential gain is in measuring the actual rational
exponent in (BE.1), not in inserting only a uniform archimedean
bound. A positive density of primitive-root primes does not
give a constant pointed quotient at one of those primes, nor
does it supply a bound for this exponent after the cover and
its decks have been defined. Passing to a constant-deck field
$\mathbf F_{Q^d}$ replaces the square-root term by $Q^{d/2}$.
Before that passage Frobenius is semilinear for the deck action,
so the constant cyclic quotient and character-polynomial
argument above cannot simply be used over the original field.
No uniform control of that field degree or actual exponent is
proved here; no proof through Hooley follows from A34 alone.

\paragraph{A realized exponent-two witness of degree nineteen.}
The rank-one curve of A27 gives a strict example for the
actual-exponent criterion. Write
\[
 k_0=\mathbf F_{64}=\mathbf F_2[z]/(z^6+z+1),\qquad
 u=z^5+z^4+z^3+z,
 \qquad u^2+u+1=0.
\]
In the reciprocal coordinates $X=1/t$, $Y=X^2y$, its equation
and its rational Weierstrass base point are
\[
 C:\quad Y^2+X^2Y=uX^5+X^3+uX,
              \qquad W=\infty.
 \tag{BE.14}
\]
Thus $W$ is the simple-pole point of the original A27 model.
The two elliptic Frobenius factors of A27, with traces $1,-2$
over $\mathbf F_4$, have third power sums $-11,16$.
Consequently
\[
 \#C(k_0)=60,\qquad
 \#J_C(k_0)=(65+11)(65-16)=3724=19\cdot196.
 \tag{BE.15}
\]
All field elements in the following certificate are encoded
by their binary coefficients in the displayed polynomial basis.
The reduced Mumford divisor
\[
 \Delta=\bigl(X^2+54X+8,\;49X+50\bigr)
 \tag{BE.16}
\]
has exact order nineteen. Multiplication by $196$ defines
the surjective quotient $b:J_C(k_0)\to\mathbf Z/19$ by
$[196]A=b(A)\Delta$. The occupancies of the rational points
under $R\mapsto b([R-W])$, in residue order $0,\ldots,18$, are
\[
 (6,2,2,2,5,2,5,2,2,5,5,2,2,5,2,5,2,2,2).
 \tag{BE.17}
\]
These finite statements are checked by two separately
implemented divisor calculations: binary-linear interpolation
and extended-GCD Cantor composition. Both retain every point
image, the entire nineteen-element subgroup, and the exact
finite-field and cyclotomic arithmetic. The second calculation
uses cyclic length-nineteen convolution, independently of the
first calculation's polynomial-remainder multiplication.

Pull back the Lang quotient belonging to $\ker b$ along the
Abel map $R\mapsto[R-W]$. The connectedness and split-fibre
argument of A12 gives a geometrically connected constant
cyclic cover $D\to C$ of degree nineteen, with a rational
point above $W$. Let $B$ be its whole Prym, of dimension
eighteen. Since $W$ is Weierstrass, $\kappa_W=0$, and a
faithful character has
\[
 L_\chi(T)=1+cT+64T^2,
      \qquad c=\sum_{j=0}^{18}N_j\zeta_{19}^j.
 \tag{BE.18}
\]
The same two calculations give
\[
 \operatorname{ord}(\bar c)=7,
 \qquad 1+c^7\equiv2\pmod4,
 \qquad 1+c^7\equiv6\pmod8
 \quad\text{in }\mathbb Z[\zeta_{19}].
 \tag{BE.19}
\]
Here $\bar c$ is reduced modulo two. The order calculation
uses $\operatorname{ord}_{19}(2)=18$, so $\Phi_{19}$ is
irreducible modulo two; in particular the local coefficient
ring $\mathcal O=\mathbb Z_2[\zeta_{19}]$ is unramified of
degree eighteen. Only the congruence modulo four is needed
for the exponent conclusion.

Set $k=\mathbf F_{64^7}$. To justify the extension step
without assuming a unit-root comparison, let
$L_\chi(T)=(1-\alpha T)(1-\beta T)$ and put
$s_j=\alpha^j+\beta^j$. Then
\[
 s_0=2,\quad s_1=-c,\quad
 s_j=-c\,s_{j-1}-64s_{j-2}.
 \tag{BE.20}
\]
Modulo four this gives $s_7\equiv(-c)^7$.
The character factor over $k$ is
$1-s_7T+64^7T^2$, so (BE.19) gives
\[
 L_{\chi,k}(1)\equiv1+c^7\equiv2\pmod4.
\]
Formula (BE.8), applied over this stated field, now proves
\[
 B(k)[2^\infty]\simeq\mathcal O/2\mathcal O
                  \simeq(\mathbf Z/2)^{18},
 \qquad M=2,
 \qquad v_2\bigl(\#B(k)\bigr)=18.
 \tag{BE.21}
\]
Thus the relative class number is even, whereas
$19>(2\cdot2-1)^2=9$. A34 separates the section of $W$
from every distinct rational point of $C_k$ using this same
cover. The rational-unit criterion rejects this particular
cover over $k$, so the asserted strict extension is realized.
This does not assert that every other cover over $k$ fails
the unit criterion, or that a uniform supply of such
bounded-exponent covers has been proved.


\paragraph{Exact behaviour under every constant extension.}
Keep this particular degree-nineteen cover over $\mathbf F_{64}$
and its Prym $B$.  For $m\geq1$, let $M_m$ be the exponent of
$B(\mathbf F_{64^m})[2^\infty]$, with exponent one for the zero
group.  Then
\[
 M_m=
 \begin{cases}
  1,&7\nmid m,\\
  2,&7\mid m\text{ and }m\text{ is odd},\\
  2^{v_2(m)+2},&14\mid m.
 \end{cases}                                                   \tag{BE.22}
\]
Consequently this same cover satisfies A34's sufficient inequality
\[
                    19>(2M_m-1)^2
                    \quad\Longleftrightarrow\quad14\nmid m.   \tag{BE.23}
\]
For every field in the successful cases it separates the section
of $W$ from every distinct rational point.  Failure of this
inequality when $14\mid m$ does not assert failure of point
separation by this cover or by other covers.

\begin{proof}
Use the already checked character polynomial
$L_\chi(T)=1+cT+64T^2$ and the complete unramified ring
$\mathcal O=\mathbf Z_2[\zeta_{19}]$ of degree eighteen.
By (BE.19), $c$ is a unit.  The reduction of
$X^2+cX+64$ has the two distinct roots $0$ and $\bar c$.
Hensel's lemma for complete local rings
\cite[Lemma~10.153.9, Tag~\texttt{04GM}]{Stacks04GM}
gives a unit root $u\in\mathcal O$; its other root is
$w=64/u$.  The root equation implies
\[
 u\equiv-c\pmod{64\mathcal O},\qquad
 \operatorname{ord}(\bar u)=7,
 \qquad u^7\equiv3\pmod{8\mathcal O}.                          \tag{BE.24}
\]
For the last congruence use $1+c^7\equiv6\pmod8$ from
(BE.19), so $(-c)^7\equiv3\pmod8$.

The deck action is already constant over $\mathbf F_{64}$.
Thus on the same faithful character space, extension of
constants to degree $m$ raises both Frobenius roots to their
$m$th powers.  The value of the resulting character factor is
\[
           L_{\chi,m}(1)=(1-u^m)(1-w^m).
                                                               \tag{BE.25}
\]
Since $v_2(w)=6$, the second factor is a unit.  Formula
(BE.8) applies over each stated extension and gives
\[
 \begin{gathered}
 B(\mathbf F_{64^m})[2^\infty]\simeq
      \mathcal O/2^{e_m}\mathcal O,
 \qquad e_m=v_2(1-u^m),\\
 v_2\bigl(\#B(\mathbf F_{64^m})\bigr)=18e_m.
 \end{gathered}
                                                               \tag{BE.26}
\]
No identification of a crystalline or etale Frobenius
eigenvalue is needed beyond the roots of the established
character polynomial and the finite-module proof of (BE.8).

If $7\nmid m$, (BE.24) gives $e_m=0$.
Otherwise write $m=7j$ and put $a=u^7$.
The congruence $a\equiv3\pmod8$ gives
$v_2(a-1)=1$ and $v_2(a+1)=2$.
For odd $j$, the factor $1+a+\cdots+a^{j-1}$ reduces to
one modulo two, so $v_2(a^j-1)=1$.
For even $j=2^t b$ with $t\geq1$ and $b$ odd, first
$v_2(a^2-1)=3$.  Each subsequent squaring increases this
valuation by exactly one: the previous even power is one
modulo eight, so its sum with one has valuation one.
The final odd power $b$ does not change the valuation,
by the same geometric-sum argument.  Hence
$e_m=t+2=v_2(m)+2$.  This proves (BE.22).

In its first two cases $M_m$ is respectively one or two,
so the right side of A34's inequality is respectively one
or nine.  If $14\mid m$, then $M_m\geq8$ and that right
side is at least $225$.  This proves (BE.23).
\end{proof}


\subsection{The two-primary theta locus is rational over a small field}
\label{sec:two-primary-theta-field}

\paragraph{Finding A35 (a finite constant-field bound).}
Let $C/\mathbf F_q$ be a smooth proper geometrically connected
genus-two curve in characteristic two. Put $X=C^{(1)}$, let $B$ be
the locally exact differential line bundle on $X$, and write
\[
 J=J_X,\qquad
 \Theta=\{L\in J:h^0(X,B\otimes L)>0\}
       =\{\mathcal O_X(R)\otimes B^{-1}:R\in X\}.
 \tag{T2F.1}
\]
Thus $B^{\otimes2}\simeq\omega_X$, with the same Frobenius-twist
normalization as A11 and A24. Let $f\in\{0,1,2\}$ be the
geometric $2$-rank, let $F$ denote $q$-power Frobenius, and let
$d$ be the order of its action on the finite geometric etale group
$J[4](\overline{\mathbf F}_q)\simeq(\mathbf Z/4)^f$; take $d=1$
if $f=0$. Then
\[
 \Theta(\overline{\mathbf F}_q)\cap J[2^\infty]
              \subset J(\mathbf F_{q^d}),
 \qquad
 \begin{cases}
 d=1,&f=0,\\
 d\in\{1,2\},&f=1,\\
 d\in\{1,2,3,4,6\},&f=2.
 \end{cases}
 \tag{T2F.2}
\]
Here $J[2^\infty]$ denotes geometric points of two-power order,
not the full nonreduced group scheme. In particular the rank-one
locus is rational over $\mathbf F_{q^2}$. In every rank it is
rational over $\mathbf F_{q^{12}}$; for an ordinary curve the
sharper field in (T2F.2) has degree at most six, but need not be a
subfield of $\mathbf F_{q^6}$.

The proof uses the two-torsion theta intersection and an elementary
matrix squaring argument. The general bound of two on a nonzero
theta translate is the genus-two Riemann--Roch argument already
used in A19 and A24. Recovered successor A.95 identifies the
rank-one theta characteristic explicitly in its $(3,1)$
Artin--Schreier model \cite[A.95]{ASrun}; no such model or rational
Weierstrass point is assumed here.

\begin{proof}
\emph{1. Nonzero two-torsion translates meet theta only in
two-torsion points.}
Work over $\overline{\mathbf F}_q$. The separable hyperelliptic
double map of $X$ has exactly $f+1$ distinct branch points, by
the Artin--Schreier $2$-rank formula
\cite[Lemma~2.6]{PriesZhu2012}. Write their points on $X$ as
$W_0,\ldots,W_f$. Each fibre divisor at such a point is $2W_i$,
so $2W_i\sim K_X$.

Every difference $[W_i-W_j]$ with $i\ne j$ is a nonzero
two-torsion point. It is nonzero because the Abel map of a curve
of positive genus is injective on geometric points. If $f=1$
this supplies the unique nonzero point of $J[2]$. If $f=2$ the
three unordered pairs supply three distinct points: equality
between two pairs sharing one index would make the other two
Weierstrass points equal in the Abel map, since signs do not
matter for two-torsion. These are all the nonzero points of
$J[2]\simeq(\mathbf Z/2)^f$.

Fix $T=[W_i-W_j]\ne0$. If
$z=\mathcal O_X(R)\otimes B^{-1}\in\Theta$ and
$z+T=\mathcal O_X(S)\otimes B^{-1}\in\Theta$, then
\[
                        R+W_i\sim S+W_j.
 \tag{T2F.3}
\]
If these two effective divisors are equal, their distinct fixed
points $W_i,W_j$ force $R=W_j$ and $S=W_i$. If they are
different, their degree-two class has at least two sections.
Riemann--Roch in genus two then forces that class to be canonical.
The canonical pencil is the hyperelliptic pencil; its fibre
through $W_i$ is $2W_i$. Thus $R=W_i$ and $S=W_j$.
Conversely both displayed possibilities do occur. Consequently
\[
 \Theta\cap(\Theta+T)
   =\{\mathcal O_X(W_i)\otimes B^{-1},
       \mathcal O_X(W_j)\otimes B^{-1}\}
                  \subset J[2]
 \quad\hbox{on geometric points}.
 \tag{T2F.4}
\]
The last inclusion follows from $2W_i\sim K_X\sim2B$.
This formulation uses the actual Raynaud theta divisor. On an
ordinary curve $B$ is non-effective, and replacing $B$ by a
Weierstrass point without accounting for the translation would
give a different theta divisor.

\emph{2. Frobenius cannot move a two-primary theta point.}
For $f=0$ the only geometric two-primary point is zero, so the
claim is immediate. Suppose $f>0$ and identify the geometric
etale Tate module with $\mathbf Z_2^f$. Its associated torsion
group is $(\mathbf Q_2/\mathbf Z_2)^f$, with the compatible
Frobenius action. By the definition of $d$, the matrix of $F^d$
has the form
\[
                        U=I+A,\qquad A\in4M_f(\mathbf Z_2).
 \tag{T2F.5}
\]
Suppose $z\in\Theta$ has exact order $2^a$ and is not fixed by
$U$. Write $z=2^{-a}v$ modulo $\mathbf Z_2^f$, with $v$ a
primitive integral vector. The minimum valuation of the entries
of $Av$ is a finite integer $t<a$, and $t\ge2$.

Set $A_j=U^{2^j}-I$. Induction using
\[
                         A_{j+1}=2A_j+A_j^2
 \tag{T2F.6}
\]
shows both
\[
 A_j\in2^{j+2}M_f(\mathbf Z_2),\qquad
                         v_2(A_jv)=t+j.
 \tag{T2F.7}
\]
Indeed the first assertion is immediate at $j=0$ and is preserved
by (T2F.6). For the second, $2A_jv$ has valuation $t+j+1$, while
$A_j^2v$ has valuation at least
$(j+2)+(t+j)=t+2j+2\ge t+j+2$. Hence their leading terms
cannot cancel, proving the asserted equality at $j+1$.

Choose $j=a-t-1\ge0$. Equation (T2F.7) says that
\[
                    T:=F^{d2^j}z-z
 \tag{T2F.8}
\]
is a nonzero point of exact order two. The line bundle $B$ and
its theta divisor are defined over $\mathbf F_q$, so both $z$
and $z+T$ lie on $\Theta$. By (T2F.4), $z\in J[2]$. But $U$
acts identically on $J[4]$, hence on $J[2]$, contradicting the
choice of $z$. Every two-primary theta point is therefore fixed
by $F^d$, which proves the inclusion in (T2F.2).

\emph{3. The small extension degrees.}
For rank one the automorphism group of the geometric four-torsion
is $(\mathbf Z/4)^\times$, of order two. For rank two, reduction
\[
                  \mathrm{GL}_2(\mathbf Z/4)
                          \longrightarrow\mathrm{GL}_2(\mathbf F_2)
\]
has a kernel of exponent two: $(I+2H)^2=I$ modulo four.
The target has element orders one, two, and three. Therefore
every element upstairs has order dividing two, four, or six.
This gives exactly the asserted list of possible $d$ and proves
all the field bounds.
\end{proof}

\paragraph{The precise finite residual set.}
For $R\in X(\overline{\mathbf F}_q)$ put
$z_R=\mathcal O_X(R)\otimes B^{-1}$. Its canonical class is
\[
                    \kappa_R=[K_X-2R]=-2z_R.
 \tag{T2F.9}
\]
Thus $\kappa_R$ is two-primary if and only if $z_R$ is
two-primary. By (T2F.2) and injectivity of the Abel map defined
by $B$, every such $R$ belongs to $X(\mathbf F_{q^d})$.
An exact enumeration over that one finite field therefore
captures this entire geometric residual set, even when arbitrary
constant extensions are allowed afterwards. The corresponding
statement for the original curve follows by the Frobenius twist.
Also $[R-\iota R]=2z_R$, so the same bound applies to those
hyperelliptic-pair differences that are two-primary.

This is a finite candidate statement, not a separation theorem.
It gives no bound of this kind for arbitrary two-primary
differences $[R-S]$, and it does not treat canonical classes
with a nonzero odd-primary part supported on primes at which
the order of two is odd. Those cases remain distinct from the
purely two-primary residual set considered here. No GRH or
prime-distribution assertion enters the proof.

\paragraph{The ordinary real-multiplication targets need only degree six.}
Suppose in addition that $C$ is ordinary and its Jacobian has an
action of $\mathcal O_5=\mathbf Z[(1+\sqrt5)/2]$ defined over
$\mathbf F_q$, as do the ordinary reflection targets in A33.
The action carries over to $J_X$ by Frobenius twist. Then
\[
 \Theta(\overline{\mathbf F}_q)\cap J[2^\infty]
                     \subset J(\mathbf F_{q^6}),
 \qquad
 \{R:\kappa_R\text{ is two-primary}\}
                     \subset X(\mathbf F_{q^6}).
 \tag{T2F.10}
\]
Here is the additional argument. The polynomial $T^2+T-1$ is
irreducible modulo two, so
$\mathcal O=\mathcal O_5\otimes\mathbf Z_2$ is the ring of
integers of the unramified quadratic extension of $\mathbf Q_2$.
The geometric etale Tate module has rank two over $\mathbf Z_2$.
It is a finite torsion-free $\mathcal O$-module: a nonzero
$\mathcal O$-annihilator would, on taking its norm, give a
nonzero $\mathbf Z_2$-annihilator. Since $\mathcal O$ is a
discrete valuation ring, the module is free, necessarily of
rank one over $\mathcal O$.

Frobenius commutes with this action because the endomorphisms are
defined over $\mathbf F_q$. It therefore acts by a unit of
$\mathcal O$. Reduction of $(\mathcal O/4)^\times$ modulo two
has quotient $\mathbf F_4^\times$, of order three, and kernel
$1+2\mathcal O$ modulo four, of exponent two. Thus
$(\mathcal O/4)^\times$ has exponent six, and $F^6$ acts
identically on $J[4]$. Applying the proof of (T2F.2) with this
power gives (T2F.10). This confines the purely two-primary case
on the reduced targets to one explicit finite field; it still
does not force a canonical class to have a good odd-primary
component or separate any of these remaining points.


\subsection{An ordinary reflection target need not carry a five-primary canonical detector}

\paragraph{Finding A36 (an actual failure of the forced-five shortcut).}
There is a constant, pointed cyclic etale cover of degree five of the
rank-one curve in A27 whose reflection quotient $E/\mathbf F_{256}$
is an ordinary genus-two curve with the Rosati-invariant
$\mathcal O_5$ action of A33 and
\[
 \#J_E(\mathbf F_{256})=288^2=2^{10}3^4.
 \tag{NF.1}
\]
Consequently every class $\kappa_R=[K_E-2R]$ with
$R\in E(\mathbf F_{256})$ has zero five-primary component.
Each such fixed class continues to have zero five-primary component
after every constant extension. The same is true of the two
correspondence constants described below. This disproves an inference
that the quintic transfer or its integral real multiplication must
force a five-primary canonical detector. It does not exclude a
three-primary detector, another useful odd prime on another target,
or five-primary classes belonging to points that first appear over
larger fields.

\begin{proof}
\emph{1. The finite quotient and its character polynomial.}
Use the field and reciprocal model already used in A32:
\[
 \begin{split}
 k&=\mathbf F_2[z]/(z^8+z^4+z^3+z+1),\qquad u=188,\quad u^2+u+1=0,\\
 C/k&:\quad Y^2+X^2Y=uX^5+X^3+uX,\qquad W=\infty.
 \end{split}
 \tag{NF.2}
\]
The integer encoding records binary coefficients in increasing powers
of $z$. The two elliptic Frobenius factors from A27 have fourth power
sums $17,-16$. Thus
\[
 \begin{aligned}
 \#C(k)=256,
 \qquad \#J_C(k)&=(257-17)(257+16)=65520,\\
                       65520&=5\cdot13104,
                       \qquad 5\nmid13104.
 \end{aligned}
 \tag{NF.3}
\]
The reduced Mumford class
\[
 \Delta=((224,80,1),(180,59))
 \tag{NF.4}
\]
has exact order five; coefficient lists in (NF.4) are in increasing
degree. Define a surjective homomorphism
$b:J_C(k)\to\mathbf Z/5$ by
$[13104]A=b(A)\Delta$. In residue order $0,1,2,3,4$, the
occupancies of $R\mapsto b([R-W])$ are
\[
                         (76,45,45,45,45).
 \tag{NF.5}
\]
The retained finite certificates specify the complete five-element
subgroup and every one of the $256$ point images. The extended-GCD
Cantor calculation independently reproduces the binary-linear
interpolation calculation, using no interpolation arithmetic.
Both preserve the curve equation during divisor reduction; the
invariant and reduction formulas are those explained in A32.

The Lang quotient for $b$, pulled back along the Abel map based at
$W$, gives a geometrically connected constant cyclic cover
$\pi:D\to C$ of degree five with a rational lift $P$ of $W$.
Indeed the Jacobian isogeny has connected domain and kernel
$\mathbf Z/5$, and the Abel map induces the required surjection
on the geometric abelian prime-to-two fundamental group, as in A12.
Its fibre above the identity, hence the fibre above $W$, is split.
For a faithful character $\chi(1)=\zeta_5$, the genus-two
Riemann--Roch calculation in recovered A1 B.2 gives
\cite[B.2]{A1run}
\[
 L_\chi(T)=1+c_\chi T+256T^2,
 \qquad
 c_\chi=76+45\sum_{i=1}^4\zeta_5^i=31.
 \tag{NF.6}
\]
Here the top coefficient is $256$ because $W$ is Weierstrass, so
$K_C\sim2W$. The same formula holds for all four faithful
characters. For the whole new factor
$\mathcal P=\ker(\operatorname{Nm}_\pi)^0$, character
factorization therefore gives
\[
                H_{\mathcal P}(T)=(1+31T+256T^2)^4.
 \tag{NF.7}
\]
Its reduction modulo two has degree four, the dimension of
$\mathcal P$, so $\mathcal P$ is ordinary. This uses the
standard Frobenius-polynomial and geometric rank comparison
already specified in A32
\cite[Propositions~2.2.4--2.2.5]{KatzCongruence}, not a
rational-class-number inference.

\emph{2. A reflection defined over the stated field.}
The hyperelliptic involution of $C$ lifts to reflections on $D$,
as in A33. Among these lifts there is a unique reflection $\tau$
fixing the chosen $k$-rational point $P$ above $W$: on the split
five-point fibre, the reflections act by $j\mapsto a-j$.
Galois preserves the cover and $P$, so this unique lift is
$k$-defined. Its quotient $v:D\to E=D/\langle\tau\rangle$
and the deck generator $\sigma$ are consequently defined over $k$.
The complete genus and isogeny calculation of A33 gives
\[
 \begin{aligned}
 g(E)&=2,\qquad f_E=2,\qquad \mathcal P\sim_k J_E\times J_E,\\
 T&=v_*\sigma v^*,\qquad T^2+T-1=0\quad\hbox{on }J_E.
 \end{aligned}
 \tag{NF.8}
\]
The embedding generated by $T$ is integral and fixed by Rosati.
Because the isogeny in (NF.8) is defined over $k$,
$H_E(T)^2=H_{\mathcal P}(T)$. Both sides have constant
coefficient one, so in the integral polynomial ring (NF.7) implies
\[
 H_E(T)=(1+31T+256T^2)^2
      =1+62T+1473T^2+15872T^3+65536T^4.
 \tag{NF.9}
\]
Evaluation at one proves (NF.1). Every $k$-rational canonical
class is an element of this finite group. Its order is therefore
prime to five, and the order of a fixed geometric point of an
abelian variety does not change upon extending its field of
definition. This proves the asserted persistent scope.
\end{proof}

\paragraph{The exact inhomogeneous correspondence identity.}
The degree-zero relation in (NF.8) is not the whole relation on
the Picard group. Work on $X=E^{(1)}$, carrying the correspondence
and the quotient map over by Frobenius twist; all the data remain
defined over $k$. Write $B$ for the locally exact differential line
bundle on $X$ attached to $E\to X$, the actual Raynaud theta
characteristic of degree one. Let $F\in\operatorname{Pic}^5(X)$
be the fibre class of the degree-five map
$X\to C^{(1)}/\langle\iota\rangle\simeq\mathbf P^1$.
Use $T$ also for the degree-doubling correspondence on the full
Picard group. For every $L\in\operatorname{Pic}^d(X)$ one has
\[
                       (T^2+T-1)L=dF.
 \tag{NF.10}
\]
Here the right side denotes addition of divisor classes, not a
degree-zero endomorphism.

For completeness, put $T_2=v_*\sigma^2v^*$. The exact
pullback--norm identity $v^*v_*=1+\tau$ and
$\sigma\tau\sigma=\tau$ give
$T^2=T_2+2$ on the Picard group. At a point $R=v(Q)$,
the divisor $R+T(R)+T_2(R)$ is the degree-five fibre over the
image of $\pi(Q)$ in the hyperelliptic projective line. On an
unramified fibre its five terms are $v(\sigma^jQ)$ for
$j=0,\ldots,4$. At a branch fibre the repeated images have
exactly their multiplicities: one point has multiplicity one and
the other two have multiplicity two. Thus the identity holds
without dividing by two. Extending additively to divisors and
then to line-bundle classes proves (NF.10).

Put
\[
             a=T(B)-2B\in J_X(k),
             \qquad b_0=F-5B\in J_X(k).
 \tag{NF.11}
\]
Substituting $L=B$ in (NF.10), and using $T(B)=2B+a$,
gives the exact identities
\[
       b_0=(T+3)a,
       \qquad (2-T)b_0=5a,
       \qquad (2-T)(T+3)=5\quad\hbox{on }J_X.
 \tag{NF.12}
\]
In particular the two constants have the same support away from
five: both $2-T$ and $T+3$ are invertible on every finite
$\ell$-primary group for $\ell\ne5$, by the last identity.
These linear relations alone do not produce a nonzero five-primary
part. In the actual example above the Frobenius twist has the same
number of $k$-rational Jacobian points as $E$, and (NF.1) shows
that the five-primary parts of both $a$ and $b_0$ are zero.
They remain zero for these fixed constants after all extensions.

Thus the proposed inhomogeneous constants do not repair the
forced-five inference. The good-prime criterion A29 still has a
separate possible three-primary application on this target, and
no absence of such a detector is asserted here.


\subsection{Finite-support avoidance persists with real multiplication by
\texorpdfstring{$\mathcal O_5$}{O5}}
\label{sec:rm5-avoidance}

\paragraph{Finding A37 (the ordinary real-multiplication locus).}
Put $F=\mathbf Q(\sqrt5)$ and
$\mathcal O=\mathcal O_5=\mathbf Z[(1+\sqrt5)/2]$.
Consider smooth genus-two curves in characteristic two together with
a specified embedding $\mathcal O\hookrightarrow\operatorname{End}(J_C)$
fixed by the Rosati involution of the canonical principal polarization.
On the ordinary Hilbert--Blumenthal moduli locus of these objects the
following assertions hold.

\begin{enumerate}
\item For its geometric generic curve $C$, the Raynaud theta support on
$J=\operatorname{Pic}^0(C^{(1)})$ contains no nonzero point of odd order.
\item For every finite set $S$ of odd primes there is an ordinary such
curve over a finite field $\mathbf F_Q$ with
\[
 \Theta(\overline{\mathbf F}_Q)\cap
 \bigl(J(\overline{\mathbf F}_Q)\{S\}\setminus\{0\}\bigr)
                         =\varnothing.                       \tag{RM.1}
\]
The chosen $\mathcal O$ action and the curve are defined over the same
field $\mathbf F_Q$.  The notation $J\{S\}$ permits all powers and all
products of primes in $S$.
\end{enumerate}

The finite-field examples can be chosen in every prescribed nonempty
open of the ordinary locus with its specified real-multiplication
embedding and principal polarization.  Thus they are Zariski dense
there.  In particular neither $\{3,5\}$ nor any other fixed finite odd
prime support becomes a universal theta-witness supply after the
ordinary real-multiplication reduction of A33.  These assertions are
unconditional; no bound for $Q$ is claimed.

\begin{proof}
\emph{1. The precise moduli input, including characteristic two.}
For an odd integer $h\geq3$, let $H_h$ be the characteristic-two
geometric fibre of the Hilbert--Blumenthal scheme
$M_{(\mathcal O,\mathcal O^+),h}$, with a fixed primitive $h$th root of
unity and full $\mathcal O$-linear symplectic level structure.
Here $\mathcal O^+$ denotes the totally positive elements.
The polarization type $L=\mathcal O$ is exactly an
$\mathcal O$-linear principal polarization: see Deligne--Pappas,
Section~2.6, printed p.~65 \cite{DelignePappas1994}.
Their Theorem~2.2, printed p.~64, gives smoothness since
$2\nmid h\operatorname{disc}(F)=5h$.
Yu's Theorem~1.1, explicitly attributed to Rapoport and
Deligne--Pappas, gives geometric irreducibility for every full level
prime to the residue characteristic \cite[Theorem~1.1, author
manuscript p.~2]{Yu2009HB}.  Thus $H_h$ is smooth and geometrically
irreducible.  The hypotheses require that the level be odd;
they do not require that it be prime to five.

The genus-two Torelli map is an immersion over $\operatorname{Spec}
\mathbf Z$, including characteristic two, by Landesman,
Theorem~1.2 \cite{Landesman2021Torelli}.  In characteristic two its image
is the complement of the product locus $\mathcal A_{1,1}$; the exact
statement is Cl\'ery--van der Geer, Section~4, printed p.~5211
\cite{CleryVanderGeer2022}.  Consequently its pullback to the
smooth, reduced scheme $H_h$ is the open genus-two Jacobian locus,
and this open carries a universal genus-two curve.  More explicitly,
an immersion with this open underlying image becomes the open
subscheme after this base change: its defining ideal has radical
zero on a reduced scheme and is therefore zero.
All genus-two curves are hyperelliptic, so Landesman's
$\mathcal H_2$ is the genus-two curve stack here.

Intersect with the ordinary open and denote the resulting open by
$U_h$.  It is nonempty: the A13 curve, with the integral
Rosati-invariant action proved in the supplement to A33, gives a
geometric point, and a full level-$h$ structure can be chosen over
the algebraic closure.  There is therefore a universal smooth proper
curve $\mathcal C\to U_h$ with its principally polarized Jacobian
$\mathcal A$ and its specified $\mathcal O$ action.  This construction
produces the curve before any finite-field Frobenius specialization;
no extension of the eventual specialization field is needed to
recover a curve from its Jacobian.

\emph{2. Full joint finite-level monodromy and the prime five.}
Fix an odd integer $n>1$ and choose the auxiliary level $h\geq3$
prime to $n$.  Work first over $\overline{\mathbf F}_2$, with compatible
choices of roots of unity.  The forgetful map
\[
                      H_{hn}\longrightarrow H_h             \tag{RM.2}
\]
is the torsor of full $\mathcal O$-linear symplectic level-$n$
structures.  Its group is
\[
                         G_n=\operatorname{SL}_2(\mathcal O/n\mathcal O).
                                                               \tag{RM.3}
\]
Here is the pairing convention, which matters when $5\mid n$.
Let $\mathfrak d^{-1}$ be the inverse different of $F/\mathbf Q$.
The principal polarization determines a perfect alternating
$\mathcal O/n\mathcal O$-bilinear pairing with values in
$\mathfrak d^{-1}\otimes\mu_n$; see Deligne--Pappas,
Sections~2.11--2.13, printed pp.~66--67
\cite{DelignePappas1994}.  After choosing a generator of this
rank-one module, the form is the standard determinant form on a
free rank-two module.  Its $\mathcal O$-linear automorphism group
is (RM.3).  Equivalently, since
$\mathfrak d^{-1}=\frac1{\sqrt5}\mathcal O$, the underlying
$\mathbf Z/n$ pairing is obtained from
\[
 (x,y)\longmapsto
   \operatorname{Tr}_{F/\mathbf Q}
       \bigl((x_1y_2-x_2y_1)/\sqrt5\bigr)\pmod n.             \tag{RM.4}
\]
The inverse-different trace duality makes this form perfect even
at five.  Using the trace pairing on $\mathcal O$ without the
inverse different would not justify that assertion.

Both schemes in (RM.2) are geometrically irreducible by the cited
theorem.  The inverse image of $U_h$, or of any nonempty open in
$U_h$, is a nonempty open of $H_{hn}$ and hence is irreducible.
Therefore the geometric monodromy of $\mathcal A[n]$ on every
such open is all of $G_n$.  This is a joint statement for the
entire integer $n$, so no independence assertion between prime
levels is being assumed.

Put $\mathcal X=\mathcal C^{(1)}$ and
$\mathcal J=\operatorname{Pic}^0(\mathcal X/U_h)$.
Relative Frobenius identifies the odd-torsion local systems of
$\mathcal A$ and $\mathcal J$ by an $\mathcal O$-linear
symplectic similitude of multiplier two.  Conjugation by this
similitude preserves the determinant-one group.  Thus the same
full $G_n$ monodromy holds for $\mathcal J[n]$.

\emph{3. Every nonzero local orbit contains an odd translation cycle.}
Write the finite ring as a product of local principal ideal rings
\[
 \mathcal O/n\mathcal O=\prod_{\mathfrak p\mid n}
                 R_{\mathfrak p},\qquad
 R_{\mathfrak p}=\mathcal O/\mathfrak p^{a_{\mathfrak p}}.
                                                               \tag{RM.5}
\]
Let $z$ be a nonzero element of $(\mathcal O/n\mathcal O)^2$.
Choose a factor in which its component is nonzero, write that
factor as $R=\mathcal O/\mathfrak p^a$, and let $\ell$ be its
odd residue characteristic.  Choose a uniformizer $\varpi$ of
the local ring.  A matrix in $\operatorname{SL}_2(R)$ sends
this component to $\varpi^v e_1$ for some $0\leq v<a$:
divide both coordinates by their largest common power of
$\varpi$, complete the resulting unimodular column to a
determinant-one basis, and invert that basis matrix.

After this change of basis, let $T$ have component
$\varpi^{a-1}e_2$ in the selected factor and zero in every
other factor.  It is nonzero and has additive order $\ell$.
For $j\in\mathbf Z/\ell$, the lower unipotent matrix with
lower-left entry $j\varpi^{a-1-v}$ sends the selected
component of $z$ to
\[
                    \varpi^v e_1+j\varpi^{a-1}e_2.             \tag{RM.6}
\]
Act as the identity in all remaining factors.  It follows that
the $G_n$ orbit of $z$ contains the entire cycle
$z+\langle T\rangle$, consisting of $\ell\geq3$ distinct points.
This proof covers split and inert primes.  It also covers five:
there $\mathcal O/5^r\mathcal O$ is a local ring of length
$2r$ in its uniformizer, and the same argument applies without
any assumption that the ring be reduced.

\emph{4. Generic theta avoidance.}
Let $K$ be the function field of $U_h$ over
$\overline{\mathbf F}_2$ and let $\Theta$ be the Raynaud theta
support of the generic curve.  It is defined over $K$.
If it contained a nonzero point $z\in J[n](\overline K)$,
Galois invariance and Step~2 would put the entire $G_n$ orbit
of $z$ in $\Theta$.  Step~3 then supplies a nonzero $T$ of
odd prime order $\ell$ such that
\[
                   z+\langle T\rangle
                         \subset\Theta\cap(\Theta+T).         \tag{RM.7}
\]
For a genus-two translated Abel curve the intersection with
a distinct nonzero translate has at most two geometric points,
as proved in A19.  Equation (RM.7) contradicts $\ell\geq3$.
This proves absence of all nonzero $n$-torsion at the geometric
generic point for each fixed odd $n$.

The auxiliary level can be chosen prime to each integer under
consideration.  Forgetting this finite auxiliary level does not
change the geometric generic curve or whether its theta contains
a geometric torsion point.  Thus the conclusion holds for every
odd $n$ on the geometric generic curve of the original
real-multiplication locus.  It is not an assertion that a
single open set excludes infinitely many different levels.

\emph{5. A finite avoidance open and an exact Frobenius class.}
Let $S$ be nonempty, put $N=\prod_{\ell\in S}\ell$, and choose
$h$ prime to $N$.  Begin in the inverse image of any prescribed
nonempty ordinary moduli open.  On its universal curve the
relative sequence
\[
 0\longrightarrow\mathcal O_{\mathcal X}
 \longrightarrow F_*\mathcal O_{\mathcal C}
 \longrightarrow\mathcal B\longrightarrow0                 \tag{RM.8}
\]
has a line-bundle quotient of fibre degree one.  As in A22,
the image of the proper relative Abel map
$R\mapsto\mathcal O(R)\otimes\mathcal B^{-1}$ is precisely
the relative Raynaud theta support.  Its intersection with
the finite etale scheme $\mathcal J[N]\setminus\{0\}$ has
closed image on the base.  Step~4 shows that this image
misses the generic point.  Remove it and choose a nonempty
smooth affine open $U$.  Every geometric fibre over $U$ has
\[
                         \Theta_t\cap
                    (\mathcal J_t[N]\setminus\{0\})
                              =\varnothing.                  \tag{RM.9}
\]

Descend the family and this open to a finite field
$k_0=\mathbf F_{q_0}$ large enough that all the data are
defined, $U$ is geometrically integral, and
$q_0\equiv1\pmod{hN^2}$.  The full
$\mathcal O$-linear symplectic level-$N^2$ torsor
\[
                      V\longrightarrow U,
 \qquad G=\operatorname{SL}_2(\mathcal O/N^2\mathcal O),         \tag{RM.10}
\]
is finite etale and geometrically integral by Step~2.
For every $\ell\in S$, set
\[
 u_\ell=1+\ell\pmod{\ell^2},\qquad
 \gamma_\ell=
   \operatorname{diag}(u_\ell,u_\ell^{-1})
       \in\operatorname{SL}_2(\mathcal O/\ell^2\mathcal O).
                                                               \tag{RM.11}
\]
These are $\mathcal O$-linear matrices, including at five.
Each has order $\ell$; by the Chinese remainder theorem
$\gamma=(\gamma_\ell)_{\ell\in S}\in G$ has order $N$.
On the underlying rank-four $\mathbf Z/\ell^2$ module,
\[
 \gamma_\ell\equiv I\pmod\ell,
 \qquad (\gamma_\ell-I)/\ell\pmod\ell
             \text{ is invertible}.                          \tag{RM.12}
\]
Indeed on the rank-two $\mathcal O/\ell$ module the latter
map is $\operatorname{diag}(1,-1)$.  Both entries remain
units when $\mathcal O/5$ is nonreduced.

For completeness, the finite-field realization is exactly the
affine twisting argument of A22, with the group in (RM.10).
Let $\mathcal F$ be arithmetic $q_0$-power Frobenius on $V$.
The action $v\mapsto\mathcal F(v)\gamma^{-1}$ defines a
finite descent datum over $\mathbf F_{q_0^N}/k_0$.
Since $V$ is affine it is effective by the Stacks Project,
Lemma~35.37.1, Tag~\texttt{0245} \cite{Stacks0245}.
The resulting twist is geometrically integral.  Lang--Weil,
in the form of Poonen, Theorem~7.7.1(ii)--(iii), printed
p.~222 \cite{Poonen2017RPV}, supplies a rational point over
$\mathbf F_{q_0^m}$ for every sufficiently large $m$.
Choose $m\equiv1\pmod N$.  Its underlying frame satisfies
\[
              \mathcal F^m(v)=v\gamma^m=v\gamma.               \tag{RM.13}
\]
It lies over $t\in U(\mathbf F_Q)$, $Q=q_0^m$, and the
universal family gives the required ordinary curve and
its specified real multiplication over this exact field.

\emph{6. From a finite level to every exponent and product.}
Let $J$ be the Frobenius-twisted Jacobian of this fibre and
$\Phi$ its $Q$-power Frobenius.  Lifting the torsion frame
to its $\ell$-adic Tate module, (RM.12)--(RM.13) give
\[
             \Phi-I=\ell A_\ell,
             \qquad A_\ell\in\operatorname{GL}_4(\mathbf Z_\ell).
                                                               \tag{RM.14}
\]
Hence, on $(\mathbf Q_\ell/\mathbf Z_\ell)^4$,
$\ker(\Phi-I)=\ker[\ell]$, and therefore
\[
 J(\mathbf F_Q)[\ell^\infty]=J[\ell],\qquad
                       J(\mathbf F_Q)\{S\}=J[N].              \tag{RM.15}
\]
In particular Frobenius fixes every $J[\ell]$ for
$\ell\in S$.  Finding~A19 with $d=1$ now puts every
geometric $S$-primary theta point in this rational group.
Equation (RM.9) excludes its nonzero elements and proves
(RM.1).  If $S$ is empty the assertion is immediate.

The construction began in an arbitrary prescribed nonempty
open, and all subsequent operations preserve a nonempty
open before specialization.  This proves the density
assertion as well.
\end{proof}

\paragraph{Scope of the obstruction.}
A37 rules out a uniform finite odd prime support even inside
the ordinary Rosati-$\mathcal O_5$ locus singled out by A33.
It does not exclude theta witnesses whose prime factors
depend on the individual curve.  It also makes no claim
that the particular explicit curves in A13 or A36 avoid
an arbitrary prescribed support; those curves provide
nonemptiness of the family, whereas the examples in
(RM.1) are obtained by the stated specialization.


\subsection{Ordinary real multiplication does not force a good canonical prime}

\paragraph{Finding A38 (failure of a universal canonical forcing statement).}
Let $C/\mathbf F_2$ be the ordinary genus-two curve of A13,
\[
                   y^2+y=x+x^{-1}+(x+1)^{-1}.
 \tag{BC.1}
\]
Its Jacobian has the integral Rosati-invariant $\mathcal O_5$
action proved in the supplement to A33. Among its $129$
$\mathbf F_{128}$-rational points, exactly three are Weierstrass.
For every one of the other $126$ points $R$,
\[
                  \operatorname{ord}([K_C-2R])=4159,
             \qquad \operatorname{ord}_{4159}(2)=2079
                       \quad\hbox{is odd}.
 \tag{BC.2}
\]
Consequently none of these canonical classes has a nonzero
$r$-primary component for any odd prime $r$ with even
$\operatorname{ord}_r(2)$. This remains true for each of these
fixed points after every constant extension. Thus integral
ordinary real multiplication by $\mathcal O_5$, even with a
non-Weierstrass prescribed point, does not force the sufficient
canonical-prime hypothesis of A29.

\begin{proof}
The three simple Artin--Schreier poles $0,1,\infty$ give
genus two and ordinarity; their ramification points are the
three rational Weierstrass points. The checked polynomial
in A13 is
\[
                         P_2(Z)=Z^4-Z^2+4.
 \tag{BC.3}
\]
Its roots occur in opposite pairs. Every odd power sum is
therefore zero, and the Frobenius point-count formula gives
$\#C(\mathbf F_{128})=128+1=129$
\cite[Theorem~11.1]{MilneJV}.
For completeness, let $\beta_1,\beta_2$ be the two roots of
$Z^2-Z+4$ and put $t_m=\beta_1^m+\beta_2^m$.
The recurrence $t_0=2$, $t_1=1$, $t_m=t_{m-1}-4t_{m-2}$
gives
\[
               (t_0,\ldots,t_7)=(2,1,-7,-11,17,61,-7,-251).
\]
Since seven is odd, pairing the opposite eigenvalues in
the Jacobian point-count product
\cite[Theorem~19.1(b)]{MilneAV} gives
\[
 \#J_C(\mathbf F_{128})
    =(1-\beta_1^7)(1-\beta_2^7)
    =4^7+1-t_7=16636=4\cdot4159.
 \tag{BC.4}
\]
Trial division through $\lfloor\sqrt{4159}\rfloor=64$
proves that $4159$ is prime. The exact residues
\[
 \begin{gathered}
 2079=3^3\cdot7\cdot11,\qquad
 2^{2079}\equiv1,\\
 2^{693}\equiv2554,\qquad
 2^{297}\equiv1245,\qquad
 2^{189}\equiv102\pmod{4159}
 \end{gathered}
 \tag{BC.5}
\]
certify the exact multiplicative order in (BC.2).
These are the arithmetic inputs already present in A13;
the retained additional certificate independently rechecks
them and directly enumerates all $126$ nonbranch points.

Write $G=J_C(\mathbf F_{128})$. The three distinct pairwise
differences of the rational Weierstrass points are all the
nonzero geometric points of $J_C[2]$, because the curve has
two-rank two. Thus $G[2]$ has order four and is
$(\mathbf Z/2)^2$. By (BC.4) it is already the entire
two-primary subgroup of $G$, which is therefore killed by two.
Choose a rational Weierstrass point $W$. Since $K_C\sim2W$,
\[
                         \kappa_R=[K_C-2R]=2[W-R].
 \tag{BC.6}
\]
Its two-primary part vanishes. If $R$ is non-Weierstrass,
this class is nonzero: otherwise $2R\sim K_C$ would be a
fibre of the canonical hyperelliptic pencil, forcing $R$ to
be a ramification point. Consequently its order is the
prime $4159$, proving (BC.2). Extension of constants preserves
the order of each fixed geometric divisor class, so it cannot
introduce another primary component.

An explicit example is available without a Jacobian-divisor
search. Let
\[
 k=\mathbf F_2[a]/(a^7+a+1),\qquad
 R=(a+1,\ a^5+a^4+a^2+a).
 \tag{BC.7}
\]
The finite certificate verifies the field, the curve equation,
and that this point is nonbranch; in binary coefficient
encoding its coordinates are $(3,54)$. Hence its canonical
class also has the exact order in (BC.2).
\end{proof}

\paragraph{What this does and does not decide after the quintic transfer.}
The example refutes the universal assertion that every ordinary
Rosati-$\mathcal O_5$ pointed genus-two curve has an A29
canonical detector. The example alone does not show that any of the
displayed points is the particular image of a prescribed split
point in a quintic transfer. A40 below supplies that inverse
dihedral construction with its pointing and field of definition.
A43 later supplies the stronger qualifications for the special
A13 inverse family: the sources have two-rank one and both
Weierstrass points have only bad canonical reflection images.
Those qualifications are proved there from the special Frobenius
identities, rather than inferred from real multiplication alone.

There is also no point-section counterexample here. A13 already
supplies a geometric theta witness of order $4159$ because
$\langle2,-1\rangle=(\mathbf Z/4159)^\times$, so A11
provides point separation after the indicated pointing extension.
The ordinary canonical criterion and that geometric theta
criterion have different hypotheses. Since (BC.1) is defined
over $\mathbf F_2$, its Frobenius twist is the same curve;
no identification over a nonperfect field is being assumed.


\subsection{Two actual pointed successes of the rational quintic transfer}
\label{sec:pointed-quintic-success}

\paragraph{Finding A39 (three-primary detectors on an explicit reflection quotient).}
Keep the field $k=\mathbf F_{256}$, the curve $C$, and the
degree-five cover in A36. Write $W_s=\infty$ and
$W_t=(0,0)$ in its reciprocal model. There is a reflection
quotient defined over this same field for which the image $R_s$
of its fixed lift of $W_s$ is Weierstrass, while the image $R_t$
of its fixed lift of $W_t$ has a nonzero three-primary canonical
class. The other two images of the fibre over $W_s$ also have
nonzero three-primary canonical classes. Consequently A29 can
be applied over this stated field to separate each of the two
original Weierstrass points, using the same quintic cover and
an additional cyclic cover of degree three or nine of its
reflection quotient. The resulting etale cover of $C$ has
degree fifteen or forty-five. These are statements about this
fixed example, not a universal pointing theorem.

\begin{proof}
\emph{1. An explicit pointed Kummer realization.}
Use the binary field encoding of A36 and put
\[
 \begin{aligned}
 U(X)&=X^2+80X+224,\qquad V_0(X)=59X+180,\\
 A(X)&=X^5+43X^4+39X^3+106X^2+188,\\
 B(X)&=13X^2+13X+80,\qquad F=A+BY.
 \end{aligned}
 \tag{RPT.1}
\]
The exact polynomial identity in the retained certificate is
\[
 F\,\iota(F)=A^2+ABX^2+B^2(uX^5+X^3+uX)=U^5.
 \tag{RPT.2}
\]
Moreover $F$ vanishes to order five along the degree-two divisor
$D_0$ represented by $(U,V_0)$, and its conjugate is nonzero
there. Indeed $F$ is zero modulo $(U,Y-V_0)$, while
$\gcd(U,BX^2)=1$. Its conjugate $F+BX^2$ is therefore a
unit along $D_0$, so the norm identity assigns all five zeros
there. The polynomial $U$ is separable.
Since $A$ is monic of degree five and $BY$ has smaller pole order
at $W_s$, one obtains
$\operatorname{div}(F)=5D_0-10W_s$.

Thus $D$ can be written $z^5=F$, with deck generator
$\sigma z=12z$, where $12$ has exact order five in $k$.
The class $D_0-2W_s$ has exact order five by A36, so this
Kummer cover is connected and etale. It is the pointed cover
of A36 up to a deck-generator change: $J_C(k)[5]$ has order
five, the field contains $\mu_5$, and the split-fibre
normalization fixes the constant twist.

The reflection is
\[
 \tau(X,Y,z)=(X,Y+X^2,U/z).
 \tag{RPT.3}
\]
Let $P$ be the point above $W_s$ specified by $z/X\to1$.
Let $Q$ be the point above $W_t$ with $z=92=\sqrt{224}$.
The identities $92^5=188$ and (RPT.2) show that $Q$ lies on
$D$. Both $P$ and $Q$ are $k$-rational and fixed by $\tau$.
Put $v:D\to E=D/\langle\tau\rangle$ and
\[
 R_s=v(P),\quad R_t=v(Q),\quad
 R_1=v(\sigma P),\quad R_2=v(\sigma^2P),\quad A_s=R_1+R_2.
 \tag{RPT.4}
\]
The five points above $W_s$ have images $R_s,R_1,R_1,R_2,R_2$.

\emph{2. The canonical divisor and the two tested classes.}
The lower jumps at $W_s,W_t$ are one and three. The completed
local rings used in A33 therefore give
$\operatorname{Diff}(v)=2P+4Q$.
Since $K_C\sim2W_s$ and $\pi$ is etale,
\[
             v^*K_E\sim2\sum_{j=1}^4\sigma^jP-4Q
                         =v^*\bigl(2A_s-2R_t\bigr).
 \tag{RPT.5}
\]
Pullback by $v$ is injective on geometric degree-zero divisor
classes. Indeed if $v^*D_1=\operatorname{div}(g)$, then
$\tau g=cg$ for a constant with $c^2=1$, hence $c=1$ in
characteristic two. The function descends to $E$, and $D_1$
is principal. Consequently $K_E\sim2(A_s-R_t)$.

Write $B_0=A_s-R_t$, a degree-one theta characteristic, and
$Z_R=[R-B_0]$. It is not necessary to identify $B_0$ with the
Raynaud characteristic for this canonical-class computation.
Then
\[
 \begin{aligned}
 \kappa_R&=-2Z_R,\\
 v^*Z_{R_s}&=2P+2Q-\sum_{j=1}^4\sigma^jP,\\
 v^*Z_{R_t}&=4Q-\sum_{j=1}^4\sigma^jP.
 \end{aligned}
 \tag{RPT.6}
\]

\emph{3. The finite Riemann--Roch certificate.}
Here is the complete divisor-space interpretation of the
retained finite matrix. A function in
$\mathcal L(32\pi^*W_s)$ decomposes into its five deck
characters as
\[
           f=\sum_{i=0}^4\frac{z^i}{U^i}
                              \bigl(A_i(X)+B_i(X)Y\bigr),
 \tag{RPT.7}
\]
where
\[
 \deg A_i\le16+i,\quad \deg B_i\le13+i,\quad
 A_i+B_i\bar V_i\equiv0\pmod{U^i}.
 \tag{RPT.8}
\]
The polynomial $\bar V_i$ is the unique curve-equation root
modulo $U^i$ reducing to $V_0+X^2$ modulo $U$.
The congruence removes poles on the conjugate divisor
$\iota D_0$; poles of the coefficient along $D_0$ are canceled
by the zeros of $z^i$. These are all finite places where a
denominator can occur. At infinity $X,Y$ have pole orders two
and five, and their two parities prevent cancellation between
the leading terms of $A_i$ and $B_iY$. This proves precisely
the degree bounds. Conversely multiplication by $U^i$ clears
all finite poles of each deck coefficient, so the affine
coordinate ring $k[X,Y]$ proves that (RPT.7)--(RPT.8) include
every function in this Riemann--Roch space.

For each $i$ an explicit basis consists of
\[
 \begin{gathered}
 (A_i,B_i)=(U^iX^j,0),\quad 0\le j\le16-i,\\
 (A_i,B_i)=
   (\operatorname{rem}_{U^i}(X^j\bar V_i),X^j),
                         \quad0\le j\le13+i.
 \end{gathered}
 \tag{RPT.9}
\]
There are $31$ independent functions for each character,
hence $155$ in total.

At $P$ use the uniformizer $t=Y/X^3$, put $S=1/X=t^2a(t)$,
and write $z=XZ(t)$. Their unique power series are fixed by
\[
 \begin{aligned}
 1+(t+u)a+t^4a^3+ut^8a^5&=0, & a(0)&=u^{-1},\\
 Z^5&=\sum_j A_jS^{5-j}
                   +t\sum_j B_jS^{2-j}, & Z(0)&=1.
 \end{aligned}
 \tag{RPT.10}
\]
In the second line $A_j,B_j$ are the coefficients of (RPT.1).
Both derivatives with respect to the unknown series have
nonzero constant term. At $Q$ use $t=Y$, with
\[
 t^2+X^2t=uX+X^3+uX^5,\quad X(0)=0,\quad
 z^5=A(X)+B(X)t,\quad z(0)=92.
 \tag{RPT.11}
\]
Again the derivatives are units, so these local expansions
are unique. The certificate retains them through degree
$127$ and checks all defining identities to that precision.

A function with divisor
$32(4Q-\sum_{j=1}^4\sigma^jP)$ would belong to (RPT.7), be
regular at $P$, and vanish to order at least $128$ at $Q$.
Regularity at $P$ gives the first $32$ coefficients of
$t^{32}f$ equal to zero; the condition at $Q$ gives another
$128$ coefficients. On the basis (RPT.9) these are the
retained $160$ homogeneous linear conditions in $155$
unknowns. Their matrix has rank $155$ over $\mathbf F_{256}$.
The independent certificate reconstructs the matrix from the
retained residual identities and checks rank $1240$ after
expanding each field entry into its eight-by-eight binary
multiplication matrix. Therefore no such nonzero function
exists, and $32Z_{R_t}\ne0$.

There is also a short positive certificate for $R_s$.
In the expression $\sum_i c_i z^i(A_i+B_iY)/U^i$, take
the following coefficient lists in increasing degree:
\[
\begin{array}{c|c|l|l}
 i&c_i&A_i&B_i\\\hline
 1&225&(224,80,1)&0\\
 2&1&(225,12,162,224)&1\\
 3&13&(224,189,202,174,189)&(13,1)\\
 4&1&(177,0,61,175,79,225)&(236,1,1).
\end{array}
 \tag{RPT.12}
\]
Each summand is regular off $\pi^{-1}(W_s)$ and has pole
order at most two at each point of that fibre. The expansions
(RPT.10)--(RPT.11) give a zero of order four at $P$ and a zero
of order four at $Q$, both with leading coefficient $189$.
The function is nonzero, since its first deck-character
component is $225z$. Its remaining four possible poles have
total degree at most eight, equal to the displayed zero
degree. Thus its divisor is exactly
$4P+4Q-2\sum_{j=1}^4\sigma^jP$.
By (RPT.6) and injectivity of pullback, $2Z_{R_s}=0$ and
$\kappa_{R_s}=0$. Therefore $R_s$ is Weierstrass.

\emph{4. The actual three-primary detectors.}
The rational group has order $2^{10}3^4$ by A36. Its integral
$\mathcal O_5$ action makes its two-primary part a module over
the unramified quadratic ring $\mathcal O_5\otimes\mathbf Z_2$.
At most one nonzero cyclic module summand can occur, since
each contributes two dimensions to the group killed by two
and the geometric two-rank is two. Its order $2^{10}$
therefore gives exponent $32$. Thus $32Z_{R_t}\ne0$ proves
that $Z_{R_t}$, and hence $\kappa_{R_t}$, has a nonzero
three-primary component.

For the other fibre put $Z_i=[R_i-B_0]$ and
$a_0=T(B_0)-2B_0$. The divisor relations are
$T(R_s)=2R_1$ and $T(R_1)=R_s+R_2$.
On the three-primary part $Z_{R_s}=0$, so these relations
give
\[
 a_0=2Z_1,\qquad Z_2=(T+2)Z_1,\qquad
                  Z_{R_t}=Z_1+Z_2=(T+3)Z_1.
 \tag{RPT.13}
\]
The operator $T+3$ is invertible on this part because
$(2-T)(T+3)=5$, and $T+2$ is invertible because
$(1-T)(T+2)=1$. Therefore both $\kappa_{R_1}$ and
$\kappa_{R_2}$ have nonzero three-primary components.

All these points and maps are defined over $k$. Since
$\operatorname{ord}_3(2)=2$, A29 applies to $R_t,R_1,R_2$.
The ring $\mathcal O_5\otimes\mathbf Z_3$ is also unramified
quadratic. A finite module over it of order $3^4$ has exponent
at most nine, so the cyclic detector furnished by A29 has
degree three or nine. Its pullback along the degree-two map
$v$ is connected, since these degrees are coprime, and is
etale over $D$ with the chosen rational lift split.

Finally a source Weierstrass point $W$ needs only one useful
image of a lift. A lifted collision pushed through $v$
becomes equal images by A29. The two possible points in that
quotient fibre are interchanged by $\tau$, whose images on
$C$ are interchanged by the hyperelliptic involution. Since
$W$ is fixed by that involution, both possibilities lie above
$W$. Use $Q$ and $R_t$ for $W_t$, and use $\sigma P$ and
$R_1$ for $W_s$. In both cases the original point collision
is impossible. The etale composite degree is $5\cdot3$ or
$5\cdot9$, as claimed. The composite need not be Galois
over $C$: for $m\in\{3,9\}$ its geometric normal-closure
group embeds in $C_m^5\rtimes C_5$, so has odd order and
the cover belongs to the pro-odd framework.
\end{proof}


\subsection{An inverse quintic construction and actual bad transferred images}

\paragraph{Finding A40 (inverse dihedral construction in characteristic two).}
Let $E$ be an ordinary smooth projective genus-two curve over an
algebraically closed field $k$ of characteristic two. Suppose its
canonically polarized Jacobian $J$ has an integral Rosati-symmetric
endomorphism
\[
                    T^2+T-1=0,\qquad T^\dagger=T.
 \tag{ID.1}
\]
There is a smooth connected genus-six curve $D$ with an action of
the dihedral group of order ten such that
\[
 \begin{gathered}
 D/\langle\sigma\rangle=C,\quad g(C)=2,\quad
 D\longrightarrow C\ \hbox{is etale of degree five},\\
 D/\langle\tau_1\rangle\simeq E\simeq D/\langle\tau_2\rangle.
 \end{gathered}
 \tag{ID.2}
\]
where $\sigma=\tau_1\tau_2$ has order five. Each $\tau_i$
descends to the hyperelliptic involution of $C$.

If the initial data are over a finite field, the construction and
all these maps are defined after a finite constant extension.
Every prescribed geometric point $R$ of $E$ is then the image of
a rational lift of a point $P$ on $C$ for a constant, $P$-split
cyclic degree-five cover, after one further finite extension.
No assertion that this $C$ has two-rank one, or that this $P$ is
Weierstrass, is part of the statement.

\paragraph{Correspondence background.}
Kumar--Mukamel describe the passage from a Jacobian endomorphism
to an effective correspondence and give a genus-six $(2,2)$
real-multiplication example over a number field
\cite[Sections~2--3, pp.~5--7]{KumarMukamel2016}.
The proof here establishes integrality and generic smoothness
for each fixed ordinary characteristic-two curve, and then
proves the unramified dihedral and pointed finite-field
assertions. No priority claim for the general correspondence
method is made.

\begin{proof}
\emph{1. The residual theta correspondence and its exact numerical class.}
Choose a degree-one theta characteristic $B$ on $E$, so $2B\sim K_E$,
and write
\[
             z_P=[P-B],\qquad \Theta=\{z_P:P\in E\}\subset J.
 \tag{ID.3}
\]
This is a smooth symmetric theta divisor: $z_{\iota P}=-z_P$,
where $\iota$ is the hyperelliptic involution. The construction
does not require this choice to be the Raynaud characteristic.
Equation (ID.1) makes $T$ an automorphism, with inverse $T+1$.

We first record the intersection numbers, including their
normalization. The characteristic polynomial of $T$ on a
four-dimensional rational Tate module is $(X^2+X-1)^2$.
Indeed it has rational coefficients and its only possible
eigenvalues are the two conjugate roots of the irreducible
quadratic in (ID.1), with equal multiplicities. Thus
$\operatorname{Tr}(T)=-2$, and
\[
 \operatorname{Tr}(T^\dagger T)=6,
 \qquad
 \operatorname{Tr}((T^{-1})^\dagger T^{-1})=6,
 \qquad \Theta\cdot T\Theta=3.
 \tag{ID.4}
\]
Here $\Theta^2=2$ and the last equality is the Rosati trace
formula $\Theta\cdot a^*\Theta=\operatorname{Tr}(a^\dagger a)/2$,
applied to $a=T^{-1}$; see Milne, Theorem~17.3 and its proof
\cite[pp.~35--36]{MilneAV}. In particular $T\Theta$ cannot
be a translate of $\Theta$, since their self-intersections
are two whereas their mutual intersection is three.

For $t\in E$ define
\[
 F_t:E\times E\longrightarrow J,
       \qquad F_t(P,Q)=T(z_P-z_t)-z_Q,
 \qquad V_t=\{t\}\times E.
 \tag{ID.5}
\]
The map $F_t$ is finite of degree three. Its fibres are the
intersections of $T\Theta$ with translates of $\Theta$, so
they are finite by (ID.4); their generic intersection degree
is three. It is also generically separable: the tangent
directions $dT(T_P E)$ and $T_Q E$ are generically distinct,
because the canonical Gauss map of a genus-two Abel curve
is nonconstant. Being finite between smooth surfaces, $F_t$
is flat by the regular/Cohen--Macaulay flatness criterion
\cite[Lemma~10.128.1, Tag~\texttt{00R4}]{Stacks00R4}.
Put
\[
                  D_t=F_t^*\Theta-V_t.
 \tag{ID.6}
\]
This is an effective divisor. The component $V_t$ occurs in
$F_t^*\Theta$ with multiplicity one: along its generic point
the normal derivative is the pairing of the nonzero vector
$dT(T_tE)$ with the conormal of $\Theta$ at $-z_Q$, and is
not identically zero as $Q$ varies.

There is no other vertical component. Such a component over
$P$ would give $\Theta+T(z_P-z_t)=\Theta$, whose translation
is zero because a principal theta divisor has trivial
translation stabilizer. Then $P=t$. There is no horizontal
component either, since that would make $T\Theta$ a translate
of $\Theta$. The two projection degrees of $D_t$ are both two:
the original degrees of $F_t^*\Theta$ are two over the first
factor and three over the second, and subtracting $V_t$
decreases only the second degree.

The same construction determines the sign of its Jacobian
correspondence. For $P\ne t$, its two points in the second
factor, counted with their multiplicities, form the unique
effective divisor
\[
                   Q+R\sim K_E+T([P-t]).
 \tag{ID.7}
\]
Indeed $F_t(P,Q)=z_R$ is precisely this equality, and the
degree-two line bundle has a unique effective divisor when
$T([P-t])\ne0$, by Riemann--Roch. This also identifies the
intersection multiplicities: away from the canonical class
the Abel map $E^{(2)}\to\operatorname{Pic}^2(E)$ is an
isomorphism. Therefore the induced homomorphism from the
first Jacobian to the second is $+T$, not $-T$.

Finally, finite flat pullback of intersection numbers gives
$(F_t^*\Theta)^2=3\Theta^2=6$; also
$F_t^*\Theta\cdot V_t=2$ and $V_t^2=0$.
Consequently
\[
 D_t^2=2,\qquad
 D_t\cdot K_{E\times E}=8,\qquad p_a(D_t)=6.
 \tag{ID.8}
\]
The pullback formula is Milne, Lemma~8.3
\cite[p.~13]{MilneAV}; the last equality is adjunction on
the smooth product surface.

\emph{2. The residual divisor is integral.}
Every irreducible component has positive degree over both
factors. A divisor of bidegree $(2,2)$ which is reducible or
nonreduced must therefore be the sum of two graphs of
automorphisms $\varphi,\psi$ of $E$, allowing the two graphs
to coincide. A component of bidegree $(1,1)$ is a graph
because its normalization maps with degree one to each
smooth curve $E$.

Let $u,v$ be the induced automorphisms of $J$. They preserve
the canonical polarization, so $u^\dagger=u^{-1}$ and
$v^\dagger=v^{-1}$. By (ID.7), $T=u+v$. With
$w=u^{-1}v$ it follows that
\[
              T^2=T^\dagger T=2+w+w^{-1},
       \qquad w+w^{-1}=-1-T.
 \tag{ID.9}
\]
The latter element satisfies $S^2+S-1=0$. Multiplication
by $w^2$ therefore gives
\[
                  1+w+w^2+w^3+w^4=0.
 \tag{ID.10}
\]
Hence $w$ has exact order five; it cannot be the identity
since multiplication by five is not zero on $J$.
Automorphisms of a genus-at-least-two curve act faithfully
on its Jacobian. This follows, for example, from the
uniqueness argument in Milne's Theorem~12.1(b)
\cite[pp.~38--39]{MilneJV}.
Thus $\varphi^{-1}\psi$ would be an automorphism of $E$
of order five.

This is impossible for an ordinary genus-two curve in
characteristic two. Its separable hyperelliptic map has
exactly three geometric branch points. Every automorphism
descends to the projective line and permutes these three
points; the kernel is the hyperelliptic group of order two,
and the induced group embeds in $S_3$. Its order divides
twelve. This contradiction proves that $D_t$ is integral
and reduced, for every $t$.

\emph{3. All but finitely many parameters give a smooth curve.}
Here the characteristic-two argument is essential. Let
\[
                \gamma:E\longrightarrow\mathbf P(T_0J)
 \tag{ID.11}
\]
send $P$ to the tangent line of its Abel image translated
to the origin. Pullback of invariant differentials identifies
this map with the canonical map of $E$. Thus $\gamma$ is the
separable hyperelliptic double map, and its geometric fibres
are $\{P,\iota P\}$. Its ramification points are precisely
the three Weierstrass points. Let $\overline T$ be the
projective automorphism induced by the invertible differential
$dT$.

Consider the reduced critical curve
\[
 Z=\{(P,Q):\gamma(Q)=\overline T\gamma(P)\}
             \subset E\times E,
 \qquad h(P,Q)=Tz_P-2z_Q.
 \tag{ID.12}
\]
On the normalization of every irreducible component $Z_0$,
both projections to $E$ are nonconstant and generically
separable. Indeed this is a fibre product of two finite
separable double maps, and the assertion holds off their
finite branch sets. Since $d[2]=0$ in characteristic two,
\[
                         dh=dT\,dz_P\ne0
                         \quad\hbox{generically on }Z_0.
 \tag{ID.13}
\]
In particular $h$ is nonconstant on each such component.

No image $h(Z_0)$ is contained in $T\Theta$. Otherwise
there would be a nonconstant map $s:Z_0\to E$ with
$Tz_s=h(P,Q)$. Differentiating and using (ID.13) gives
$\gamma(s)=\gamma(P)$ generically. The two points in a
canonical fibre imply $s=P$ or $s=\iota P$ on this
irreducible component. The first possibility would make
$2z_Q=0$ identically, contradicting the nonconstant second
projection. The second gives
$2(Tz_P-z_Q)=0$. The map $Tz_P-z_Q$ would then be constant,
because its image lies in the finite group scheme $J[2]$
and the source is reduced and connected. Surjectivity of
both projections would make $T\Theta$ a translate of
$\Theta$, contrary to (ID.4). Thus each curve $h(Z_0)$
has only finitely many intersections with $T\Theta$.

Suppose $(P,Q)$ with $P\ne t$ is singular on $F_t^*\Theta$,
and write $F_t(P,Q)=z_R$. The differential criterion for
this pullback of a smooth divisor is
\[
                    \gamma(R)=\gamma(Q)=\overline T\gamma(P).
 \tag{ID.14}
\]
Thus $R=Q$ or $R=\iota Q$. The latter would give
$T(z_P-z_t)=0$, hence $P=t$, which was excluded. In the
former case $(P,Q)\in Z$ and $h(P,Q)=Tz_t$.
The previous paragraph shows that this can happen for only
finitely many parameters $t$. Outside that finite set,
$D_t$ is smooth away from $V_t$.

It remains to check the intersections with $V_t$, rather
than simply discard this fixed component. Near $(t,Q)$
take a local parameter $a$ on the first factor. A local
equation of $F_t^*\Theta$ has the form $aG(a,Q)$, and $G$
defines the residual divisor. Up to a local unit,
$G(0,Q)$ is the pairing of the conormal of $\Theta$ at
$-z_Q$ with $dT(T_tE)$. Its zeros are exactly
\[
                         \gamma(Q)=\overline T\gamma(t).
 \tag{ID.15}
\]
At a non-Weierstrass $Q$ this zero is simple, since
$\gamma$ is etale there. Hence the derivative of $G$ in
the second-factor parameter is nonzero and the residual
curve is smooth at this point. Only finitely many $t$
have $\overline T\gamma(t)$ equal to a branch value of
$\gamma$ (at most six geometric points). Excluding them
as well proves that $D=D_t$ is smooth for every $t$ in a
nonempty open subset of $E$. By (ID.8), it has genus six.

\emph{4. The two involutions give the unramified dihedral cover.}
Let $v_1,v_2:D\to E$ be the projections. Both have degree
two and are separable. A purely inseparable degree-two map
between smooth curves over this perfect field would be a
relative Frobenius followed by an isomorphism and would
preserve the genus, whereas here the genera are six and two.
Let $\tau_i$ be the involution of $v_i$, and put
\[
 U_i=v_i^*:J\longrightarrow J_D,
       \qquad N_i=v_{i,*}:J_D\longrightarrow J.
 \tag{ID.16}
\]
Pullback and norm are adjoint for the canonical principal
polarizations. Thus (ID.7), Rosati symmetry, and the degree
two identities give
\[
 N_iU_i=2,\quad U_iN_i=1+\tau_i,\quad
 N_2U_1=T=N_1U_2.
 \tag{ID.17}
\]
The map $U=(U_1,U_2):J^2\to J_D$ has finite kernel:
its adjoint composite is
\[
 NU=\begin{pmatrix}2&T\\T&2\end{pmatrix},
 \qquad
 (4-T^2)=(3+T),\qquad (3+T)(2-T)=5.
 \tag{ID.18}
\]
The adjugate shows that this matrix is an isogeny, since
multiplication by five is an isogeny. Its image $A$ has
dimension four. Its polarized complementary abelian
subvariety $B_D=\ker(N_1,N_2)^0$ has dimension two, and
$J_D$ is isogenous to $A\times B_D$.

Under the isogeny $J^2\to A$, the two involutions and their
product are represented by
\[
 \tau_1=\begin{pmatrix}1&T\\0&-1\end{pmatrix},\quad
 \tau_2=\begin{pmatrix}-1&0\\T&1\end{pmatrix},\quad
 \sigma=\tau_1\tau_2=
            \begin{pmatrix}-T&T\\-T&-1\end{pmatrix}.
 \tag{ID.19}
\]
The last matrix has determinant one and trace $-1-T$.
Therefore $\sigma+\sigma^{-1}=-1-T$ and, exactly as in
(ID.10), $\Phi_5(\sigma)=0$ on $A$.
Also $\sigma^2$ has matrix
$\left(\begin{smallmatrix}0&-1\\1&T\end{smallmatrix}\right)$,
so $N_1\sigma^2U_1=T$: choosing the squared deck generator
recovers the originally specified real-multiplication
endomorphism, rather than its conjugate $-1-T$.
On $B_D$, both $\tau_i$ act as $-1$ by (ID.17), so
$\sigma$ acts as $1$. It follows that $\sigma^5=1$ on
$J_D$ and that its action is not the identity. Faithfulness
of the action of curve automorphisms on the Jacobian gives
an automorphism of $D$ of exact order five. Together with
$\tau_1$ it generates the dihedral group, since
$\tau_1\sigma\tau_1=\sigma^{-1}$.

The $\sigma$-invariant part of the rational Tate module of
$J_D$ has dimension four: it is the part belonging to
$B_D$, with none from $A$ because $\Phi_5(1)=5\ne0$.
For the quotient $C=D/\langle\sigma\rangle$, pullback
identifies its Jacobian up to isogeny with this invariant
abelian subvariety; thus $g(C)=2$.
Riemann--Hurwitz now gives
\[
                 2g(D)-2=10=5(2g(C)-2),
 \tag{ID.20}
\]
so the different is zero and the degree-five map is etale.
The quotients by the two reflections are exactly the two
projection targets $E$. On $J_C\sim B_D$, either reflection
acts as $-1$. Its quotient has genus zero, so it is the
hyperelliptic involution of $C$. This proves (ID.2).

\emph{5. Finite-field descent and an arbitrary prescribed image.}
Suppose $E$ and the original data are over a finite field.
The theta characteristic, a point $t$ in the nonempty
smoothness open, the residual divisor and all its maps have
finitely many algebraic coefficients. They are defined over
one finite constant extension $k_1$; the two degree-two
deck involutions and their product are then defined over
that same field. Given any $R\in E(\overline{k_1})$, choose
$P'\in v_1^{-1}(R)$ and a finite extension $k_2/k_1$ over
which $R$ and $P'$ are rational. Put $P=\pi(P')\in C(k_2)$.
The entire fibre of the etale map $\pi:D\to C$ above $P$
is $\{\sigma^jP':0\le j<5\}$, so it is split over $k_2$.
Its deck group is constant. This proves the pointed
assertion without identifying $R$ with the auxiliary
parameter $t$ and without prescribing the rank of $C$.
\end{proof}

\paragraph{Consequences for the canonical forcing attempt.}
Apply A40 to the ordinary integral-$\mathcal O_5$ curve of
A13 and A38. Choose one of its $126$ non-Weierstrass
$\mathbf F_{128}$-rational points, for example (BC.7).
After the finite extensions just constructed, this is an
actual image of a split lift in a quintic transfer with an
ordinary reflection quotient. Its fixed canonical class
still has exact order $4159$, and
$\operatorname{ord}_{4159}(2)=2079$ is odd. Hence even among
actual prescribed images of quintic transfers from genus-two
curves, the A29 canonical-prime hypothesis need not hold.

This does not assert that both reflection images associated
to one lift fail, that every possible quintic cover fails,
or that the rank-one-source/Weierstrass-point subcase has
been refuted. It closes the broad image-membership gap left
in A38. The geometric order-$4159$ theta witness of A13
still provides point separation; no section collision is
constructed here. The proof is unconditional and does not
use a prime-density assertion.

The next finding A41 supplies the stronger simultaneous failure at all
five reflection images for a fixed-field member of this inverse family.
It still leaves a suitable choice of a different cover open.


\subsection{All reflection images can fail the canonical-prime condition}

\paragraph{Finding A41 (a fixed-field simultaneous failure of the sufficient detector).}
Let $E/\mathbf F_{128}$ be the ordinary curve of A13 and A38.
There is a smooth genus-two curve $C/\mathbf F_{128}$ and a
constant etale cyclic cover $\pi:D\to C$ of degree five,
with ordinary reflection quotient $E$, together with a split
rational point $P\in C(\mathbf F_{128})$, having the following
property. For every rational lift of $P$ to $D$ and every
one of the five dihedral reflection quotients, its image
$R$ has
\[
                  \operatorname{ord}([K_E-2R])=4159,
             \qquad \operatorname{ord}_{4159}(2)=2079
                         \quad\hbox{odd}.
 \tag{AI.1}
\]
Here the quotient is identified with $E$ by the maps of A40;
the order assertion is independent of that identification.
Thus no choice of lift or reflection within this same
degree-five cover supplies the sufficient canonical-prime
hypothesis of A29. The assertion persists for these fixed
images after every extension of constants.

The proof supplies at least $95$ rational parameters for
the A40 inverse family over this field. Each resulting
cover has at least eleven distinct rational split source
points with the property just stated. These are existence
and counting statements in an explicit algebraic family;
no selected parameter or covering equation is claimed to
have been individually enumerated.

\begin{proof}
\emph{1. At most thirty-four parameters are excluded in A40.}
Use the notation of A40 over an algebraically closed field.
Let $F_1=\{P_0\}\times E$ and $F_2=E\times\{Q_0\}$ be
the two fibre classes on $E\times E$. The critical curve
\[
 Z=\{(P,Q):\gamma(Q)=\overline T\gamma(P)\},
         \qquad h(P,Q)=Tz_P-2z_Q
 \tag{AI.2}
\]
is the pullback of the diagonal in $\mathbf P^1\times\mathbf P^1$.
Both canonical maps have degree two, so its numerical
divisor class is
\[
                             [Z]=2F_1+2F_2.
 \tag{AI.3}
\]
This is also the class of its reduced curve used in A40.
Indeed its projections are generically separable, so the
Cartier divisor is generically reduced; a Cartier divisor
on this smooth surface is Cohen--Macaulay and has no
embedded associated components. It is therefore reduced.

Put $H=h^*\mathcal O_J(T\Theta)$. Its restrictions to the
two fibres have degrees
\[
                         \deg(H|_{F_1})=12,
                \qquad \deg(H|_{F_2})=2.
 \tag{AI.4}
\]
For the first equality, varying $Q$ gives a translate of
$-[2]\Theta$, so pullback by $[2]$ multiplies the divisor
class by four \cite[Corollary~6.6 and Remark~12.13,
pp.~9 and~24]{MilneAV}; (ID.4) then gives
$4(\Theta\cdot T\Theta)=12$.
For the second, varying $P$ gives a translate of $T\Theta$,
whose self-intersection is two. Translation does not alter
these degrees. Hence
\[
                             Z\cdot H=28.
 \tag{AI.5}
\]

Step~3 of A40 proves that no component of $Z$ maps into
$T\Theta$ under $h$. The pulled-back effective divisor
therefore meets $Z$ properly, and the sum of its local
intersection lengths is exactly $28$. In particular there
are at most $28$ geometric intersection points. Each
off-vertical singularity of the residual curve $D_t$ gives
one of these points with $h(P,Q)=Tz_t$, and $Tz_t$
determines $t$ uniquely. Thus at most $28$ parameters can
produce such a singularity.

The other excluded parameters in A40 satisfy
$\overline T\gamma(t)=\gamma(W)$ for a Weierstrass point
$W$. There are three branch values, each with at most two
preimages under the degree-two map $\overline T\gamma$.
This excludes at most six additional parameters. Consequently
\[
        \#\{t\in E(\bar k):D_t\hbox{ is not smooth}\}\le34.
 \tag{AI.6}
\]
The integrality of every residual curve was already proved
in A40, so no additional reducibility exclusions are needed.

\emph{2. The inverse covers already exist over $\mathbf F_{128}$.}
For the A13 curve, choose $B=\mathcal O_E(W)$ at one of
its three $\mathbf F_2$-rational Weierstrass points. This
is a theta characteristic and makes the symmetric theta
embedding in A40 defined over $\mathbf F_2$.
The integral Rosati endomorphism $T$ is also defined over
$\mathbf F_2$, by the checked supplement to A33.
Thus the entire family $D_t$, its two projection maps,
and its dihedral data are defined over the field of $t$.
Finding A38 gives $\#E(\mathbf F_{128})=129$.
By (AI.6), at least
\[
                              129-34=95
 \tag{AI.7}
\]
of these rational parameters give smooth genus-six curves
$D_t$ over $\mathbf F_{128}$. For any such $t$, A40 gives
a genus-two quotient $C_t$ and a constant etale cyclic
degree-five map over this same field. No further extension
is required to define its deck group or reflections.

The pullback, norm, and complementary-factor construction
of A40 are all defined over $\mathbf F_{128}$. They give
the isogeny over that field
\[
                      J_{D_t}\sim J_E^2\times J_{C_t}.
 \tag{AI.8}
\]
The Frobenius trace of $E$ over this field is zero, again
by A38. Taking traces in (AI.8) and applying the curve
point-count formula therefore gives
\[
 \begin{gathered}
 \#D_t(\mathbf F_{128})=\#C_t(\mathbf F_{128})
       \ge129-4\sqrt{128}>83,\\
                    \#D_t(\mathbf F_{128})\ge84.
 \end{gathered}
 \tag{AI.9}
\]
Here the inequality is the genus-two Weil bound
\cite[Theorem~11.1]{MilneJV}. It applies to $C_t$ and
does not require an explicit equation or a point-count
computation for that individual curve.

\emph{3. All five images can be simultaneously non-Weierstrass.}
Write $v:D_t\to E$ for its first projection and let
$\sigma$ generate its constant degree-five deck group.
The five maps
\[
                         v\sigma^j:D_t\longrightarrow E,
                              \qquad 0\le j<5,
 \tag{AI.10}
\]
are the five reflection quotient maps, using the
$\mathbf F_{128}$-defined identifications with $E$.
Each has degree two. The union of the inverse images of
the three Weierstrass points under all five maps has at
most
\[
                               5\cdot3\cdot2=30
 \tag{AI.11}
\]
geometric points, and hence at most that many rational
points. By (AI.9), at least $54$ rational points $P'$ of
$D_t$ have all five images non-Weierstrass. For each of
them A38 proves (AI.1) for every image.

This good set is invariant under $\sigma$, since applying
$\sigma$ merely permutes the five maps. The degree-five
action is free. Thus its size, already at least $54$, is
a multiple of five and is at least $55$. It supplies at
least eleven distinct points $P=\pi(P')$ of $C_t$.
For each, the entire fibre $\{\sigma^jP'\}$ is rational,
so $P$ is split. Changing the lift permutes the same five
images. The five reflections are the conjugates
$\sigma^{-j}\tau\sigma^j$, and thus (AI.10) covers every
reflection choice. A different isomorphism of a quotient
with $E$ preserves the order of its canonical divisor
class. Finally, constant extension preserves the exact
order of each of these fixed classes. This proves every
assertion.
\end{proof}

\paragraph{The remaining choice of cover.}
A41 rules out forcing the A29 canonical-prime hypothesis
by changing only the lift or reflection of a fixed
quintic transfer. It leaves open the possibility of
choosing a different cyclic cover for the prescribed
source point, or using a geometric theta witness or a
different rational-Prym criterion. A43 below determines the
two-rank of these sources to be one and proves that both
Weierstrass points are rational and split with the same failure
at every reflection image. The order-$4159$
geometric theta witness on $E$ from A13 is compatible
with all these canonical failures, so no failure of
point-section separation follows.


\subsection{A fixed Lang cover counts prescribed primary components}
\label{sec:fixed-lang-density}

\paragraph{Finding A42 (an explicit finite-primary equidistribution refinement).}
Let $X/k$ be a smooth projective geometrically connected curve of
genus $g\ge1$ over $k=\mathbf F_q$, of characteristic $p$.
Let $B$ be a degree-one line bundle on $X$ defined over $k$, and put
\[
 J=\operatorname{Pic}^0(X),\qquad
 i_B(R)=\mathcal O_X(R)\otimes B^{-1}.
 \tag{PE.1}
\]
No rational effective representative of $B$, or rational point of
$X$, is assumed. Fix a finite nonempty set $S$ of rational primes,
which may contain $p$. For a torsion abelian group, a subscript $S$
denotes the sum of its primary components at primes in $S$. Put
\[
 G_S=J(k)_S,\qquad d=\#G_S,\qquad b=d(g-1)+1.
 \tag{PE.2}
\]
Let $\Phi$ be $q$-power Frobenius. For each $\ell\in S$, let
$h_\ell$ be the prime-to-$\ell$ part of the finite order of its
action on the \emph{geometric point group}
$J[\ell](\overline k)$; set $h_\ell=1$ if that group is zero.
Define
\[
 M=\operatorname{lcm}_{\ell\in S}(\ell h_\ell),\qquad
 \mathcal I=\{n\ge1:\gcd(n,M)=1\}.
 \tag{PE.3}
\]
For every $n\in\mathcal I$ there is equality of actual subgroups
of $J(\overline k)$,
\[
                    J(\mathbf F_{q^n})_S=G_S.
 \tag{PE.4}
\]
For each fixed $c\in G_S$ there is a smooth projective
geometrically connected curve $Y_c/k$, independent of $n$, and a
constant etale cover $Y_c\to X$ of degree $d$, such that
\[
 \begin{aligned}
 N_{c,n}&:=\#\{R\in X(\mathbf F_{q^n}):(i_B(R))_S=c\},\\
 dN_{c,n}&=\#Y_c(\mathbf F_{q^n}),\qquad n\in\mathcal I,\\
 \left|N_{c,n}-\frac{q^n+1}{d}\right|
             &\le\frac{2b}{d}q^{n/2}.
 \end{aligned}
 \tag{PE.5}
\]
In particular, along $n\in\mathcal I$ tending to infinity,
\[
                \frac{N_{c,n}}{\#X(\mathbf F_{q^n})}
                           \longrightarrow\frac1d.
 \tag{PE.6}
\]
Every prescribed component $c$ therefore occurs on infinitely
many points of the Abel curve, hence on a Zariski-dense subset.
The progression $n\equiv1\pmod M$ supplies arbitrarily large
admissible extension degrees without a theorem about primes.

\paragraph{Source credit and scope.}
The connected Lang quotient is the construction of Poonen,
\emph{Multiples of subvarieties in algebraic groups over finite
fields}, Lemma~5.1 (author PDF, p.~6). His proof of Theorem~1.4 on
the same page already controls primary subgroups by restricting
extension degrees; Lemmas~6.1--6.2, p.~7, give related finite-primary
extension bounds \cite{Poonen2005}. Milne's
\emph{Jacobian Varieties}, Proposition~9.1, p.~30, includes the
characteristic-primary etale part; its proof is completed on
pp.~31--32 \cite{MilneJV}. A12 already used the same connectedness
argument. We give the full fixed-cover, norm-scaling, and explicit
count refinement here, without a priority claim for the underlying
Lang or torsion-density argument.

\begin{proof}
\emph{1. Extension degrees that preserve every selected primary subgroup.}
For $n\ge1$ write
\[
 V_n(T)=1+T+\cdots+T^{n-1},\qquad
              \Phi^n-1=(\Phi-1)V_n(\Phi).
 \tag{PE.7}
\]
Fix $\ell\in S$ and $n\in\mathcal I$. On the finite
$\mathbf F_\ell$-vector space $J[\ell](\overline k)$,
$\Phi$ is invertible. Every eigenvalue $\lambda$ over
$\overline{\mathbf F}_\ell$ has order dividing $h_\ell$.
If $\lambda\ne1$, then $\lambda^n\ne1$ because
$\gcd(n,h_\ell)=1$, and
\[
 V_n(\lambda)=\frac{\lambda^n-1}{\lambda-1}\ne0.
\]
For $\lambda=1$ the value is $n\ne0$ in $\mathbf F_\ell$.
Thus $V_n(\Phi)$ is invertible on $J[\ell](\overline k)$,
also when $\Phi$ is not semisimple modulo $\ell$.

It is then injective on every finite group
$J[\ell^a](\overline k)$: a nonzero kernel would contain an
element of order $\ell$, contradicting the preceding conclusion.
It is therefore an automorphism on each such group. If
$x\in J(\overline k)$ is $\ell$-primary and $\Phi^n x=x$,
(PE.7) gives $V_n(\Phi)(\Phi-1)x=0$, hence $\Phi x=x$.
The converse is immediate. This proves (PE.4), prime by prime.
Only geometric torsion points are used; no assertion that the
group scheme $J[p]$ is etale is made.

This also proves the expected integral statement on the etale
$\ell$-adic Tate module: the same polynomial is an automorphism
because its reduction modulo $\ell$ is invertible. The finite-group
argument above suffices and avoids a choice of Tate-module basis.
If desired, one may replace $M$ by the explicit larger integer
\[
        M_0=\prod_{\ell\in S}
                   \left(\ell\prod_{j=1}^{2g}(\ell^j-1)\right).
 \tag{PE.8}
\]
Indeed the prime-to-$\ell$ part of the order of any element of
$\operatorname{GL}_r(\mathbf F_\ell)$, with $r\le2g$, divides
$\prod_{j=1}^{2g}(\ell^j-1)$. The characteristic-primary geometric
point group has rank at most $g$, so it satisfies this bound too.

\emph{2. One Lang quotient over the fixed ground field.}
Write $K=J(k)$ and let $H=K_{S^c}$ be its complementary primary
subgroup. The homomorphism $\Phi-1:J\to J$ is an etale isogeny:
its differential is $-1$ and its geometric kernel is the finite
constant group $K$. Quotienting by the constant subgroup $H$
gives
\[
 \begin{gathered}
 J\xrightarrow{\pi_H} A=J/H\xrightarrow{\psi}J,\\
       \psi\pi_H=\Phi-1,\qquad\ker\psi=K/H\simeq G_S.
 \end{gathered}
 \tag{PE.9}
\]
All maps and the indicated group identification are over $k$.
For $c\in G_S\subset J(k)$ form
\[
       Y_c=X\times_{i_B-c,\,J,\,\psi} A.
 \tag{PE.10}
\]
It is a constant etale $G_S$-torsor over $X$ of degree $d$.

\emph{3. Geometric connectedness, with no exclusion of $p$.}
Here is the universal-property proof, also showing that the
translation in (PE.10) causes no difficulty. Work over
$\overline k$ and choose $P_0\in X(\overline k)$.
If $Y_c$ were disconnected, its geometric monodromy would be a
proper subgroup $A_0\subset G_S$. Quotienting the isogeny by
$A_0$ would give a nontrivial intermediate etale isogeny
$\psi':A'\to J$ such that $i_B-c$ lifts to a morphism
$v:X\to A'$. The trivial monodromy of the quotient torsor is
exactly what supplies this lift. Subtracting $v(P_0)$ gives a
pointed lift of the usual Abel map $i_{P_0}$, since
\[
 \psi'\bigl(v(R)-v(P_0)\bigr)=i_B(R)-i_B(P_0)=[R-P_0].
\]
The Jacobian universal property supplies a homomorphism
$u:J\to A'$ with $\psi'u=\operatorname{id}_J$
\cite[Proposition~6.1, pp.~19--20]{MilneJV}.
The homomorphism $u$ is an isogeny, since it has finite kernel
and its image has the full dimension. Multiplicativity of degree
contradicts $\deg\psi'>1$. Thus $Y_c$ is geometrically connected.
It is smooth and projective, and unramified Riemann--Hurwitz gives
\[
                              g(Y_c)=d(g-1)+1=b.
 \tag{PE.11}
\]
This argument applies unchanged when $p$ divides $d$.

\emph{4. The exact norm and component detected by a split fibre.}
Let $k_n=\mathbf F_{q^n}$ and $R\in X(k_n)$, and put
$x=i_B(R)$. Choose $y\in J(\overline k)$ such that
$(\Phi-1)y=x-c$. The image $z=\pi_H(y)$ lies over $R$ in
(PE.10). Define
\[
 N_n=1+\Phi+\cdots+\Phi^{n-1}:J(k_n)\longrightarrow J(k).
\]
The target assertion follows from
$(\Phi-1)N_nx=(\Phi^n-1)x=0$. Direct telescoping gives
\[
 \begin{aligned}
 \Phi^n z-z
   &=\pi_H\bigl((\Phi^n-1)y\bigr)
     =\pi_H\bigl(N_nx-nc\bigr),\\
 \operatorname{pr}_{G_S}(N_nx-nc)
   &=n(x_S-c),\qquad n\in\mathcal I.
 \end{aligned}
 \tag{PE.12}
\]
For the second equality, (PE.4) places $x_S$ in $G_S$, on which
$\Phi$ acts trivially. The complementary component of $x$ has
order prime to every prime in $S$, so its image under $N_n$ has
the same property and projects to zero in $G_S$.
Multiplication by $n$ is an automorphism of $G_S$ because
$n$ is prime to every prime in $S$.
Consequently the fibre is split over $k_n$ if and only if
$x_S=c$. Since the deck group is constant, such a fibre has
exactly $d$ rational points, and every other fibre has none.
This proves the exact equality in (PE.5).

The first line of (PE.12) holds for every $n$; the second and the
split-fibre equivalence use the specified extension condition.
Thus no compatibility between different base-field Lang maps is
assumed. A single quotient (PE.9) and a fixed cover $Y_c/k$ are
being counted throughout. In particular, $n$ need not be one
modulo the exponent of $G_S$.

\emph{5. Counting, density, and descent.}
The Weil bound for the geometrically connected curve $Y_c$ gives
\[
       \bigl|\#Y_c(k_n)-(q^n+1)\bigr|\le2bq^{n/2}
\]
\cite[Theorem~11.1]{MilneJV}. Dividing by $d$ proves (PE.5).
Applying the same bound to $X$ proves (PE.6).
The numerator in the lower bound tends to infinity along
$n\equiv1\pmod M$, so the union of the counted point sets is
infinite. Since $i_B$ is an Abel embedding, it is Zariski dense
in its curve image. For a finite search, any admissible $n$ with
\[
                  q^n+1-2bq^{n/2}>d
 \tag{PE.13}
\]
gives at least two points in the desired primary coset. For
example $q^n>(2b+d)^2$ is sufficient. There is at most one point
with $i_B(R)=0$, so (PE.13) also guarantees a nonzero torsion
class when $c=0$.

Every object used in the construction is defined over the
specified finite field. If one first enlarges that field to
rationalize a desired finite torsion component or other input,
the theorem applies with that new $q$, $G_S$, and $M$. There is
no subsequent need to change the ground field or the Lang cover
as $n$ varies.
\end{proof}

\paragraph{Pure odd theta torsion and dense canonical-good points.}
Now let $C/\mathbf F_q$ be any smooth projective geometrically
connected genus-two curve in characteristic two. Set
$X=C^{(1)}$ and let $B$ be its degree-one locally exact
differential line bundle, so that $B^{\otimes2}\simeq\omega_X$.
The actual Raynaud theta support is
\[
                   \Theta=i_B(X)\subset J_X.
 \tag{PE.14}
\]
This line bundle and embedding are defined over $\mathbf F_q$;
no rational Weierstrass point is required. Taking $S=\{2\}$ and
$c=0$ in (PE.5) proves that theta points of odd order have
limiting proportion $1/\#J_X(\mathbf F_q)[2^\infty]$ along the
specified extension degrees. There are infinitely many nonzero
such points, for every $2$-rank and without a simplicity
assumption on the Jacobian. If the geometric $2$-torsion is
already rational, one may take every odd $n$; otherwise the
additional condition in (PE.3) excludes the finite set of
relevant unit-order prime divisors.

More generally, fix an odd prime $r$ with
$\operatorname{ord}_r(2)$ even; one may always choose $r=3$.
After a finite extension $k_0$ choose a nonzero
$c_r\in J_X[r](k_0)$. Apply the theorem over $k_0$ with
$S=\{2,r\}$ and $c=c_r$. It gives a Zariski-dense set of
points $R\in X(\overline k)$, with the explicit positive
proportion (PE.6) along the admissible extensions of $k_0$, for
which
\[
 \begin{aligned}
 (i_B(R))_2&=0,\qquad (i_B(R))_r=c_r\ne0,\\
 \kappa_R&=[K_X-2R]=-2i_B(R),\qquad
                  (\kappa_R)_r=-2c_r\ne0.
 \end{aligned}
 \tag{PE.15}
\]
At each such rational point, over the finite field in which it
was counted, A29 supplies an allowable pointed cyclic
$r$-power cover and separates its point section from those of
the other rational points. Frobenius twisting by the automorphism
of the finite ground field transports this statement, including
its rational points and pointed covers, to $C$.

\paragraph{What this does and does not settle.}
The characteristic-primary part of (PE.9) is an auxiliary
counting cover. It is \emph{not} a pro-prime-to-characteristic
neighbourhood and is never used as one. The allowable cover in
the last paragraph is the separate odd-degree construction of
A29. Pure odd theta torsion is now available quantitatively on
each fixed finite-field genus-two curve, and useful canonical
points are dense on each target. This does not force any image
of an arbitrarily prescribed source point, among the finitely
many images in a chosen dihedral cover, to be one of them;
A41 exhibits a fixed cover where all those images fail.
Nor does (PE.5) force the \emph{whole order} of a theta point
to be prime, a prime power, or an odd integer $m$ with
$\langle2,-1\rangle=(\mathbf Z/m)^\times$. Its unprescribed
primary components can contain additional primes. Consequently
the universal prescribed-cover choice and primitive-orbit
selection gates remain open; no previously open row is marked
closed by this density refinement.


\subsection{Rank-one sources with two bad Weierstrass transfer fibres}

\paragraph{Finding A43 (the rank-one and Weierstrass qualifications do not repair the fixed-cover detector).}
Let $E/\mathbf F_2$ be the ordinary curve of A13 and A38, and
base change it to $k=\mathbf F_{128}$. In the inverse family
of A40 there are at least $95$ parameters $t\in E(k)$ for
which the smooth cyclic degree-five cover
$\pi_t:D_t\to C_t$ has the following properties:
\[
 \begin{gathered}
       g(C_t)=2,\qquad f(C_t)=1,\\
 W_t^{(1)},W_t^{(2)}\in C_t(k)\text{ are both Weierstrass points,}\\
 \text{both fibres are split in }D_t,\\
 \operatorname{ord}([K_E-2R])=4159,\\
 \text{for every lift of either point and every reflection image }R.
 \end{gathered}
 \tag{RWI.1}
\]
Since $\operatorname{ord}_{4159}(2)=2079$ is odd, none of
these images satisfies the sufficient hypothesis of A29.
Their exact canonical orders persist after every extension
of constants. Thus changing the lift or reflection of this
fixed cover does not force the canonical detector even for
a rank-one source and either prescribed Weierstrass point.
The assertion does not quantify over other cyclic covers
of that source point.

\begin{proof}
\emph{1. The special Frobenius and automorphism identities.}
Write $F$ for the geometric-point map obtained by squaring
the coordinates on $E$, and write $\pi$ for the corresponding
Frobenius endomorphism of $J=J_E$. Put $V=\pi^\dagger$.
The checked A33 supplement supplies the integral,
Rosati-symmetric endomorphism $T$ with
\[
 \pi V=2,\qquad \pi+V=2T+1,\qquad T^2+T=1.
 \tag{RWI.2}
\]
All these endomorphisms are defined over $\mathbf F_2$ and
commute. The element
\[
 u=T-\pi
 \quad\text{satisfies}\quad
 u^\dagger u=1,\qquad u^2+u+1=0.
 \tag{RWI.3}
\]
Indeed $u^\dagger=T-V$, and both identities follow by
substitution in (RWI.2). It is an integral polarized
automorphism of exact order three. The genus-two Torelli
uniqueness statement and descent give an automorphism
$\rho:E\to E$ over $\mathbf F_2$ inducing $u$; its order
is three by faithfulness of the curve action on its
Jacobian \cite[Theorem~12.1 and Corollary~12.2,
pp.~38--39]{MilneJV}.

Use the actual Raynaud theta characteristic $B$ on $E$;
the Frobenius twist is identified with $E$ over
$\mathbf F_2$. It is defined over $\mathbf F_2$, is preserved
by $\rho$, and is non-effective because $E$ is ordinary.
Put $z_P=[P-B]$ and $\Theta=\{z_P:P\in E\}$.
Thus $0\notin\Theta$, $z_{\iota P}=-z_P$,
$\pi z_P=z_{F(P)}$, and $u z_P=z_{\rho(P)}$.
The three Weierstrass classes are exactly the three
nonzero geometric points of $J[2]$ on $\Theta$, as in A35.
They are all $\mathbf F_2$-rational. Since $T$ is an
automorphism, the intersection number from (ID.4) implies
\[
             \Theta\cap T\Theta
                =J[2](\bar k)\setminus\{0\}
             \quad\text{as a set of three points}.
 \tag{RWI.4}
\]
Each already belongs to both divisors, whose intersection
number is three.

Changing the theta characteristic does not change the
residual family. Intrinsically, for $P\ne t$ its fibre is
$Q+R\sim K_E+T([P-t])$, as in (ID.7); equivalently, if
two theta characteristics differ by $a\in J[2]$, the two
maps and theta divisors in (ID.5)--(ID.6) are translated
by the same $a$. Hence the family used here is the one
counted in A41.

Differentiating (RWI.2) gives $dV=1$, since $d\pi=0$
and $d[2]=0$. In particular $V$ is etale. Its degree is
four and its kernel is all four geometric points of
$J[2]$: $\pi$ acts identically on those points and
$V\pi=[2]$. Also $d\rho=du=dT$.
Let $\gamma:E\to\mathbf P(T_0J)$ be the canonical Gauss
map. In an $\mathbf F_2$ coordinate its three branch
values are exactly $\mathbf P^1(\mathbf F_2)$, because
the branch values of the displayed A13 model are
$0,1,\infty$. The projective action of $dT=d\rho$
permutes those three values.

\emph{2. The normalized double-point locus.}
The image
\[
                         \Gamma=V\Theta\subset J
 \tag{RWI.5}
\]
has normalization $E$. Indeed $V$ is etale, and the only
identifications on $\Theta$ by its nonzero kernel points
are pairs of Weierstrass classes, by the theta-intersection
calculation in A35. All three Weierstrass classes map to
zero. Their tangent lines are distinct and $dV=1$, so
$\Gamma$ has three smooth transverse branches at zero
and is smooth away from zero. Numerically $\Gamma\equiv2\Theta$:
$\pi_*\Theta=2\Theta$ and $[2]_*\Theta=4\Theta$, while
$[2]=V\pi$ and $V|_\Theta$ is birational. The pushforward
identity follows from the multiplication pullback and projection
formulas \cite[Corollary~6.6 and Lemma~8.3]{MilneAV}.
Thus $\Gamma^2=8$, whereas $(T\Theta)^2=2$, so $\Gamma$
cannot coincide with a translate of $T\Theta$.

The condition that a degree-two effective divisor is a
double point is the image of $Q\mapsto 2z_Q$.
Because $[2]=V\pi$ and $F:E\to E$ is a radicial degree-two
map, its reduced image is precisely $\Gamma$.
For non-Weierstrass $t$ put
$H_t=T(\Theta-z_t)$. The total intersection number is six.
For example the Rosati trace formula gives
\[
 \begin{aligned}
 H_t\cdot\Gamma
 &=\Theta\cdot V^*(T\Theta)\\
 &=\tfrac12\operatorname{Tr}
       \bigl((T^{-1}V)^\dagger(T^{-1}V)\bigr)\\
 &=\tfrac12\operatorname{Tr}\bigl(2(T+2)\bigr)=6.
 \end{aligned}
 \tag{RWI.6}
\]
Here $T^{-2}=T+2$ and the rational Tate trace of $T$ is
$-2$. At zero the curve $H_t$ has tangent
$dT\gamma(t)$, which is not one of the three branch
lines of $\Gamma$, since $t$ is not Weierstrass.
Thus zero contributes exactly three to (RWI.6).

We determine the remaining intersection multiplicities.
On $E\times E$ use the normalized map
\[
                    \Psi(P,Q)=Tz_P-Vz_Q.
 \tag{RWI.7}
\]
Away from zero, a tangency of $H_t$ and $\Gamma$ is
characterized by $\gamma(Q)=dT\gamma(P)$, since
$dV=1$. Its two critical graphs are
$Q=\rho(P)$ and $Q=\iota\rho(P)$.
The identities in (RWI.2)--(RWI.3) give
\[
 \begin{aligned}
 T-Vu&=-u^2\pi^2,\\
 T+Vu&=-T\pi.
 \end{aligned}
 \tag{RWI.8}
\]
On the first graph, $\Psi(P,Q)=-u^2\pi^2z_P$ lies in
$\Theta$. It cannot equal $Tz_t$ for non-Weierstrass $t$,
by (RWI.4). On the second graph it equals $-T\pi z_P$.
Consequently there is exactly one possible nonzero
point of tangency, specified by
\[
                   t=\iota F(P),\qquad Q=\iota\rho(P).
 \tag{RWI.9}
\]
These points are non-Weierstrass and hence $Vz_Q\ne0$.

The tangency has multiplicity exactly two, not three.
At the point in (RWI.9), the differential of $\Psi$ has
one-dimensional image $dT\gamma(P)$. The tangent of
$T\Theta$ at $Tz_t$ is $dT\gamma(F(P))$.
These lines differ: equality would make $\gamma(P)$
fixed by Frobenius on $\mathbf P^1$, hence a branch
value and $P$ Weierstrass. The inverse image
$\Psi^{-1}(T\Theta)$ is therefore a smooth reduced
curve locally at this point. It contains the smooth
second critical graph, so locally equals that graph.
Its map to $T\Theta$ is, by (RWI.9), the radicial
Frobenius map of degree two. Thus its local fibre has
length exactly two. Since $V\Theta$ is smooth here
and its normalization is an isomorphism locally, this
is exactly the local intersection length of $H_t$
with $\Gamma$.

After the origin contribution three and this contribution
two, (RWI.6) leaves exactly one further simple nonzero
intersection. For a smooth $D_t$, these two nonzero
intersections are exactly the ramification points of
$v_1:D_t\to E$. Off $P=t$, this follows directly from
the unique effective pair in (ID.7): its two points
coincide precisely on the double-point locus. Over
$P=t$, the residual fibre is characterized by
$\gamma(Q)=dT\gamma(t)$; it consists of two distinct
points because this value is not a branch value.
There is therefore no further ramification there.

\emph{3. The source has two-rank one.}
The degree-two Galois cover $v_1$ has exactly two
geometric ramification points. Deuring--Shafarevich gives
\[
                 f(D_t)-1=2(f(E)-1)+2=4,
                      \qquad f(D_t)=5
 \tag{RWI.10}
\]
\cite[Corollary~1.8, pp.~35--36]{Crew1984}.
The isogeny $J_{D_t}\sim J_E^2\times J_{C_t}$ from
A40 and invariance and additivity of two-rank under
isogeny give $5=4+f(C_t)$, so $f(C_t)=1$.
The excluded vertical-branch parameter set in A41
contains the three Weierstrass parameters, since
$dT$ permutes their Gauss values. Its entire excluded
set still has at most $34$ points. Thus at least $95$
parameters in $E(\mathbf F_{128})$ are both smooth and
non-Weierstrass, and all have this rank-one source.

\emph{4. A reflection-fixed lift with all five images bad.}
Fix such a parameter $t$, and take
$P=\iota F^{-1}(t)\in E(k)$. Set
\[
 \begin{gathered}
 x=z_P,\qquad z_t=-\pi x,\qquad
 y=-u\pi^{-1}x=z_{\iota\rho F^{-1}(P)},\\
 x'=u^2\pi^{-1}x=z_{\rho^2F^{-1}(P)}.
 \end{gathered}
 \tag{RWI.11}
\]
The inverse $\pi^{-1}$ here means inverse Frobenius
on geometric points; on $k$-points it is $\pi^6$.
All three displayed classes are points of $\Theta(k)$
and correspond to non-Weierstrass points. They are
pairwise distinct. For example $x=y$ would imply
$Tx=0$, $x=x'$ would imply $(T+1)x=0$, and $x'=y$
would imply $x=0$. These are impossible because $T$
and $T+1$ are automorphisms and $0\notin\Theta$.

None of their points equals the auxiliary parameter
$t$. For $x=z_t$, (RWI.9) and $Vz_Q\ne0$ already
exclude the equality. If $x'=z_t$, then
$\pi^2x=-u^2x$; applying the seventh power and using
$\pi^7x=x$ gives $x=-u^2x$, hence $ux=0$.
If $y=z_t$, then $\pi^2x=ux$; the seventh power
implies $ux=x$ and then $\pi^2x=x$. The corresponding
point of $E$ is in $E(\mathbf F_4)$, whose only points
are the three Weierstrass points by A13. Again this
contradicts the choice of $P$.

For two theta classes $a,b$, abbreviate
$F_t(a,b)=Ta-b-Tz_t$. Direct use of (RWI.2)--(RWI.3)
gives
\[
 \begin{gathered}
 F_t(x,y)=y,\qquad F_t(x',y)=x',\qquad
 F_t(y,x')=x. 
 \end{gathered}
 \tag{RWI.12}
\]
All first coordinates are different from $t$, so the
three pairs in (RWI.12) belong to the residual curve,
not just to its removed vertical component. Let
$a_0=(x,y)\in D_t(k)$, identifying theta classes with
their unique points of $E$. The first equality makes
$a_0$ fixed by $\tau_1$, the involution of $v_1$.
The other two, and uniqueness of the other point in
each degree-two fibre, give the five distinct points
\[
 \begin{gathered}
 a_0=(x,y),\quad \sigma a_0=(x',x'),\quad
 \sigma^2a_0=(y,x),\\
 \sigma^3a_0=(y,x'),\quad\sigma^4a_0=(x',y),
             \qquad \sigma=\tau_1\tau_2.
 \end{gathered}
 \tag{RWI.13}
\]
For instance $\tau_2(x,y)=(x',y)$ and
$\tau_1(x',y)=(x',x')$; the next step uses $F_t(y,x')=x$.
The last two orbit points also follow directly from
$\tau_1\sigma^j a_0=\sigma^{-j}a_0$, since
$\tau_1a_0=a_0$, by swapping the first-projection fibre.

Put $W_t^{(1)}=\pi_t(a_0)$. It is rational and its entire
cyclic degree-five fibre is the rational orbit
(RWI.13), so it is split. Since $\tau_1a_0=a_0$
and $\tau_1$ descends to the hyperelliptic involution
of $C_t$, the point $W_t^{(1)}$ is Weierstrass.
The five reflection images are the first coordinates
of (RWI.13), namely $x,x',y,y,x'$.

\emph{5. The other Weierstrass fibre has the same property.}
Put $a=z_t$ and define three further theta points
\[
       x_2=-\pi a,\qquad y_2=u\pi^{-1}a,
                       \qquad x'_2=u^2\pi^{-1}a.
 \tag{RWI.14}
\]
They are respectively the classes of
$\iota F(t)$, $\rho F^{-1}(t)$ and
$\rho^2F^{-1}(t)$, so all are non-Weierstrass and
$k$-rational. They are pairwise distinct. Equality
of $x_2$ with $y_2$ or $x'_2$ would give
$\pi^2a=-u^j a$ for $j=1$ or $2$. Raising to the
seventh power shows $\pi^2a=a$, putting $t$ in
$E(\mathbf F_4)$, which is impossible. Equality
$y_2=x'_2$ would give $ua=a$, hence $3a=0$ by
(RWI.3). But $\#J(k)=4\cdot4159$ by A38, so this
would force $a=0$, again impossible.

None of these three points is $t$. Equality $x_2=a$
implies $\pi^2a=a$ and is excluded as above.
Equality $y_2=a$ implies $\pi a=ua$; taking seventh
powers gives $ua=a$, excluded by the same absence
of rational three-torsion. Finally $x'_2=a$
implies $(\pi-u^2)a=(T+1)a=0$, contradicting
invertibility of $T+1$.

The same ring calculations give
\[
 F_t(x_2,y_2)=y_2,\qquad
 F_t(x'_2,y_2)=x'_2,\qquad
 F_t(y_2,x'_2)=x_2.
 \tag{RWI.15}
\]
Consequently $a_2=(x_2,y_2)$ is another rational
$\tau_1$-fixed point of $D_t$, and its orbit is
exactly (RWI.13) with every variable replaced by
its subscript-two version. Its source point
$W_t^{(2)}=\pi_t(a_2)$ is therefore rational,
Weierstrass and split, with all five reflection
images non-Weierstrass.

The two fixed lifts are distinct: their first
coordinates are $\iota F^{-1}(t)$ and $\iota F(t)$,
whose equality would again put $t$ in $E(\mathbf F_4)$.
A reflection has at most one fixed point in a free
$C_5$ fibre. Indeed if it fixes $b$ and $\sigma^j b$,
then inversion of the deck group gives
$\sigma^{2j}b=b$, forcing $j=0$. Hence
$W_t^{(1)}\ne W_t^{(2)}$.
The hyperelliptic cover of a rank-one genus-two curve
has exactly two geometric branch points, by the same
Deuring--Shafarevich formula. These are therefore both
Weierstrass points of $C_t$.

All images in both orbits are non-Weierstrass points
of $E(k)$, so A38 gives exact canonical order $4159$
for each. Changing the lift only permutes its orbit,
and changing the reflection runs through the same
five images. An isomorphism of a quotient with $E$
preserves canonical-class order. Finally constant
extension preserves the order of each fixed geometric
class. This proves (RWI.1) and the stated detector failure.
\end{proof}


\paragraph{Supplement to A43: the order-three mechanism is special.}
The automorphism used in A43 cannot be supplied on every ordinary
real-multiplication target.  More precisely, let $C$ be an ordinary
genus-two curve over an algebraically closed field of characteristic
two, put $X=C^{(1)}$, and let $B$ be the locally exact differential
line bundle for $F:C\to X$.  If an automorphism $\rho$ of $C$ of
order three induces $u$ on $J_C$ with
\[
                         u^2+u+1=0,
\]
then the actual Raynaud theta support of $C$ contains a point of
exact order three.  Consequently there are ordinary finite-field
genus-two curves with the specified integral Rosati-$\mathcal O_5$
action which admit no such geometric automorphism, even after
extending constants.

\begin{proof}
Write $\rho_1=\rho^{(1)}$ and $u_1=u^{(1)}$ for the induced
automorphisms of $X$ and $J_X$.  The quotient
$X\to X/\langle\rho_1\rangle$ is a separable degree-three map
of smooth proper curves.  Its ramification is tame because the
characteristic is two.  If $\rho_1$ had no geometric fixed point,
this quotient would be etale: a nontrivial stabilizer in its
prime-order group would be the whole group.  Riemann--Hurwitz
would then give
\[
                    2=3\bigl(2g(X/\langle\rho_1\rangle)-2\bigr),
\]
which is impossible for an integer genus.  Here the exact input is
the smooth-proper separable formula and its tame different term,
Stacks Project, Lemmas~53.12.2 and~53.12.4, Tags
\texttt{0C1D} and \texttt{0C1F}; in the etale case the different
is zero \cite{StacksRH}.  Thus a fixed point $R\in X$ exists.

The required invariance is that of the actual Raynaud line bundle.
Its defining sequence is
\[
 0\longrightarrow\mathcal O_X\longrightarrow F_*\mathcal O_C
                     \longrightarrow B\longrightarrow0
\]
\cite[Section~8.1, p.~472]{OzmanPries2019}.
The natural identity $F\rho=\rho_1F$ identifies this sequence
with its pullback by $\rho_1$: both horizontal maps in the
corresponding square are isomorphisms, and their pullbacks
identify the displayed inclusion of structure sheaves.
Taking the quotient therefore gives $\rho_1^*B\simeq B$.
The induced pushforward action also fixes the class of $B$.
Hence
\[
 z=[\mathcal O_X(R)\otimes B^{-1}]
     \in\Theta_C,\qquad u_1z=z,
     \qquad 3z=(u_1^2+u_1+1)z=0.
\]
Ordinarity gives $H^0(X,B)=0$, equivalently
$0\notin\Theta_C$; see the first paragraph of
\cite[Section~8.2, p.~472]{OzmanPries2019}.
Thus $z\ne0$, and its order is exactly three.

Now apply Finding~A37(2) with $S=\{3\}$.
It supplies an ordinary finite-field genus-two curve with its
specified integral Rosati-$\mathcal O_5$ action and with no
nonzero geometric $3$-primary point on its Raynaud theta support.
The conclusion just proved excludes the stated automorphism
over the algebraic closure, so no constant extension restores it.
In particular the polarized endomorphism of A43 cannot exist on
this curve: the genus-two Torelli theorem realizes a
polarization-preserving Jacobian automorphism by a curve
automorphism, absorbing the possible sign by the hyperelliptic
involution, and its faithful action on the Jacobian preserves
the order \cite[Theorem~12.1]{MilneJV}.
\end{proof}

The extra relation used in A43 is
$\pi+V=2T+1$ for the specific $\mathbf F_2$ model, in addition
to $\pi V=2$ and $T^2+T=1$.  For a general finite-field
target whose specified real multiplication is defined over the
ground field, $T$ commutes with the $q$-Frobenius $\pi$ and
with $V=\pi^\dagger$.  Real multiplication alone then gives only
\[
 (T-\pi)^\dagger(T-\pi)=T^2-T(\pi+V)+q,
 \qquad
 (T-\pi)^2+(T-\pi)+1
                =\pi^2-(2T+1)\pi+2.
\]
Neither expression has the special value required in A43 just
from the real-multiplication relation.  The preceding argument
shows that the missing automorphism is actually absent on some
targets, rather than merely absent from the proof.
This limits that particular inverse-family calculation; it
does not disprove another rank-one realization, a different
choice of cyclic cover, or point-section separation.


\subsection{A characteristic-power tower has only finitely many eligible primes}
\label{sec:pro-two-tower-obstruction}

\paragraph{Finding A44 (a finite obstruction to the controlled-tower attempt).}
Let $A/\mathbf F_q$ be an abelian surface, with $q$ a power of two.
Write its Frobenius polynomial as
\[
 \begin{gathered}
 P_A(T)=T^4-aT^3+bT^2-qaT+q^2,\qquad a,b\in\mathbf Z,\\
 K=\bigcup_{m\ge0}\mathbf F_{q^{2^m}}.
 \end{gathered}
 \tag{CT.1}
\]
Call an odd prime $\ell$ eligible if
$\langle2,-1\rangle=(\mathbf Z/\ell)^\times$.
Define the following nonzero integers:
\[
 \begin{aligned}
 D_\epsilon&=(\epsilon b-q^2-1)(q+\epsilon)^2-qa^2,
                         &&\epsilon\in\{1,-1\},\\
 L_A&=\left|\#A(\mathbf F_{q^{16}})
                    (q^4-1)D_1D_{-1}\right|.
 \end{aligned}
 \tag{CT.2}
\]
Every eligible prime occurring in the order of a point of $A(K)$
divides $L_A$. Consequently there is a finite, computable set
\[
 S_A=\{\ell:\ell\text{ is odd},\ \ell\mid L_A,\
                 \langle2,-1\rangle=(\mathbf Z/\ell)^\times\}
 \tag{CT.3}
\]
containing every such prime. Put
\[
 e=\max_{\ell\in S_A}v_2\!\left(
                   \prod_{j=1}^{4}(\ell^j-1)\right),
 \tag{CT.4}
\]
with $e=0$ if $S_A$ is empty. Then
\[
          A(K)_{S_A}=A(\mathbf F_{q^{2^e}})_{S_A}.
 \tag{CT.5}
\]
In particular, all points of $A(K)$ whose odd order $n>1$ satisfies
$\langle2,-1\rangle=(\mathbf Z/n)^\times$ belong to the finite
group on the right of (CT.5). This is an unconditional statement.

For a genus-two curve in characteristic two, the actual Raynaud
theta support nevertheless contains infinitely many nonzero
odd-order points over this same field $K$. All but finitely many
of them therefore have a prime factor in their order which is
individually ineligible. Thus abundance in this constant-field
tower cannot by itself supply an unbounded family of the orders
required by the single-orbit primitive bridge.

\paragraph{Known inputs and exact scope.}
We use the Frobenius polynomial, its point-count formula and its
Weil absolute values \cite[Proposition~12.9 and Theorem~19.1]{MilneAV}, and the
prime-to-characteristic torsion module
\cite[Theorem~8.2 and Remark~8.4]{MilneAV}.
The finite-primary stability in a tower is consistent with
Poonen's Lemma~6.2 \cite{Poonen2005}; the new calculation here
isolates a finite set of eligible primes in the characteristic-two
tower. The argument does not treat towers of odd auxiliary-prime
degree. Nor does finiteness say that the remaining finite set of
theta witnesses is empty: the example below is nonempty.

\begin{proof}
\emph{1. The integers in (CT.2) do not vanish.}
All complex roots of $P_A$ have absolute value $\sqrt q>1$.
If $D_\epsilon=0$, put $s=a/(1+\epsilon q)$. Direct coefficient
comparison would give a factorization over $\mathbf Q$,
\[
 P_A(T)=(T^2-sT+\epsilon)
                 (T^2-\epsilon qsT+\epsilon q^2).
 \tag{CT.6}
\]
The two roots of the first factor would be roots of $P_A$,
so their product would have absolute value $q$. Its constant
term instead has absolute value one. This is a contradiction.
The other factors in $L_A$ are positive nonzero integers.

\emph{2. A prime first appearing late supplies a cyclotomic eigenvalue.}
Let an eligible prime $\ell$ divide $\#A(\mathbf F_{q^{2^m}})$,
with $m$ least. If $m\le4$, then
$\ell\mid\#A(\mathbf F_{q^{16}})$ and it is already in (CT.2).
Suppose $m\ge5$. On the four-dimensional space
$A[\ell](\overline{\mathbf F}_q)$, Frobenius $\Phi$ has an
eigenvalue $\lambda$ of exact multiplicative order $2^m$.
Indeed the nonzero kernel of $\Phi^{2^m}-1$ is semisimple for
$\Phi$, since $T^{2^m}-1$ is separable modulo $\ell$.
If all eigenvalues in that kernel had order dividing $2^{m-1}$,
there would be nonzero rational $\ell$-torsion at the preceding
level, contrary to minimality. The minimal polynomial of
$\lambda$ over $\mathbf F_\ell$ divides $P_A\bmod\ell$; hence
\[
       [\mathbf F_\ell(\lambda):\mathbf F_\ell]
               =\operatorname{ord}_{2^m}(\ell)\le4.
 \tag{CT.7}
\]
This degree is a power of two, so it is $1$, $2$ or $4$.
No semisimplicity of the full Frobenius matrix is assumed.

Eligibility excludes $\ell\equiv1\pmod8$: in that residue class
both $2$ and $-1$ are squares modulo $\ell$, so they cannot
generate the full unit group. Here the assertion for $2$ follows
directly from Gauss's lemma: for $\ell=8u+1$ the number of
$j$ with $1\le j\le(\ell-1)/2$ and $2j>\ell/2$ is $2u$.
If $\ell\equiv3$ or $5\pmod8$, then
$v_2(\ell^4-1)=4$, contradicting (CT.7) when $m\ge5$.
Consequently
\[
                              \ell\equiv7\pmod8.
 \tag{CT.8}
\]
The degree in (CT.7) cannot be one, because then
$\ell\equiv1\pmod{2^m}$.

\emph{3. Degree two forces one of two fixed divisibilities.}
Suppose $\lambda$ has degree two, and write
$t=v_2(\ell+1)$; by (CT.8), $t\ge3$ and $v_2(\ell-1)=1$.
The condition $2^m\mid\ell^2-1$ implies $t\ge m-1$.
Therefore
\[
          \lambda^{\ell+1}=\epsilon\in\{1,-1\},
       \qquad
          \operatorname{min}_\lambda(T)=T^2-sT+\epsilon.
 \tag{CT.9}
\]
The Frobenius polynomial is $q$-reciprocal, so its roots also
include $q/\lambda$ and $q/\lambda^\ell$. If these two roots
overlap the first pair, then either $q=\epsilon$ in
$\mathbf F_\ell$, or $q=\lambda^2$.
In the latter case $\lambda^2$ is in $\mathbf F_\ell$ and
$\lambda\notin\mathbf F_\ell$, so
$\lambda^\ell=-\lambda$ and $\epsilon=-q$.
Thus every overlap case gives $\ell\mid q^2-1$, already a factor
of $L_A$.

In the remaining case all four roots are distinct and exhaust
$P_A\bmod\ell$. The factorization (CT.6) holds modulo $\ell$.
Its coefficients give
\[
        a=(1+\epsilon q)s,\qquad
        b=\epsilon(1+q^2+qs^2).
 \tag{CT.10}
\]
Eliminating $s$ without division gives $\ell\mid D_\epsilon$.
This proves the required assertion in degree two.

\emph{4. Degree four forces another fixed divisibility.}
If $\lambda$ has degree four, then
$2^m\mid\ell^4-1$ but $2^m\nmid\ell^2-1$.
Since $v_2(\ell^2+1)=1$, this forces
$v_2(\ell+1)=m-2$. Hence
\[
 v_2(1+\ell+\ell^2+\ell^3)
       =v_2((1+\ell)(1+\ell^2))=m-1.
\]
The norm of $\lambda$ to $\mathbf F_\ell$ is therefore $-1$.
Its degree-four minimal polynomial is $P_A\bmod\ell$, whose
constant term is $q^2$. We conclude that
$\ell\mid q^2+1$, again a factor of $L_A$.
Together with Steps~2--3 this proves (CT.3).

\emph{5. Arbitrary primary exponents are already in one finite field.}
Let $\ell\in S_A$ and let $x\in A(K)$ have order $\ell^r$.
Its field degree over $\mathbf F_q$ is a power of two, and is
the length of its Frobenius orbit. This divides the order of
the Frobenius matrix in
$\operatorname{GL}_4(\mathbf Z/\ell^r)$ and hence the order of
that finite group. Reduction modulo $\ell$ has an $\ell$-group
kernel. As $\ell$ is odd, the two-part of its order is exactly
that of $\operatorname{GL}_4(\mathbf F_\ell)$, namely
\[
        2^{\,v_2(\prod_{j=1}^4(\ell^j-1))}.
 \tag{CT.11}
\]
The orbit degree thus divides $2^e$. Adding the finitely many
primary components proves (CT.5). Reduction of the unit group
modulo any prime divisor of an odd $n$ is surjective; therefore
the full signed-unit condition at $n$ implies eligibility at
every such prime. This proves the finite-field assertion for
all the stated orders, with no bound on their exponents.

\emph{6. Odd theta points in the tower are still unbounded.}
We prove this last statement so that the obstruction does not
merely come from a lack of odd torsion. First, for any abelian
surface as in (CT.1), put
\[
 d_m=\#A(\mathbf F_{q^{2^m}})[2^\infty].
\]
There are $m_0\ge0$ and an integer $0\le t\le2$ such that
\[
                     d_m=d_{m_0}2^{t(m-m_0)}
                         \qquad(m\ge m_0).
 \tag{CT.12}
\]
To see this, put all four roots $\alpha_i$ of $P_A$ in a finite
extension of $\mathbf Q_2$, with valuation normalized by $v(2)=1$.
The point-count product gives
$v_2(d_m)=\sum_i v(1-\alpha_i^{2^m})$.
The roots are integral. A nonunit contributes zero. A unit whose
residue is not one also contributes zero: its residue has odd
order, which cannot divide $2^m$.
For a unit with residue one, repeated squaring eventually gives
$v(\alpha_i^{2^{m_0}}-1)>1$. Explicitly, while the positive
valuation is less than one, squaring doubles it; at valuation
one the next valuation is at least two. Thereafter
\[
  v(u^2-1)=v(u-1)+1\quad\text{when }v(u-1)>1.
\]
No root is a root of unity, by its complex absolute value.
Thus each such unit contributes an affine function of $m$ with
slope one for all sufficiently large $m$.
The number of unit roots is at most two: the reciprocal pairing
$\alpha\mapsto q/\alpha$ sends them injectively to nonunit
roots, with multiplicities. Summing proves (CT.12).

Now take $A=J_X$, where $X=C^{(1)}$ and
$\Theta=i_B(X)$ is the actual genus-two Raynaud theta support.
At the individual field of size $Q_m=q^{2^m}$, use the Lang
quotient whose deck group is the rational two-primary subgroup.
The argument of A12, or Steps~2--5 of A42 at this single field,
gives
\[
 \begin{aligned}
 O_m&=\#\bigl(\Theta(\mathbf F_{Q_m})
                  \cap J_X(\mathbf F_{Q_m})_{\rm odd}\bigr),\\
 O_m&\ge\frac{Q_m+1}{d_m}
                      -2\left(1+\frac1{d_m}\right)\sqrt{Q_m}.
 \end{aligned}
 \tag{CT.13}
\]
Here the counting covers may vary with $m$; A42's fixed-primary
extension hypothesis is not being asserted. The bound (CT.12)
is polynomial in $2^m$, whereas $Q_m$ is exponential in $2^m$.
The right side of (CT.13) tends to infinity. There is at most
one zero theta class, so this gives infinitely many nonzero
odd theta points over $K$.
The group of points of $J_X(K)$ supported on eligible primes
is finite by (CT.3)--(CT.5). Every other odd theta point must
therefore have an ineligible prime factor in its order.
\end{proof}

\paragraph{A finite candidate set can contain useful points.}
For the A13 curve
\[
 y^2+y=x+x^{-1}+(x+1)^{-1}
\]
over $\mathbf F_2$, $P_A(T)=T^4-T^2+4$, so $a=0$, $b=-1$,
$D_1=-54$ and $D_{-1}=-4$. The established quadratic recurrence
for $4$-power Frobenius gives
\[
 \#A(\mathbf F_{2^{16}})=65760^2
                    =2^{10}3^2 5^2 137^2.
 \tag{CT.14}
\]
Since $137\equiv1\pmod8$, (CT.3) gives $S_A=\{3,5\}$.
In particular the order-$4159$ witness of A13 lies outside this
constant-field tower; it was constructed over an odd-degree
extension.

There are in fact order-three theta points in the tower.
Put $h=X^2+X$ and use the polynomial model
\[
 Y^2+hY=X^5+X^3+X^2+X=:f.
\]
Let $W_0,W_1,W_\infty$ be its three Weierstrass points.
The curve and its Frobenius twist have this same
$\mathbf F_2$-model. The Raynaud theta characteristic is
$B=W_0+W_1-W_\infty$.
Here are the details of this identification. The three simple
Artin--Schreier poles make the genus-two curve ordinary, so
$J[2](\overline{\mathbf F}_2)$ has four elements and there are
exactly four geometric theta characteristics. The three $W_i$
are distinct effective ones. Since $2W_i\sim K_C$, the displayed
$B$ is another theta characteristic. It differs from $W_0$
and $W_1$ because distinct degree-one point divisors are not
linearly equivalent. If $B\sim W_\infty$, then
$W_0+W_1\sim2W_\infty$ would be a canonical hyperelliptic fibre,
impossible for two distinct branch points. Thus $B$ is the
fourth characteristic. Any effective characteristic is a point
$R$ with $2R\sim K_C$, hence a Weierstrass point. Consequently
$B$ is ineffective. The locally exact differential
characteristic is ineffective on an ordinary curve, by the
Frobenius cohomology sequence, so it is precisely this $B$.
All three $W_i$ are rational over $\mathbf F_2$.

Choose $r\in\mathbf F_4$ with $r^2+r+1=0$, and
$a_0\in\mathbf F_{16}$ with $a_0^2+a_0=r^2$.
Then $R=(r,a_0)$ lies on the curve, and direct multiplication gives
\[
 \begin{aligned}
 F&=Y+a_0h,\\
 F\,\iota(F)&=f+(a_0^2+a_0)h^2
                          =h(X+r)^3,\\
 \operatorname{div}(F)&=3R+W_0+W_1-5W_\infty.
 \end{aligned}
 \tag{CT.15}
\]
The conjugate value at $R$ is one, since $h(r)=1$; the pole of
$F$ at infinity has order five. These facts and the norm identity
justify the full divisor formula. Since
$\operatorname{div}(h)=2W_0+2W_1-4W_\infty$, it follows that
\[
                    \operatorname{div}(F/h^2)=3(R-B).
 \tag{CT.16}
\]
The class $R-B$ is nonzero because $B$ is ineffective.
This gives four exact order-three theta points over
$\mathbf F_{16}$, one for each of the two choices of $r$ and
two choices of $a_0$. Thus the finite set in A44 need not be
empty, and a finite member can already suffice for A11.
The proof does not assert that these are all eligible theta
points of this curve.

\paragraph{Status of the attempted main-route application.}
The computation closes the specific proposed inference from
unbounded odd theta torsion in a characteristic-power constant
tower to an unbounded supply of signed-unit-compatible orders.
It leaves a finite, decidable candidate set for that tower,
which can contain successful witnesses as (CT.15) shows.
It does not settle the existence of a witness in that set for
every curve, and it does not rule out an odd auxiliary-prime
tower, a different cover, or another primitive-factor mechanism.
No previously open universal prescribed-point row is marked
closed by this obstruction.


% Checked necessary-condition supplement to Finding A44.
\paragraph{Supplement to Finding A44: odd auxiliary-prime towers.}

This note continues the controlled-tower test after Finding~A44.
The pro-two conclusion does not extend by the same argument to an
odd auxiliary-prime tower.  There is, however, a precise reduction
of the orders which that method would have to produce.  This is
an analysis of the remaining gate, not a proof that the required
orders occur.

Let $A/\mathbf F_q$ be an abelian surface, where $q$ is a power
of two, and write its Frobenius polynomial as
\[
 P_A(T)=T^4-aT^3+bT^2-qaT+q^2.
\]
Put
\[
 D_+=(b-q^2-1)(q+1)^2-qa^2\ne0.
 \tag{OT.1}
\]
The nonvanishing is the argument already proved in A44: equality
would give a factor $T^2-sT+1$ of $P_A$ over $\mathbf Q$, with
$s=a/(q+1)$, although the product of the absolute values of two
Weil roots is $q>1$.
Fix an odd prime $r$, and put
\[
 K_r=\bigcup_{m\ge0}\mathbf F_{q^{r^m}},\qquad
 S_0=\{\hbox{odd primes dividing }
           r\,\#A(\mathbf F_q)(q-1)D_+\}.
 \tag{OT.2}
\]
Then every odd prime $\ell\notin S_0$ occurring in the order of a
point of $A(K_r)$ satisfies $\ell\equiv1\pmod r$.
More precisely, if $m\ge1$ is the first index for which
$\ell\mid\#A(\mathbf F_{q^{r^m}})$, then
\[
 \ell\equiv1\pmod{r^m},
 \qquad P_A(T)\bmod\ell
 \text{ has a root in }\mathbf F_\ell
 \text{ of exact order }r^m.
 \tag{OT.3}
\]

\begin{proof}
Since $\ell\ne r,2$, the polynomial $T^{r^m}-1$ is separable
modulo $\ell$.  The first-appearance argument from A44 therefore
gives a Frobenius eigenvalue $\lambda$ on $A[\ell]$ of exact
order $r^m$.  Its degree over $\mathbf F_\ell$ is
\[
                 d=\operatorname{ord}_{r^m}(\ell)\le4.
\]
The multiset of roots of $P_A\bmod\ell$ is stable under the
involution $\alpha\mapsto q/\alpha$.  Its distinct Frobenius
orbits are either equal or disjoint, including when the
characteristic polynomial has repeated factors.

Degree three is impossible.  The orbit of three roots of the
irreducible factor containing $\lambda$ and its reciprocal orbit
cannot be disjoint inside a polynomial of degree four.  They must
be the same orbit.  The reciprocal involution commutes with its
three-cycle Frobenius permutation.  The centralizer of that
three-cycle on its three elements is cyclic of order three, so
the commuting involution is the identity.  This would give
$\lambda^2=q$, which has degree at most two over $\mathbf F_\ell$,
a contradiction.

If $d=2$, the cyclicity of
$(\mathbf Z/r^m\mathbf Z)^\times$ gives
$\ell\equiv-1\pmod{r^m}$, so
$\lambda^{\ell+1}=1$.  Its irreducible quadratic is therefore
$g=T^2-sT+1$ for some $s\in\mathbf F_\ell$.
If its reciprocal orbit coincides with its orbit, the reciprocal
involution cannot be the identity: that would put
$\lambda^2=q$ in $\mathbf F_\ell$, and since $\lambda$ has odd
order, $\lambda$ itself would belong to $\mathbf F_\ell$.
It must be the nontrivial Frobenius permutation, and consequently
$q=\lambda^{\ell+1}=1$ modulo $\ell$.
Otherwise the two quadratic orbits are disjoint and exhaust the
four roots, giving
\[
 P_A(T)=(T^2-sT+1)(T^2-qsT+q^2)\pmod\ell.
\]
Comparing coefficients without division gives $\ell\mid D_+$.
Both possibilities are excluded by $\ell\notin S_0$.

If $d=4$, its orbit exhausts all four roots.  The reciprocal
involution commutes with the four-cycle Frobenius permutation.
It is not the identity, again because that would give
$\lambda^2=q$ and degree at most two.  It is therefore the square
of that permutation.  Since an element of order four in the
cyclic group $(\mathbf Z/r^m\mathbf Z)^\times$ squares to $-1$,
\[
 q=\lambda^{\ell^2+1}=1\pmod\ell,
\]
which is also excluded.  Only $d=1$ remains, proving (OT.3).
For $r=3$ the degree-four case cannot arise; the same proof
otherwise applies without a restriction $r\ge5$.
\end{proof}

The whole-order condition is substantially stronger than
eligibility of its individual prime divisors.  Call an odd
integer $n>1$ compatible when
\[
                   \langle2,-1\rangle
                      =(\mathbf Z/n\mathbf Z)^\times.
 \tag{OT.4}
\]
If a point of $A(K_r)$ has compatible order $n$, at most one
prime divisor of $n$ lies outside $S_0$.  If such a prime occurs,
then $v_r(n)\le1$.  In particular, the sparse cofactor supplied by
Poonen's tower proof cannot retain two new primary components
and still satisfy (OT.4).

\begin{proof}
The Sylow-$r$ subgroup of a group generated by $2$ and $-1$ is
cyclic: the subgroup generated by $2$ has index at most two,
and $r$ is odd.  Each distinct prime $\ell\notin S_0$ dividing
$n$ supplies a nontrivial Sylow-$r$ factor in
$(\mathbf Z/\ell^{v_\ell(n)}\mathbf Z)^\times$, since
$r\mid\ell-1$.  Two such factors would give a noncyclic product.
If $r^2\mid n$, the factor for $r^{v_r(n)}$ supplies another
nontrivial Sylow-$r$ factor, proving the second assertion.
The same argument also excludes any other cofactor prime $s$
with $r\mid s-1$ when a prime outside $S_0$ already occurs.
\end{proof}

There is an explicit finite list of possible cofactors.  Let
\[
 e=\max_{s\in S_0\setminus\{r\}}
       v_r\!\left(\prod_{j=1}^4(s^j-1)\right),
 \qquad
 N_0=\#\bigl(A(\mathbf F_{q^{r^e}})\{S_0\setminus\{r\}\}\bigr),
 \tag{OT.5}
\]
where an empty maximum is zero, and the final notation denotes
the order of the indicated primary subgroup.  Then
\[
 A(K_r)\{S_0\setminus\{r\}\}
  =A(\mathbf F_{q^{r^e}})\{S_0\setminus\{r\}\}.
 \tag{OT.6}
\]
Indeed, a point of $s$-power order, $s\ne r$, has a Frobenius
orbit of $r$-power length dividing the order of
$\mathrm{GL}_4(\mathbf Z/s^j\mathbf Z)$ for some $j$.
Its $r$-part is bounded by the product in (OT.5), independently
of $j$, because the reduction kernel is an $s$-group.  This is
Poonen's complete Lemma~6.1 argument and the finite-subgroup
conclusion of Lemma~6.2, with the field bound made explicit
\cite[Lemmas~6.1--6.2]{Poonen2005}.

Also, (OT.4) forces $n$ to have at most two distinct odd prime
factors: its unit group has $2$-rank equal to their number,
whereas two generators give rank at most two.  Thus any compatible
$n$ containing a prime outside $S_0$ has the form
\[
 n=c\ell^a,\quad a\ge1,\quad \ell\notin S_0,\qquad
 c\in\mathcal M,
 \tag{OT.7}
\]
where the following possibly redundant set is finite and
computable:
\[
 \mathcal M=\{1,r\}\ \cup
 \{s^j:s\in S_0\setminus\{r\},\ r\nmid s-1,
                         \ 1\le j\le v_s(N_0)\}.
 \tag{OT.8}
\]
One must still test (OT.4) at the actual integer in (OT.7).
The inequalities and the list are necessary conditions, not
sufficient conditions for a theta incidence.

For a genus-two curve in characteristic two, take
$A=J_{C^{(1)}}$ and the actual Raynaud support
$\Theta_C=C^{(1)}-B_C$.  Finding~A19 makes its subset of points
supported on the fixed odd set $S_0$ finite, with an effective
finite-field enumeration bound.  Consequently an unbounded
compatible theta supply from this odd tower would require
infinitely many incidences of the form (OT.7), with $c$ in the
fixed finite list.  The pure-prime-power case $c=1$ remains among
them.  This conclusion does not say that no such incidences
exist or that an arbitrary individual curve lacks a witness.

\paragraph{Exact location of the unresolved inference.}
In Poonen's proof of Lemma~6.4, a large finite-prime initial set
is removed and the remaining subgroup $H$ has uniformly bounded
index; a translated Lang pullback then produces a point in
$a+H$.  This controls the total size of the removed primary
part, but allows the point's residual component to use several
new primes.  In the tower above, two such primes already
contradict (OT.4).  Replacing $H$ by the union of its individual
primary subgroups would impose the needed restriction, but
that union need not be a subgroup and does not define the same
Lang quotient.  The bounded-index argument therefore stops at
this step.  The density-zero conclusion of Poonen's
Lemma~6.3 does not remove this issue
\cite[Lemmas~6.3--6.6]{Poonen2005}.

Nor is the surviving congruence itself an arithmetic obstruction
to eligibility: for $r=3$, the primes $7$ and $13$ are both
$1\bmod3$ and satisfy $\langle2,-1\rangle=U_\ell$.
Their product $91$ does not satisfy that equality, since its
unit group has a Sylow-three subgroup $C_3\times C_3$.
These are assertions about finite unit groups only; they do not
assert Frobenius eigenvalues or theta points on a chosen curve.

This probe sharpens the residual selection gate for the
controlled-tower construction.  It supplies neither a
prime-incidence theorem nor a resolution of a previously open
universal prescribed-point row.  No density statement about
Artin primes alone is used to infer incidence.


\subsection{The actual-exponent criterion before deck constantization}
\label{sec:nonconstant-deck-exponent}

\paragraph{Finding A45 (a geometrically cyclic pointed cover).}
Let $C/k$ be a smooth projective geometrically connected
curve of genus $g\ge2$ over a finite field of characteristic $p$,
and let $P\in C(k)$.
Let $\pi:D\to C$ be a geometrically connected finite etale cover
defined over $k$, such that $D_{\bar k}\to C_{\bar k}$ is cyclic
of odd degree $n>1$ prime to $p$. Assume that a point $P'\in D(k)$ lies over $P$.
The geometric deck group need not be constant over $k$. Put
\[
 B=\ker(\operatorname{Nm}:J_D\longrightarrow J_C)^0,
 \qquad M=\exp\bigl(B(k)[p^\infty]\bigr),
 \tag{ND.1}
\]
where the trivial group has exponent one. If
\[
                         n>(2M-1)^2,
 \tag{ND.2}
\]
then any distinct $Q\in C(k)$ whose point section collides with
that of $P$ forces a separable map
\[
 C_{\bar k}\longrightarrow E,\qquad
 \deg(C_{\bar k}/E)=p^a,\quad p\le p^a\le g-1,
 \tag{ND.2a}
\]
to an elliptic curve, totally ramified at both $P$ and $Q$.
Consequently (ND.2) separates $P$ from every distinct rational
point whenever $2\le g\le p$. In any genus it does so if no
map of the type (ND.2a) is totally ramified at $P$.
Only the rational point $P'$ is required; the entire fibre above
$P$ need not split over $k$. The exponent in (ND.1) is measured
over the original field $k$.

This extends A34 to nonconstant geometric deck groups and to
every characteristic, with the stated higher-genus obstruction. Whether it improves the numerical test for a specified
cover depends on the actual exponents before and after extending
constants. The primitive-conductor descent statement below records
an important case where those exponents are equal. No numerical
strictness is inferred just from the change of formulation.

\begin{proof}
\emph{1. Project the rational collision class before applying a deck map.}
Suppose that a distinct $Q\in C(k)$ has the same point section as
$P$. The pointed-cover implication supplies $Q'\in D(k)$ above
$Q$ with
\[
                       \delta=[Q'-P']\in J_D(k)[p^\infty].
 \tag{ND.3}
\]
This implication is the one used in A01 and A34; the geometric
monodromy of this cover is the cyclic group of order prime to $p$, so this is an
allowed prime-to-$p$ neighbourhood even though its arithmetic
Frobenius acts nontrivially on that group.

Choose a power $N$ of $p$ killing $\delta$ and an integer $r$
with $rn\equiv1\pmod N$. Define
\[
                b=\delta-r\pi^*\operatorname{Nm}(\delta).
 \tag{ND.4}
\]
All maps in this formula are defined over $k$, so $b$ is rational
over $k$ and is killed by $N$. Also $\operatorname{Nm}(b)=0$.
We verify that it lies in the connected norm kernel, without
assuming that the whole norm kernel is connected.

Over $\bar k$, choose a generator $\sigma$ of the geometric deck
group and write $S=1+\sigma+\cdots+\sigma^{n-1}$.
The norm--pullback identity is $\pi^*\operatorname{Nm}=S$.
One has $\operatorname{im}(\sigma-1)=B$: the image is connected
and lies in the norm kernel, whereas $n-S$ is divisible by
$\sigma-1$ in the integral group ring and restricts to $[n]$ on
$B$. Surjectivity of $[n]:B\to B$ proves the reverse inclusion.
Since $rn$ acts as the identity on the $N$-torsion class $S\delta$,
\[
             nb=(n-S)\delta\in B(\bar k),\qquad b=r(nb)\in B(\bar k).
 \tag{ND.5}
\]
Together with rationality, this proves $b\in B(k)[p^\infty]$.
The pullback term in (ND.4) is fixed by every geometric deck map.
Consequently
\[
              (\sigma-1)\delta=(\sigma-1)b,
              \qquad M(\sigma-1)\delta=0.
 \tag{ND.6}
\]
The class on the left need not itself be $k$-rational; its stated
annihilator follows from the rational class $b$, and requires no
extension of constants.

Let $m$ be its exact order, with $m=1$ for the zero class. Then
$m$ is a power of $p$ and $m\mid M$. On $D_{\bar k}$ there is a
function $f$ with
\[
 \operatorname{div}(f)
       =m(\sigma Q'+P'-Q'-\sigma P'),\qquad \deg f=d=2m.
 \tag{ND.7}
\]
The four points are distinct: the two base fibres are disjoint
and the deck action is free. If $m>1$ and $f$ were a $p$th power,
one $p$th of this divisor would be principal, contradicting the exact
order of $(\sigma-1)\delta$. If $m=1$, the simple poles already
exclude a $p$th power. Thus $f$ is separable over the perfect field
$\bar k$.

\emph{2. A small intermediate field.}
Set $K=\bar k(D)$ and
\[
 F=\bar k(\gamma f:\gamma\in\langle\sigma\rangle),
 \qquad e=[K:F],
\]
and let $Z$ be the smooth projective curve with function field
$F$. Poonen's Corollary~2.3, with Proposition~2.1 and Lemma~2.2,
applies because $F$ is generated by degree-$d$ functions; its
controlled-field integer is $d/e$. It gives
\[
                    e\mid d,\qquad g(Z)\le(d/e-1)^2.
 \tag{ND.8}
\]
The exact published locator is author p.~3
\cite{Poonen2007Gonality}. These extensions lie inside the
separable extension $K/\bar k(f)$. Since $D\to C$ is etale,
$g(D)=n(g-1)+1$. Inequality (ND.2) implies $n>(d-1)^2$, so $e=1$
is impossible. Hence $e>1$ divides $2p^b$ for some $b\ge0$, so
$\gcd(e,n)=1$. The resulting
map $v:D_{\bar k}\to Z$ is separable of degree $e$, and
$f=h\circ v$ for a function $h$ on $Z$.

\emph{3. The geometric deck group acts freely on $Z$.}
The field $F$ is stable under the group. Its action on $Z$ is
faithful, since a kernel would have order dividing both $n$ and
$e$. If a nontrivial stabilizer existed, it would contain a
subgroup of odd prime order $\ell$. This subgroup acts freely
on $D$ and preserves the fibre divisor over its fixed point
in $Z$. All its orbits in that divisor have length $\ell$ and
constant multiplicities, forcing $\ell\mid e$, a contradiction.
Thus the action is free. Riemann--Hurwitz gives
\[
                    g(Z)-1=n\bigl(g(Z/\langle\sigma\rangle)-1\bigr).
\]
The bound $0\le g(Z)<n$ forces both genera to be one. Therefore
\[
                   \deg\operatorname{Diff}(v)=2g(D)-2=2n(g-1).
 \tag{ND.9}
\]
This is the same geometric conjugate-field construction used in
A34 and recovered successor A.53 \cite[A.53]{ASrun}.

\emph{4. The two positive fibres force a small power of $p$.}
Suppose first that $P'$ and $\sigma Q'$ have the same image
under $v$. That fibre contains exactly these two positive
points: any other point would be an additional zero of $f$.
If the zero of $h$ has multiplicity $s$, both indices are
$m/s=e/2$. Equivariance gives the same index $e/2$ at $Q'$
and $\sigma P'$. Their pole images must coincide as well:
otherwise each would form a singleton degree-$e$ fibre and
have index $e$, a contradiction. Writing $x=v(P')$ and
$y=v(Q')$, these two fibre equalities give
$x=\sigma y$ and $y=\sigma x$. Thus $\sigma^2x=x$, contrary
to the free action, since $n>1$ is odd.

Therefore $P'$ and $\sigma Q'$ have distinct images in $Z$.
Each is the sole point in its zero fibre, and is totally
ramified with index $e$. Since its multiplicity in $f$ is $m$,
one has $e\mid m$. Thus $e$ is a nontrivial power of $p$.
The two deck orbits are disjoint because they lie above distinct
points of $C$, and each has $n$ elements. Every point in both
orbits has ramification index $e$ and different exponent at
least $e$. Indeed, for a local parameter expression
$u=a_et^e+\cdots$ the leading derivative vanishes because
$p\mid e$, so $\operatorname{ord}_t(du/dt)\ge e$; this derivative
is nonzero by separability. Together with (ND.9) this gives
\[
                   2ne\le2n(g-1),\qquad p\le e\le g-1.
 \tag{ND.10}
\]
For $g\le p$ this is already impossible, proving the stated
unconditional range.

\emph{5. Identify the exact elliptic obstruction in any genus.}
Let $E=Z/\langle\sigma\rangle$; it is a smooth genus-one
curve over $\bar k$, hence an elliptic curve after choosing
an origin. The fields satisfy
\[
 \bar k(E)=F^{\langle\sigma\rangle},\qquad
 \bar k(C)=K^{\langle\sigma\rangle},\qquad
 K=F\,K^{\langle\sigma\rangle}.
 \tag{ND.11}
\]
For the last equality, the subgroup of the Galois group of
$K/K^{\langle\sigma\rangle}$ fixing the compositum is the
kernel of its faithful action on $F$, and is therefore trivial.
Since $[F:\bar k(E)]=n$ and $[K:F]=e$, one obtains
$[\bar k(C):\bar k(E)]=e$. This extension is separable:
$K/F$ is separable and $F/\bar k(E)$ is etale. Thus it defines
a separable degree-$e$ map $C_{\bar k}\to E$.

The finite etale pullback $C_{\bar k}\times_E Z$ has function
field $K$ by (ND.11), so is connected and smooth and is precisely
$D_{\bar k}$. Ramification indices in $D_{\bar k}\to Z$ are
therefore those of $C_{\bar k}\to E$, because the other base
change map $Z\to E$ is etale. The total ramification at $P'$
and $\sigma Q'$ descends to total ramification at $P$ and $Q$.
Together with (ND.10), this proves the obstruction (ND.2a)
and all asserted sufficient cases.
\end{proof}

\paragraph{Integral semilinear descent for a primitive conductor.}
Let $p$ be a prime, let $n>1$ be prime to $p$, and let
$a\in(\mathbf Z/n)^\times$ have order
$f$.  Use the full cyclotomic algebra
\[
 \mathcal O_n=\mathbf Z_p[Z]/\Phi_n(Z),\qquad
 \tau(Z)=Z^a,\qquad R_n=\mathcal O_n^{\langle\tau\rangle}.
 \tag{SD.1}
\]
This algebra can be a product of unramified discrete valuation
rings.  It is not to be replaced by one chosen factor unless
that factor is preserved by $\tau$.
Let $M$ be a finite $\mathcal O_n$-module killed by a power of
$p$, with an additive automorphism $F$ satisfying
\[
 F(cm)=\tau(c)F(m),\qquad F^f=1.
 \tag{SD.2}
\]
Then the natural map and its exponent consequence are
\[
 \mathcal O_n\otimes_{R_n}M^F\ \xrightarrow{\ \sim\ } M,
 \quad c\otimes m\longmapsto cm,
 \qquad \exp(M)=\exp(M^F).
 \tag{SD.3}
\]
The exponent of the zero group is one.  No freeness of $M$,
and no assumption that $p$ is prime to $f$, is required.

\begin{proof}
Since $p\nmid n$, $Z^n-1$ has distinct roots modulo $p$.
Choose a finite unramified extension $S/\mathbf Z_p$ in which
these roots all split; one obtains it by lifting the distinct
roots over $\mathbf F_{p^h}$ with $n\mid p^h-1$.
The resulting faithfully flat base change identifies
\[
 \mathcal O_n\otimes_{\mathbf Z_p}S
                    \simeq\prod_{j\in(\mathbf Z/n)^\times} S.
 \tag{SD.4}
\]
The automorphism $\tau$ acts by multiplication by $a$ on the
primitive-root indices.  Every orbit has exactly $f$ elements:
$a^ij=j$ with $j$ a unit implies $a^i=1$ modulo $n$.
Taking invariants commutes with this flat base change, since
invariants are the kernel of $\tau-1$ on the finite free
$\mathbf Z_p$-module $\mathcal O_n$.  Hence
\[
 R_n\otimes_{\mathbf Z_p}S
       \simeq\prod_{(\mathbf Z/n)^\times/\langle a\rangle}S.
 \tag{SD.5}
\]
It follows, by faithfully flat descent, that
$\mathcal O_n/R_n$ is finite etale and faithfully flat of
rank $f$.  More explicitly, the map
\[
 \mathcal O_n\otimes_{R_n}\mathcal O_n
       \longrightarrow\prod_{i=0}^{f-1}\mathcal O_n,
 \qquad c\otimes d\longmapsto(c\tau^i(d))_i
 \tag{SD.6}
\]
is an isomorphism: after the faithfully flat base change $S$
it is the bijection of coordinates supplied by a free cyclic
permutation on each orbit.  This is the finite etale
$\mathbf Z/f$-torsor property, including when $p$ divides $f$.

The powers of $F$ in (SD.2) therefore give an actual descent
datum for the $\mathcal O_n$-module $M$ along $R_n\to\mathcal O_n$.
The cocycle condition is exactly $F^f=1$ and the semilinear
composition law.  Faithfully flat descent for modules, in its
unrestricted module form, identifies the descended module with
$M^F$ and gives (SD.3); the exact primary statement is the
Stacks Project, Proposition~35.3.9, Tag~\texttt{023N}
\cite{Stacks023N}.  In particular this step applies to finite
torsion modules and does not average by $1/f$.
If $p^e$ kills $M^F$, it kills the displayed tensor product and
therefore $M$.  Conversely $M^F$ is a subgroup of $M$.
This proves the equality of exponents.
\end{proof}

\paragraph{Application to geometric cyclic covers.}
Let $\pi:D\to C$ be a geometrically connected etale cover of
smooth proper curves over a finite field $k$ of characteristic
$p$, geometrically Galois with cyclic deck group of odd order
$n$ prime to $p$.  Choose a geometric generator $\sigma$ and let $F$ be
arithmetic Frobenius on geometric points.  Write
\[
                 F\sigma F^{-1}=\sigma^a,
       \qquad f=\operatorname{ord}_n(a),\qquad k_s=\mathbf F_{q^s}.
 \tag{SD.7}
\]
Thus $k_f$ is the minimal constant extension defining every
deck transformation.  Put
$B=\ker(\operatorname{Nm}_\pi:J_D\to J_C)^0$.
Because $n$ is a unit in $\mathbf Z_p$, the group algebra
decomposes into its cyclotomic factors:
\[
 \mathbf Z_p[\langle\sigma\rangle]
       \simeq\prod_{d\mid n}\mathcal O_d.
 \tag{SD.8}
\]
Indeed the reductions of the cyclotomic factors of $Z^n-1$
are pairwise coprime.  If $e_d$ is the resulting idempotent,
write
\[
 \mathscr T_d=e_d\bigl(B(\overline k)[p^\infty]\bigr),
              \qquad d\mid n,\quad d>1.
 \tag{SD.9}
\]
These are the primitive conductor-$d$ torsion modules.  They
are stable under $F$, since raising a primitive root to $a$
preserves its conductor.  The trivial conductor is absent
from $B[p^\infty]$, and
$B(\overline k)[p^\infty]=\bigoplus_{d\mid n,\ d>1}\mathscr T_d$.
One can see this directly from
$\pi^*\operatorname{Nm}_\pi=\sum_{i=0}^{n-1}\sigma^i$:
the average of this sum projects onto conductor one, and
$[n]$ is invertible on all $p$-primary groups.

These modules can also be realized by actual subvarieties,
without replacing a factor by an arbitrary isogenous variety.
Fourier orthogonality on the $n$th roots of unity gives
\[
 \begin{aligned}
 p_d&=ne_d=\sum_{j=0}^{n-1}c_d(j)\sigma^j,
                       &p_d^2&=np_d,\\
 c_d(j)&=\sum_{\substack{\zeta^d=1\\\operatorname{ord}(\zeta)=d}}
                         \zeta^{-j}\in\mathbf Z.
 \end{aligned}
 \tag{SD.9a}
\]
The coefficients are integers because they are rational
algebraic integers.  The projector is fixed by Frobenius,
which permutes the primitive characters.  Thus
$B_d=\operatorname{im}(p_d:B\to B)$ is defined over $k$.
On $B_d$ the map $p_d$ is $[n]$; since $p$ is prime to $n$ this shows
$B_d(\overline k)[p^\infty]=\mathscr T_d$.
Indeed both inclusions follow by applying $p_d$ and the
inverse of $[n]$ on the finite cyclic group generated by a
$p$-primary point.  Hence the following exponent statements
refer to these canonical actual conductor subvarieties.

For each $d>1$ put $f_d=\operatorname{ord}_d(a)$ and
\[
                      M_d=\mathscr T_d^{F^{f_d}}.
\]
This is finite, since it is a subgroup of $B(k_{f_d})$.
On $\mathcal O_d$, $F$ is semilinear for the automorphism
$Z\mapsto Z^a$ of order $f_d$, and $F^{f_d}=1$ on $M_d$.
The preceding lemma gives the exact equality
\[
       \exp\bigl(\mathscr T_d^{F^{f_d}}\bigr)
                    =\exp\bigl(\mathscr T_d^F\bigr).
 \tag{SD.10}
\]
Thus a primitive conductor factor has the same rational
$p$-primary exponent over $k$ as over its own minimal deck
constantization field.  In particular, when $n$ is prime,
the whole Prym has this equality between $k$ and $k_f$.
It also holds for the primitive conductor-$n$ factor of any
odd degree, including a prime power.

For the whole Prym of a composite-degree cover one must retain
the distinction between $f_d$ and $f$.  If $f_d<f$, the module
$\mathscr T_d^{F^f}$ is acted on by $F$ with $F^f=1$, but
$F^{f_d}$ need not be the identity.  The coefficient action
has already become trivial at $f_d$; its remaining action on
the module is an ordinary linear action.  Consequently
(SD.2) with exponent $f_d$ is not a hypothesis on this larger
module.  What descent identifies is exactly
\[
 \bigl(\mathscr T_d^{F^f}\bigr)^{F^{f_d}}
        =\mathscr T_d^{F^{f_d}},
 \tag{SD.11}
\]
whose exponent is given by (SD.10).  It gives no equality
between the exponent of $\mathscr T_d^{F^f}$ and that of
$\mathscr T_d^F$.  The full exponent is the maximum of the
conductor exponents, so proper-conductor growth is precisely
the part not excluded by this argument.


\paragraph{An actual improvement before the deck field: degree thirty-nine.}
Let $k=\mathbf F_4=\mathbf F_2(u)$, with $u^2+u+1=0$, and put
\[
 C/k:\quad Y^2+X^2Y=uX^5+uX^4+X^3+uX,
                       \qquad W=\infty.
 \tag{ND.12}
\]
There is a geometrically connected etale cyclic cover
$D\to C$ of degree $39$, with a $k$-point over $W$, whose
minimal deck field is $K=\mathbf F_{256}$. For its whole Prym
$B$ the exact exponents are
\[
       \exp B(k)[2^\infty]=2,
             \qquad \exp B(K)[2^\infty]=16.
 \tag{ND.13}
\]
Thus A45 applies over $k$, since $39>9$, whereas the actual
constant-deck test of A34 for this same cover over $K$ fails,
since $39\le961$. The whole new factor is ordinary of dimension
$38$, and its primitive conductor-$39$ factor already has
nonzero rational two-primary torsion over $k$. This is a
strict improvement for a specified pointed cover and field;
it asserts no absence of other successful covers.

\begin{proof}
\emph{1. The twist and the exact stable quotient.}
The curve (ND.12) is the quadratic twist of A27's model
$C_0:Y^2+X^2Y=uX^5+X^3+uX$. Adding $uX^4$ adds the constant
$u$ to the Artin--Schreier equation, and
$\operatorname{Tr}_{\mathbf F_4/\mathbf F_2}(u)=1$.
The genus is two and the geometric two-rank is one, as for
$C_0$. The point $W$ is rational and $K_C\sim2W$.

Use the field and integer encoding of A32:
\[
 K=\mathbf F_2[z]/(z^8+z^4+z^3+z+1),
       \quad u=188,\quad b=92,
       \quad b^2+b=u,\quad b^4+b=1.
 \tag{ND.14}
\]
The isomorphism to $C_0$ over $K$ is
$(X,Y)\mapsto(X,Y+bX^2)$. It fixes $W$.
The checked A27 polynomial gives $\#J_C(k)=18$, and A32 gives
$\#J_C(K)=65520=2^4\cdot3^2\cdot5\cdot7\cdot13$.
Write $G=J_C(K)$, identified with $J_{C_0}(K)$ by this map.
The exact reduced Mumford representatives
\[
 \Delta_9=(X^2+189,\,189X+80),\qquad
 \Delta_{13}=(X^2+80X+237,\,106X+207)
 \tag{ND.15}
\]
on $C_0$ have orders nine and thirteen. On the twist,
$\Delta_9=(X^2+u+1,(u+1)X+u)$ is already $k$-rational.
The finite computation verifies that these generate the entire
respective primary subgroups of $G$ and that the original
four-power Frobenius acts on them by $8$ modulo nine and $5$
modulo thirteen. The twist changes that Frobenius to its
negative on the Jacobian; the two actions are therefore $1$
modulo nine and $8$ modulo thirteen.

Define $a_9$ and $a_{13}$ by
\[
 [7280]Q=a_9(Q)\Delta_9,\qquad
 [5040]Q=a_{13}(Q)\Delta_{13},
\]
and let $a:G\to\mathbf Z/39$ be given by
$a\bmod3=a_9\bmod3$ and $a\bmod13=a_{13}$.
These are homomorphisms, and $a$ is surjective: multiplication
by the displayed cofactors kills the other primary factors
and is invertible on the selected primary subgroup.
For the twist Frobenius $\Phi$,
\[
                a(\Phi Q)=34a(Q)\pmod{39},
                        \operatorname{ord}_{39}(34)=4.
 \tag{ND.16}
\]
The action on the conductor-three quotient is trivial.
Its actions modulo thirteen and modulo thirty-nine both have
order four.

The finite etale isogeny $\Phi^4-1:J_C\to J_C$ is defined
over $k$ and has geometric kernel $G$.
Its subgroup $H=\ker a$ is Frobenius stable by (ND.16).
The quotient isogeny $J_C/H\to J_C$ is therefore defined
over $k$, has cyclic geometric kernel $G/H$ of order $39$,
and has deck Frobenius $34$. Pull it back through
$C\to J_C$, $R\mapsto[R-W]$. The result is geometrically
connected by Milne's full Abel--Jacobi fundamental-group
statement, Proposition~9.1 \cite{MilneJV}; the isogeny cover
is connected and etale. The zero point gives a $k$-point
over $W$. This constructs $D$ and proves that its minimal
deck field is precisely $K$.

\emph{2. Complete finite character data.}
Over $K$ this is the constant Lang quotient of $G$ by $H$.
For $c_j=\#\{R\in C(K):a([R-W])=j\}$, exact enumeration gives
\[
\begin{aligned}
 (c_0,\ldots,c_{12})&=(14,7,6,6,6,7,4,6,7,4,6,6,5),\\
 (c_{13},\ldots,c_{25})&=(15,7,6,6,6,5,6,6,5,6,6,6,7),\\
 (c_{26},\ldots,c_{38})&=(15,5,6,6,4,7,6,4,7,6,6,6,7).
\end{aligned}
 \tag{ND.17}
\]
The sum is $256$. Reducing the labels modulo three gives
$(74,91,91)$; reducing modulo thirteen gives exactly the A32
table (OR.9). The four rational points of (ND.12) over $k$ are
\[
              W,\quad(0,0),\quad(u,u+1),\quad(u,0).
\]
Their conductor-three labels are respectively $0,0,1,2$.
These are also their labels for the conductor-three cover
over $k$: extending constants from $k$ to $K$ multiplies a
degree-one Frobenius label by four, which is one modulo three.

For any nontrivial character $\chi$ of the indicated cyclic
quotient, the genus-two character polynomial at $W$ is
\[
         L_\chi(T)=1+c_\chi T+qT^2,
             \qquad c_\chi=\sum_R\chi(a([R-W])).
 \tag{ND.18}
\]
Here $q$ is the size of the field over which the deck group is
constant; this is the exact formula already proved in the
recovered A1 B.2 and used in A29 and A32. For conductor three
it gives $c_\chi=1$ over $k$, and $c_\chi=-17$ over $K$.
Thus $L_\chi(1)$ is respectively $6$ and $240$.
Since two is primitive modulo three, the integral module
argument (BE.8) applies to this genus-two degree-three cover.
It gives
\[
              \exp B_3(k)[2^\infty]=2,
                 \qquad \exp B_3(K)[2^\infty]=16.
 \tag{ND.19}
\]
Here the degree-thirteen quotient $D\to D_3$ identifies the
degree-three Prym's two-primary module with the actual $B_3$:
its pullback and norm have composite $[13]$ in both directions
on this conductor, and $[13]$ is invertible on every two-primary
group.
For conductor thirteen, the A32 character values give odd
$\#B_{13}(K)$, hence no rational two-primary torsion over
either field.

\emph{3. Exact exponent of the conductor-thirty-nine factor.}
Let $Z$ denote a primitive thirty-ninth root formally, and
reduce $257+\sum c_jZ^j$ modulo $\Phi_{39}(Z)$. The result is
\[
\begin{aligned}
 L_{39}(1)=2\bigl(&128+Z+Z^4-Z^6+Z^7-Z^9-Z^{12}\\
                &+Z^{14}+Z^{17}-Z^{18}+Z^{20}-Z^{21}\bigr).
\end{aligned}
 \tag{ND.20}
\]
Write $\beta$ for the reduction modulo two of the expression
in parentheses. It is a unit modulo $\overline\Phi_{39}$.
For a direct coefficient certificate, the polynomials
\[
\begin{aligned}
 U&=Z^{20}+Z^{17}+Z^9+Z^7+Z^6+Z^4,\\
 V&=Z^{17}+Z^{14}+Z^{11}+Z^{10}+Z^9+Z^8
                         +Z^6+Z^5+Z^2+Z+1
\end{aligned}
\]
satisfy $U\beta+V\overline\Phi_{39}=1$ in $\mathbf F_2[Z]$.
Thus in the full finite etale algebra
$\mathcal O_{39}=\mathbf Z_2[Z]/\Phi_{39}$, the value (ND.20)
is two times a unit. Its norm has two-adic valuation
$\varphi(39)=24$. Character factorization and the norm--pullback
isogeny therefore give
\[
                       v_2\bigl(\#B_{39}(K)\bigr)=24.
 \tag{ND.21}
\]

Moreover (ND.20) says $c_\chi\equiv1\pmod2$ for every
primitive character. Each corresponding coherent $H^1$
summand has dimension one, because the base genus is two.
As in A32, the reduced character polynomial is the determinant
of linear $256$-power Frobenius on that summand. It is $1+T$;
thus the scalar is one on every summand. The functorial
Artin--Schreier stable-part comparison makes $B_{39}$ ordinary
of dimension $24$ and all its geometric etale two-torsion
rational over $K$. The exact cited comparison inputs are
Deligne, Theorem~2.2(b) and formula~(2.4.1), with the hypotheses
of Section~2.1 \cite{DeligneModP}, and Katz,
Propositions~2.2.4--2.2.5 \cite{KatzCongruence}, as specified
in A32. The deck group is constant over $K$, so its distinct
odd-order characters split the coherent space; there is no
uncontrolled Jordan block between these one-dimensional lines.

Consequently $B_{39}(K)[2]$ already has order $2^{24}$.
By (ND.21) this is the entire two-primary group, proving
\[
                        \exp B_{39}(K)[2^\infty]=2.
 \tag{ND.22}
\]
Primitive semilinear descent, (SD.10), gives the same exponent
over $k$, since the conductor-thirty-nine action in (ND.16)
has order four. The conductor idempotents have odd denominators
and split the two-primary group of the actual whole Prym into
its conductor-$3$, $13$, and $39$ components. Taking the maximum
of (ND.19), (ND.22), and the conductor-thirteen exponent one
proves (ND.13). Each of the three geometric components is
ordinary, of dimensions $2$, $12$, and $24$, respectively;
hence $B$ is ordinary of dimension $38$. The cover itself has
genus $40$ and two-rank $39$.

Finally $g=2=p$ and $39>(2\cdot2-1)^2$, so A45 separates
the section of $W$ over the original field. The comparison
over the full deck field uses the actual larger exponent
$16$ and fails the inequality. This proves the stated strict
comparison for this cover.
\end{proof}

\paragraph{Finite computation scope.}
The table calculation \texttt{probe\_A27\_c39.py} uses the
retained, independently checked A32 point list and its binary
linear-interpolation Jacobian arithmetic. It recomputes the
three-primary projections, their cyclic order-nine structure,
the Frobenius actions, the twist's four rational points, the
new thirty-nine labels, the cyclotomic reductions, and the
displayed Bezout identity. Its JSON retains all $256$ point
images and the exact data above. This use of retained arithmetic
is not described as a newly independent Jacobian implementation.


The separate checker \texttt{verify\_A45\_conductor39\_cantor.py}
uses the previously checked extended-GCD Cantor arithmetic,
imports none of the linear-interpolation producer, and freshly
enumerates every point and both primary projections. It verifies
all $256$ new quotient labels, both complete cyclic subgroups,
the twisted Frobenius permutation of every point, the four
$k$-points, and the two exact cyclotomic unit identities.
Both implementations pass and agree on every point label.
A further rerun of that checker reproduces its JSON exactly;
that rerun is not counted as a third independent implementation.


\subsection{The integral sign removes the primitive-factor exception}
\label{sec:primitive-integral-sign}

\paragraph{Finding A46 (the rational primitive criterion in all characteristics).}
Let $C/k$ be a smooth projective geometrically connected curve of
genus $g\ge2$ over a finite field of characteristic $p$.
Fix $P\in C(k)$. Suppose that, over a finite extension $L/k$,
there is a geometrically connected finite etale cover of degree $n$
\[
               D\longrightarrow C_L
\]
which is geometrically Galois with cyclic deck group, where
$n>1$ is odd and $p\nmid n$, and a point $P'\in D(L)$ above $P$.
The geometric deck group need not be constant over $L$.
Let $P_n\subset J_D$ be the actual primitive conductor-$n$
abelian subvariety, defined over $L$ by the integral projector
specified below. If
\[
                         P_n(L)[p^\infty]=0,
 \tag{IS.1}
\]
then the point section of $P$ is distinct from that of every
different point of $C(k)$ in the geometrically pro-prime-to-$p$
arithmetic fundamental group. In fact the same criterion over
$L$ separates $P$ from every distinct point of $C(L)$.

Consequently, if a geometrically cyclic etale cover of odd degree
prime to $p$ has primitive factor of geometric $p$-rank zero,
then it gives the point-section separation conclusion for the
base curve after the usual descent and pointing of the cover.
This geometric sufficient condition applies in every genus
$g\ge2$ and every positive characteristic.
No hypothesis about an elliptic quotient or ordinarity is needed
in either assertion.

\paragraph{Prior result and exact strengthening.}
The recovered \texttt{ANABELIAN\_SUCCESSOR} final row A.14,
working identifier AC.11, supplies the composite-projector
argument with an exclusion of a separable double map of the base
to a genus-one curve. Findings A15 and A18 retained the complete
small-field argument, and removed that exclusion in specified
geometric cases. The new step here uses the signs of the
\emph{integral principal divisor} to show that the degree-two
outcome itself is impossible. It removes the rational-only
qualification in A15/A18 and does not use their
Deuring--Shafarevich computation. We give the full proof, including
the descent of the actual primitive subvariety, independently of
any nonconstant-deck statement elsewhere in this report.

\begin{proof}
\emph{1. The canonical integral primitive projector.}
Over an algebraic closure of $L$, choose a generator $\sigma$
of the cyclic deck group $G$. In $\mathbf Q[G]$, let $e_n$
be the idempotent which is one on the faithful characters and
zero on the characters of proper conductor. Explicitly,
\[
 e_n=\frac1n\sum_{a=0}^{n-1}c_n(a)\sigma^a,
 \qquad
 c_n(a)=\sum_{d\mid\gcd(n,a)}d\,\mu(n/d)\in\mathbf Z.
 \tag{IS.2}
\]
Here the formula for $a=0$ uses $\gcd(n,0)=n$.
To verify the formula, write $\zeta$ for a primitive complex
$n$th root. Inclusion--exclusion on the units modulo $n$ gives
\[
 \sum_{u\in(\mathbf Z/n)^\times}\zeta^{-au}
       =\sum_{d\mid\gcd(n,a)}d\,\mu(n/d)=c_n(a).
\]
The finite Fourier sum therefore makes $e_n$ the stated
idempotent. Put
\[
        p_n=ne_n\in\mathbf Z[G],\qquad
        P_n=\operatorname{im}(p_n:J_D\longrightarrow J_D).
 \tag{IS.3}
\]
The coefficients satisfy $c_n(ua)=c_n(a)$ for every unit $u$
modulo $n$. Arithmetic Galois acts on $G$ by such automorphisms,
so $p_n$ is an endomorphism defined over $L$ and its image is an
abelian subvariety defined over $L$. This is an actual image,
not an unspecified isogenous representative.

Choose the geometric group-ring operator
\[
                 A=\prod_{\ell\mid n}(1-\sigma^{n/\ell}),
 \tag{IS.4}
\]
where the product is over distinct prime divisors. On the
character $\sigma\mapsto\zeta^i$, it acts by
$\prod_{\ell\mid n}(1-\zeta^{i n/\ell})$, nonzero exactly
when $\gcd(i,n)=1$. Hence
\[
                  Ae_n=A,\qquad Ap_n=nA
                       \quad\hbox{in }\mathbf Z[G].
 \tag{IS.5}
\]
The rational identity is integral because $\mathbf Z[G]$
is torsion-free. The images of $A$ and $p_n$ over the algebraic
closure are the same abelian subvariety: $A$ vanishes on every
proper conductor and is invertible on the faithful factor of
$\mathbf Q[G]$. Thus (IS.3) agrees with the primitive factor
used in the constant-deck arguments.

\emph{2. A collision gives a principal reduced projector divisor.}
Suppose that $Q\ne P$ had the same point section, and restrict
the collision to $L$. The covering formulation gives a point
$Q'\in D(L)$ above $Q$, with the same lifted section as $P'$,
and therefore
\[
                        \delta=[Q'-P']
                              \in J_D(L)[p^\infty].
 \tag{IS.6}
\]
For completeness, a finite etale cover which is geometrically
cyclic of degree prime to $p$ is a permitted neighbourhood of
the geometrically pro-prime-to-$p$ group. Equality of sections
lifts to the same pointed neighbourhood. The subgroup topology
is compatible: taking the core of a geometric prime-to-$p$
quotient of the normal index-$n$ subgroup gives a quotient
whose order divides a product of prime-to-$p$ orders and $n$.
The lifted sections have equal prime-to-$p$ abelian Kummer
classes. On a finite-field Jacobian the kernel of that Kummer
map on rational points is exactly the $p$-primary subgroup,
which gives (IS.6). This is the same covering/Kummer implication
used in A1 and the recovered B.13--B.14 arguments.

Now $p_n\delta$ belongs to $P_n(L)[p^\infty]$, so (IS.1)
makes it zero. Equivalently, multiplication by $n$ is an
automorphism on the finite $p$-primary group and
\[
 b=e_n\delta:=[n]^{-1}(p_n\delta)
                      \in P_n(L)[p^\infty]
 \tag{IS.7}
\]
is well-defined and zero. Equation (IS.5) gives
$nA\delta=Ap_n\delta=0$. Since $A\delta$ is $p$-primary and
$p\nmid n$, it follows that
\[
                               A\delta=0.
 \tag{IS.8}
\]
We now use (IS.8) over the algebraic closure. Although $A$ need
not descend to $L$, no constant-field extension has been used
to preserve or replace the rational condition (IS.1).

Put $s=\omega(n)$ and $d=2^s$. The $2^s$ subset products of the
elements $\sigma^{n/\ell}$ are distinct: their factors have
orders in distinct prime-primary components of the cyclic
group. Thus the expansion of $A$ has exactly $d$ coefficients,
each $1$ or $-1$. Its augmentation is zero, so half have each
sign. The two free deck orbits of $P'$ and $Q'$ are disjoint,
because their images on $C$ are distinct. Consequently
$A(Q'-P')$ is a nonzero integral divisor with exactly $d$
simple zeros and $d$ simple poles. Equation (IS.8) makes it
principal. Over $\overline L$ choose $f$ with
\[
                       \operatorname{div}(f)=A(Q'-P').
 \tag{IS.9}
\]
It has degree $d$ and is separable in every characteristic:
a $p$th power has all divisor multiplicities divisible by $p$,
whereas (IS.9) has simple zeros and poles.

\emph{3. The field generated by the deck conjugates.}
Write
\[
 K=\overline L(D),\quad
 M=\overline L(\gamma f:\gamma\in G),\quad e=[K:M].
\]
Every conjugate has degree $d$. The Castelnuovo--Severi
consequence in Poonen, Proposition~2.1, Lemma~2.2 and
Corollary~2.3, gives
\[
                    e\mid d,
                    \qquad g(M)\le(d/e-1)^2.
 \tag{IS.10}
\]
The exact applicability is Corollary~2.3: over the perfect field
$\overline L$, the subfield $M$ is generated by finitely many
elements all of degree $d$, so it is $d$-controlled in the
definition preceding Lemma~2.2. Thus there is a positive integer
$u$ with $[K:M]=d/u$ and $g(M)\le(u-1)^2$; substituting
$u=d/e$ gives (IS.10)
\cite[Proposition~2.1, Lemma~2.2 and Corollary~2.3]{Poonen2007Gonality}.

For every odd $n>1$,
\[
                         n>(2^{\omega(n)}-1)^2=(d-1)^2.
 \tag{IS.11}
\]
For $s=1,2$ use $n\ge3,15$, respectively. For $s\ge3$,
$n\ge3\cdot5^{s-1}>4^s>(2^s-1)^2$; the middle inequality
holds at $s=3$ and its ratio increases with $s$.
Since the original cover is etale,
$g(D)=n(g(C)-1)+1\ge n+1$.
Thus (IS.10)--(IS.11) exclude $e=1$, and $e>1$ is a power of two.

Let $Z$ be the smooth projective curve with function field $M$,
and let $v:D_{\overline L}\to Z$ be the corresponding morphism
of degree $e$. The group $G$ preserves $M$ and acts faithfully
on $Z$: a kernel would have order dividing both $n$ and $e$.
This action is free as well. Indeed, a nontrivial stabilizer
contains an element of some odd prime order $\ell\mid n$.
That element acts freely on $D$ and preserves the fibre divisor
$v^*z$ above the point it fixes on $Z$. Its orbits in that fibre
have length $\ell$, with constant multiplicities along each
orbit, forcing $\ell\mid e$. This contradicts $e$ being a
power of two.

\emph{4. A reduced fibre and its integral signs give the contradiction.}
Write $f=h\circ v$ with $h\in\overline L(Z)$.
Choose a zero $z$ of $h$. The fibre divisor $v^*z$ is reduced:
every multiplicity in it, multiplied by the positive order of
$h$ at $z$, is a zero multiplicity of $f$, which is one.
It therefore consists of $e$ distinct geometric points, all
lying in the two original deck orbits of $P',Q'$.
Each such orbit maps injectively to $Z$, because the induced
$G$-action there is free. Hence $e\le2$, and so $e=2$.

The extension $K/M$ is separable, since it is intermediate in
the separable extension $K/\overline L(f)$. Its quadratic
involution $\tau$ is normalized by $G$, which preserves $M$.
The group of order two has no nontrivial automorphism, so
$\tau$ commutes with $G$.
The reduced fibre just chosen contains one point of each
original orbit. Its two points are exchanged by $\tau$;
therefore $\tau$ exchanges the two entire $G$-orbits.
For some $j$ modulo $n$,
\[
                  \tau P'=\sigma^jQ',\qquad
                  \tau Q'=\sigma^{-j}P'.
 \tag{IS.12}
\]
The second equality follows from the first, $\tau^2=1$,
and commutation with $G$.

Since $f\in M$, its divisor is $\tau$-invariant. Applying
$\tau$ to (IS.9), and comparing its coefficients on the free
$P'$-orbit, gives
\[
                         A\sigma^{-j}=-A
                                \quad\hbox{in }\mathbf Z[G].
 \tag{IS.13}
\]
This comparison is literal: the map from $\mathbf Z[G]$
to the integral divisors supported on that orbit,
$\sum a_\gamma\gamma\mapsto\sum a_\gamma\gamma P'$,
is injective. If $t$ is the order of $\sigma^j$, then $t$
is odd. Iterating (IS.13) gives
\[
                           A=(-1)^tA=-A.
\]
But $\mathbf Z[G]$ is torsion-free and $A\ne0$, as its
subset-product expansion already shows. This is impossible.
The argument uses integer divisor coefficients, so the two
signs remain distinct also in characteristic two.

The assumed collision is therefore excluded. In particular,
there is no need to calculate the genus of $Z$, to descend an
elliptic quotient, or to impose an additional condition on it.

\emph{5. The geometric consequence and its constant-field scope.}
A primitive factor of geometric $p$-rank zero has no nonzero
geometric points of $p$-power order. Thus its rational
$p$-primary subgroup vanishes over every finite constant
extension. Descend the geometric cyclic cover and its deck
action to a finite field, and extend constants once more to
rationalize a chosen point above the prescribed base point.
Condition (IS.1) then holds, so the result just proved applies.
For an initially rational criterion, in contrast, the proof
uses its specified field $L$ throughout the arithmetic step;
it does not enlarge that field to make the deck group constant.
\end{proof}

\paragraph{Effect on the earlier gates.}
A15's rational primitive-factor criterion no longer needs
rank one or exclusion of an elliptic double quotient. A18's
geometric criterion extends to every characteristic for odd
cyclic degrees prime to that characteristic. In characteristic
two and genus two, the already proved actual-theta orbit
coverage test still supplies geometric primitive rank zero,
and hence remains a sufficient input. The proof above does
not assert existence of such a cover or such theta incidence
for every prescribed curve. The closure here is the earlier
\emph{rational-only primitive-factor qualification}, not the
universal incidence gate.


\subsection{A bounded exponent on the actual primitive factor}
\label{sec:bounded-primitive-exponent}

\paragraph{Finding A47 (bounded primitive torsion and its exact obstruction).}
Let $C/k$ be a smooth projective geometrically connected curve
of genus $g\ge2$ over a finite field of characteristic $p$.
Let $P\in C(k)$, and let $D\to C$ be a geometrically connected
finite etale cover of odd degree $n>1$ prime to $p$, geometrically
Galois with cyclic group. Assume that $P'\in D(k)$ lies above $P$.
The geometric deck group need not be constant over $k$.
Let $P_n\subset J_D$ be the actual primitive conductor-$n$
subvariety defined by the integral projector (IS.2)--(IS.3), and put
\[
 s=\omega(n),\qquad c=2^s,\qquad
 M_n=\exp\bigl(P_n(k)[p^\infty]\bigr),
 \tag{BP.1}
\]
with exponent one for the zero group. Suppose that
\[
                         n>(cM_n-1)^2.
 \tag{BP.2}
\]
All section statements below concern the geometrically
pro-prime-to-$p$ arithmetic fundamental group.
Then a collision between the point sections of $P$ and a distinct
$Q\in C(k)$ forces a separable map to an elliptic curve
\[
 C_{\bar k}\longrightarrow E,
 \qquad \deg(C_{\bar k}/E)=p^a=e>1,
 \qquad e\mid M_n,\quad e\le g-1,
 \tag{BP.3}
\]
totally ramified at both $P$ and $Q$. Moreover the given geometric
cover $D_{\bar k}\to C_{\bar k}$ is the pullback of an etale
cyclic cover $Z\to E$ of degree $n$ between genus-one curves.
Consequently (BP.2) separates the section of $P$ from all other
rational point sections whenever $2\le g\le p$, or whenever
there is no map (BP.3) totally ramified at $P$.

If $M_n=1$, the conclusion is separation in every genus
$g\ge2$: (BP.3) is then impossible, and (BP.2) is automatic
by (IS.11). Thus the rational primitive criterion A46 is
recovered, including its nonconstant-deck scope. No existence
of a cover satisfying (BP.2) on every prescribed curve is asserted.

\begin{proof}
\emph{1. The principal projector divisor and its exact order.}
Suppose that the sections of $P$ and $Q\ne P$ collide.
The same pointed-cover and finite-field Kummer implication as
in A46 gives $Q'\in D(k)$ with
\[
                 \delta=[Q'-P']\in J_D(k)[p^\infty].
 \tag{BP.4}
\]
Choose a geometric generator $\sigma$ and put
\[
 A=\prod_{\ell\mid n}(1-\sigma^{n/\ell}),\qquad
 p_n=ne_n\in\mathbf Z[\langle\sigma\rangle],
 \tag{BP.5}
\]
where $e_n$ is the faithful-character projector. The actual
endomorphism $p_n$ is defined over $k$, and $Ap_n=nA$ by
(IS.2)--(IS.5). Hence
\[
 M_np_n\delta=0,\qquad nM_nA\delta=0,
                         \qquad M_nA\delta=0.
 \tag{BP.6}
\]
The last implication uses that $A\delta$ is $p$-primary and
$p\nmid n$. In particular the geometric operator $A$ is used
only after a rational bound has been obtained on the actual
subvariety over $k$; the field has not been enlarged.

Let $m$ be the exact order of $A\delta$, with $m=1$ for zero.
Then $m$ is a power of $p$ and $m\mid M_n$. The expansion of
$A$ has $c=2^s$ distinct coefficients, each $1$ or $-1$, half
of each sign, as in A46. The two deck orbits of $P'$ and $Q'$
are disjoint and free. There is therefore a function $f$ on
$D_{\bar k}$ with
\[
 \operatorname{div}(f)=mA(Q'-P'),\qquad \deg f=d=cm.
 \tag{BP.7}
\]
Every nonzero divisor coefficient is $m$ or $-m$.
The function is separable. For $m>1$, a $p$th power would
make $(m/p)A(Q'-P')$ principal, contradicting exact order;
for $m=1$, the simple poles exclude a $p$th power directly.

\emph{2. The controlled field and free cyclic action.}
Set
\[
 K=\bar k(D),\qquad
 F=\bar k(\gamma f:\gamma\in\langle\sigma\rangle),
                  \qquad e=[K:F].
\]
Let $Z$ be the smooth projective curve with field $F$ and
$v:D_{\bar k}\to Z$ the associated map. The field is generated
by finitely many degree-$d$ elements over the perfect field
$\bar k$. Poonen's Corollary~2.3 and its definition of
$d$-controlled therefore give exactly
\[
                 e\mid d,\qquad g(Z)\le(d/e-1)^2
 \tag{BP.8}
\]
\cite[Proposition~2.1, Lemma~2.2 and Corollary~2.3]{Poonen2007Gonality}.
Since $g(D)=n(g-1)+1$ and (BP.2) implies $n>(d-1)^2$,
one has $e>1$. Moreover $e\mid2^sp^b$ for some $b$, so
$\gcd(e,n)=1$. The extension $K/F$ is separable, being
intermediate in $K/\bar k(f)$.

The cyclic group preserves $F$ and acts faithfully on $Z$:
the order of a kernel divides both $e$ and $n$.
Its action is free. If a point stabilizer contained an element
of prime order $\ell\mid n$, its free action on $D$ would
partition the fibre divisor above that point into orbits of
length $\ell$ with equal multiplicities on each orbit.
This would give $\ell\mid e$, a contradiction.
Consequently the quotient $E=Z/\langle\sigma\rangle$ is etale
of degree $n$. The bound $0\le g(Z)<n$ and
\[
                         g(Z)-1=n(g(E)-1)
\]
force $g(Z)=g(E)=1$. In particular
\[
                 \deg\operatorname{Diff}(v)=2n(g-1).
 \tag{BP.9}
\]

\emph{3. Integral signs exclude identical image orbits.}
The maps from each of the two original deck orbits into $Z$
are injective, because its cyclic action is free. Their images
are either identical or disjoint.
Suppose first that they are identical, and choose $j$ with
\[
                         v(P')=\sigma^jv(Q').
\]
Let $r_P,r_Q$ be the ramification indices of $v$ at $P',Q'$;
they are positive and constant along the respective deck orbits.
Write $f=h\circ v$. On the image orbit based at $v(Q')$, put
\[
 H=\sum_{i=0}^{n-1}
       \operatorname{ord}_{\sigma^iv(Q')}(h)\,\sigma^i
                         \in\mathbf Z[\langle\sigma\rangle].
\]
Comparing the literal integer coefficients of (BP.7) on the
two free orbits gives
\[
           mA=r_QH,\qquad -mA=r_P\sigma^{-j}H.
\]
Eliminating $H$ and cancelling the positive integer $m$ in
the torsion-free integral group ring yields
\[
                        r_QA=-r_PA\sigma^{-j}.
 \tag{BP.10}
\]
If $t$ is the order of $\sigma^j$, then $t$ is odd. Iteration
gives $r_Q^tA=-r_P^tA$, which is impossible: $r_Q^t+r_P^t$
is positive and $A\ne0$. This argument does not assume a
quadratic map, equal indices, or reduced fibres. Its signs
are integer divisor signs, also in characteristic two.

\emph{4. Disjoint image orbits force the wild elliptic obstruction.}
The images must therefore be disjoint. Take any point in the
support of (BP.7). A fibre above its image can contain only
points in that divisor's support, since it is a zero or pole
fibre of $h$. It contains at most one point from either original
orbit, by injectivity, and cannot contain one from the other
orbit, by disjointness. It is consequently a singleton fibre
of total ramification index $e$. Its divisor multiplicity
$m$ is a multiple of $e$, so
\[
                         e\mid m\mid M_n.
 \tag{BP.11}
\]
Since $e>1$, it is a nontrivial power of $p$.
The identity coefficient of $A$ is one; in particular $P'$
and $Q'$ themselves belong to the support. Equivariance makes
every point of their two disjoint orbits totally ramified
with index $e$.

At such a point, a local parameter expression
$u=a_et^e+\cdots$ has nonzero derivative by separability,
but its degree-$e$ leading derivative vanishes since $p\mid e$.
Thus the different exponent is at least $e$; this is the local
different formula in the Riemann--Hurwitz derivation
\cite[Lemma~53.12.4, Tag~\texttt{0C1F}]{StacksRH}.
The two orbits have $2n$ points in total, and (BP.9) gives
\[
                         2ne\le2n(g-1),
                    \qquad p\le e\le g-1.
 \tag{BP.12}
\]
This is impossible for $2\le g\le p$, and is also impossible
when $M_n=1$ by (BP.11).

\emph{5. The obstruction descends to the given base curve.}
Let $G=\langle\sigma\rangle$.
The field $K/K^G$ is Galois with group $G$. Its subgroup fixing
$FK^G$ is the kernel of the faithful action on $F$, so
\[
             K=FK^G,\qquad \bar k(C)=K^G,
                 \qquad \bar k(E)=F^G.
 \tag{BP.13}
\]
It follows that $[\bar k(C):\bar k(E)]=e$. This extension is
separable, since $K/F$ and $F/F^G$ are separable. Choosing an
origin on the genus-one curve $E$ makes it an elliptic curve.
The etale fibre product $C_{\bar k}\times_E Z$ is connected
and has function field $K$, hence is $D_{\bar k}$.
Ramification indices are unchanged by the etale base change
$Z\to E$. The total ramification at $P',Q'$ therefore gives
total ramification at $P,Q$ in the degree-$e$ map $C_{\bar k}\to E$.
This proves every assertion in (BP.3) and the theorem.
\end{proof}

\paragraph{Why the remaining configuration is not a sign contradiction.}
On an elliptic curve a fixed-point-free finite cyclic action
consists of translations. Indeed an automorphism with a nontrivial
origin-preserving part has a fixed point, since one minus that
part is a surjective nonzero endomorphism. Translations act
trivially on $\operatorname{Pic}^0$. Thus, in the configuration
of (BP.3), $A(y-x)$ is principal for any two points $x,y$ of
$Z$. If $x=v(P')$ and $y=v(Q')$, total ramification gives
\[
                   v^*A(y-x)=eA(Q'-P').
 \tag{BP.14}
\]
The remaining positive-power principal divisor is therefore
compatible with this elliptic pullback geometry. This explains
why the proof retains (BP.3) when $M_n>1$; it does not assert
that such a divisor arises from a full point-section collision,
or that an arbitrary elliptic pullback satisfies (BP.2).

\paragraph{The precise remaining witness requirement.}
For the characteristic-two genus-two part of U01 and
\texttt{Q-HOOLEY}, including the ordinary real-multiplication
targets of A33, A46's sufficient input $M_n=1$ can now be
replaced by a pointed cover satisfying
\[
                 M_n<\frac{\sqrt n+1}{2^{\omega(n)}}.
 \tag{BP.15}
\]
The primitive factor may have positive geometric two-rank and
nonzero rational two-primary torsion. This changes the sufficient
input to the point-section route; it does not solve the literal
primitive-rank-zero existence problem or the rational-unit
witness questions in the recovered A1 rows B.24, B.5 and B.32.

Compared with A45, the whole-Prym exponent $M$ has been replaced
by $M_n$, at the cost of replacing $2$ by $2^{\omega(n)}$ in
the bound. The primitive test has a smaller threshold when
$M>2^{\omega(n)-1}M_n$; in other cases the two numerical tests
need not imply one another. For prime-power degree the coefficient
is still two, and only the top primitive conductor is measured.
No general estimate supplying (BP.15) for every prescribed curve
or point follows from this proof or from Hooley's prime-density
theorem. In particular, arbitrary constant extension is not a
way to improve a rational torsion exponent.


\paragraph{The coarse Weil estimate still does not supply A47.}
This is a limitation of a proposed estimate, not a lower bound
for the actual exponent. Work in genus two and characteristic
two. First suppose that the cyclic deck group of a pointed
cover of degree $n>1$ is constant over $\mathbf F_q$, and put
\[
 B=(\sqrt q+1)^2,\qquad f=\operatorname{ord}_n(2),
                         \qquad r=\varphi(n)/f\ge1.
 \tag{BPW.1}
\]
The pointed abelian-cover classification, equivalently the
Abel-map fundamental-group statement and the finite-field
Lang quotient, makes its cyclic group a quotient of
$J_C(\mathbf F_q)$ \cite[Proposition~9.1]{MilneJV}.
The Weil bound therefore gives
\[
                  n\mid\#J_C(\mathbf F_q),\qquad n\le B^2
 \tag{BPW.2}
\]
\cite[Theorem~19.1]{MilneAV}.

The primitive factor has dimension $\varphi(n)$: the
degree-$d$ quotient has genus $d+1$, and the cyclic conductor
decomposition of its Prym dimensions is
$d-1=\sum_{a\mid d,\ a>1}\varphi(a)$.
Since the deck maps are rational, $P_n(\mathbf F_q)[2^\infty]$
is a module over
$\mathcal O_n=\mathbf Z_2[Z]/\Phi_n(Z)$.
This is a product of unramified discrete valuation rings of
common residue degree $f$. If a component contains an element
of order $M_n=2^a$, its cyclic module has annihilator exactly
$(2^a)$ and cardinality $2^{fa}=M_n^f$.
Taking a component of maximal exponent and applying the Weil
bound to the actual primitive abelian subvariety yields
\[
 M_n^f\le\#P_n(\mathbf F_q)\le B^{\varphi(n)},
                         \qquad M_n\le B^r.
 \tag{BPW.3}
\]
The unrounded substitution of (BPW.3) in (BP.2) would require
$n>(2^{\omega(n)}B^r-1)^2$. Its right side exceeds $B^2$,
since $2^{\omega(n)}\ge2$, $r\ge1$, and $B>1$.
It is incompatible with (BPW.2).

Rounding down to the largest power of two allowed by (BPW.3)
does not rescue this estimate. For $q\ge8$, that rounded
upper bound is at least $q$, because $q$ is itself a power
of two and $B^r\ge B>q$. But
\[
           2q-1-B=q-2\sqrt q-2>0\qquad(q\ge8).
 \tag{BPW.4}
\]
The inequality holds at $q=8$ since $6>4\sqrt2$, and the
expression increases for $q\ge8$. For $q=2$, the largest
power of two at most $B=3+2\sqrt2$ is four, and
$2\cdot4-1=7>B$. For $q=4$, it is eight since $B=9$,
and $2\cdot8-1=15>B$.
For larger $r$ the rounded upper bound can only increase.
Thus even the rounded substitution in (BP.2) always asks for
$n>B^2$, contradicting the available bound.

For a nonconstant geometric deck group, first pass to its own
minimal deck field. The actual primitive exponent is unchanged
by (SD.10); the pointed cover then satisfies the constant-deck
hypotheses just used. No full cyclotomic action on the rational
group before that passage is assumed. The same estimate fails
over the deck field as well.

This leaves the exact finite-field primitive exponent as a
missing input. Geometric rank information or local character
values can make it much smaller than (BPW.3); the argument
does not exclude such covers. Hooley's candidate-order density
does not supply that additional information. This supplements
A34's earlier coarse-bound obstruction and is not a new
universal nonexistence or main-route supply result.


\subsection{Good unpointed twists and the prescribed-point obstruction}
\label{sec:twists-pointing-obstruction}

\paragraph{Finding A48 (the exact limit of the constant-twist supply argument).}
Let $C/k$ be a smooth projective geometrically connected curve of
genus $g\ge2$ over a finite field of characteristic $p$.
Let $D\to C$ be a geometrically connected finite etale cyclic
cover of odd degree $n>1$ prime to $p$, with constant deck group
$G=\langle\sigma\rangle$ over $k$. Let $P_n$ be its actual
primitive conductor-$n$ abelian subvariety, as in A46, and put
\[
                       f_n=\dim_{\mathbf F_p}P_n[p](\bar k).
\]
Among the $n$ constant twists $D_t\to C$, indexed by
$t\in\mathbf Z/n$, at least
\[
                n-f_n\ \ge\ n-(g-1)\varphi(n)
 \tag{TW.1}
\]
have primitive factor with no nonzero rational $p$-power
torsion. In genus two the lower bound is at least
$n-\varphi(n)>0$, for every such degree $n$.
The conclusion concerns \emph{unpointed} covers.

For a specified $P\in C(k)$, exactly one of these twists is
$P$-split. If that pointed twist has nonzero rational primitive
$p$-primary torsion, then every other twist has nonzero rational
primitive $p$-primary torsion over every constant extension in
which its fibre above $P$ acquires a rational point.
The same no-improvement assertion holds for twists of a fixed
geometrically cyclic cover with nonconstant deck group, starting
from its pointed model over $k$.
Thus good unpointed twists, followed by a pointing extension,
do not by themselves supply the witness required by A46.

\begin{proof}
\emph{1. Bound the primitive character multiplicities.}
Over $\bar k$, the coherent eigensheaves of the cyclic cover are
nontrivial degree-zero line bundles for every nontrivial
character. Riemann--Roch gives
\[
             \dim_{\bar k}H^1(C_{\bar k},L_\chi)=g-1.
 \tag{TW.2}
\]
The functorial Artin--Schreier stable-part comparison embeds
$H^1_{\mathrm{et}}(D_{\bar k},\mathbf F_p)\otimes\bar k$
as the Frobenius-stable part of $H^1(D_{\bar k},\mathcal O_D)$.
We use the same precise comparison as A32: Katz,
Proposition~2.0, Corollary~2.1 and Propositions~2.2.4--2.2.5
\cite{KatzCongruence}. These statements apply to proper schemes
in every cohomological degree, with no restriction on genus or
ordinarity. The comparison is compatible with deck automorphisms.
The functorial identification of $H^1_{\mathrm{et}}(D_{\bar k},
\mathbf F_p)$ with the dual of $T_pJ_D/pT_pJ_D$, and
$T_pJ_D/pT_pJ_D\simeq J_D[p](\bar k)$, shows that passing to
the $p$-torsion representation only inverts characters. Therefore,
in
\[
 V=P_n[p](\bar k),\qquad
 V\otimes_{\mathbf F_p}\overline{\mathbf F}_p
                    =\bigoplus_{\chi\ {\rm faithful}}V_\chi,
\]
each multiplicity $m_\chi=\dim V_\chi$ is at most $g-1$.
In particular $f_n=\sum_\chi m_\chi\le(g-1)\varphi(n)$.
The sum is a direct sum because $p\nmid n$; no ordinary-rank
hypothesis has been imposed.

\emph{2. Count the twists with a primitive fixed vector.}
Let $F$ denote arithmetic Frobenius acting on geometric points.
It commutes with $G$, since the deck transformations are defined
over $k$. A constant twist $D_t$ has Frobenius descent action
\[
                           F_t=\sigma^t F.
 \tag{TW.3}
\]
These $n$ twists exist by finite Galois descent: the constant
cocycle sends Frobenius to $\sigma^t$ and is trivial after a
finite constant extension. They are distinct as forms over $C$,
whose full geometric automorphism group over $C$ is $G$.
Their primitive abelian subvarieties have the same geometric
$p$-torsion space $V$ and the twisted action (TW.3).

On $V_\chi$, let $\alpha$ run through the eigenvalues of $F$,
counted with algebraic multiplicity. The action of $F_t$ has
an eigenvalue one on this summand exactly when
\[
                         \chi(\sigma)^t\alpha=1
 \tag{TW.4}
\]
for at least one such $\alpha$. Since $\chi$ is faithful,
$t\mapsto\chi(\sigma)^t$ is injective. Each $\alpha$ therefore
forbids at most one value of $t$. This remains true when the
matrix of $F$ has a nontrivial Jordan block: a fixed vector
exists exactly when one is an eigenvalue. Summing this union
bound over the faithful characters gives at most $f_n$ bad
twists. A finite rational $p$-primary group is nonzero if and
only if it has an element of order $p$. Thus the complementary
twists have zero rational primitive $p$-primary group, proving
(TW.1).

In particular every geometric cyclic cover of a genus-two
curve, after descending it and making its deck group constant,
has a twist with this unpointed rational vanishing property.
This conclusion uses neither an Artin-prime statement nor GRH.
It makes no assertion about the rational fibre above a
prescribed base point.

\emph{3. The point chooses one twist.}
Choose a geometric point $R$ of the fibre above $P$, and write
$F(R)=\sigma^a R$. Because $F$ commutes with $G$, its action
on this entire fibre is translation by $a$. The action for
$D_t$ is translation by $a+t$. It has a fixed point if and
only if $a+t=0$ modulo $n$, in which case every point of the
fibre is rational. Thus the unique $P$-split twist is
$D_P=D_{-a}$, with
\[
                              F_P=\sigma^{-a}F.
\]
Over $k_m=\mathbf F_{q^m}$, the fibre of $D_t$ above $P$
has a rational point exactly when $m(a+t)=0$ modulo $n$.
In that situation
\[
 F_t^m=\sigma^{m(a+t)}F_P^m=F_P^m.
 \tag{TW.5}
\]
The same geometric cover with these equal descent actions is
the same form over $k_m$. Its actual primitive subvariety is
identified as well. Consequently
\[
             P_{n,P}(k)[p^\infty]
                      \hookrightarrow P_{n,t}(k_m)[p^\infty].
 \tag{TW.6}
\]
Any nonzero torsion in the pointed model therefore returns as
soon as the other twist can be used as a pointed cover.
This proves the stated obstruction.

\emph{4. The same cancellation for nonconstant deck groups.}
Start now with a geometrically cyclic pointed cover $D_P/k$
whose deck group is not required to be constant. Fix its chosen
point $R$ over $P$ as the origin in the geometric fibre, and
write
\[
 F_P\sigma F_P^{-1}=\sigma^a,\qquad a\in(\mathbf Z/n)^\times.
\]
For a twist represented by $F_t=\sigma^t F_P$, define
$N_m(a)=1+a+\cdots+a^{m-1}$ in $\mathbf Z/n$. On the
geometric fibre, the $m$th Frobenius power acts by
\[
                         x\longmapsto a^m x+N_m(a)t.
 \tag{TW.7}
\]
If this fibre has a rational point over $k_m$, there is an $x$
with $N_m(a)t=(1-a^m)x$. Conjugating the full descent action
by the geometric deck translation $\sigma^x$ then gives
\[
 \sigma^{-x}F_t^m\sigma^x
       =\sigma^{N_m(a)t-(1-a^m)x}F_P^m=F_P^m.
 \tag{TW.8}
\]
Thus the two covers are isomorphic over $C_{k_m}$, and their
canonical primitive subvarieties are isomorphic over $k_m$.
The nonzero rational $p$-primary subgroup of the pointed model,
if present, persists under this isomorphism and the constant
extension. This proves the final assertion. Unlike (TW.1),
no count of nonconstant twists is being asserted.
\end{proof}

\paragraph{What remains for a prescribed-point witness.}
For a fixed geometric cyclic cover and a prescribed point,
twisting can relocate a Frobenius eigenvalue one while
simultaneously relocating the rational fibre. Equations
(TW.5)--(TW.8) show why that operation cannot create the
pointed vanishing condition from a model which fails it.
This does not rule out changing the geometric cyclic cover,
its conductor, or the base curve. Nor does it rule out direct
Fourier verification of a successful pointed cover, as in the
retained examples.

The good-twist bound is therefore a positive unpointed supply
statement with an exact pointing obstruction. It does not
close the universal rational primitive-factor existence gate.
In particular, an average nonvanishing assertion, or a positive
density of arithmetic candidate primes under number-field
GRH, cannot be inserted into this argument as proof that the
unique prescribed-point twist is among the good ones.
A further condition controlling that particular twist, or a
construction varying the geometric cover while keeping the
pointing, is still needed.


\paragraph{Why pullback iteration does not improve an existing primitive factor.}
Let $X$ be a smooth projective geometrically connected curve
of genus at least two. Let $D\to X$ be a geometrically connected
etale cover of odd degree $n>1$ prime to $p$, geometrically
Galois with cyclic group, and let $a:Y\to X$ be a finite etale
cover of degree $d$ prime to $p$.
Assume that all curves and maps are defined over the same finite
field $k$ of characteristic $p$, and that
\[
                         D'=D\times_XY
 \tag{PI.1}
\]
is geometrically connected. Then $D'\to Y$ has the same
geometric cyclic degree $n$. Let $P_n$ and $P'_n$ be the
actual primitive conductor-$n$ subvarieties in $J_D,J_{D'}$.
For every finite extension $k'/k$, pullback gives an injection
\[
             P_n(k')[p^\infty]\lhook\joinrel\longrightarrow
                              P'_n(k')[p^\infty],
 \qquad
             f(P_n)\le f(P'_n),
 \tag{PI.2}
\]
where $f$ denotes geometric $p$-rank. In particular the old
rational exponent divides the new one.

\begin{proof}
Let $h:D'\to D$ be the projection, of degree $d$.
The deck groups are naturally identified by base change;
write their corresponding generators as $\sigma,\sigma'$.
Pullback and norm are equivariant for these generators.
For the canonical integral projectors $p_n=ne_n$ and
$p'_n=ne'_n$ of (IS.2)--(IS.3), this gives
\[
 h^*p_n=p'_nh^*,\qquad h_*p'_n=p_nh_*.
 \tag{PI.3}
\]
Consequently $h^*$ maps the actual image $P_n$ into the actual
image $P'_n$, and $h_*$ maps $P'_n$ back into $P_n$.
These maps are defined over $k$, even if the geometric deck
groups themselves are not constant. Their composite on $P_n$ is
\[
                          h_*h^*=[d].
 \tag{PI.4}
\]
Since $p\nmid d$, multiplication by $d$ is invertible on every
$p$-primary torsion group. Thus $h^*$ is injective on those
groups over $\bar k$ and over every $k'$. Restricting to
geometric $p$-torsion proves the $p$-rank inequality, and
the rational injections prove the exponent divisibility.
No arbitrary isogenous representative of a primitive factor
is substituted.
\end{proof}

Thus a nonzero old rational primitive $p$-primary group cannot
become zero by this construction. If the same-$n$ numerical
inequality of A47 already fails for the old exponent, it still
fails for the new exponent. An additional constant extension
cannot improve either conclusion, since old rational torsion
remains rational. This concerns the vanishing and exponent
inputs to A46/A47; it does not forbid a newly chosen cyclic
cover of $Y$ from satisfying them.

Connectedness and retention of the conductor are essential.
If the fibre product is disconnected and one selects a
component, the cyclic degree or conductor can drop; no
same-conductor assertion is made in that situation.
The pointing also remains a separate hypothesis when applying
a section criterion: rational lifts in both given covers
produce a rational point of the fibre product, but the
injection alone does not supply those lifts.

\paragraph{Raynaud's simultaneous theta incidences are not nilpotence.}
Raynaud's Theorem~2, printed p.~76, gives a prime-to-$p$
solvable Galois cover $Y\to X$ and a finite representation
factoring through its group which has no theta divisor.
This is a geometric existence statement over an algebraic
closure; finite-field descent and pointing may require a
finite extension of the initial field.
Proposition~1(5), printed p.~75, then gives
\[
 a^{(1)*}J_{X^{(1)}}\ \subset\ \operatorname{Supp}(\Theta_Y),
 \tag{PI.5}
\]
where the theta divisor is that of locally exact differentials
\cite[Proposition~1(5), Theorem~2]{RaynaudNonordinary2000}.
In particular all the multiples of a pulled-back torsion class
lie on this theta divisor, provided the class retains its order.
The theorem and its construction are stronger than a single
nonordinary-cover assertion.

Nevertheless, if an original cyclic cover has positive
primitive $p$-rank, (PI.2) makes the connected pulled-back
cover's primitive $p$-rank positive as well. The simultaneous
incidences therefore cannot be interpreted as a statement that
the primitive Frobenius operator is nilpotent. In genus
$g(Y)>2$, a nontrivial character has coherent cohomology
dimension $g(Y)-1>1$. Singularity of every arrow around a
Frobenius orbit does not imply nilpotence of its product:
the matrices $\operatorname{diag}(0,1,\ldots,1)$ already show
the linear-algebra distinction. This matrix observation is a
logical comparison, not a claimed geometric counterexample.

The underlying obstruction to a decreasing-rank induction is
also direct. For any allowed degree-$d$ pullback, trace gives
\[
 H^1(X,\mathcal O_X)\ \lhook\joinrel\longrightarrow\
 H^1(Y,\mathcal O_Y),\qquad
                 d^{-1}\operatorname{Tr}\circ a^*=1.
 \tag{PI.6}
\]
These maps commute with Frobenius. Hence an existing stable
Frobenius subspace persists under further allowed covers.
A fresh nonordinary representation can add a defect in a
larger space, but these cited results supply no decreasing
stable dimension or zero stable part of a new primitive factor.

Finally, a quotient obtained only from a subgroup of the
prime-to-$p$ deck group cannot supply an elliptic intermediate
curve over the original base: every such intermediate cover
$Z\to X$ is etale and
$g(Z)=1+[Z:X](g(X)-1)\ge2$.
An additional non-deck quotient or correspondence would require
new geometric input. This records the exact limit of the
pullback-iteration route, not a nonexistence theorem for fresh
witnesses upstairs and not a universal main-route advance.


\paragraph{A geometric instance with every primitive arrow singular.}
The preceding distinction has an actual finite-field realization.
There are a finite extension $k/\mathbf F_{256}$, a smooth
projective curve $Y/k$, and a connected cyclic etale cover
$D'\to Y$ of degree thirteen such that every primitive
Frobenius character arrow is singular, but the actual primitive
factor has positive geometric two-rank. More precisely, one
can arrange
\[
 \begin{gathered}
 d=5\cdot3^a,\qquad 13\le a\le65,\qquad g(Y)=d+1,\\
 \dim P'_{13}=12d,\qquad 12\le f(P'_{13})\le12(d-1).
 \end{gathered}
 \tag{PIA.1}
\]
There is a geometric cyclic class $\mathcal M$ of exact order
thirteen on $Y^{(1)}$, defining the cover under relative
Frobenius, with
\[
 H^0\bigl(Y^{(1)},B_Y\otimes\mathcal M^i\bigr)\ne0
                           \qquad(1\le i\le12).
 \tag{PIA.2}
\]
Thus all primitive classes meet the actual locally exact
differential theta divisor, yet the primitive Frobenius operator
is not nilpotent. This is an existence corollary, not an
explicit equation or an exact computation of the new two-rank.

\begin{proof}
\emph{1. Verify the degree support in the cited construction.}
Work first over $\overline{\mathbf F}_2$ with the fixed genus-two
curve $C$ and cyclic thirteen-cover $D\to C$ of A32.
Its primitive factor $P_{13}$ is ordinary of dimension twelve,
so $f(P_{13})=12$.

For genus two and characteristic two, Raynaud's proof of
Theorem~2 makes the explicit choices $m=5$ and $n=3$
on printed p.~86. The first stage is a cyclic etale cover
$C'\to C$ of degree five; its genus is six.
Corollary~18 applies on $C'$ with the usual principal
polarization: $3\mid6$ and $6/3=2\le5=(6-1)/(2-1)$.
In the construction preceding Theorem~17, printed p.~84,
the resulting etale Galois cover $Z\to C'$ has group
\[
 1\longrightarrow\mu_3\longrightarrow H
            \longrightarrow J_{C'}[3]\longrightarrow1,
                         \qquad |H|=3^{13}.
 \tag{PIA.3}
\]
It supplies a representation $\rho'$ of $\pi_1(C')$ factoring
through $H$ with no theta divisor
\cite[pp.~84--86, Theorem~17, Corollary~18 and proof of Theorem~2]{RaynaudNonordinary2000}.

Take the Galois closure $Y\to C$ of $Z\to C$ geometrically.
Over $C'$ its function field is the compositum of the five
conjugates of the $H$-extension. Restriction embeds its Galois
group over $C'$ into $H^5$, and maps it onto $H$ for the
original conjugate. Hence that group has order $3^a$ with
$13\le a\le65$, and $[Y:C]=5\cdot3^a=d$.
In particular thirteen does not divide this degree.
This argument verifies the support after Galois closure;
it does not assume that the first Heisenberg cover was
already Galois over $C$.

\emph{2. Retain the theta-containment conclusion.}
The induced representation
$\rho=\operatorname{Ind}_{\pi_1(C')}^{\pi_1(C)}\rho'$
factors through the Galois group of $Y/C$: its kernel
contains the core of the subgroup defining $Z$.
By Raynaud's Proposition~1(4), a representation without a
theta divisor still has none after this induction.
Proposition~1(5) therefore gives, for $\alpha:Y\to C$,
\[
       \alpha^{(1)*}J_{C^{(1)}}
                         \subset\operatorname{Supp}(\Theta_Y)
 \tag{PIA.4}
\]
\cite[Proposition~1(4)--(5), p.~75]{RaynaudNonordinary2000}.

Choose the exact-order-thirteen line bundle $\mathcal L$
on $C^{(1)}$ defining the cyclic cover corresponding to
$D\to C$ under the relative-Frobenius equivalence of etale
covers, and put $\mathcal M=\alpha^{(1)*}\mathcal L$.
Norm composed with pullback is multiplication by $d$.
Since $13\nmid d$, $\mathcal M$ still has order thirteen.
Its relative-Frobenius pullback defines
$D'=D\times_CY\to Y$. Every $\mathcal M^i$, including all
twelve primitive powers, belongs to the image in (PIA.4).
The defining support property of the actual theta divisor
gives (PIA.2).

\emph{3. Positive stable rank survives.}
The covers $D/C$ and $Y/C$ have coprime geometric degrees
thirteen and $d$. They are linearly disjoint, so $D'$ is
geometrically connected and remains cyclic of degree thirteen
over $Y$. The actual-projector injection (PI.2) applies to
$D'\to D$, whose degree $d$ is odd. It gives
\[
                            f(P'_{13})\ge f(P_{13})=12.
 \tag{PIA.5}
\]
Riemann--Hurwitz gives $g(Y)=d+1$ and $g(D')=13d+1$;
since thirteen is prime, the primitive factor is the whole
Prym and has dimension $12d$.

For each $1\le i\le12$, the locally exact differential
sequence twisted by $\mathcal M^i$ identifies the kernel of
the corresponding relative-Frobenius map on coherent $H^1$
with $H^0(B_Y\otimes\mathcal M^i)$. The degree-zero line
bundles at either end have no global sections because they
are nontrivial of order thirteen. Each character space has
dimension $g(Y)-1=d$, and (PIA.2) makes every such map
singular. This is the character interpretation in
Pop--Sa\"idi, Section~1.3, printed p.~112 \cite{PopSaidi}.
The primitive coherent space is their direct sum, of dimension
$12d$; its Frobenius kernel has dimension at least twelve.
Its stable rank is consequently at most $12d-12$.
Together with (PIA.5) this proves (PIA.1), including the
failure of nilpotence despite every primitive arrow being
singular.

\emph{4. Descend to a finite field.}
All the curves, maps and the finitely many required deck
transformations over $\overline{\mathbf F}_2$ descend to
a common finite extension $k$ of $\mathbf F_{256}$.
Geometric connectedness, degrees, theta incidences and
geometric two-ranks are unchanged by this descent.
If a rational pointing is desired, a further finite extension
rationalizes a chosen point of $Y$ above A32's split point
and hence a point of the fibre product. No original-field
rational-unit assertion is made after such extensions.
\end{proof}

This corollary strengthens the linear-algebra illustration to
actual curves over a finite field. It does not supply a new
primitive-unit witness, and no main-gate advance is inferred
from it. The earlier A04 already refuted the weaker logical
inference; the present construction records a more specific
higher-genus failure with all primitive theta arrows present.


\subsection{An exact Fourier criterion on an odd quotient}
\label{sec:odd-quotient-fourier}

\paragraph{Finding A49 (all conductors of a prescribed odd quotient).}
Let $C/k$ be a smooth projective geometrically connected genus-two
curve over $k=\mathbf F_q$ of characteristic $p$, and let $P\in C(k)$.
Put $G=J_C(k)$, $N=\#C(k)$ and $\kappa_P=[K_C-2P]$.
Let
\[
 b:G\twoheadrightarrow H,\qquad
 n=\exp H>1,\qquad n\text{ odd},\qquad p\nmid n
 \tag{OQ.1}
\]
be a finite abelian quotient satisfying
\[
 b(\kappa_P)=0,\qquad
       \langle p,-1\rangle=(\mathbf Z/n)^\times.
 \tag{OQ.2}
\]
For $h\in H$, let $[h]$ denote the corresponding basis element of
$\mathbf F_p[H]$, and define
\[
 F_H=[0]+\sum_{R\in C(k)}[b([R-P])],\qquad
 U_H=\frac{N+1}{|H|}\sum_{h\in H}[h].
 \tag{OQ.3}
\]
The fraction is taken in $\mathbf F_p$; its denominator is invertible.
Then the following are equivalent:
\begin{enumerate}
\item $F_H\ne U_H$ in $\mathbf F_p[H]$.
\item There is a cyclic quotient $c:H\twoheadrightarrow\mathbf Z/m$,
      with $m>1$ dividing $n$, whose constant $P$-split Lang cover
      $D\to C$ over $k$ has actual primitive conductor-$m$ factor
      $\mathcal P_m$ satisfying
      \[
                  p\nmid\#\mathcal P_m(k).
       \tag{OQ.4}
      \]
\end{enumerate}
The cover in the second assertion is geometrically connected.
When these equivalent conditions hold, A46 separates the point section of the prescribed $P$
from that of every distinct point of $C(k)$.
The conclusion concerns the primitive factor of this cover;
proper-conductor factors may have rational $p$-primary torsion.

In particular, the sufficient numerical conditions
\[
                 p\nmid(N+1),\qquad |H|>N
 \tag{OQ.5}
\]
imply (OQ.4) for some such cyclic quotient. In characteristic two,
the first condition is simply that $N$ is even. It therefore holds
for every genus-two curve of two-rank one by A26.
At a rational Weierstrass point $P$, the canonical condition in
(OQ.2) holds for every quotient, since $K_C\sim2P$.

\paragraph{Recovered inputs and the scope of the refinement.}
The Fourier identity and empty-coset argument are already present
in recovered successor A.10 and A.35, and A26 isolates their
even-$N$ consequence for a prime quotient on which $p$ is primitive
\cite[A.10, A.35]{ASrun}. The genus-two character polynomial is
recovered A1 B.2; its canonical-class symmetry is the one used in
recovered successor A.26 and A29 \cite[B.2]{A1run}\cite[A.26]{ASrun}.
Here the quotient may have higher prime powers or several prime
divisors. Condition (OQ.2) uses the canonical symmetry to join the
$p$-power and inverse character orbits, and A46 permits selecting
only the resulting primitive conductor. The proof below establishes
the exact equivalence and its original-field pointing; it makes no
priority claim for the underlying Fourier or Lang constructions.

\begin{proof}
\emph{1. Every cyclic quotient gives a cover over the original field.}
Fix a cyclic quotient $c:H\to A=\mathbf Z/m$, with $m>1$, and put
$a=c\circ b$ and $K=\ker a\subset G$.
Let $\Phi$ be $q$-power Frobenius on $J=J_C$. The Lang isogeny
$\Phi-1$ is etale with constant kernel $G$, and factors over $k$ as
\[
             J\longrightarrow J/K\xrightarrow{\psi_a}J,
                    \qquad \ker\psi_a=G/K=A.
 \tag{OQ.6}
\]
Pull it back by $i_P:C\to J$, $R\mapsto[R-P]$.
The result $D_a=C\times_J(J/K)$ is an etale cover with constant
deck group $A$. The Abel--Jacobi connected-pullback theorem
\cite[Proposition~9.1]{MilneJV}, as used in A12, makes it
geometrically connected. Its fibre over $P$ is the constant kernel
of $\psi_a$, containing the origin. Thus it is $P$-split over $k$
itself. No extension to rationalize a point or deck transformation
has been made.

\emph{2. The exact primitive rational test.}
For a nontrivial characteristic-zero character $\chi$ of $A$, use
the Lang Frobenius convention of recovered A1 B.2. Its exact
unramified genus-two polynomial and middle-coefficient symmetry are
\[
\begin{aligned}
 L_P(\chi,T)&=1+S_\chi T+q\chi(a(\kappa_P))T^2,\\
 S_\chi&=\sum_{R\in C(k)}\chi(a([R-P])),\\
 S_\chi&=\chi(a(\kappa_P))S_{\chi^{-1}}.
\end{aligned}
 \tag{OQ.7}
\]
These are precisely the Riemann--Roch identities proved there
and in (PC.6)--(PC.7). Reversing the Lang convention inverts all
characters and changes none of the tests below. Here
$a(\kappa_P)=0$ by (OQ.2), so $S_\chi=S_{\chi^{-1}}$.

Define $\mathcal P_m$ to be the actual image of the integral
canonical projector $m e_m$ on $J_{D_a}$, as in A46; $e_m$ is
the rational idempotent supported on the characters of exact
order $m$. It is defined over $k$. The deck-character decomposition
of the zeta numerator, or equivalently its action on a Tate module
of auxiliary characteristic prime to $pm$, gives the exact identity
\[
       \#\mathcal P_m(k)
       =\prod_{\substack{\chi\in\widehat A\\\operatorname{ord}\chi=m}}
                         L_P(\chi,1).
 \tag{OQ.8}
\]
The image has exactly these rational character summands; finite-field
point counts depend on the Frobenius polynomial, so this equality
uses the actual image and does not require a chosen isogenous
representative. All factors on the right are algebraic integers.

Fix a prime above $p$ in a cyclotomic field containing the $n$th
roots of unity. Reduction identifies their character group with
the $n$-torsion roots of unity in $\overline{\mathbf F}_p$, since
$p\nmid n$. By (OQ.7),
\[
        \overline{L_P(\chi,1)}
          =1+\sum_{R\in C(k)}\bar\chi(a([R-P]))
          =(\bar\chi\circ c)(F_H).
 \tag{OQ.9}
\]
Consequently (OQ.4) is equivalent to nonvanishing of (OQ.9)
for every faithful character of $A$. Indeed the product (OQ.8)
is a rational integer, and its reduction is nonzero exactly
when every factor has nonzero reduction at this fixed prime.

\emph{3. One character controls its entire conductor.}
Reduction of units $(\mathbf Z/n)^\times\to(\mathbf Z/m)^\times$
is surjective for every $m\mid n$. Thus (OQ.2) implies
\[
             \langle p,-1\rangle=(\mathbf Z/m)^\times.
 \tag{OQ.10}
\]
For any character $\eta:H\to\overline{\mathbf F}_p^\times$,
the coefficients of $F_H$ lie in $\mathbf F_p$, giving
\[
                 \eta^p(F_H)=\eta(F_H)^p.
 \tag{OQ.11}
\]
The hyperelliptic involution is defined over $k$ and satisfies
$[\iota R-P]=\kappa_P-[R-P]$.
As $b(\kappa_P)=0$, its permutation of $C(k)$ also gives
\[
                         \eta^{-1}(F_H)=\eta(F_H).
 \tag{OQ.12}
\]
Hence vanishing, and also nonvanishing, is preserved by the
$p$-power and inverse character operations. If $\eta$ has
exact order $m$, all generators of its cyclic character subgroup
are reached using (OQ.10). Thus a single nonzero value supplies
nonzero values for every faithful character of the quotient
$H/\ker\eta$ of order $m$.

\emph{4. Fourier inversion and the numerical corollary.}
The finite Fourier transform on $H$ is invertible over
$\overline{\mathbf F}_p$, because $p\nmid|H|$.
The trivial character evaluates $F_H$ to $N+1$.
All its nontrivial character values vanish if and only if
\[
                   F_H=\frac{N+1}{|H|}\sum_{h\in H}[h]=U_H.
 \tag{OQ.13}
\]
If $F_H\ne U_H$, choose a nontrivial character $\eta$ with
$\eta(F_H)\ne0$. Its image has order $m>1$ dividing $n$.
Step 3 and (OQ.9) make every faithful factor of its cyclic
quotient nonzero; Step 2 proves (OQ.4), and Step 1 supplies
the required original-field cover. Conversely, a quotient
satisfying (OQ.4) has a faithful character with nonzero value
on $F_H$, so (OQ.13) cannot hold. This proves the equivalence.

Under (OQ.5), the coefficient $(N+1)/|H|$ is nonzero.
There are at most $N$ occupied Abel-image cosets, and the zero
coset is occupied by $P$. Since $|H|>N$, an empty nonzero coset
has coefficient zero in $F_H$ and a nonzero coefficient in
$U_H$. Thus $F_H\ne U_H$. Finally (OQ.4) says
$\mathcal P_m(k)[p^\infty]=0$, so the point-separation
conclusion follows directly from A46, with its odd cyclic
degree and prescribed rational lift hypotheses all satisfied.
\end{proof}

\paragraph{Two retained comparisons over $\mathbf F_4$.}
On A27's curve
$y^2+y=ux^3+x+u/x$, $u^2+u+1=0$, take either rational
Weierstrass point $W$. The retained counts are
$\#C(\mathbf F_4)=6$ and $\#J_C(\mathbf F_4)=28$.
Its seven-primary quotient is $H=\mathbf Z/7$.
The powers of two modulo seven are $\{1,2,4\}$, and their
negatives give the other three units. Thus (OQ.2) and (OQ.5)
hold and A49 gives a constant $W$-split degree-seven primitive
rational witness over $\mathbf F_4$. The primitive-root
hypothesis of A26 fails at seven. A31 already supplies a
geometric degree-seven witness for this same curve; the
comparison concerns the reach of the sufficient criteria.

For the quadratic twist used in A45, namely
\[
          C:\ Y^2+X^2Y=uX^5+uX^4+X^3+uX,
                     \qquad W=\infty,
 \tag{OQ.14}
\]
the retained data (ND.12)--(ND.19) give $N=4$,
$\#J_C(\mathbf F_4)=18$, and a rational class of exact
order nine. Hence its full three-primary quotient is
$H=\mathbf Z/9$. The powers of two modulo nine are
$1,2,4,8,7,5$, so the unit condition holds; $\kappa_W=0$
and $9>4$. The only possible nontrivial cyclic conductors
are three and nine. The unique degree-three quotient,
up to an automorphism of its deck group, has
$L_\chi(T)=1+T+4T^2$ for each nontrivial character.
Its primitive factor has order $6^2$, and therefore fails
(OQ.4). A49 consequently supplies a degree-nine quotient
whose primitive conductor-nine factor has odd rational order.
Here the prime-level quotient has size $3\le4$, while the
full primary quotient has size $9>4$.

One can check that last odd order exactly from the same retained
point data, without a new curve search. The two Weierstrass
points have label zero in $H$; the other two points have labels
$a,-a$ with $3\nmid a$, because their reductions modulo three
are $1,2$ in (ND.17)--(ND.19). For a primitive ninth root $z$,
the faithful values at $T=1$ are $7+z^j+z^{-j}$, for
$j\in(\mathbf Z/9)^\times$. The three distinct values
$s=z^j+z^{-j}$ satisfy $s^3-3s+1=0$, each occurring twice.
Therefore
\[
 \#\mathcal P_9(\mathbf F_4)
     =\left(7^3-3\cdot7-1\right)^2
     =321^2=103041.
 \tag{OQ.15}
\]
The proper conductor-three factor still has even rational order;
no assertion about the full relative class number being odd
is made. The accompanying finite algebra checker verifies the
unit orbits and the ninth-root product. Its input curve counts
and labels are the separately retained A27/A45 certificates,
not a new enumeration in that checker.

\paragraph{The remaining supply question.}
This is a test on a quotient which must actually exist over the
displayed field. A49 does not guarantee a quotient satisfying
(OQ.2), nonuniformity in (OQ.3), or the sufficient bound (OQ.5)
for every curve and prescribed point. In particular a nonzero
Fourier value without (OQ.2) need not control a whole faithful
conductor. The canonical-kernel requirement is also explicit;
A29 addresses a different situation, where a suitable primary
canonical component is nonzero. The two comparisons above
already have other witnesses in this run, so they do not prove
that the full previously unresolved universal witness gate has
shrunk. No consequence about this geometric supply is inferred
from Artin-prime density or from number-field GRH.


\paragraph{Nonuniformity is unavoidable if the conductor is left unrestricted.}
The counting part of A49 can always be arranged after allowing
arbitrary odd conductors. Precisely, fix a smooth projective
geometrically connected genus-two curve $C/k$ over a finite field
$k=\mathbf F_q$, a point $P\in C(k)$ and a finite set $S$ of primes containing
$2$, the characteristic $p$, and every prime dividing the order
of $\kappa_P$ in $J_C(k)$.
There are arbitrarily large extensions $k_j/k$ and quotients
\[
 H_j=J_C(k_j)_{S^c}
\]
which kill $\kappa_P$ and for which the augmentation element
$F_{H_j}$ of (OQ.3) is nonuniform.
Consequently there is a constant connected $P$-split cyclic cover
over $k_j$, of odd conductor supported outside $S$, with at least
one faithful character $L$-value nonzero modulo $p$.
In fact its entire signed $p$-orbit is nonzero. No assertion that
all faithful characters are nonzero is made.

\begin{proof}
Use the explicit integer $M$ of (PE.3) in A42, so that
$J_C(\mathbf F_{q^j})_S=J_C(k)_S$ whenever $\gcd(j,M)=1$.
Put $B=\#J_C(k)_S$ and $Q=q^j$ on this infinite set of degrees;
the progression $j\equiv1\pmod M$ already supplies an infinite set.
The projection to $H_j$ has kernel of the fixed size $B$ and
kills $\kappa_P$. The Weil bounds \cite[Theorem~19.1]{MilneAV} give
\[
 |H_j|\ge\frac{(\sqrt Q-1)^4}{B},\qquad
 Q+1-4\sqrt Q\le N_j:=\#C(\mathbf F_Q)
                         \le Q+1+4\sqrt Q.
\]
For all sufficiently large admissible $j$, these imply
$|H_j|>N_j$ and $N_j>2(B-1)+1$.

For every nonzero Jacobian class there are at most two ordered
representations as a difference of two distinct geometric points
of a genus-two curve; this is the difference bound used in recovered
successor A.10 and A.35 \cite[A.10, A.35]{ASrun}.
Hence at most $2(B-1)$ rational points
share an $H_j$-coset with another rational point: the total number
of ordered distinct pairs in the same coset is at most that bound.
There are therefore at least two singleton cosets. At most one
is zero, so a nonzero singleton coset has coefficient one in
$F_{H_j}$. Since $|H_j|>N_j$ and $P$ occupies zero, there is
also an empty nonzero coset, whose coefficient is zero.
Thus $F_{H_j}$ is not a constant multiple of the sum of all
group elements, in every characteristic and with no parity
condition on $N_j$.

Fourier inversion now gives a nontrivial character with nonzero
value. Its cyclic image has order supported outside $S$.
The original-field Lang construction and character formula of
A49 identify this value with the corresponding faithful
$L_P(\chi,1)$ modulo a fixed prime above $p$.
Since the quotient kills $\kappa_P$, $p$-power and inversion
preserve nonvanishing. This proves precisely the signed-orbit
assertion. Iterating after adjoining the primes in the selected
conductor to $S$ also gives pairwise coprime conductors, without
making those conductors prime or signed-primitive.
\end{proof}

This combines the already checked A42 primary-stabilization
argument with the recovered difference bound and Fourier identity;
it is not a new claim about Lang construction or torsion density.
It answers the nonuniformity-only supply question positively.
It does not give the unit-group hypothesis of A49: an exponent
or selected conductor can have several signed $p$-orbits, and
the other orbits can still contain vanishing factors. The exact
remaining step is to control an entire faithful conductor,
or to bound the rational primitive exponent as in A47.
The statement therefore does not close the existing universal
geometric-witness row or automatically count as main-gate movement.


\subsection{Integral operator limits for a partial character packet}
\label{sec:partial-packet-operators}

\paragraph{Finding A50 (integral limits of a partial-packet projector).}
The unrestricted-quotient supplement to A49 produces one good
signed $p$-orbit of character values. A $p$-adic projector onto
that orbit need not be an integral operator of small coefficient
sum. We give an exact norm bound and a finite module example
which make this distinction precise. Neither assertion constructs
a curve or proves failure of any actual prescribed-point cover.

\paragraph{A quantitative bound for approximating a partial projector.}
Let $n>1$ be prime to a rational prime $p$, put
$U=(\mathbf Z/n)^\times$ and $r=\varphi(n)$, and let
$\varnothing\ne\Omega\subsetneq U$ be a union of orbits of
multiplication by $p$. Fix a primitive $n$th root $\zeta$ in
$\overline{\mathbf Q}_p$. For
\[
 A=\sum_{i=0}^{n-1}a_i\sigma^i\in\mathbf Z[C_n],\qquad
                     d=\sum_i|a_i|,
 \tag{PP.1}
\]
suppose $A(\zeta^u)$ is a $p$-adic unit for at least one
$u\in\Omega$, and, for an integer $a\ge1$,
\[
                    v_p(A(\zeta^u))\ge a
                         \quad(u\in U\setminus\Omega).
 \tag{PP.2}
\]
The normalization is $v_p(p)=1$. Then
\[
                  d\ \ge\ p^{a(r-|\Omega|)/r}.
 \tag{PP.3}
\]
In particular, integral approximations modulo $p^a$ to the
$p$-adic primitive-character projector supported on $\Omega$
cannot have uniformly bounded coefficient sum as $a\to\infty$.

\begin{proof}
The algebraic integer $\alpha=A(\zeta_n)$ in
$\mathbf Q(\zeta_n)$ is nonzero: otherwise every primitive
evaluation would vanish, contrary to the unit assumption.
Its field norm is therefore a nonzero rational integer.
The product over $u\in U$ is that norm. All its factors are
$p$-adically integral, and (PP.2) gives
\[
 v_p\!\left(N_{\mathbf Q(\zeta_n)/\mathbf Q}(\alpha)\right)
                         \ge a(r-|\Omega|).
 \tag{PP.4}
\]
The equality between the product and norm uses all primitive
embeddings at the one fixed $p$-adic embedding; no place is
discarded. Since $p\nmid n$, $p$ is unramified in this cyclotomic
field and the normalization agrees with the usual valuation on
the integer norm. At every complex embedding,
$|A(\zeta_n^u)|\le d$. Hence
\[
 p^{a(r-|\Omega|)}
       \le\left|N_{\mathbf Q(\zeta_n)/\mathbf Q}(\alpha)\right|
       \le d^r,
\]
which proves (PP.3). A projector approximation is one on the
good characters and zero on the others modulo $p^a$, so it
satisfies the stated assumptions. Exact integral isolation of
a nonempty proper part of one rational cyclotomic block is
also impossible: vanishing at one primitive characteristic-zero
root forces divisibility by $\Phi_n$, hence vanishing at all.
\end{proof}

This is a lower bound conditional on the displayed divisibility
requirements, not an estimate of actual geometric torsion
exponents. It leaves open the existence of a useful small
operator for a particular cover and particular collision class.

\paragraph{A complete length-thirty-one spectral calculation.}
Let $G=C_{31}=\langle\sigma\rangle$, let
$K=\mathbf F_{32}=\mathbf F_2[z]/(z^5+z^2+1)$, and put
$\zeta=z$. Its multiplicative order is thirty-one. Define
\[
 \Omega=\{\pm2^j\bmod31:0\le j<5\},\qquad |\Omega|=10.
 \tag{PP.5}
\]
In the binary group algebra $\mathbf F_2[G]$, the entire
subspace whose Fourier transform is supported on $\Omega$
consists of the words
\[
 c(a,b)_i=\operatorname{Tr}_{K/\mathbf F_2}
             (a\zeta^i+b\zeta^{-i}),
       \qquad 0\le i<31,\quad (a,b)\in K^2.
 \tag{PP.6}
\]
It has dimension ten, and its minimum nonzero Hamming weight
is ten. The complete weight distribution is
\[
\begin{array}{c|rrrrrrr}
\text{weight}&0&10&12&14&16&18&20\\
\hline
\text{number of words}&1&31&155&310&217&155&155.
\end{array}
 \tag{PP.7}
\]

\begin{proof}
Expand the trace as a sum of the five $2$-power conjugates.
For the Fourier convention
$\widehat c(k)=\sum_{i=0}^{30}c_i\zeta^{ki}$, character
orthogonality gives
\[
 \widehat{c(a,b)}(-2^j)=a^{2^j},\qquad
 \widehat{c(a,b)}(2^j)=b^{2^j},
 \tag{PP.8}
\]
and zero outside $\Omega$. The two orbits are disjoint:
two has order five modulo thirty-one, and its powers do not
contain minus one. In particular (PP.8) recovers $a,b$ and
proves injectivity of the ten-dimensional binary parameter
space. Conversely a binary word satisfies
$\widehat c(2k)=\widehat c(k)^2$.
On each of the two five-element orbits, one arbitrary element
of $K$ determines all the Fourier values. Thus the entire
binary spectral subspace has dimension at most ten, and
(PP.6) is all of it, not just a subcode.

The finite assertion (PP.7) is verified by enumeration of all
$32^2=1024$ pairs in the retained certificate. It uses exact
arithmetic in the displayed polynomial quotient, checks that
the thirty-one powers of $z$ are all nonzero elements, and
retains every parameter pair, word and weight. The sum of
the displayed multiplicities is $1024$ and every word is distinct.
The pair $a=b=1$ has weight ten. A separate implementation
using five-coefficient polynomial convolution reproduces
the entire distribution and checks the Fourier support of
all ten generators. This is finite exhaustive verification,
not a claimed formal proof-assistant verification.
\end{proof}

\paragraph{A finite module with no sufficiently small annihilator.}
There is a cyclic $\mathbf Z/8[G]$-module $W$, with a generator
$\delta$, such that its good packet $\Omega$ is absent, but
every nonzero $A\in\mathbf Z[G]$ with $A\delta=0$ satisfies
\[
                         \|A\|_1\ge8.
 \tag{PP.9}
\]
Consequently no such nonzero annihilator meets
$31>(\|A\|_1-1)^2$, the small-degree comparison used in
the genus-two integral-operator argument.

\begin{proof}
Let $e_0\in\mathbf F_2[G]$ be the idempotent with Fourier
value one on $\Omega$ and zero elsewhere. It is the word
$c(1,1)$ in (PP.6). Lift its coefficients to integers and
apply the polynomial $f(t)=3t^2-2t^3$ twice in
$R=(\mathbf Z/8)[G]$. Since
\[
 f(t)^2-f(t)=(t^2-t)^2(4t^2-4t-3),
\]
the resulting element $e$ is idempotent modulo eight and
reduces to $e_0$. Define
\[
                        W=(1-e)R,\qquad\delta=1-e.
 \tag{PP.10}
\]
It is cyclic, and multiplication by $e$ kills it. Its reduction
modulo two contains exactly the complement of $\Omega$,
including the trivial character; its rank over $\mathbf Z/8$
is twenty-one. Thus there is a whole missing good signed
packet, not merely one missing character.

Suppose $A\ne0$ and $A\delta=0$.
Write $A=2^bB$ with at least one coefficient of $B$ odd.
If $b\ge3$, then $\|A\|_1\ge8$.
If $b<3$, the equality
$2^bB(1-e)=0$ modulo eight implies, on reduction modulo two,
\[
                         \bar B(1-e_0)=0.
\]
The nonzero binary word $\bar B$ therefore belongs to the
entire spectral subspace of (PP.6), so it has at least ten
nonzero coefficients. Hence
$\|A\|_1=2^b\|B\|_1\ge2^b\cdot10\ge10$.
This proves (PP.9); the bound eight is attained by the scalar
operator $8$, so the unrestricted minimum is exactly eight.
The degree comparison would require $\|A\|_1\le6$, and
therefore fails for every nonzero annihilator in this example.
\end{proof}

For an integral group-ring operator and two disjoint free
deck orbits, the positive part of $A(Q'-P')$ has degree
$\|A\|_1$. Thus (PP.9) addresses the precise coefficient
quantity in that construction. At exponent two, by contrast,
the scalar two already annihilates the entire abstract module;
the binary weight bound alone would only obstruct a lift
which remains nonzero modulo two. Using the exponent-eight
module in (PP.10) avoids that mistaken inference.

This is an abstract finite-module counterexample to an
unqualified small-integral-operator deduction from one good
packet. It is not asserted to be the rational torsion module
of an actual Jacobian, or to arise from a collision of point
sections. It does not refute geometric supply, exclude other
operators on a particular curve, or prove that every successful
argument must use a full faithful conductor. The norm bound
and finite example specify why a $p$-adic packet projector by
itself does not close the remaining step of the A49 supply note.


\subsection{Limits of norm compatibility and tower descent}
\label{sec:norm-compatibility-towers}

\paragraph{Finding A51 (compatible classes and rational generator bounds).}
Let $C/k$ be smooth, projective and geometrically connected,
of genus $g\ge2$, where $k=\mathbf F_q$ has characteristic $p$.
Suppose hypothetically that distinct $P,Q\in C(k)$ have equal
geometrically pro-prime-to-$p$ point sections. The covering
implication and independent uniqueness argument of recovered
A1, A.31, give compatible rational lifts on every permitted
pointed cover and classes
\[
       \delta_D=[Q_D-P_D]\in J_D(k)[p^\infty]
       \tag{NC.1}
\]
compatible under norms \cite[claim A.31]{A1run}.
The refinement and point-count conclusions below concern these
actual classes. Separately, a Frobenius-stable rank-one
$\mathbf Z_\ell$ quotient, when it exists, gives the stated
uniform bound on the number of rational $p$-primary generators.
The abstract examples prove limits of those algebraic
implications; none is asserted to be a geometric tower or
a point-section collision. No universal witness is supplied.

\begin{proof}
\emph{1. Exact refinement and order formulas.}
Let $h:D'\to D$ be a finite etale refinement of degree $d$
prime to $p$, with the compatible rational point lifts.
Write $\delta'=\delta_{D'}$ and $\delta=\delta_D$.
Pushforward of the point divisors and pullback of a divisor
give
\[
                     h_*\delta'=\delta,
                     \qquad h_*h^*=[d].
       \tag{NC.2}
\]
The second identity follows on each closed point from the
degree-$d$ fibre formula, and hence on divisor classes.
Since $[d]$ is invertible on $p$-primary groups, the map
$d^{-1}h^*$ is a section of $h_*$ on these groups. Therefore
\[
\begin{aligned}
 J_{D'}(k)[p^\infty]
   &=d^{-1}h^*J_D(k)[p^\infty]\oplus K_h,\\
 K_h&=\ker\bigl(h_*:J_{D'}(k)[p^\infty]
                         \longrightarrow J_D(k)[p^\infty]\bigr),\\
 \delta'&=d^{-1}h^*\delta+b,\qquad b\in K_h.
\end{aligned}
       \tag{NC.3}
\]
Indeed subtracting $d^{-1}h^*h_*$ from the identity gives
the second summand, and applying $h_*$ proves that the
intersection is zero. If $m,m'$ are the exact orders of
$\delta,\delta'$, respectively, this direct sum proves
\[
           m\mid m',\qquad
           m'=\operatorname{lcm}\bigl(m,\operatorname{ord}(b)\bigr).
       \tag{NC.4}
\]
These formulas preserve the old order and leave the new
norm-kernel component. They do not assert that every possible
$b$ is represented by compatible point lifts.

\emph{2. The principal-function norm relation.}
A geometric function with divisor $m(Q_D-P_D)$ descends,
after multiplication by a scalar, to $f\in k(D)^\times$.
Its Galois conjugates differ by constants, and Hilbert90
removes that scalar cocycle
\cite[Theorem~59.24.1, Tag~\texttt{03P8}]{Stacks03P8}.
Likewise choose $f'\in k(D')^\times$ with divisor
$m'(Q_{D'}-P_{D'})$.
The functions are separable: if one were a $p$th power over
the perfect field, its divisor would contradict the exact
order; for order one, its simple poles already exclude that
possibility. The divisor-of-norm formula gives
\[
                 N_h(f')=c\,f^{m'/m},\qquad c\in k^\times
       \tag{NC.5}
\]
\cite[Lemma~42.18.1, Tag~\texttt{02RT}]{Stacks02RT}.
Here the rational points have residue degree one, so their
pushforwards have precisely the displayed multiplicities.
The quotient of the two sides has zero divisor on the
proper geometrically connected curve, hence is constant.

A $p$th-power norm does not, by itself, force a $p$th-power
function upstairs. To check this directly, choose $d>1$
prime to $p$ and a finite field containing a nontrivial
$d$th root $\zeta$. In $k(z)/k(t)$ with $t=z^d$,
\[
                    N\!\left(\frac{z-1}{z-\zeta}\right)=1.
       \tag{NC.6}
\]
The function has a simple zero and pole and is not a $p$th
power. This is only an algebraic counterexample to the
unrestricted norm inference: the map is ramified, and its
zero and pole have the same image. It is not an example of
the pointed etale-cover configuration in (NC.1).

\emph{3. The split-point estimate is a lower bound.}
Suppose $D\to C$ has degree $n$, is geometrically Galois,
and has constant deck group over $k$. A rational lift of
$P$ makes its whole fibre rational; the lift of $Q$ does
the same for its disjoint fibre. Thus $\#D(k)\ge2n$.
The function $f$ of Step2 has degree $m$. Its zero and pole
fibres contain one point each, and each other rational value
has at most $m$ rational preimages. Consequently
\[
     2n\le\#D(k)\le(q-1)m+2,
     \qquad
     m\ge\left\lceil\frac{2n-2}{q-1}\right\rceil.
       \tag{NC.7}
\]
This estimate concerns the particular point-difference order.
The constant-deck hypothesis is essential for the $2n$
rational points. No existence of arbitrarily large constant
pointed covers over this fixed field is asserted, and
enlarging the field changes the bound.

\emph{4. Proper-quotient norms miss the faithful component.}
Let $D_i\to C$ be constant cyclic covers of coprime odd
degrees $n_i>1$ prime to $p$, both with rational lifts of $P$.
Their fibre product $D$ is geometrically connected: the two
geometric Galois fields have intersection of degree dividing
both $n_1,n_2$. Its group is
$G=C_{n_1}\times C_{n_2}\simeq C_n$, $n=n_1n_2$.
For $h_i:D\to D_i$, set $K_i=\ker(G\to C_{n_i})$.
The pullback of the pushed-forward divisor is its orbit sum,
so
\[
 h_i^*h_{i*}=\sum_{\gamma\in K_i}\gamma,
 \qquad
 \left(\sum_{\gamma\in K_i}\gamma\right)p_n=0,
 \qquad p_n=ne_n.
       \tag{NC.8}
\]
The second identity holds in $\mathbf Z[G]$: a faithful
character is nontrivial on $K_i$, so its character sum is
zero, while $p_n$ kills all nonfaithful characters.
Character evaluation is faithful on $\mathbf Q[G]$.
For the actual primitive subvariety $P_n=\operatorname{im}p_n$
of A46, $p_n$ acts as $[n]$. Since $n$ and $\deg h_i$ are
prime to $p$, (NC.2) and (NC.8) show
\[
                       h_{i*}\bigl(P_n(k)[p^\infty]\bigr)=0.
       \tag{NC.9}
\]
Thus both old norm constraints are insensitive to the full
faithful product-conductor component. The same proof applies
to every proper cyclic quotient of a fixed cyclic cover.
A good packet downstairs constrains its inflated characters,
not the faithful characters killed by that norm. This is a
statement about actual Jacobian maps. It does not assert that
an arbitrary undetected torsion point can still be represented
by the compatible two-point divisor, or that no different
argument can combine the packet information.

\emph{5. The minimal deck field in a rank-one tower.}
Fix an odd prime $\ell\ne p$. Suppose a Frobenius-stable
rank-one $\mathbf Z_\ell$ quotient of
$\pi_1(C_{\bar k})^{\rm ab}$ is given. Let its Frobenius
multiplier be $a\in\mathbf Z_\ell^\times$, using arithmetic
Frobenius on geometric points. The pointed section gives
the associated compatible geometrically cyclic covers
$D_r\to C$ of degree $\ell^r$ over $k$.
The Abel--Jacobi identification with the Tate module
\cite[Proposition~9.1]{MilneJV}
\cite[Remark~15.3]{MilneAV}
shows that $a$ is a Frobenius eigenvalue on $T_\ell J_C$.
The Weil bound on its complex absolute values makes it
a non-root of unity \cite[Theorem~19.1]{MilneAV}.
Set
\[
 t=\operatorname{ord}_{\ell}(a),\qquad
 s=v_\ell(a^t-1),\qquad
 f_r=\operatorname{ord}_{\ell^r}(a).
       \tag{NC.10}
\]
Then $t\mid\ell-1$, $1\le s<\infty$, and
\[
                         f_r=t\ell^{\max(0,r-s)}.
       \tag{NC.11}
\]
To prove this formula, write $a^t=1+\ell^s u$ with $u$ a
unit. A power prime to $\ell$ preserves the valuation $s$;
an $\ell$th power raises it by exactly one. Both assertions
follow from binomial expansion, since $\ell$ is odd and
the first nonconstant term has strictly smaller valuation
than the remaining terms. Iteration gives
$v_\ell((a^t)^h-1)=s+v_\ell(h)$.
Any exponent yielding one modulo $\ell$ is a multiple of
$t$, proving (NC.11). This is exactly the minimal extension
degree defining all geometric deck transformations of $D_r$.
The assertion is conditional on the stable rank-one quotient;
it does not claim that every chosen $\ell$ supplies one.

\emph{6. A uniform rational generator bound, with no exponent bound.}
Let $B_r$ be the actual primitive conductor-$\ell^r$
subvariety of $J_{D_r}$ and let $k_j=\mathbf F_{q^j}$.
Use the full finite etale algebra, with its indicated action,
\[
\begin{aligned}
 \mathcal O_r&=\mathbf Z_p[Z]/\Phi_{\ell^r}(Z),\\
 \tau_r(Z)&=Z^a,\qquad R_r=\mathcal O_r^{\tau_r},\\
 U_r&=B_r(k_{f_r})[p^\infty],\qquad
 N_r=B_r(k)[p^\infty].
\end{aligned}
       \tag{NC.12}
\]
The integral descent result (SD.3)--(SD.10) in A45 applies
to $U_r$: its Frobenius is $\tau_r$-semilinear and has
$f_r$th power one. It gives
\[
                 U_r\simeq\mathcal O_r\otimes_{R_r}N_r.
       \tag{NC.13}
\]
Its known descent input is effective faithfully flat descent
for arbitrary modules, not averaging by $1/f_r$
\cite[Proposition~35.3.9, Tag~\texttt{023N}]{Stacks023N}.
The finite etale $R_r$-algebra $\mathcal O_r$ has constant
rank $f_r$. Since $R_r$ is a product of unramified discrete
valuation rings, this module is free of rank $f_r$, component
by component. Hence (NC.13) is, additively and noncanonically,
a direct sum of $f_r$ copies of $N_r$, and
\[
\begin{aligned}
 \#U_r&=(\#N_r)^{f_r},\qquad \exp U_r=\exp N_r,\\
 \dim_{\mathbf F_p}U_r[p]&=f_r\dim_{\mathbf F_p}N_r[p].
\end{aligned}
       \tag{NC.14}
\]
The dimension of $B_r$ is $(g-1)\varphi(\ell^r)$: for a
prime-power cover it is the new factor between $D_r$ and
$D_{r-1}$, so this follows by subtracting their etale
Riemann--Hurwitz genus formulas. The geometric $p$-rank
of an abelian variety is at most its dimension
\cite[Remark~8.5, the bound $r\le g$]{MilneAV}.
Thus (NC.14) gives
\[
\begin{aligned}
 \dim_{\mathbf F_p}B_r(k)[p]
 &\le \frac{(g-1)\varphi(\ell^r)}{f_r}\\
 &=\frac{(g-1)(\ell-1)\ell^{s-1}}{t}
 \qquad(r\ge s).
\end{aligned}
       \tag{NC.15}
\]
This is a uniform bound on the number of rational
$p$-primary generators.

For every $e\ge1$, however, the \emph{abstract, nongeometric}
module $\mathcal O_r/p^e\mathcal O_r$ with Frobenius
$\tau_r$ satisfies exactly the same semilinear descent
hypotheses. Its invariants are $R_r/p^eR_r$, by that descent,
and have exponent $p^e$. Thus the displayed descent identities
and generator bound do not by themselves supply an exponent
bound or vanishing. No such module is asserted to be a
Jacobian factor or part of a geometric tower.

\emph{7. Finite invariant groups also allow arbitrary abstract exponents.}
There is a compatible \emph{abstract, nongeometric} example
with one fixed group action. Let $G=\mathbf Z_\ell$, choose
arbitrary integers $e_r\ge1$, and put
\[
 A=\bigoplus_{r\ge1}\mathcal O_r/p^{e_r}\mathcal O_r,
 \qquad H_m=\ell^m\mathbf Z_\ell.
       \tag{NC.16}
\]
Here a topological generator of $G$ acts on the $r$th
summand by multiplication by $Z$, not by $\tau_r$.
Every element has finite support and an open stabilizer,
so this action on the discrete group is continuous.
For $r\le m$ the summand is fixed by $H_m$.
For $r>m$, the element $Z^{\ell^m}-1$ is a unit in
$\mathcal O_r$: its value at every primitive $\ell^r$th
root in characteristic $p$ is nonzero, since $p\ne\ell$.
It is also a unit modulo $p^{e_r}$, so that summand has no
invariants. Consequently
\[
\begin{aligned}
 A^{H_m}&=\bigoplus_{1\le r\le m}
                 \mathcal O_r/p^{e_r}\mathcal O_r,\\
 \exp(A^{H_m})&=p^{\max_{1\le r\le m}e_r}\qquad(m\ge1).
\end{aligned}
       \tag{NC.17}
\]
Every open-subgroup invariant group is finite, while these
exponents can grow arbitrarily. This is solely an algebraic
limitation; no curve, tower of Jacobians or compatible point
divisors realizing $A$ are claimed.
\end{proof}

The unresolved covering-route input is still an upper bound
on a suitable projected \emph{point-difference} order, or a
geometric argument forcing its faithful component to vanish
from the signed-orbit constraints. Norm compatibility supplies
neither conclusion through the formulas proved here. This
records exact consequences and limits, without turning the
absence of a theorem in a bounded source search into a theorem
of nonexistence, and without a claim of main-gate reduction.


\paragraph{Supplement to Finding A51: finite matrices and the missing rationality premise.}
Two elementary statements distinguish a finite matrix over a completed
group algebra from a fixed rational Frobenius determinant. The first is
an abstract module construction, with no asserted geometric realization.
The second gives a conditional sufficient premise for the exponent
criterion of A47. Neither supplies an actual curve tower or a pointed
primitive factor.

\emph{1. Completed coefficients, including semilinear Frobenius, permit
arbitrary large conductor exponents.}
Let $p\ne\ell$ be primes and put
\[
 \mathcal R=\varprojlim_n\mathbf Z_p[C_{\ell^n}],
 \qquad \gamma_n\text{ a compatible generator of }C_{\ell^n}.
 \tag{FX.1}
\]
The transition maps are induced by the cyclic-group quotients.
For $a\ge0$ define the component of $e_a\in\mathcal R$ at level $n$ by
\[
 e_a^{(n)}=
 \begin{cases}
 1,&n\le a,\\[1mm]
 \displaystyle\ell^{-(n-a)}
       \sum_{j=0}^{\ell^{n-a}-1}\gamma_n^{j\ell^a},&n>a.
 \end{cases}
 \tag{FX.2}
\]
These averages are compatible under the quotient maps. On a character
of exact order $\ell^b$, their values are $1$ for $b\le a$ and $0$
otherwise. Splitting the finite cyclic group over a finite unramified
extension of $\mathbf Z_p$ therefore proves
$e_ae_c=e_{\min(a,c)}$. Thus
\[
 \varepsilon_0=e_0,\qquad \varepsilon_a=e_a-e_{a-1}\ (a\ge1)
 \tag{FX.3}
\]
are orthogonal conductor idempotents. At level $n$ they vanish for
$a>n$ and sum to $1$ for $0\le a\le n$.

Choose any integers $b_a\ge2$. The expression
\[
 B=\sum_{a\ge0}p^{b_a}\varepsilon_a\in\mathcal R
 \tag{FX.4}
\]
is well defined, since it is a finite sum at each level. Multiplication
by $1-B$ is an automorphism of the free rank-one $\mathcal R$-module.
For every exact-conductor-$\ell^a$ character its value is
$1-p^{b_a}$, a $p$-adic unit. The corresponding cokernel of
$1-(1-B)$ has exact exponent $p^{b_a}$.

This remains an obstruction to a bound when the deck action is
nonconstant. Fix $u\in\mathbf Z_\ell^\times$, acting on $\mathcal R$
by $\gamma\mapsto\gamma^u$. It fixes each $\varepsilon_a$.
The automorphism
\[
           \mathcal F=(1-B)u
 \tag{FX.5}
\]
is semilinear over $\mathcal R$ and linear over $\mathbf Z_p$.
After a finite unramified coefficient extension splitting the
level-$a$ cyclic group, its exact-conductor characters form $u$-orbits
of length $d_a=\operatorname{ord}_{\ell^a}(u)$ for $a\ge1$
(and $d_0=1$). On each orbit
$\mathcal F$ is a cyclic permutation times $c_a=1-p^{b_a}$.
The relations for the additive cokernel of $1-\mathcal F$ eliminate
all but one basis vector and leave the relation
$1-c_a^{d_a}$. Consequently the exact-conductor additive cokernel
has exponent
\[
          p^{b_a+v_p(d_a)}.
 \tag{FX.6}
\]
Here is the valuation check, including $p=2$. If $v_p(x)\ge2$,
then a power prime to $p$ preserves the valuation of $(1+x)^m-1$:
the linear term has valuation $v_p(x)$ and all later terms have
larger valuation. Raising $1+x$ to the $p$th power instead increases
that valuation by one, because its linear term has valuation
$v_p(x)+1$ and every other term has larger valuation. Apply these
two observations to $x=-p^{b_a}$ and factor $d_a$ into its $p$-part
and prime-to-$p$ part. Finite unramified scalar extension preserves
the elementary-divisor exponents, proving (FX.6) over $\mathbf Z_p$
as well. Since the $b_a$ are arbitrary, finite matrix size and this
semilinear structure alone allow arbitrarily rapid exponent growth.
No Weil-weight, polarization, Jacobian, or point-section condition
is claimed for these modules. This does not contradict the actual
descent identities of A45 or the conditional bounds in A51.

\emph{2. A fixed nonzero rational function has eventually constant
cross-prime valuation.}
Let $K/\mathbf Q_p$ be finite, with residue field
$\mathbf F_{p^s}$, uniformizer $\varpi$, and valuation
$v_\varpi(\varpi)=1$. For nonzero $f(X)\in K(X)$ write
\[
             f(X)=\varpi^\mu G(X)/H(X),
 \tag{FX.7}
\]
where $G,H$ are integral Laurent polynomials, each with a unit
coefficient. Clearing monomials does not change the valuations
at roots of unity; let $D$ bound the degrees of their nonzero
reductions after doing so. If $\zeta$ has exact order $\ell^a$ and
\[
                  \operatorname{ord}_{\ell^a}(p^s)>D,
 \tag{FX.8}
\]
then
\[
                        v_\varpi(f(\zeta))=\mu.
 \tag{FX.9}
\]
Indeed adjoining prime-to-$p$ roots of unity is unramified, and
reduction preserves their exact orders. The minimal polynomial of
$\bar\zeta$ over $\mathbf F_{p^s}$ has degree
$\operatorname{ord}_{\ell^a}(p^s)$. It cannot divide either nonzero
reduction of degree at most $D$, so both evaluated polynomials are
units. Finally $\operatorname{ord}_{\ell^a}(p^s)\to\infty$:
otherwise arbitrarily large $\ell^a$ would divide one of the
finitely many nonzero integers $p^{sj}-1$ with $j$ bounded.
This proves (FX.9) for every sufficiently large $a$ and every
primitive root of that level.

In particular, suppose an actual fixed integral Laurent matrix
$M(X)$ has nonzero determinant $f(X)=\det(1-M(X))$.
At high conductor the sum of the nonnegative elementary-divisor
valuations of $1-M(\zeta)$ is the constant $\mu$ in (FX.9).
Its cokernel therefore has uniformly bounded $p$-primary exponent.
Only an additional identification of the actual rational primitive
groups with such evaluated cokernels, or subgroups of finite direct
sums of them, would give a uniform $M_n$ for A47. In genus two and
characteristic two, along odd prime powers $n=\ell^a$, a uniform
$M_n$ would eventually satisfy $n>(2M_n-1)^2$.
No such geometric identification is obtained here.

\emph{3. Exact primary scope and the unresolved premise.}
Finite presentation of the full proper etale fundamental group is
known \cite[Theorem~1.1]{ESS2022}
\cite[Theorem~1.1 and Corollary~1.2]{Shusterman2022}, but the
presentations are profinite, not asserted to have discrete finite-word
relators or finite-word arithmetic Frobenius images. The familiar
surface-word presentation is for the prime-to-$p$ quotient
\cite[p.~4]{ESS2022}; its open subgroups have zero continuous
$H^1(-,\mathbf F_p)$ and do not retain the full curve's
characteristic-$p$ etale cohomology. Arithmetic Fox constructions in
the inspected number-field pro-$\ell$ setting have completed
power-series coefficients, including explicit infinite Milnor sums
\cite[Sections~4.1--4.3, Proposition~4.3.4]{KMT2017}.

The geometric skew-abelian tower paper of Asvin~G. states a finite
matrix description over a $\mathbf Z_\ell$ power-series ring.
Its return matrix is the product
around the whole growing Frobenius character orbit
\cite[Definition~8, Lemma~17, Section~3.1, equation~(6)]{Asvin2023}.
The congruences are at $\ell$, and its polynomial truncation uses
$\zeta_{\ell^a}-1$ being $\ell$-adically small
\cite[Theorems~23,26 and Section~3.3, p.~2169]{Asvin2023}.
That element is a $p$-adic unit for $p\ne\ell$. Neither this
truncation nor a determinant of a growing orbit product supplies
the single fixed rational function required in (FX.7).
There is also an integral issue in its geometric argument:
Lemma~11 claims that the natural pullback on integral $H^1$
is a direct factor for every finite curve map, and its proof
infers surjectivity on Tate modules from a surjective map of
abelian varieties. For
$[\ell]:E\to E$ on an elliptic curve, both induced maps are
multiplication by $\ell$ and have cokernel
$(\mathbf Z/\ell\mathbf Z)^2$. Thus that inference, and the
claimed natural-image saturation, fail integrally. We do not
use its geometric integral finite-free premise as an established
input. The abstract Section~3 coefficient-prime scope and the
two independent algebra proofs above do not depend on that premise.
Sinnott's rational-measure theorem likewise assumes an actual fixed
rational Fourier function and an appropriate unit-group support
condition \cite[Section~1.16 and Theorem~3.2]{Sinnott1987}.
No identification with that transform is provided by finite
profinite presentation or matrix size. These are completed
algebraic checks on the proposed route, not a geometric
counterexample or a closure of the primitive-factor supply gate.


\paragraph{Supplement to A51: full deck transitivity in characteristic two.}
Let $C/k$ be a smooth projective geometrically connected genus-two
curve over a finite field of characteristic two. Let $D\to C$ be
a geometrically connected etale cover of odd degree $n>1$,
geometrically Galois with cyclic group generated by $\sigma$.
The deck group need not be constant. Write
\[
 F\sigma F^{-1}=\sigma^a,\qquad
 f=\operatorname{ord}_n(a)=\varphi(n),
 \tag{TD.1}
\]
where $F$ is arithmetic Frobenius on geometric points.
Thus $a$ acts transitively on all faithful characters.
Let $B=\operatorname{im}(ne_n:J_D\to J_D)$ be the actual
primitive conductor-$n$ subvariety. Then
\[
\begin{aligned}
 &f(B)=0\quad\hbox{or}\quad f(B)=\dim B=\varphi(n),\\
 &B(k)[2]=0\quad\Longleftrightarrow\quad f(B)=0,\\
 &f(B)=\varphi(n)\quad\Longrightarrow\quad
                         \dim_{\mathbf F_2}B(k)[2]=1.
\end{aligned}
 \tag{TD.2}
\]
Here $f(B)$ denotes geometric two-rank. In particular, in this
fully transitive case the rational primitive-unit test of A46
is equivalent to zero geometric primitive two-rank. Full
arithmetic transitivity does not remove the rational two-torsion
of an ordinary primitive factor.

\begin{proof}
Use the full cyclotomic algebra and its reduction
\[
 \mathcal O=\mathbf Z_2[Z]/\Phi_n(Z),\qquad
 A=\mathcal O/2\mathcal O,\qquad \tau(Z)=Z^a.
 \tag{TD.3}
\]
Since $n$ is odd, $A$ is a finite product of finite fields.
After extending to $\overline{\mathbf F}_2$ its factors are
indexed by all primitive $n$th roots. Condition (TD.1) makes
$\tau$ transitive on those indices, so
\[
              \dim_{\mathbf F_2}A=f,\qquad A^\tau=\mathbf F_2.
 \tag{TD.4}
\]
In particular it permutes the field factors of $A$ transitively.

The actual geometric group $V=B[2](\bar k)$ is an $A$-module
by (SD.8)--(SD.9a), and $F$ is a bijective
$\tau$-semilinear map on $V$.
Semilinearity makes the ranks of $V$ over the field factors
of $A$ equal. Consequently $V$ is free of some rank $r$ over
$A$ and $\dim_{\mathbf F_2}V=rf$.
The genus-two unramified character formula gives
$\dim B=\varphi(n)=f$: each nontrivial character has degree
two in the zeta numerator, and the actual primitive image
retains exactly the faithful characters
\cite[B.2]{A1run}; see also (OQ.7)--(OQ.8).
The two-rank bound $f(B)\le\dim B$
\cite[Remark~8.5, the bound $r\le g$]{MilneAV}
therefore forces $r=0$ or $r=1$.

If $r=0$, there is no geometric or rational two-torsion.
If $r=1$, choose an $A$-basis of $V$. In that basis
$F=c\tau$ for a unit $c\in A^\times$, and hence
\[
 F^f=c\tau(c)\cdots\tau^{f-1}(c)
                         \in(A^\tau)^\times=\mathbf F_2^\times=
                         \{1\}.
 \tag{TD.5}
\]
The finite etale torsor descent already proved in A45,
(SD.3)--(SD.6), now applies to this module killed by two.
It gives
\[
                  V\simeq A\otimes_{\mathbf F_2}V^F.
 \tag{TD.6}
\]
Equivalently, this is the norm-one Hilbert90 calculation for
the full finite etale product algebra. No hypothesis that
$A$ is a field, and no division by $f$, is needed.
Comparing dimensions in (TD.6) gives
$\dim_{\mathbf F_2}V^F=1$.
Since $V^F=B(k)[2]$, all assertions of (TD.2) follow.
The descent input is the same unrestricted module theorem
\cite[Proposition~35.3.9, Tag~\texttt{023N}]{Stacks023N}.
\end{proof}

The determinant makes the obstruction equally explicit.
Put $q=\#k$. Over an auxiliary-adic splitting field the
faithful summands of the Tate module have dimension two,
and $F$ cyclically permutes all $f$ of them.
Transport a basis successively by $F$ around this cycle;
the successive block maps are identities except for the last,
which is the return matrix $R$ on the first summand.
Eliminating these successive blocks in $1-TF$ gives
$\det(1-TF)=\det(1-T^fR)$. Therefore the
actual zeta numerator of $B$ has the form
\[
                       H_B(T)=1+S T^f+q^fT^{2f},\qquad S\in\mathbf Z.
 \tag{TD.7}
\]
The integrality and top coefficient are the Frobenius
polynomial identities for the $f$-dimensional abelian variety
$B$ \cite[Theorem~19.1]{MilneAV}.
In the ordinary case, (TD.2) makes
$\#B(k)=H_B(1)=1+S+q^f$ even, so $S$ is odd and
\[
                        \overline{H_B(T)}=1+T^f.
 \tag{TD.8}
\]
Thus the reduced return scalar equals one, rather than being
forced away from one by polarization or canonical symmetry.
The rational two-primary group has one generator and its
exponent is exactly
$2^{v_2(1+S+q^f)}$; no bound on that positive valuation has
been obtained. In the zero-two-rank case the rational group
is zero and $S$ is even. These are statements about the
actual primitive factor, not abstract module realizations.

\paragraph{Why the unrestricted quotient supply does not target this tower.}
In the rank-one $\ell$-tower of (NC.10)--(NC.11), with
$n_r=\ell^r$ and minimal deck field $K_r=k_{f_r}$,
one has $f_r=t\ell^{r-s}$ for large $r$.
Over $K_r$, the prescribed pointed cover has a constant
cyclic Lang quotient of size $n_r$.
Set $Q_r=q^{f_r}$. The genus-two Weil bounds give
\cite[Theorem~11.1]{MilneJV}
\[
\begin{aligned}
 \#C(K_r)&\ge Q_r+1-4\sqrt{Q_r}>n_r
                                      &&\hbox{for all large }r,\\
 \#\ker\bigl(J_C(K_r)\twoheadrightarrow C_{n_r}\bigr)
   &\ge\frac{(\sqrt{Q_r}-1)^4}{n_r}\longrightarrow\infty.
\end{aligned}
 \tag{TD.9}
\]
The kernel bound uses the Jacobian Weil bound
\cite[Theorem~19.1]{MilneAV}.
Thus the empty-coset inequality $|H|>\#C(K_r)$ and the fixed
small kernel used in A49's unrestricted-quotient supply proof
are unavailable for this prescribed tower quotient.
That proof may still supply a unit character on another,
much larger quotient, with an uncontrolled conductor.

Norms from a common refinement constrain the inflated
characters and kill its faithful component, as (NC.8)--(NC.9)
show. Moreover (NC.15) bounds generators of each individual
primitive factor, not the sum of the new factors on all
levels or independent covers. Neither statement puts the
supplied unit orbit in the remaining component of the same
compatible two-point class. In the fully transitive ordinary
case, (TD.2)--(TD.8) specifically show that the presumed
missing rational two-torsion is present. The unresolved
input remains a useful zero-rank cover, a controlled actual
return valuation, or an additional pointed geometric
argument. No universal existence conclusion follows here.


\paragraph{Supplement to Finding A51: fields of quintic iteration and available canonical detectors.}
Let $k=\mathbf F_q$ have characteristic two, let
$\Sigma=\{2,3,5,13\}$, and set
\[
 k_\Sigma=\bigcup_{\substack{m\ge1\;:\text{ every prime}\,
                              \text{dividing }m\text{ lies in }\Sigma}}
                              \mathbf F_{q^m}.
 \tag{FI.1}
\]
The pointed geometric quintic-cover and reflection choices of A33,
iterated from a genus-two curve and a rational point over $k$, have
representatives over $k_\Sigma$. Nevertheless, on every fixed
genus-two target defined over a finite subfield of $k_\Sigma$,
there are infinitely many points in this same field extension whose
canonical class has nonzero $3$-primary component. These two assertions
do not force the prescribed iterated reflection images to meet those
points.

\begin{proof}
\emph{1. One uniform degree support for the geometric choices.}
Let $(C,P)$ be defined over a finite field $K\subset k_\Sigma$.
Its principal polarization gives a nondegenerate alternating Weil
pairing on $J_C[5]$, and Frobenius acts as a similitude, with
multiplier $\#K$ modulo $5$. Consequently all $5$-torsion becomes
rational over an extension of degree dividing
\[
 \begin{aligned}
 G_5&=\#\operatorname{GSp}_4(\mathbf F_5)
      =(5-1)5^4(5^2-1)(5^4-1)\\
    &=37{,}440{,}000=2^9\,3^2\,5^4\,13.
 \end{aligned}
 \tag{FI.2}
\]
For completeness, symplectic bases in dimension four over
$\mathbf F_t$ are counted by choosing a nonzero first vector
in $t^4-1$ ways, its pairing-one partner in $t^3$ ways, then
a nonzero vector in their two-dimensional orthogonal complement
in $t^2-1$ ways and its partner in $t$ ways. This gives
$\#\operatorname{Sp}_4(\mathbf F_t)=t^4(t^2-1)(t^4-1)$;
the surjective multiplier adds the factor $t-1$.
Lagrange's theorem gives the divisibility of the Frobenius order.

Over that extension every geometric cyclic quotient of degree five is
Frobenius invariant with constant deck group. The section at $P$
extends the geometric quotient of the abelianized fundamental group
to an arithmetic quotient by sending that section to zero. This
constructs its constant, $P$-split model; all five lifts of $P$
are rational. Equivalently, among the constant cyclic twists the
$P$-split choice is the unique one, as in A48. Thus specifying a
geometric cover and a geometric lift of $P$ adds no obstruction
to representing the pointed isomorphism class in this field.

The hyperelliptic involution of $C$ is defined over $K$ and acts
as inversion on the cyclic quotient. It has five lifts to the
cover, all involutions, as in A33. Write one as $\tau$ and a deck
generator as $\sigma$. Because the deck group is now constant,
Frobenius acts on the lift set by
\[
                  \tau\longmapsto\sigma^b\tau
                  \qquad(b\in\mathbf Z/5\mathbf Z).
 \tag{FI.3}
\]
This is a translation of that five-element torsor, so an extension
of degree dividing five defines every reflection. Their genus-two
quotients and all images of the five rational lifts are defined
there too. The local quotient/genus argument is the one already
proved in A33.

Thus one step can be represented after a degree dividing
\[
                  B=5G_5=187{,}200{,}000.
 \tag{FI.4}
\]
Inductively every depth-$j$ choice has a representative over
$\mathbf F_{q^{B^j}}$: every intermediate relative degree divides
$B$, and the same field contains all such relative subextensions.
This asserts a common prime support, not a single finite field
for all depths. It concerns geometric pointed choices represented
by their split models, not arbitrary preselected equations or
arbitrary forms over unrelated fields. If another constant twist
is first specified, its five-element fibre torsor is killed by
an extension of degree dividing five, which has the same prime
support.

\emph{2. Canonical $3$-primary detectors already occur densely inside
that constant extension.}
Fix any genus-two target $E/K$ with finite $K\subset k_\Sigma$.
An extension inside $k_\Sigma$ rationalizes $J_E[3]$, since
\[
 \#\operatorname{GSp}_4(\mathbf F_3)
       =103{,}680=2^8\,3^4\,5.
 \tag{FI.5}
\]
After a further extension of degree a power of five, the Weil
bound gives a rational point $O\in E(K)$; enlarge $K$ accordingly.
Put $Q=\#K$, $G=J_E(K)[3^\infty]$, and $d=\#G\ge3^4$.
If $\Phi$ is its Frobenius, it acts as the identity on $J_E[3]$.
For $n=5^m$ the polynomial
\[
  V_n(\Phi)=1+\Phi+\cdots+\Phi^{n-1},\qquad
  \Phi^n-1=(\Phi-1)V_n(\Phi)
 \tag{FI.6}
\]
is invertible on $3$-power torsion: modulo three it is $nI$,
and a nonzero kernel at a higher finite level would contain an
element of order three. It follows that, as actual subgroups,
\[
                  J_E(\mathbf F_{Q^{5^m}})[3^\infty]=G.
 \tag{FI.7}
\]
Apply the already proved fixed Lang-cover count A42 with the
singleton prime set $\{3\}$ and the Abel map $R\mapsto[R-O]$.
Here its modulus is $3$, so every $n=5^m$ is admissible. For
every $h\in G$ it gives
\[
 \left|\#\{R\in E(\mathbf F_{Q^n}):[R-O]_3=h\}
                 -\frac{Q^n+1}{d}\right|
       \le\frac{2(d+1)}d Q^{n/2}.
 \tag{FI.8}
\]
This uses the connected Lang quotient of degree $d$ and genus
$d+1$, with its exact norm scaling; no unproved density
principle is substituted for A42.

For the canonical class $\kappa_R=[\omega_E(-2R)]$ one has
\[
             (\kappa_R)_3=(\kappa_O)_3-2[R-O]_3.
 \tag{FI.9}
\]
Exactly one $h\in G$ makes the right side zero, since
multiplication by two is invertible on $G$. Each of the other
$d-1$ components occurs on unboundedly many points by (FI.8).
All these points belong to $E(k_\Sigma)$ and have a nonzero
canonical $3$-primary component. Their relative density along
these fields tends to $(d-1)/d$, and in particular they are
Zariski dense. Since $\operatorname{ord}_3(2)=2$, A29 applies
to each such pointed target and supplies its prime-power
canonical detector.
\end{proof}

The pro-two or single-auxiliary-prime constant-tower restrictions
in A44 concern a fixed Jacobian and its Frobenius polynomial;
they do not uniformly restrict this changing family of reflection
Jacobians. The preceding calculation gives the sharper relevant
endpoint: the common field support is compatible with abundant
good canonical points on every fixed target. It still does not
force any prescribed finite set of reflection images, or the
full iterated orbit from a specified point, to contain one.
Neither a Hecke-orbit theorem nor the Lang count establishes
that missing incidence. No main-gate reduction is claimed.


\paragraph{Supplement to Finding A51: weights and principal polarization do not force a small exponent.}
Fix $q=p^e$. For every integer $f\ge1$ there exists an ordinary,
principally polarized abelian variety $A_f/\mathbf F_q$ of dimension
$f$ such that
\[
 \begin{aligned}
 H_{A_f}(T)&=1-T^f+q^fT^{2f},\\
 A_f(\mathbf F_q)&\simeq\mathbf Z/q^f\mathbf Z,\\
 \dim_{\mathbf F_p}A_f(\mathbf F_q)[p]&=1,
 \qquad v_p\bigl(H_{A_f}(1)\bigr)=ef.
 \end{aligned}
 \tag{WV.1}
\]
Thus Weil weights, integrality, principal polarization, the displayed
cyclic Frobenius-polynomial shape, and one rational $p$-primary
generator do not by themselves give the small exponent required in
A47. This assertion concerns actual abelian varieties. No cyclic
deck action, Jacobian-cover realization, or compatibility in a
single geometric tower is asserted.

\begin{proof}
Put $Q=q^f$. Waterhouse's elliptic isogeny-class theorem gives
an elliptic curve $E/\mathbf F_Q$ with trace one: the trace is
prime to $p$ and $1<2\sqrt Q$. This is its ordinary case
\cite[Theorem~4.1(1), pp.~536--538; Corollary~4.4, p.~540]{Waterhouse1969AV}.
Consequently
\[
       H_E(U)=1-U+QU^2,
       \qquad \#E(\mathbf F_Q)=Q.
 \tag{WV.2}
\]
The finite group $E(\mathbf F_Q)$ is therefore a nonzero
$p$-group. An ordinary elliptic curve has exactly $p$ geometric
points killed by $p$, so this rational group has exactly one
$p$-primary generator. The elementary classification of finite
abelian groups makes it cyclic of order $Q$.

Take the Weil restriction
\[
               A_f=\operatorname{Res}_{\mathbf F_Q/\mathbf F_q}E.
 \tag{WV.3}
\]
It is an abelian variety of dimension $f$, geometrically the
product of the $f$ Frobenius conjugates of $E$, and is ordinary.
The canonical degree-one polarization of $E$ gives a principal
polarization of $A_f$: after base change it is the product of
those canonical polarizations. Preservation of the principal
polarization, including its descent, is the Weil-restriction
result in \cite[Proposition~1.22, pp.~27--28]{Diem2001Weil}.
The defining functor gives the equality of groups
\[
                  A_f(\mathbf F_q)=E(\mathbf F_Q).
 \tag{WV.4}
\]
On an auxiliary-adic Tate module, Frobenius has the induced
representation from $E$; the characteristic-polynomial formula
for Weil restriction is
\[
                  H_{A_f}(T)=H_E(T^f)
                              =1-T^f+q^fT^{2f}
 \tag{WV.5}
\]
\cite[Proposition~1.21, p.~26]{Diem2001Weil}.
Combining (WV.2), (WV.4), and (WV.5) proves every assertion
of (WV.1), in particular the exact exponent $q^f$ rather than
merely its order as an upper bound.

The Weil condition can also be read directly from the monic
polynomial
\[
                    X^{2f}-X^f+q^f.
 \tag{WV.6}
\]
If $\alpha$ is any root, then $\alpha^f$ is a root of
$Y^2-Y+q^f$. The latter two roots are complex conjugates of
absolute value $q^{f/2}$, so every $\alpha$ has absolute value
$\sqrt q$. The coefficients are integral and the reciprocal
identity for a dimension-$f$ abelian variety holds exactly.
Thus the large valuation is not produced by abandoning the
Weil constraints: the middle coefficient is simply $-1$,
and $1+(-1)$ cancels at $T=1$.
\end{proof}

For comparison with a putative transitive prime-power cover,
set $n=\ell^r$ and $f=\varphi(n)$, with $\ell$ odd and
$\ell\ne p$. These abelian varieties have the same polynomial
shape and rational generator number occurring in (TD.7), but
their exponents $M=q^{\varphi(n)}$ grow exponentially in $n$.
In particular they do not satisfy $n>(2M-1)^2$ for large $n$.
This defeats a deduction from only the listed numerical and
polarization constraints. It does not supply the integral
$\mathbf Z[\zeta_n]$ deck action, proper-conductor norm maps,
cover point-count congruences, or a fixed underlying genus-two
curve. The varieties $E$ may vary with $f$. Those additional
geometric and tower constraints remain possible sources of an
actual bound; no failure of such a bound for genuine curve
towers is proved here.


\subsection{Small integral annihilators and a fixed-field lattice bound}
\label{sec:small-integral-annihilator}

\paragraph{Finding A52 (small integral annihilators and a fixed-field lattice bound).}
Let $C/k$ be a smooth projective geometrically connected curve of genus
$g\ge2$ over a finite field of characteristic $p$. Let $D\to C$ be
geometrically connected and etale, geometrically Galois with an odd cyclic
group $G=\langle\sigma\rangle$ of order $n$ prime to $p$. The deck group
need not be constant over $k$. Suppose that a collision of distinct
point sections at $P,Q\in C(k)$ gives compatible lifts
$P',Q'\in D(k)$ and
\[
                  \delta=[Q'-P']\in J_D(k)[p^\infty].
\tag{IA.1}
\]
Assume that there is an integral operator $A\in\mathbf Z[G]$ such that
\[
 A\delta=0,\qquad A(\zeta_n)\ne0,\qquad
 n>(\|A\|_1-1)^2,
\tag{IA.2}
\]
where $\zeta_n$ is a primitive complex $n$th root and the norm is the
sum of the absolute values of the coefficients in the degree-$<n$
representative. Then the given geometric cover is the pullback of an
etale cyclic degree-$n$ cover between genus-one curves along a separable
map
\[
 C_{\bar k}\longrightarrow E,\qquad
 \deg(C_{\bar k}/E)=p^a=e>1,\qquad e\le g-1,
\tag{IA.3}
\]
totally ramified at $P$ and $Q$. In particular (IA.2) excludes the
collision whenever $2\le g\le p$, or whenever there is no map (IA.3)
totally ramified at $P$.

The criterion permits arbitrary integer coefficients and does not
assume that the operator is defined over $k$ or that its degree is
coprime to $n$. It is a consequence of an operator on the actual class
(IA.1), not of a freely chosen torsion module.

\begin{proof}
Choose, among all integral annihilators of $\delta$ with nonzero faithful
evaluation, an operator $B=\sum b_i\sigma^i$ of least positive coefficient
norm $d$. Thus $d\le\|A\|_1$. The greatest common divisor of its nonzero
coefficients is a power of $p$. Indeed, if a prime $\ell\ne p$ divided
every coefficient, then $B=\ell B_1$ and $\ell B_1\delta=0$ would imply
$B_1\delta=0$, since $B_1\delta$ is $p$-primary. This would contradict
minimality, and faithful evaluation would remain nonzero.

The two deck orbits of $P'$ and $Q'$ are free and disjoint. Hence a
principal function $f$ with
\[
                  \operatorname{div}(f)=B(Q'-P')
\tag{IA.4}
\]
has degree exactly $d$. It is separable: if $f$ were a $p$th power,
every $b_i$ would be divisible by $p$, and the divisor of its $p$th root
would show that $B/p$ is a smaller annihilator with nonzero faithful
evaluation. Over the perfect field $\bar k$, inseparability of a
nonconstant function is equivalent to being a $p$th power.

Put $K=\bar k(D)$, let $F$ be generated by all $G$-translates of $f$,
and let $v:D_{\bar k}\to Z$ be the corresponding map of smooth curves.
Write $e=[K:F]$. The extension is separable. Poonen's controlled-field
result, applied to the finitely many degree-$d$ generators, gives
\[
                e\mid d,\qquad g(Z)\le(d/e-1)^2
\tag{IA.5}
\]
\cite[Proposition~2.1, Lemma~2.2 and Corollary~2.3]{Poonen2007Gonality}.
Since $g(D)=n(g-1)+1$ and $n>(d-1)^2$, one has $e>1$.

The group $G$ preserves $F$, so $v$ is equivariant. We first prove
injectivity on the two relevant orbits without assuming that $G$ acts
freely everywhere on $Z$. Put $x=v(P')$, and suppose $\sigma^j x=x$.
The ramification index $r_P$ of $v$ is constant along the $P'$ orbit.
Writing $f=h\circ v$, coefficient comparison in (IA.4) gives
\[
 -b_i/r_P=\operatorname{ord}_{\sigma^i x}(h)
        =\operatorname{ord}_{\sigma^{i+j}x}(h)=-b_{i+j}/r_P
\tag{IA.6}
\]
for every $i$, including zero coefficients. Thus $\sigma^jB=B$.
Evaluation at $\zeta_n$, where $B(\zeta_n)\ne0$, forces $j=0$ modulo
$n$. The same argument applies to $v(Q')$. In particular the action of
$G$ on $Z$ is faithful, and both image orbits have exactly $n$ points.

These two image orbits cannot coincide. If they did, the coefficient
comparison used in A47, with the positive ramification indices $r_P,r_Q$,
would give, for some $j$,
\[
                         r_QB=-r_PB\sigma^j.
\tag{IA.7}
\]
Iterating this equality through the odd order $u$ of $\sigma^j$ gives
$(r_Q^u+r_P^u)B=0$ in the torsion-free integral group ring, which is
impossible. The signs here are integer divisor signs in every
characteristic.

Choose any nonzero coefficient of $B$. A fibre of $v$ through its
corresponding zero or pole of $f$ contains only points in the support of
(IA.4), because it is a zero or pole fibre of $h$. Injectivity on each
image orbit and disjointness of the two image orbits show that this
fibre consists of a single point, of ramification index $e$. The same
holds for every nonzero coefficient, and equivariance then makes all
$2n$ points in the two original orbits totally ramified. In particular
\[
                           e\mid\gcd_i(b_i).
\tag{IA.8}
\]
The preceding minimality argument now shows that $e$ is a nontrivial
power of $p$, and hence $\gcd(e,n)=1$.

This last coprimality makes the $G$-action on all of $Z$ free. Otherwise
a point stabilizer would contain an element of prime order $\ell\mid n$.
Its free action on $D$ would partition the fibre divisor into
$\ell$-element orbits of equal multiplicities, implying $\ell\mid e$.
Thus $Z\to E=Z/G$ is etale of degree $n$. Since $g(Z)<n$, the identity
$g(Z)-1=n(g(E)-1)$ gives $g(Z)=g(E)=1$.

At each of the $2n$ totally ramified points the different exponent is
at least $e$: the leading derivative of a local parameter of order
$e$ vanishes because $p\mid e$, while separability makes a later
derivative nonzero. The local different formula and Riemann--Hurwitz
therefore give
\[
                  2ne\le\deg\operatorname{Diff}(v)=2n(g-1)
\tag{IA.9}
\]
\cite[Lemma~53.12.4, Tag~\texttt{0C1F}]{StacksRH}.
This proves $e\le g-1$.

Finally, faithfulness of the action on $F$ gives $K=FK^G$ inside the
Galois extension $K/K^G$. Consequently $C_{\bar k}\to E$ has degree
$e$, and the connected etale fibre product $C_{\bar k}\times_E Z$ has
function field $K$. It is $D_{\bar k}$. Ramification is unchanged by
the etale base change, so the map is totally ramified at $P,Q$.
This proves (IA.3).
\end{proof}

\paragraph{The lattice bound in one fixed prime-power tower.}
Suppose now that the geometric deck groups form a compatible
$\mathbf Z_\ell$ tower, where $\ell$ is odd and $\ell\ne p$, over the
fixed finite field $k$. Let the fixed deck Frobenius multiplier be
$a\in\mathbf Z_\ell^\times$, assumed not to be a root of unity. Put
\[
 t=\operatorname{ord}_{\mathbf F_\ell^\times}(a),\qquad
 s=v_\ell(a^t-1),\qquad
 E=\mathbf Q(\zeta_{\ell^s})^{\langle a\rangle},\qquad
 h=[E:\mathbf Q]=\varphi(\ell^s)/t.
\tag{IA.10}
\]
Here the action in the definition of $E$ is reduction modulo $\ell^s$.
For every $r\ge s$, put $n=\ell^r$ and
$\mathcal O_r=\mathbf Z_p[Z]/\Phi_n(Z)$. The minimal deck field has
degree $f_r=t\ell^{r-s}$. Its deck action $\tau$ has invariants
\[
 R=\mathcal O_r^\tau=\mathcal O_E\otimes\mathbf Z_p,
                      \qquad \operatorname{rank}_{\mathbf Z_p}R=h.
\tag{IA.11}
\]
These identifications are compatible with the embeddings from level
$\ell^s$. They follow because $\langle a\bmod\ell^r\rangle$ contains
the whole kernel of reduction to level $\ell^s$, as in (NC.11).
Invariants of the integral cyclotomic ring are the integers of the
fixed field, and flat tensoring with $\mathbf Z_p$ preserves these
invariants.

Fix a $\mathbf Z$-basis $\theta_1,\ldots,\theta_h$ of $\mathcal O_E$.
Choose integral polynomial representatives
$\theta_i=C_i(\zeta_{\ell^s})$ and put
\[
                         C_E=\sum_{i=1}^h\|C_i\|_1.
\tag{IA.12}
\]
This constant is fixed once and for all and is independent of $r$.
Let $P_n$ be the actual primitive factor of $J_{D_r}$, and define
\[
 \begin{aligned}
 N_r&=P_n(k)[p^\infty],\\
 M_r&=\exp(N_r),\\
 W_r&=[R:\operatorname{Ann}_R(N_r)].
 \end{aligned}
\tag{IA.13}
\]
The invariant ring acts on $N_r$ by actual deck endomorphisms and
semilinear descent, as in A45. Use $M_r=W_r=1$ for the zero group.
If the cover has a rational lift of the prescribed point, then
\[
 n>\left(2\min\{M_r,C_E\lfloor W_r^{1/h}\rfloor\}-1\right)^2
\tag{IA.14}
\]
excludes a collision when $2\le g\le p$, and otherwise gives the
elliptic obstruction (IA.3).

\begin{proof}
Let $I=\mathcal O_E\cap\operatorname{Ann}_R(N_r)$. The index of $I$
in $\mathcal O_E$ is exactly $W_r$: the finite quotient is $p$-primary,
and the dense map $\mathcal O_E\to R$ is surjective modulo the open
annihilator ideal. Set $L=\lfloor W_r^{1/h}\rfloor$.
The $(L+1)^h>W_r$ elements
\[
                 \sum_{i=1}^h u_i\theta_i,\qquad 0\le u_i\le L,
\tag{IA.15}
\]
contain two in the same coset modulo $I$. Their nonzero difference
$c\in I$ has a polynomial representative at level $n$ of coefficient
norm at most $C_E L$: substitute $Z^{\ell^{r-s}}$ into the fixed
$C_i$, and use coefficient differences of absolute value at most $L$.
Because $E$ is a field, this nonzero $c$ has nonzero faithful
cyclotomic evaluation.

For any hypothetical collision put $b=e_n\delta\in N_r$, where
$e_n$ is the actual primitive projector on $p$-primary groups. This
is well defined since $p\nmid n$, and it is rational since the
integral endomorphism $ne_n$ is rational. The integral operator
\[
                  A=(1-\sigma^{\ell^{r-1}})c
\tag{IA.16}
\]
vanishes on every proper conductor, kills $b$, and thus kills all of
$\delta$. Its faithful evaluation is nonzero, and
$\|A\|_1\le2C_E L$. The alternative operator
$M_r(1-\sigma^{\ell^{r-1}})$ has norm $2M_r$ and the same properties.
Apply the preceding criterion to the better bound. This proves
(IA.14).
\end{proof}

\paragraph{What the bound does and does not supply.}
The same lattice argument for the annihilator of one actual class
$b\in N_r$ replaces $W_r$ by $\#(Rb)$. The uniform condition (IA.14)
uses the annihilator of the whole actual rational group so that it
does not depend on a hypothetical collision.
Writing $R=\prod_vR_v$, with unramified residue degrees $h_v$ and
component exponents $p^{e_v}$, gives
\[
                       W_r=p^{\sum_v h_ve_v},\qquad
                       M_r=p^{\max_v e_v}.
\tag{IA.17}
\]
Thus the lattice expression can be smaller than the maximum-exponent
expression when component exponents differ sufficiently. No example
outside the union of the previously proved A45--A47 tests is asserted
here, and no universal estimate for $W_r$ is established.

In the fully transitive case $E=\mathbf Q$, choose the basis $1$.
Then $h=C_E=1$ and $W_r=M_r$, so (IA.14) is exactly the previous
prime-power inequality $n>(2M_r-1)^2$. More generally a bounded
number of rational generators does not bound the $e_v$. The valid
module $R/p^uR$ has $W_r=p^{uh}$ and $W_r^{1/h}=p^u$, with no saving
over its exponent. This last module is an algebraic illustration,
not a claimed Jacobian-cover realization. The variable-quotient
single-packet supply of A49 does not bound the remaining component
exponents or target a prescribed tower. Hence the new sufficient
integral-lattice condition leaves an arithmetic witness-supply
requirement.


\paragraph{Supplement to Finding A52: the norm lower bound and the trace boundary.}
In a compatible cyclic $\ell$-power tower as above, write
$\delta_r=[Q_r-P_r]$, let $b_r=e_{\ell^r}\delta_r$ be its actual
rational primitive projection, and put
\[
 \mathcal O=\mathbf Z_p[Z]/\Phi_{\ell^r}(Z),\qquad
 R=\mathcal O^\tau,\qquad f=\operatorname{ord}(\tau),\qquad
 h=\varphi(\ell^r)/f,\qquad W=\#(Rb_r).
\tag{TN.1}
\]
The rings here are the full finite-etale product algebras, and the
module is the actual primitive torsion module. Flatness in the
descent isomorphism of A45 applies to the inclusion
$Rb_r\subset P_{\ell^r}(k)[p^\infty]$ and gives
\[
                       \mathcal O b_r\simeq
                       \mathcal O\otimes_R Rb_r,
                   \qquad \#(\mathcal O b_r)=W^f.
\tag{TN.2}
\]
In particular bounded rational generator number does not bound the
size of the geometric deck-orbit module.

If $A\in\mathbf Z[C_{\ell^r}]$ annihilates $\delta_r$ and has nonzero
faithful evaluation, then
\[
                              \|A\|_1\ge W^{1/h}.
\tag{TN.3}
\]
Indeed, $\alpha=A(\zeta_{\ell^r})$ annihilates $\mathcal O b_r$.
This is a cyclic $\mathcal O$-module, and the inclusion of ideals
$(\alpha)\subset\operatorname{Ann}_{\mathcal O}(b_r)$ makes
$\mathcal O/(\alpha)$ surject onto it. The nonzero integer
cyclotomic norm therefore satisfies
\[
 W^f\mid N_{\mathbf Q(\zeta_{\ell^r})/\mathbf Q}(\alpha),
 \qquad
 \left|N_{\mathbf Q(\zeta_{\ell^r})/\mathbf Q}(\alpha)\right|
                   \le\|A\|_1^{\varphi(\ell^r)}.
\tag{TN.4}
\]
The first assertion is a divisibility assertion between powers of
$p$ and a nonzero integer, obtained from the $p$-adic index; the
second uses the triangle inequality at every complex embedding.
Together they prove (TN.3). This argument includes every component
of the product algebra, not only a selected character packet.

There is a sharper exact statement when deck Frobenius is transitive
on all faithful characters, so $R=\mathbf Z_p$. If $b_r$ has exact
order $m=p^u$, (TN.2) gives
$\mathcal O b_r\simeq\mathcal O/m\mathcal O$.
Restrict now to integral operators that vanish on \emph{every proper
conductor} and annihilate $b_r$. Their faithful component is divisible
by $m$; all other components are zero. The cyclotomic product
decomposition of $\mathbf Z_p[C_{\ell^r}]$ consequently shows that
each integral coefficient is divisible by $m$. Such an operator
has augmentation zero. Every nonzero one has norm at least $2m$,
and equality is attained by
\[
                         m(1-\sigma^{\ell^{r-1}}).
\tag{TN.5}
\]
Thus $2m$ is the exact minimum in this conductor-isolating class.
It is not a lower bound for arbitrary operators annihilating only
$b_r$: the integral trace $\Phi_{\ell^r}(\sigma)$ already kills
that primitive class with norm $\ell$.

For the full point difference, this cheap trace has a different
meaning. Represent an integral operator by its unique polynomial of
degree less than $\ell^r$. If its faithful value is zero, monic
division in $\mathbf Z[X]$ gives
\[
 A=\Phi_{\ell^r}(\sigma)B,\qquad
 \deg B<\ell^{r-1},\qquad
 \Phi_{\ell^r}(\sigma)=\sum_{j=0}^{\ell-1}
                                      \sigma^{j\ell^{r-1}}.
\tag{TN.6}
\]
For the cover map $v:D_r\to D_{r-1}$ this is the literal divisor
and class identity
\[
 A(Q_r-P_r)=v^*B(Q_{r-1}-P_{r-1}),\qquad
 A\delta_r=v^*B\delta_{r-1}.
\tag{TN.7}
\]
If the latter class is zero, then $B\delta_{r-1}=0$: pullback is
injective on $p$-primary classes because the norm of a pulled-back class
is multiplication by $\ell$ on the original class. A principal function for (TN.7) is
therefore a pullback of the older principal function, up to a
nonzero constant, over $\bar k$. This does not claim arithmetic
descent of a function when $B$ is not Frobenius-invariant.

More generally, for $r>s$ put $d_0=\ell^{r-s}$ and
\[
                 T_{r/s}=\sum_{j=0}^{d_0-1}\sigma^{j\ell^s}.
\tag{TN.8}
\]
Then $T_{r/s}\delta_r=v^*\delta_s$. If an integral operator has
zero cyclotomic components at every conductor above $\ell^s$,
successive monic division gives exactly
\[
             A=T_{r/s}B,\qquad \deg B<\ell^s,
                         \qquad\|A\|_1=d_0\|B\|_1.
\tag{TN.9}
\]
The coefficient supports of the $d_0$ shifted copies of $B$ are
disjoint, which proves the norm equality. The same norm--pullback
argument makes the principal relation descend to $D_s$. Every
deck translate of its principal function still factors through
$D_s$; hence the kernel of the action on the generated field contains the
kernel of $C_{\ell^r}\to C_{\ell^s}$. One cannot use $\ell^r$ as the faithful
group order in the controlled-field argument.

For completeness, even the formally optimistic larger-level
inequality would be stronger. For an integer $d=\|B\|_1\ge1$,
\[
 \frac{(d_0d-1)^2}{d_0}-(d-1)^2
                    =(d_0-1)d^2-1+\frac1{d_0}>0.
\tag{TN.10}
\]
Thus $\ell^r>(d_0d-1)^2$ implies
$\ell^s>(d-1)^2$. In particular, if the exact order of $\delta_s$
is $m_s$, the relation $m_sT_{r/s}$ of norm $d_0m_s$ merely
pulls back the degree-$m_s$ function at level $s$.

These assertions separate two exhaustive cases for a proposed
annihilator at the top level: a nonzero faithful value obeys the
actual orbit-module norm bound, while a zero faithful value traces
down to an older relation. Neither a bounded rational generator
count nor inflation of a trace relation supplies the missing
uniform estimate on the annihilator indices in (IA.14).


\paragraph{Supplement to Finding A52: the differential of the actual point relation.}
Let $v:D_r\to D_s$ be a finite etale refinement of degree prime
to $p$ in the compatible pointed tower of A51. Let $P_r,Q_r$
be its compatible rational lifts of distinct points, and let
$m_r,m_s$ be the exact orders of the actual $p$-primary classes
$[Q_r-P_r]$ and $[Q_s-P_s]$. Suppose that $m_r>m_s$.
Choose the rational principal functions from A51, so that
\[
 \operatorname{div}(f_i)=m_i(Q_i-P_i),\qquad
 N_v(f_r)=c f_s^{m_r/m_s},\quad c\in k^\times.
\tag{PD.1}
\]
Then the actual point relation gives a nonzero regular differential
\[
 0\ne\omega_r=\frac{df_r}{f_r}
       \in H^0(D_r,\Omega^1_{D_r/k}),\qquad
                         \operatorname{Tr}_v(\omega_r)=0.
\tag{PD.2}
\]
This conclusion is about the given point divisor, not an abstract
torsion module. It detects a new trace-zero differential at a
strict order jump; it supplies no upper bound on $m_r$ or the
annihilator index in (IA.14).

\begin{proof}
The divisibility $m_s\mid m_r$ and the norm formula are (NC.4)--(NC.5),
using the divisor-of-norm identity
\cite[Lemma~42.18.1, Tag~\texttt{02RT}]{Stacks02RT}.
Since both orders are powers of $p$ and the inequality is strict,
$p\mid m_r$ and $p\mid m_r/m_s$.
At any geometric point write $f_r=t^a u$, with $t$ a uniformizer
and $u$ a unit. Every valuation $a$ is divisible by $p$:
it is $m_r$, $-m_r$, or zero. Thus
\[
                  df_r/f_r=a\,dt/t+du/u=du/u
\tag{PD.3}
\]
is regular. This proves regularity at every point.
The differential is nonzero. Over the perfect field, $df_r=0$
would make $f_r$ a $p$th power in the function field. Its root
would have divisor $(m_r/p)(Q_r-P_r)$, contradicting exact
order. The same argument proves separability of the function
$f_r:D_r\to\mathbf P^1$.

For a finite separable extension, differentiation of the product
of the conjugates proves the identity
\[
       \operatorname{Tr}_v(df_r/f_r)
                    =dN_v(f_r)/N_v(f_r).
\tag{PD.4}
\]
One may check it after passing to a Galois closure; both sides
descend. Here the trace on differentials is defined by the
function-field trace and the etale identification of relative
differentials, so it preserves regularity. Substitution of
(PD.1) into (PD.4) gives
$(m_r/m_s)\,df_s/f_s=0$ in characteristic $p$, proving (PD.2).

The direct ramification estimate for this particular function
also gives no upper bound. Its degree is $m_r$, and its zero
and pole fibres each consist of one point, of total ramification
index $m_r$. At each, the different exponent is at least $m_r$:
in a local parameter expression the derivative of the leading
term of order $m_r$ vanishes, while separability makes a later
derivative nonzero. Consequently Riemann--Hurwitz gives only
\[
 2g(D_r)-2=-2m_r+\deg\operatorname{Diff}(f_r),
              \qquad \deg\operatorname{Diff}(f_r)\ge2m_r
\tag{PD.5}
\]
from these two forced fibres
\cite[Lemma~53.12.4, Tag~\texttt{0C1F}]{StacksRH}.
Their lower contributions cancel the degree term; no upper
bound on $m_r$ follows from this estimate. The trace-zero
identity likewise retains the new differential rather than
forcing $f_r$ to be a $p$th power. These are limits of these
specific deductions, not a geometric realization of arbitrary
orders or a nonexistence theorem for a different point-based
argument.
\end{proof}


\paragraph{Supplement to Finding A52: preserving packets and taking common covers.}
The following observations concern actual pointed covers. They do not
supply a new small integral annihilator or a universal witness.
Fix a genus-two curve $C/k$ over a finite field of characteristic $p$
and $P\in C(k)$. Suppose finitely many constant, connected,
$P$-split cyclic covers $D_i\to C$, of odd degrees prime to $p$,
have good signed $p$-packets in the sense of A49: the corresponding
character values $L_P(\chi,1)$ are units at the selected prime above $p$.
These finitely many packets can be preserved while applying the
variable-quotient supply argument of A49 again.

\begin{proof}
Let $V_i=J_{D_i}[p](\bar k)$ be the finite geometric etale
$p$-torsion vector space. Arithmetic Frobenius $F_i$ acts invertibly
on it and hence has finite order $T_i$. Choose $T$ divisible by all
$T_i$ and all the old deck-Frobenius orders; the latter are one
under the present constant-deck hypothesis. Frobenius commutes
with the deck action. If $j\equiv1\pmod T$, its action over
$k_j=\mathbf F_{q^j}$ on every $V_i$ is exactly the old action:
$F_i^j=F_i$. Thus its invariant subspace in every old character
packet is unchanged. The character-unit criterion of A49 says
precisely that this invariant subspace is zero on a good packet.
A nonzero invariant $p$-primary element would yield an invariant
element of order $p$, so this also controls the corresponding
rational $p$-primary unit test. No rational abelian subvariety
for an individual $p$-adic packet is asserted.

Enlarge the finite excluded-prime set $S$ in A49's
variable-quotient note to contain all old conductor primes, and
let $M$ be its A42 primary-stabilization modulus. The progression
$j\equiv1\pmod{\operatorname{lcm}(M,T)}$ is infinite and has
$\gcd(j,M)=1$. The fixed-kernel counting and Weil estimates of
that note therefore still supply a new good signed packet over
arbitrarily large such extensions. Its conductor can be chosen
coprime to the old ones. Repeating this argument preserves any
finite collection of the actual packets already obtained.
It does not control every faithful packet of one conductor.
\end{proof}

Here is the exact common-cover calculation. Over a common field
where the preceding packets have been preserved, take the connected
component of the fibre product through the selected rational lifts
of $P$. It is an actual connected etale cover $D\to C$ with
constant abelian deck group $H$, a subgroup of the product of the
old cyclic groups. A hypothetical point-section collision gives a
compatible rational lift of $Q$ on this same cover. Put
$\delta=[Q'-P']\in J_D(k)[p^\infty]$. For a quotient
$h_i:D\to D_i$ with kernel $K_i$, actual divisor pushforward and
pullback give
\[
 (h_i)_*\delta=\delta_i,\qquad
 h_i^*(h_i)_*=N_{K_i}:=\sum_{g\in K_i}g.
\tag{CP.1}
\]
Over a finite unramified extension of $\mathbf Z_p$ containing
the roots of unity of order dividing $\exp H$, character evaluation
therefore gives
\[
 \eta(N_{K_i})=
 \begin{cases}
 |K_i|,&\eta|_{K_i}=1,\\
 0,&\eta|_{K_i}\ne1.
 \end{cases}
\tag{CP.2}
\]
Indeed the first sum has all terms one; in the second case
multiplying the sum by a nontrivial value of the character
permutes its terms and forces it to vanish. Since $p\nmid|H|$,
these evaluations describe the full split group algebra and
its action on the actual $p$-primary group.

If an integral operator $B_i$ annihilates $\delta_i$, then
$N_{K_i}\widetilde B_i$ annihilates $\delta$, where
$\widetilde B_i$ is any lift to $\mathbf Z[H]$.
Every integral combination of these lifted relations has zero
character value on a character factoring through none of the
selected quotients. For coprime cyclic conductors this is exactly
the faithful product-conductor limitation of A51.

The genuinely noncyclic test $H=(\mathbf Z/\ell)^2$, with
$\ell\ne p$ odd, gives a more precise boundary. Include all
$\ell+1$ cyclic line quotients. Every nontrivial character $\eta$
has a unique line kernel $K_\eta$, and (CP.2) becomes
\[
 \eta(N_K)=\ell\quad(K=K_\eta),\qquad
 \eta(N_K)=0\quad(K\ne K_\eta).
\tag{CP.3}
\]
All $\eta^u$, $u\in(\mathbf Z/\ell)^\times$, have this same
kernel and form the rational cyclotomic block of that one
quotient. Its signed $p$-orbits may be smaller. Relations lifted
from every other line quotient are zero on the whole block.
Thus accumulating a good packet from each of several different
quotients does not, through these norm identities, fill a missing
signed orbit in any one block. Relabelling a quotient merely
relabells its physical characters; deck translations act within
character summands and do not permute the blocks.

These are identities of actual geometric correspondences, not a
claim that an arbitrary undetected torsion element can be a
compatible point difference. Integral combinations can cancel
coefficients, and no lower bound against every such cancellation
is claimed here. A further geometric relation using the
point-supported nature of $\delta$ could add information, but
none has been obtained in this argument. Preserving finitely many
good packets therefore leaves the universal whole-packet or
small-annihilator supply condition unresolved.


\section{Effective Galois comparison: exact analytic inputs and constants}
\label{sec:effective-galois}

Throughout this section, $D_L=|\operatorname{disc}(L)|$, $n_L=[L:\mathbb Q]$, and all logarithms are natural.  Write $\mathrm{H}_M$ for the draft assertion that GRH holds for the Hecke $L$-functions of every number field.  Every conclusion using $\mathrm{H}_M$ remains conditional on that assertion.  The trivial Hecke character gives GRH for $\zeta_L$ for every number field $L$.

\paragraph{Source coverage.}
The six supplied \texttt{SORTED\_*.md} files are summaries, with frequently truncated statements.  The later recovery of the complete A1, successor, and LM8 papers permits the row-level scope to be checked against actual proofs.  In particular, LM8 A.32 uses Serre's qualitative open-image theorem for a non-CM elliptic curve only to obtain Zariski density of the commutator image; no effective index is needed there.  LM8 D.25 uses Faltings's finiteness theorem for rational points, not the Faltings--Serre trace-comparison method.  The bounds below are explicit available refinements of finite-quotient arguments, not a claim that these two qualitative inputs were conditional.  The Gaussian and Fermat-cubic CM elliptic curves in other LM8 rows do not satisfy the non-CM hypothesis of the open-image bound below.  The row ledger records the precise finite fields in the recovered proofs.

\subsection{The finite-cover Chebotarev bound}

\paragraph{Cited theorem G-CHEB (Bach--Sorenson).}
If $L/K$ is finite Galois, $L\ne\mathbb Q$, and $C$ is a conjugacy class in $\operatorname{Gal}(L/K)$, then GRH for $\zeta_L$ and $\zeta_{\mathbb Q}$ gives an unramified degree-one prime $\mathfrak p$ of $K$ with Frobenius class $C$ and
\[
 N\mathfrak p\le B(L):=(4\log D_L+\tfrac52 n_L+5)^2.
\]
The exact reference is Bach--Sorenson, \emph{Explicit bounds for primes in residue classes}, Math. Comp. 65 (1996), 1717--1735, Theorem 5.1, p.~1729; the analytic hypotheses are specified in the last paragraph of \S2, p.~1720.  Their Theorem 3.2 and Corollary 3.3, p.~1721, make the reduction to a cyclic extension.  No Artin holomorphy conjecture is assumed.  This is already a conditional theorem of the literature; it is not a new finding of this work.

\paragraph{Finding G01 (a bounded prime outside a prescribed set).}
Let $L/\mathbb Q$ be finite Galois and let $S$ be a finite set of rational primes. Put
\[
 m=\operatorname{rad}\!\left(2\prod_{p\in S}p\right),\qquad
 \widetilde L=L(\sqrt m).
\]
For any nonempty conjugation-stable subset $C\subseteq\operatorname{Gal}(L/\mathbb Q)$, GRH for $\zeta_{\widetilde L}$ and $\zeta_{\mathbb Q}$ gives a rational prime $p\notin S\cup\{2\}$, unramified in $L$, such that
\[
 \operatorname{Frob}_p(L/\mathbb Q)\subseteq C,
 \qquad p\le B(\widetilde L).
\]
Thus $\mathrm{H}_M$ supplies this statement for arbitrary $L$. The stipulated Dirichlet GRH supplies it when $\widetilde L/\mathbb Q$ is abelian, in particular when $L/\mathbb Q$ is abelian.

\begin{proof}
The integer $m$ is even and squarefree, and $m\ge2$, so every prime dividing $m$ ramifies in the quadratic field $\mathbb Q(\sqrt m)$.  Therefore each such prime ramifies in $\widetilde L$.  Restriction $\operatorname{Gal}(\widetilde L/\mathbb Q)\to\operatorname{Gal}(L/\mathbb Q)$ is surjective.  The inverse image of $C$ is nonempty and stable under conjugation; choose one conjugacy class in it and apply G-CHEB.  The resulting prime is unramified in $\widetilde L$, hence does not divide $m$, and restriction of Frobenius proves the claim.  This argument is the squarefree-prime-avoidance construction of Mayle--Wang, Corollary~6, with $2$ explicitly included to remove a possible exceptional prime.
\end{proof}

\paragraph{Finding G02 (finite quotient trace comparison).}
Let $\rho_i:G_{\mathbb Q,S}\to\operatorname{GL}_d(R)$, $i=1,2$, be continuous representations over a finite commutative ring $R$, and let $L$ be the fixed field of the joint kernel. If their trace functions differ, they differ at a prime $p\notin S\cup\{2\}$ satisfying $p\le B(L(\sqrt m))$, under the GRH assumption in G01. The same assertion holds for any conjugation-invariant function of the joint image, including coefficients of characteristic polynomials.

\begin{proof}
The joint image is $\operatorname{Gal}(L/\mathbb Q)$.  The set on which the two trace functions differ is nonempty and stable under conjugation.  G01 applies to it.  A conjugation-invariant function is handled by precisely the same argument.
\end{proof}

\paragraph{Arbitrary number-field bases.}
For a finite Galois extension $L/K$, let $M/\mathbb Q$ be its normal closure and let $S_0$ contain the rational primes below every prime to be excluded.  With $m=\operatorname{rad}(2\prod_{p\in S_0}p)$, any prescribed class in $\operatorname{Gal}(L/K)$ occurs at a degree-one prime $\mathfrak p$ of $K$, above a rational prime outside $S_0\cup\{2\}$, with
\[
 N\mathfrak p\le B(M(\sqrt m))
\]
under GRH for $\zeta_{M(\sqrt m)}$ and $\zeta_{\mathbb Q}$.  To see this, extend a representative $g\in\operatorname{Gal}(L/K)$ to $\gamma\in\operatorname{Gal}(M/K)$ and apply G01 to its conjugacy class in $\operatorname{Gal}(M/\mathbb Q)$.  Choose a prime of $M$ whose Frobenius is precisely $\gamma$.  Since $\gamma$ fixes $K$, the prime below it in $K$ has residue degree one over $\mathbb Q$; its Frobenius in $L/K$ is $g$.  This gives an explicit form for the arbitrary-base finite affine images in the ledger, after their finite extension is fixed.  It does not bound a preceding Hilbert specialization or a depth parameter not specified by the proof.

This test compares trace functions at the chosen finite level.  Equality modulo a fixed power of a coefficient prime does not imply equality of the full adic representations.  To obtain a finite test for their full trace functions, one needs the next construction.

\subsection{A complete finite test for adic semisimplifications}

\paragraph{Finding G03 (an explicit uniform trace bound).}
Let $\mathcal O$ be the ring of integers in a finite extension $F/\mathbb Q_\ell$, with uniformizer $\varpi$ and residue field of cardinality $q$.  Let
\[
 \rho_1,\rho_2:G_{\mathbb Q,S}\longrightarrow\operatorname{GL}_d(\mathcal O)
\]
be continuous.  Set
\[
 N=2(q^{2d^2}-1),\qquad
 m=\operatorname{rad}\!\left(2\prod_{p\in S}p\right),
\]
and
\[
 \mathcal B(d,q,S)=
 \left[4\bigl((N-1)\log m+N\log N\bigr)+\tfrac52 N+5\right]^2.
 \tag{G03.1}
\]
Conditional on $\mathrm{H}_M$, equality of the exact Frobenius traces for every prime $p\nmid m$ with $p\le\mathcal B(d,q,S)$ is equivalent to equality of their trace functions on $G_{\mathbb Q,S}$.  It is therefore equivalent to isomorphism of their semisimplifications over $F$.  In particular it is equivalent to isomorphism when both representations are semisimple.

\begin{proof}
Suppose the trace functions differ.  Let $M$ be the $\mathcal O$-linear span of
\[
 \{(\rho_1(g),\rho_2(g)):g\in G_{\mathbb Q,S}\}
 \subseteq M_d(\mathcal O)\oplus M_d(\mathcal O).
\]
Products of spanning elements are again spanning elements, so $M$ is an $\mathcal O$-algebra with identity. It is a finite free $\mathcal O$-module of rank $r\le2d^2$, because it is a submodule of the finite free module on the right and $\mathcal O$ is a discrete valuation ring.  Reduction in this algebra gives a continuous finite quotient
\[
 \delta:G_{\mathbb Q,S}\longrightarrow H\subseteq(M/\varpi M)^\times.
\]
Continuity follows because $M$ is a finite module with its subspace topology and $\varpi M$ is open in $M$.  The quotient ring has $q^r$ elements; its zero element is not a unit. Hence $|H|\le q^r-1\le q^{2d^2}-1$.

Define the $\mathcal O$-linear functional
\[
 t:M\longrightarrow\mathcal O,\qquad t(A,B)=\operatorname{tr}A-\operatorname{tr}B.
\]
Its image is a nonzero ideal $\varpi^a\mathcal O$. Therefore $\lambda=\varpi^{-a}t$ is integral on $M$ and has nonzero reduction $\overline\lambda$ on $M/\varpi M$. Since the pairs $(\rho_1(g),\rho_2(g))$ span $M$, at least one has nonzero $\overline\lambda$-value. Conjugation by every such pair preserves $t$, hence also $\lambda$. Consequently
\[
 C=\{h\in H:\overline\lambda(h)\ne0\}
\]
is well defined, nonempty, and stable under conjugation. Every $g$ mapping to $C$ has unequal traces.

Let $L/\mathbb Q$ be the finite Galois extension defined by $\ker\delta$. It is unramified outside $S$, and $n_L=|H|\le q^{2d^2}-1$. Thus $\widetilde L=L(\sqrt m)$ has degree $\widetilde n\le N$ and is ramified only at primes dividing $m$. By the discriminant estimate of Serre, \emph{Quelques applications du th\'eor\`eme de densit\'e de Chebotarev}, Proposition~6 (also Mayle--Wang, Lemma~7),
\[
 \log D_{\widetilde L}
 \le(\widetilde n-1)\log\operatorname{rad}(D_{\widetilde L})
       +\widetilde n\log\widetilde n
 \le(N-1)\log m+N\log N.
\]
G01 supplies $p\nmid m$ of unequal traces with $p\le B(\widetilde L)\le\mathcal B(d,q,S)$. This proves the contrapositive of the finite-test assertion. The converse is immediate.  In characteristic zero, the Brauer--Nesbitt character criterion identifies equal trace functions with isomorphic semisimplifications.

The algebra construction is the extension to a coefficient discrete valuation ring of Mayle--Wang, Proposition~12, rather than an assumption of a known congruence depth. Its use of the quotient $M/\varpi M$ is essential: it must not be replaced by entrywise reduction of the ambient pair of matrix algebras.
\end{proof}

The exact analytic functions used by this proof are $\zeta_{\widetilde L}$ and $\zeta_{\mathbb Q}$, where $L$ is the finite quotient just constructed.  Neither the $L$-function of a positive-weight motive underlying $\rho_i$ nor Artin holomorphy is required. The finiteness and characteristic-zero character theory are independent of GRH. An executable numerical test additionally requires a way to compute the indicated exact traces; the theorem by itself does not manufacture that input.

\paragraph{Finding G04 (the precise limit of trace separation).}
Frobenius traces cannot distinguish arbitrary extension classes, even for lisse sheaves on an open subscheme of $\operatorname{Spec}\mathbb Z$.

\begin{proof}
Fix an odd prime $\ell$ and let $c(g)=\log_\ell(\chi_\ell(g)^{\ell-1})$, where $\chi_\ell$ is the cyclotomic character. This is a nonzero continuous additive homomorphism $G_{\mathbb Q,\{\ell\}}\to\mathbb Q_\ell$: its argument lies in $1+\ell\mathbb Z_\ell$, where the $\ell$-adic logarithm is an injective homomorphism, and the cyclotomic image contains this group. The two representations
\[
 \rho(g)=\begin{pmatrix}1&c(g)\\0&1\end{pmatrix},
 \qquad \rho_0(g)=\begin{pmatrix}1&0\\0&1\end{pmatrix}
\]
have trace $2$ and determinant $1$ on every element, including every unramified Frobenius. They are not isomorphic, because one has nontrivial image and the other has trivial image. Their semisimplifications are both the sum of two trivial characters.
\end{proof}

\paragraph{Finding G05 (explicit prime bounds for two LM7 predicates).}
Let $S$ be a finite set of rational primes and put
$m=\operatorname{rad}(2\prod_{p\in S}p)$.  Under the stipulated
Dirichlet GRH there is a prime $q\nmid m$ with $q\equiv2$ or $5\pmod9$
and
\[
 q\le\bigl[72\log3+24\log(4m)+35\bigr]^2. \tag{G05.1}
\]
Under $\mathrm H_M$ there is a prime $q\nmid m$ with $q\equiv4$ or
$7\pmod9$ and $3^{(q-1)/3}\not\equiv1\pmod q$, and
\[
 q\le\bigl[456\log3+72\log(4m)+95\bigr]^2. \tag{G05.2}
\]
These are quantitative refinements of the arithmetic predicates in the
recovered LM7 A.25 and A.16 registers, respectively.  Their source file
\texttt{sections/A\_agent.tex} was not recovered, so this finding does
not verify those rows' separate pro-three group presentations.

\begin{proof}
Put $L_0=\mathbb Q(\zeta_9)$ and $\alpha=\sqrt[3]{3}$.  The extension
$L_0/\mathbb Q$ is abelian of degree six.  The order
$\mathbb Z[\zeta_9]$ has discriminant of absolute value $3^9$:
for $\Phi_9(X)=X^6+X^3+1$,
\[
 \Phi_9'(\zeta_9)=3\zeta_9^2(2\zeta_3+1),\qquad
 N_{\mathbb Q(\zeta_3)/\mathbb Q}(2\zeta_3+1)=3.
\]
Taking the norm from $L_0$, whose degree over $\mathbb Q(\zeta_3)$ is
three, gives $3^9$.  Thus $D_{L_0}\le3^9$.  The two automorphisms
whose residues are $2,5\pmod9$ form the nonempty Frobenius predicate
for (G05.1).

The polynomial $X^3-3$ is Eisenstein at three.  Its real cubic field
$\mathbb Q(\alpha)$ is not Galois, so it cannot be contained in the
abelian extension $L_0$.  No other root can lie in $L_0$, either,
because $\zeta_3\in L_0$.  A cubic polynomial without a root is
irreducible; hence $L=L_0(\alpha)$ has degree eighteen over
$\mathbb Q$.  It is Galois, since $\zeta_3\in L_0$ makes it a
splitting field.  The product basis of
$\mathbb Z[\zeta_9]\mathbb Z[\alpha]$ is a basis of an order in $L$.
The cubic order has discriminant $-3^5$, and the tensor trace matrix
therefore gives
\[
 D_L\le (3^9)^3(3^5)^6=3^{57}. \tag{G05.3}
\]
For each $a=4,7$ and $j=1,2$ there is an automorphism
\[
 \zeta_9\longmapsto\zeta_9^a,\qquad
 \alpha\longmapsto\zeta_3^j\alpha.
\]
These four automorphisms form a nonempty conjugation-stable subset
of $\operatorname{Gal}(L/\mathbb Q)$.  At an unramified prime of residue
$a\pmod9$, the field $\mathbb F_q$ contains $\mu_3$, and $X^3-3$
has a root if and only if Frobenius fixes all its roots.  Thus the
four-element subset describes precisely the noncube predicate in
(G05.2).  Ordinary Chebotarev already gives density $4/18=2/9$;
the first predicate has density $2/6=1/3$.

For either $T=L_0$ or $T=L$, set $n=[T:\mathbb Q]$ and
$\widetilde T=T(\sqrt m)$.  Its degree is at most $2n$, and
\[
 \log D_{\widetilde T}\le2\log D_T+n\log(4m). \tag{G05.4}
\]
Indeed $m$ is positive, even and squarefree, so the quadratic field
has discriminant $4m$.  If it is not contained in $T$, product integral
bases give an order of discriminant $D_T^2(4m)^n$ in the degree-$2n$
compositum, whose field discriminant divides that order discriminant.
If it is contained, $\widetilde T=T$ and the displayed inequality
holds directly.  This covers the possible failure of linear disjointness.

Apply G01 and substitute (G05.4) in G-CHEB.  For $T=L_0$ we have
$[\widetilde T:\mathbb Q]\le12$ and
$\log D_{\widetilde T}\le18\log3+6\log(4m)$, giving (G05.1).
This compositum is abelian over $\mathbb Q$, so the stipulated
Dirichlet GRH supplies its Dedekind-zeta GRH.  For $T=L$ the
corresponding bounds are $36$ and $114\log3+18\log(4m)$, giving
(G05.2) under $\mathrm H_M$.  The chosen primes are unramified in
the relevant fields and avoid every divisor of $m$.  Neither use
requires Artin holomorphy.
\end{proof}

\subsection{Known explicit elliptic-curve consequences}

\paragraph{Cited result G-LOMBARDO (an unconditional bound over number fields).}
Let $E/K$ be an elliptic curve over a number field with
$\operatorname{End}_{\overline K}(E)=\mathbb Z$, let $d_K=[K:\mathbb Q]$,
and let $h_F(E)$ be its stable Faltings height in the convention of the
cited paper.  Lombardo, \emph{Bounds for Serre's open image theorem for
elliptic curves over number fields}, Algebra Number Theory 9 (2015),
2347--2395, Corollary~9.3, p.~2393 \cite{Lombardo2015}, proves
\[
 [\operatorname{GL}_2(\widehat{\mathbb Z}):\rho_{E,\infty}(G_K)]
 < \exp(10^{21483})\,d_K^{\gamma}
       \max\{1,h_F(E),\log d_K\}^{2\gamma},
 \qquad \gamma=2.4\cdot10^{10}. \tag{G-LOMBARDO}
\]
This published result is unconditional. It bounds the adelic image index;
it does not assert surjectivity or identify the exact image. We use the
stated corollary, whose isogeny input is recorded in its Theorem~2.1,
and do not substitute the improvement in Remark~9.4 that invokes the
then-unpublished input of Remarks~2.3 and~2.7. For the actual number
field $L$ in LM8 A.32, it applies directly with $E=F_a$ and $K=L$.

\paragraph{Cited results G-MW1 and G-MW2.}
Mayle--Wang, \emph{On the effective version of Serre's open image theorem}, Bull. Lond. Math. Soc. 56 (2024), 1399--1416, Theorems~1 and~2, give the following conditional bounds. For a non-CM elliptic curve $E/\mathbb Q$, its mod-$\ell$ image equals $\operatorname{GL}_2(\mathbb F_\ell)$ whenever
\[
 \ell>964\log\operatorname{rad}(2N_E)+5760.
 \tag{G-MW1}
\]
For non-CM elliptic curves $E_1,E_2/\mathbb Q$ which are not $\mathbb Q$-isogenous, some prime of good reduction for both has unequal Frobenius traces and satisfies
\[
 p\le\left(482\log\operatorname{rad}(2N_{E_1}N_{E_2})+2880\right)^2.
 \tag{G-MW2}
\]
These become conditional consequences of $\mathrm{H}_M$: the analytic input in their \S3 and the proof of Theorem~2 is G-CHEB for finite auxiliary fields, together with the squarefree quadratic compositum of Corollary~6. The open-image theorem itself and Faltings's isogeny theorem were already unconditional. These numerical bounds are not bounds for arbitrary abelian varieties, arbitrary number fields, or CM elliptic curves. In particular they do not prove Serre's uniformity conjecture.

\paragraph{Finding G06 (the actual LM8 A.32--A.33 elliptic factors).}
The factors used in these rows have the models
\[
 F_a/\mathbb Q:\quad y^2=x^3-3x+a,\qquad a=6,-6.
\]
Under $\mathrm H_M$, their mod-$\ell$ Galois image over $\mathbb Q$ is
$\operatorname{GL}_2(\mathbb F_\ell)$ whenever
\[
 \ell>964\log6+5760. \tag{G06.1}
\]
For the number field $L$ chosen in the recovered proof, the restricted
image satisfies
\[
 [\operatorname{GL}_2(\mathbb F_\ell):
       \overline\rho_{F_a,\ell}(G_L)]\le[L:\mathbb Q]
 \quad\text{for every prime satisfying (G06.1).} \tag{G06.2}
\]
Independently of $\mathrm H_M$, G-LOMBARDO bounds their adelic image
over $L$. Since the models descend to $\mathbb Q$, restriction also gives
\[
 [\operatorname{GL}_2(\widehat{\mathbb Z}):\rho_{F_a,\infty}(G_L)]
 <[L:\mathbb Q]\,\exp(10^{21483})
          \max\{1,h_F(F_a)\}^{4.8\cdot10^{10}}. \tag{G06.3}
\]

\begin{proof}
The models and their role are explicit in recovered LM8 A.32.
Their non-CM hypothesis is supplied by the ordinary-reduction
argument of LM8 A.14: for $F_6$, the traces at $5,11$ are $-1,3$,
giving the distinct imaginary quadratic Frobenius fields with
discriminants $-19,-35$. A geometric characteristic-zero CM field
would inject into both, which is impossible. The other sign is a
quadratic twist. These two point counts were rechecked directly;
no classification of CM $j$-invariants is used.

The discriminant of each displayed integral model is
$-16(4(-3)^3+27\cdot6^2)=-2^9 3^3$.
The model has good reduction outside $2,3$. Its discriminant valuations
$9$ and $3$ are smaller than $12$, so no integral Weierstrass change
can lower them: discriminant valuations of isomorphic models differ
by a multiple of $12$. Both are positive in a minimal model, so both
primes are bad and $\operatorname{rad}(2N_{F_a})=6$.
Substitute this value in G-MW1.

For any homomorphism $\rho$ and subgroup $G_L\le G_{\mathbb Q}$,
the map on coset sets is surjective, giving
$[\rho(G_{\mathbb Q}):\rho(G_L)]\le[L:\mathbb Q]$.
With the mod-$\ell$ surjectivity just obtained this proves (G06.2).
With G-LOMBARDO over $\mathbb Q$, it proves (G06.3).
Full image after restriction to $L$ is not asserted. The original
rows use only qualitative commutator Zariski density, which was
already unconditional; these formulas supply numerical refinements
for their actual elliptic factors.
\end{proof}

\paragraph{Grothendieck framework and locators.}
For $X=\operatorname{Spec}\mathbb Z[1/\prod_{p\in S}p]$, finite quotients of $G_{\mathbb Q,S}$ are connected finite \'etale covers of $X$, and continuous integral adic representations are inverse systems of finite locally constant coefficient sheaves on $X$. The precise Galois-category references are SGA~1, Expos\'e~V, Th\'eor\`eme~4.1 and \S7, together with Propositions~8.1--8.2 (the generic-point description over a normal base). G01--G03 locate a closed point of this arithmetic scheme where a prescribed finite monodromy invariant is visible.  The finite quotient is a weight-zero Artin realization.  It is not a new purity statement for the original positive-weight sheaf.  This section supplies no theorem about crystalline realizations, divided powers, or general motivic $L$-functions; those structures are not used in these proofs.  Cote~149 scans~33--36 concern Kummer point separation and prime-to-characteristic geometric fundamental groups, not the numerical constants of these number-field prime bounds.

\paragraph{Remaining row-level gate.}
To apply a numerical bound to an individual programme row, its complete proof must identify the finite quotient or pair of lisse sheaves, their coefficient ring, rank, ramification support, and the nonempty conjugacy-invariant test set. For an open-image conclusion it must also establish the relevant geometric endomorphism and subgroup hypotheses. None of these inputs can be inferred from a truncated row title alone.


\section{The cyclotomic coordinate of the finite gentle levels}
\label{sec:gt-certificates}

Write $U_N=(\mathbb Z/N\mathbb Z)^\times$, and use natural logarithms.
The word ``level'' in the supplied material can mean a matrix modulus,
a power imposed on a relative kernel, the order of a finite group, or
the modulus of the multiplier itself. These are different data. We
retain the actual quotient maps when making the comparison.

\subsection{The precise Bach input}

\paragraph{Finding T01 (the corrected uniform certificate).}
Under Dirichlet GRH, and hence under Alp\"oge's theorem as stipulated
in this problem, for every integer $N\geq3$ the set
\[
 \mathcal P_3(N)=\{p:p\text{ prime},\ p\nmid N,
                              \ p<3(\log N)^2\}
 \tag{T01}
\]
generates $U_N$. The groups $U_1,U_2$ are trivial and need no generators.
The inequality is strict, and there are no omitted small moduli.
Neither a Hecke GRH over a nonabelian field nor Masullo's claim is needed.

The coefficient $2$ attributed to unit generation in the prompt and
Fable's suggestion~7 is not the statement of Bach's cited theorem.
Bach's Theorem~2, printed page~372, treats a proper subgroup containing
every positive integer below a threshold; integers not coprime to the
modulus are allowed to be outside the subgroup. Its bound is
$2(\log N)^2$. Theorem~3 on the same page imposes coprimality and has
bound $3(\log N)^2$. The consequence for generators is stated on
printed page~373 \cite[Theorems~2--3, pp.~372--373]{Bach1990}.
The distinction is also explicitly recorded in
\cite[\S1.1]{LLS2015}.

\begin{proof}
Only the passage from Bach's integer statement to prime generators
requires proof. Let $H$ be generated by $\mathcal P_3(N)$. Every positive
integer $a<3(\log N)^2$ coprime to $N$ factors into primes in
$\mathcal P_3(N)$, so its residue lies in $H$. If $H$ were proper,
Bach's Theorem~3, applied with $x=3(\log N)^2$, would give
$x<3(\log N)^2$, a contradiction. This proves the assertion. For
$N=1,2$ the assertion is the direct trivial-group calculation.

The distinction between the two theorems cannot be discarded: in
$U_{12}$ the subgroup $\{1,5\}$ omits the positive integer $2$, but
$2$ is not a unit; the least omitted positive unit is $7$. This example
refutes the inference that a witness supplied by Theorem~2 must be a
unit. It does not refute the stronger numerical coefficient-$2$
generation statement, which is not asserted false here.
\end{proof}

The cited theorem's analytic proof has not been independently reproved
in this run, and Bach's original computation through one million moduli
has not been rerun. Its exact hypotheses, constants, strict inequalities,
and small-modulus scope were checked in the original scanned article.
The direct computations below are independent of this analytic input.

\paragraph{Finding T07 (a sharper criterion and a large-modulus consequence).}
Theorem~1.1 of Lamzouri--Li--Soundararajan, in the authors' corrected
arXiv version~v3 of 8~December~2016, applies to every proper subgroup
of $U_N$, for $N\geq3000$, with the same Dirichlet GRH input.
Define
\[
 \begin{split}
 A(N)&=\max\left(0,2\log\log N-\frac85
                    -\sum_{p\mid N}\frac{\log p}{p-1}\right),\\
 B(N)&=\max\left(0,2\log\log N+3
          +\frac{2\omega(N)(\log\log N)^2}{\log N}-2A(N)\right),
 \end{split}
\]
where $\omega(N)$ counts distinct prime divisors. A prime $\ell\nmid N$
outside the subgroup exists with
$\ell\leq(\log N+B(N))^2$
\cite[Theorem~1.1(1), corrected arXiv version~v3]{LLS2015}.
Thus the desired strict coefficient-$2$ generator bound follows
whenever $B(N)<(\sqrt2-1)\log N$, by the same subgroup argument as
T01. Part~(2) gives $\ell\leq(\log N)^2$ under the additional
condition that $N$ has no prime divisor below $(\log N)^2$.
Theorem~1.4 instead concerns a specified coset and depends on its
subgroup index; Corollary~1.1 concerns quadratic nonresidues for prime
moduli. Neither states a uniform coefficient-$2$ generator bound for
arbitrary composite moduli.

The displayed bound in Theorem~1.1(1) is not uniformly below
$2(\log N)^2$ merely for $N\geq3000$. Put
$R(N)=(\log N+B(N))^2/(\log N)^2$. Exact rational interval calculations
give
\[
 \begin{split}
 4.98743&<R(30030)<4.98744,\\
 2.02303&<R(13082761331670030)<2.02304.
 \end{split}
\]
The moduli are the products of the first six and first fourteen primes,
respectively. The executable \texttt{check\_lls\_bound\_scope.py}
certifies these intervals using the positive logarithm series of
(T05.1), also applied to rational endpoints for $\log\log N$;
its two records are in \texttt{LLS\_BOUND\_EXPRESSION\_CHECKS.tsv}.
These are evaluations of the theorem's upper-bound expression,
\emph{not} counterexamples to actual coefficient-$2$ generation.
They show why the displayed LLS theorem plus checking only $N<3000$
does not complete such a uniform claim. The authors' v3 revision notice
identifies corrections to the asymptotic Theorems~1.2--1.3; no constants
from their superseded version are used here.

Under A, the strict coefficient-$2$ generators also suffice for every
integer $N\ge\exp(300)$. This deliberately coarse threshold is a new
deduction of the displayed LLS bound, not a threshold quoted from that
paper; the remaining interval is not asserted checked.

\begin{proof}
The sufficient $B(N)$ criterion follows by the proper-subgroup argument
in T01. The two expression enclosures follow from the exact rational
positive-series estimates and complete calculations just specified.
For the final assertion put $L=\log N\ge300$ and $l=\log L$.
Since $2^{\omega(N)}\le N$ and $\log2>2/3$, the defining formula gives
$B(N)<3l^2+2l+3$. Indeed $A(N)\ge0$ and
$2\omega(N)l^2/L<3l^2$. The function
$F(L)=(3(\log L)^2+2\log L+3)/L$ has derivative
$-(3\log L-1)(\log L-1)/L^2<0$ for $L\ge300$.
Also $\log300<6$, since $e>8/3$ and $(8/3)^6>300$.
Thus $F(L)\le F(300)<123/300=41/100$ and
\[
 (L+B(N))^2<(141/100)^2L^2
             =\frac{19881}{10000}L^2<2L^2.
\]
LLS Theorem~1.1 applies since $N\ge\exp(300)>3000$ and gives a
coprime prime outside any proper subgroup below this last bound.
The same subgroup argument proves generation. The elementary lower
bound $\log2>2/3$ is, for example, the first term plus the positive
remainder in $\log2=2\sum_{j\ge0}(1/3)^{2j+1}/(2j+1)$.
\end{proof}

\subsection{What the finite certificate certifies}

\paragraph{Finding T02 (finite cyclotomic tests and their limits).}
Let $\mathcal G=\widehat{\mathrm{GT}}_0$, with multiplier
$\lambda:\mathcal G\to\widehat{\mathbb Z}^{\times}$, and let
$q:\mathcal G\twoheadrightarrow Q$ be a finite quotient. To speak of
its full multiplier modulo $N$, require an actual homomorphism
$\lambda_N:Q\to U_N$ with
\[
                    \lambda\bmod N=\lambda_N\circ q.
 \tag{T02.1}
\]
Equivalently, $\ker q\subseteq\ker(\lambda\bmod N)$. If this is not
available, the actual finite image of $(q,\lambda\bmod N)$ is a
specified refinement that does retain the multiplier; this must not
silently be identified with $Q$.

Suppose $P$ is any proved generating set of $U_N$, for instance
$\mathcal P_3(N)$ under T01, or a list certified in T05. Then:
\begin{enumerate}
\item a subgroup $J\leq Q$ has $\lambda_N(J)=U_N$ if its multiplier
image contains every residue in $P$;
\item two homomorphisms from $U_N$ agree if they agree on $P$;
\item a nontrivial character of $U_N$ is nontrivial on some member of $P$.
\end{enumerate}
These statements concern subgroups and homomorphisms already shown to
have the indicated algebraic properties. A set of candidate shadows
has not thereby been shown to be a subgroup. The certificate does not
check the two cycle equations, the pentagon, the $f$-coordinate, or
compatible lifting through all finite quotients.

\begin{proof}
Each assertion follows by writing every unit as a word in $P$. For
the third, otherwise the kernel of the character would contain all of
$U_N$. For the first, it is essential that the image be a subgroup,
so it contains all words in its specified generators. The factorization
criterion in (T02.1) is the elementary universal property of a quotient.

For the limitation, retain the supplied LM7\_PUSH E.65 theorem:
there is a continuous character $\Delta:\mathcal G\to\mathbb F_2$
vanishing on the Galois subgroup and taking value $1$ on an actual
element of $\ker\lambda$ \cite[E.58, E.65]{LM7run}. It cannot factor
through $\lambda$ or through any $\lambda\bmod N$, since a homomorphism
of that form is zero on $\ker\lambda$. Moreover $\ker\Delta$ is proper
but has the same full cyclotomic image as $\mathcal G$: it contains
$G_{\mathbb Q}$, whose cyclotomic character is onto. Thus even all of
the cyclotomic finite certificates together cannot distinguish these
two subgroups. This is a concrete counterexample within the supplied
gentle-group framework to inferring equality of groups from equality
of their multiplier images. It uses the supplied E.65 theorem as
background, not a new proof of E.58's profinite construction.
\end{proof}

For clarity, the arithmetic inclusion already gives
$\lambda(\mathcal G)=\widehat{\mathbb Z}^{\times}$, independently of
GRH. When an actual finite quotient retains $\lambda\bmod N$, its
full multiplier image is therefore already $U_N$. Bach's input gives
a bounded list of witnesses or tests, not the original surjectivity
theorem. Also, an arithmetic Frobenius at $p\nmid N$ has multiplier
$p$ modulo $N$; no claim is made that its full profinite multiplier
is the nonunit integer $p$ in $\widehat{\mathbb Z}$.

\paragraph{The scheme and fibre-functor formulation.}
Put $S_N=\operatorname{Spec}\mathbb Z[1/N]$ and
$X_N=\operatorname{Spec}\mathbb Z[1/N,\zeta_N]$. This is the connected
finite \'etale Galois cover with group $U_N$. Indeed it is the primitive
root component of the \'etale scheme of roots of $T^N-1$ after $N$ is
inverted; the cyclotomic field identifies its generic Galois group.
Arithmetic Frobenius sends $\zeta_N$ to $\zeta_N^p$. T02 is consequently
a certificate on a finite quotient of the fibre-functor automorphism
group of $S_N$, in SGA~1, Expos\'e~V, Th\'eor\`eme~4.1 and \S7
\cite{SGA1}. This is the applicable Grothendieck framework. It supplies
no additional purity, crystalline, divided-power, or motivic assertion.
For each prime $\ell$ the compatible roots of unity are the Tate module
$\mathbb Z_\ell(1)$; finite level records its cyclotomic scalar action.

\subsection{Elementary certificates for the two- and three-primary towers}

\paragraph{Finding T03 (strict coefficient-$2$ certificates without GRH).}
For integers $a,b\geq0$ and $N=2^a3^b>1$, the following prime lists generate $U_N$, and every
listed prime is coprime to $N$ and strictly less than $2(\log N)^2$:
\[
\begin{array}{c|c}
 (a,b)&\text{prime generators}\\\hline
 (1,0)&\varnothing\\
 (2,0)&3\\
 a\geq3,\ b=0&3,5\\
 a=0,\ b\geq1&2\\
 a=1,\ b\geq1&5\\
 a=2,\ b\geq1&5,7\\
 a\geq3,\ b\geq1&5,7,13
\end{array}
 \tag{T03}
\]
Also $2$ generates $U_{5^b}$ for every $b\geq1$ and satisfies the same
strict bound. These are elementary unconditional statements, requiring
neither of the analytic inputs in the prompt.

\begin{proof}
For $a\geq3$, the subgroup of units congruent to $1$ modulo $4$ has
order $2^{a-2}$ and is generated by $5$. The elementary valuation
identity
\[
                  v_2(5^{2^j}-1)=j+2
 \tag{T03.1}
\]
proves the order: the base case is $5-1=4$, and on each squaring the
factor $5^{2^j}+1$ has valuation one. Since $3\equiv-1\pmod4$, the
two primes $3,5$ generate all of $U_{2^a}$. The case $a=2$ is direct.

Both $2$ and $5$ have order $2\cdot3^{b-1}$ modulo $3^b$. To see this,
they are $-1$ modulo three, and
\[
 v_3(2^2-1)=v_3(5^2-1)=1,
 \qquad v_3(z^3-1)=v_3(z-1)+1\quad(z\equiv1\pmod3).
 \tag{T03.2}
\]
The latter identity follows by expanding $(1+3^s u)^3-1$ for
$s\geq1$ and $3\nmid u$: its first term has uniquely least
valuation $s+1$. Repeated use gives the asserted orders, equal to
$|U_{3^b}|$. This proves the first two ternary cases.

For $a=2,b\geq1$, the quotient of $U_N$ by its squares is canonically
$U_{12}$, generated by $5,7$: their products give $1,5,7,11$.
The projection of $5$ generates the ternary primary component, by (T03.2).
Thus the subgroup they generate is all of $U_N$. Here is a proof
that these two checks suffice: if $H$ is proper in a finite abelian
group $G$, then $G/H$ has a quotient of prime order; the only
possible primes dividing $|U_N|$ are two and three. Such a quotient
would contradict one of the two checks.

For $a\geq3,b\geq1$, the quotient by squares is instead $U_{24}$.
Its eight distinct elements are the products of $5,7,13$:
\[
                    1,5,7,13,11,17,19,23.
 \tag{T03.3}
\]
Reduction from $U_{2^a}$ to $U_8$ kills exactly its squares, using
(T03.1); reduction from $U_{3^b}$ to $U_3$ does the same for its
two-primary quotient, using (T03.2). This proves the assertion about
the quotient by squares. The projection of $5$ again generates the
three-primary component. The preceding prime-order-quotient argument
finishes the mixed case. Empty or absent primary components cause
no exception.

For $5^b$, $2$ has order four modulo five, and $v_5(2^4-1)=1$.
The same binomial argument, now with fifth powers, shows
$v_5(2^{4\cdot5^j}-1)=j+1$, so its order is $4\cdot5^{b-1}$.

It remains to verify the strict bounds for these entire families,
rather than just for a finite sample. The elementary series
$\log2=2(1/3+(1/3)^3/3+\cdots)$ gives $\log2>2/3$; also
$e<3$, by its factorial series, gives $\log3>1$. For the minimal
moduli $3,4,8,6,12,24$ in the nonempty rows, respectively,
these inequalities give
\[
 2(\log N)^2>2,\quad 32/9,\quad8,\quad50/9,\quad98/9,\quad18.
\]
These exceed respectively the largest proposed primes $2,3,5,5,7,13$.
The bound increases with $N$. For $5^b$ it suffices that $5>3$.
This proves strictness in every case, including all boundary cases.
\end{proof}

\subsection{Identifying the actual scalar moduli}

For the LM7 homology tower, retain the actual maps
\[
 B=\langle s,r\mid s^2=r^3=1\rangle,
 \quad t=sr,\quad
 q:B\longrightarrow\operatorname{SL}_2(\mathbb Z/12)/\{\pm I\},
\]
where $q(s)$ and $q(r)$ are the classes of
$S=\left(\begin{smallmatrix}0&1\\-1&0\end{smallmatrix}\right)$ and
$R=\left(\begin{smallmatrix}0&-1\\1&-1\end{smallmatrix}\right)$.
Set $N_B=\ker q$, $C=N_B^{\rm ab}$, and
$G_j=B/(N_B'N_B^j)$. These are the discrete source definitions of
LM7 E.1/E.4; the maps to their finite targets extend to completions.

\paragraph{Finding T04 (source-to-modulus comparison).}
In the recovered LM7\_PUSH E.1/E.4 homology tower, the cusp has order
$12j$ in $G_j=B/(N_B'N_B^j)$. The specifically reported coefficient
checks $j=2,4,8,16,3,9,27,81$, together with the retained $j=6$, therefore
have cusp scalar moduli
\[
        24,48,96,192,36,108,324,972,72.
 \tag{T04.1}
\]
For the correcting characters of LM8\_PUSH C.33--C.38,
\[
 c_{p,k}(u)=\frac{\log(u^{p-1})}{p}\pmod{p^k},
 \qquad p\geq5,\quad k\geq1,
 \tag{T04.2}
\]
the exact scalar conductor is $p^{k+1}$. At $k=1$ this is
$(u^{p-1}-1)/p\bmod p$. The T01 cutoff at that modulus is
$3(k+1)^2(\log p)^2$.

\begin{proof}
E.1 identifies $C=N_B^{\rm ab}\simeq\mathbb Z^{97}$ and records that
$c=[t^{12}]$ is primitive \cite[E.1, proof; E.4]{LM7run}.
The matrix $SR=\left(\begin{smallmatrix}1&-1\\0&1\end{smallmatrix}\right)$
has order twelve even modulo the single sign, since its diagonal
entries cannot become $-1$ modulo twelve. Thus a power $t^r$ can be trivial
in $G_j$ only if $r=12s$. Its value in the kernel $C/jC$ is then
$sc$, zero exactly when $j\mid s$. Hence $\operatorname{ord}_{G_j}(t)=12j$.
This identifies the cusp scalar, not the whole order of $G_j$.
The pure peripheral element $t^2$ has order $6j$, so the half-multiplier
$(\lambda-1)/2$ in the cycle equations is also determined by
$\lambda\bmod12j$.

For (T04.2), if $u/v\in1+p^{k+1}\mathbb Z_p$, then
$\log((u/v)^{p-1})\in p^{k+1}\mathbb Z_p$; division by $p$ makes
the difference zero modulo $p^k$. Thus the character factors through
$U_{p^{k+1}}$. On $u=1+p^k$, expansion of the logarithm gives
\[
 c_{p,k}(1+p^k)\equiv(p-1)p^{k-1}\not\equiv0\pmod{p^k}.
\]
The omitted terms have valuation at least $k$ after division by $p$,
including $k=1$ since $p\geq5$. Therefore it does not factor through
$U_{p^k}$, proving exactness. Reduction of
$\log(1+pa)/p$ modulo $p$ is $a$, which proves the $k=1$ formula.
\end{proof}

More generally, in a finite characteristic free-group target $H$ of
exponent $e$, reduction of $\lambda$ modulo $2e$ determines every
exponentiation by $\lambda$ and by $(\lambda-1)/2$ in the two cycle equations.
When $e$ is odd, reduction modulo $e$ suffices, since two is invertible.
This gives a safe finite scalar parameter once the actual target is
specified; it is not a claim that $2e$ is minimal. For example LM8
C.14 has exponent twelve and explicitly uses multiplier modulus
twenty-four, whereas LM7 D.7 has exponent $101$ in its pure target
and uses $m=0\bmod101$, equivalently $\lambda=1\bmod101$.
The redundant odd lift modulo $202$ gives the same unit group.

The complete recovered-source inventory is
\texttt{GT\_LEVEL\_SCOPE.tsv}. The following are its principal entries;
``frame'' and ``coefficient'' entries are deliberately not promoted
to unspecified scalar quotients.
\begin{center}
\begin{tabular}{p{0.24\linewidth}p{0.23\linewidth}p{0.45\linewidth}}
\hline
Source & Integer or family & Actual meaning and scope\\\hline
LM7 E.1,E.4 & $12j$ & Cusp order in the specified homology quotient.\\
LM7 E.5 & $360$ & Cusp order in $G_6\times A_5$; constructed multiplier one.\\
LM7 E.13,E.14 & $288,576$ & Explicit multiplier-one congruences in the full register.\\
LM7 E.22,E.24 & $72,144$ & Cusp order and the branches $1,73\bmod144$ over retained points.\\
LM7 D.7 & $101$ & Odd-prime multiplier parameter; the tested stratum has multiplier one.\\
LM7 E.50,E.51 & $20,28$ & Actual projective congruence quotients; symplectic tests have multiplier one.\\
LM7 E.52--E.56 & $m,2m,4m$, $(m,6)=1$ & Infinite congruence and exponent families, with specified maps.\\
LM7 E.65 & $12$ & Signed CM frame modulus on $\ker\lambda$, not a cyclotomic character.\\
LM8 C.12 & Every $N$ & Normalized implementer $\left(\begin{smallmatrix}\lambda&t\\0&1\end{smallmatrix}\right)$; determinant is $\lambda$.\\
LM8 C.14 & $24$ & Explicit multiplier modulus for the exponent-twelve $S_4$ envelope.\\
LM8 C.18 & $15$ & Exact domain character $(\lambda/5)(\lambda\bmod3)$.\\
LM8 C.33--C.38 & $p^{k+1}$ & Scalar correcting character (T04.2), not the whole corrected character.\\
LM8 C.49 & $24,60$ & Framed images restricted to multiplier one; no full arbitrary-multiplier image claimed.\\
GT\_TOWER summary & $4,12;\ m=2,3,6;\ 6m=12,18,36$ & Named matrix and power levels only; the full register is unavailable.\\\hline
\end{tabular}
\end{center}

For C.12 the determinant of an arbitrary projective matrix is defined
only modulo squares. It is the displayed unique normalized
representative that retains the full multiplier. For C.49, the orders
$96,48,2880,1440$ of the two framed/arithmetic images are not
cyclotomic moduli. Likewise $T_{24}$ and $T_{48}$ in LM7 E.59/E.62
name targets built from matrix groups of those orders; their subscripts
are not established multiplier moduli. The finite coefficient depths
in the relative nilpotent towers have the same distinction.

\subsection{Direct finite certificates and verification}

\paragraph{Finding T05 (53 exact finite checks).}
For all 53 integers in \texttt{GT\_UNIT\_CERTIFICATES.tsv}, the primes
strictly below $2(\log N)^2$ and coprime to $N$ generate $U_N$.
This assertion is unconditional and finite. A literal word in the
listed prime generators is supplied for every unit in
\texttt{GT\_UNIT\_WORDS.tsv}; there are 4,825 such words in total.
The table also verifies the coefficient-$3$ bound, though no analytic
theorem is needed for these individual computations.

\begin{proof}
The executable \texttt{verify\_gt\_units.py} constructs a subgroup by
successive multiplication and records a nonnegative prime exponent
vector for every reached residue. All arithmetic of group elements is
integer modular arithmetic. The separate read-only
\texttt{check\_unit\_certificates.py} reads the saved tables, without
performing any subgroup search, and re-evaluates each word. It verifies
that each result is a unit, no residue is repeated, and the number
of distinct residues is
$\varphi(N)=N\prod_{p\mid N}(1-1/p)$. It also compares them with
every integer residue coprime to $N$. Thus the displayed products
constitute a certificate of generation, independently of the method
used to discover them. Both executables were run successfully.

Strict bounds are certified with rational intervals, rather than
rounded logarithms. Write $N=2^h r$ with $1\leq r<2$. For rational
$a/b\in[1,2]$, set $t=(a-b)/(a+b)$ and use
\[
 2\sum_{j=0}^{K-1}\frac{t^{2j+1}}{2j+1}
 \ \leq\ \log(a/b)\ \leq\
 2\sum_{j=0}^{K-1}\frac{t^{2j+1}}{2j+1}
       +\frac{2t^{2K+1}}{(2K+1)(1-t^2)}.
 \tag{T05.1}
\]
The error bound follows by replacing every denominator of the
remaining positive terms by $2K+1$, which can only enlarge the sum,
then summing the resulting geometric series. Combining the
intervals for $h\log2$ and $\log r$, and squaring their nonnegative
endpoints, gives rational bounds for $c(\log N)^2$. The program
increases $K$ until both endpoints have the same integer part and
the lower endpoint exceeds that integer. Every proposed prime below
or equal to that integer is therefore strictly below the true bound.
This test terminated for all reported moduli at both $c=2$ and $c=3$.
The printed decimals are only a convenient display of these checks.
\end{proof}

\begin{longtable}{rrrrl}
\caption{Exact finite unit-generation certificates. Decimal bounds are informational; the cutoff column is certified with rational intervals. The lists are generating lists, not claims of minimal cardinality.}\\
\hline
$N$&$\varphi(N)$&$2(\log N)^2$ (display)&Exact cutoff&Prime generators\\\hline
\endfirsthead
\hline
$N$&$\varphi(N)$&$2(\log N)^2$ (display)&Exact cutoff&Prime generators\\\hline
\endhead
2 & 1 & 0.960906 & 0 & $\varnothing$ \\
3 & 2 & 2.413898 & 2 & $2$ \\
4 & 2 & 3.843624 & 3 & $3$ \\
5 & 4 & 5.180581 & 5 & $2$ \\
6 & 2 & 6.420804 & 6 & $5$ \\
7 & 6 & 7.573133 & 7 & $2,3$ \\
8 & 4 & 8.648154 & 8 & $3,5$ \\
9 & 6 & 9.655592 & 9 & $2$ \\
10 & 4 & 10.603796 & 10 & $3$ \\
11 & 10 & 11.499803 & 11 & $2$ \\
12 & 4 & 12.349522 & 12 & $5,7$ \\
13 & 12 & 13.157930 & 13 & $2$ \\
14 & 6 & 13.929247 & 13 & $3$ \\
15 & 8 & 14.667072 & 14 & $2,7$ \\
16 & 8 & 15.374496 & 15 & $3,5$ \\
17 & 16 & 16.054196 & 16 & $2,3$ \\
18 & 6 & 16.708498 & 16 & $5$ \\
19 & 18 & 17.339442 & 17 & $2$ \\
20 & 8 & 17.948824 & 17 & $3,11$ \\
23 & 22 & 19.662648 & 19 & $2,5$ \\
24 & 8 & 20.200052 & 20 & $5,7,13$ \\
25 & 20 & 20.722323 & 20 & $2$ \\
27 & 18 & 21.725081 & 21 & $2$ \\
28 & 12 & 22.207174 & 22 & $3,5$ \\
29 & 28 & 22.677362 & 22 & $2$ \\
30 & 8 & 23.136287 & 23 & $7,11$ \\
31 & 30 & 23.584536 & 23 & $2,3$ \\
36 & 12 & 25.683216 & 25 & $5,7$ \\
48 & 16 & 29.972395 & 29 & $5,7,13$ \\
49 & 42 & 30.292530 & 30 & $2,3$ \\
60 & 16 & 33.527315 & 33 & $7,11,13$ \\
64 & 32 & 34.592617 & 34 & $3,5$ \\
72 & 24 & 36.579746 & 36 & $5,7,13$ \\
81 & 54 & 38.622367 & 38 & $2$ \\
96 & 32 & 41.666549 & 41 & $5,7,13$ \\
101 & 100 & 42.598675 & 42 & $2$ \\
108 & 36 & 43.844706 & 43 & $5,7$ \\
120 & 32 & 45.840154 & 45 & $7,11,13,17$ \\
121 & 110 & 45.999214 & 45 & $2$ \\
144 & 48 & 49.398088 & 49 & $5,7,13$ \\
169 & 156 & 52.631722 & 52 & $2$ \\
192 & 64 & 55.282515 & 55 & $5,7,13$ \\
202 & 100 & 56.355412 & 56 & $3$ \\
288 & 96 & 64.138243 & 64 & $5,7,13$ \\
289 & 272 & 64.216783 & 64 & $2,3$ \\
324 & 108 & 66.833991 & 66 & $5,7$ \\
360 & 96 & 69.292441 & 69 & $7,11,13,17$ \\
361 & 342 & 69.357767 & 69 & $2$ \\
529 & 506 & 78.650592 & 78 & $2,5$ \\
576 & 192 & 80.800209 & 80 & $5,7,13$ \\
841 & 812 & 90.709450 & 90 & $2$ \\
961 & 930 & 94.338145 & 94 & $2,3$ \\
972 & 324 & 94.651073 & 94 & $5,7$ \\
\hline
\end{longtable}


\paragraph{Finding T06 (the remaining evidence boundary).}
The six attached \texttt{SORTED\_*.md} files are summaries, not the
complete row registers their filenames are said to accompany. The
GT\_TOWER\_CHECK summary names a register with 180 rows, but does not
contain that register or the exact definitions of every finite target.
Search and recovery in this run obtained LM8's full 304-row register
and source, and LM7's full 218-row register with only part of its proof
sources. They did not obtain GT\_TOWER\_CHECK's full register or paper.
Consequently T05 proves every listed arithmetic certificate, but is
not described as an exhaustive computation of all finite levels in
the unrecovered 180-row run. The precise remaining gate is to obtain
its finite target definitions and identify their multiplier factors;
the reusable verifier then checks any additional explicit moduli.
For the genuinely infinite families already present in the recovered
sources, T01 applies at every identified modulus, and T03 gives the
stronger unconditional coefficient-$2$ result for the indicated towers.

\begin{proof}
This is a source-coverage finding. The attached summary's own register
count is 180; its displayed statements are truncated and list only
selected rows. The recovered LM7 and LM8 CSV files have 218 and 304
rows respectively. The source inventory records, for each use, whether
the underlying proof, full statement, or only summary was available.
Thus neither a missing row nor an unknown quotient map has been
silently assigned an integer modulus. No analytic assumption can
supply those omitted definitions.
\end{proof}


\section{Audit of the available rows and their analytic hypotheses}
\label{sec:row-audit}

\paragraph{Finding R01 (exact source coverage).}
The six attached \texttt{SORTED\_*.md} files are summaries, not the
complete row lists named in them. Their statements are frequently
truncated. They report $1190$ register rows. Subsequent recovery gives
the following exact inventory; ``proof blocks'' means available numbered
source blocks, not independently verified proofs.
\begin{center}
\begin{tabular}{lrrr}
\hline
Run & Reported rows & Recovered rows & Available proof blocks\\
\hline
\texttt{ANABELIAN\_A1} &121&121&121\\
\texttt{ANABELIAN\_A1\_SUCCESSOR}&210&0&0\\
\texttt{ANABELIAN\_SUCCESSOR}&157&157&157\\
\texttt{GT\_TOWER\_CHECK}&180&0&0\\
\texttt{LM7\_PUSH}&218&218&9\\
\texttt{LM8\_PUSH}&304&304&304\\
\hline
Total&1190&800&591\\
\hline
\end{tabular}
\end{center}
The nine LM7 blocks are E.1--E.7, E.58 and E.65. The $582$ blocks in
the complete A1, successor and LM8 main sources were available for
text screening. The targeted analytic passages were read; availability
does not mean that every mathematical proof was re-proved.

The flag file has $575$ records, representing $525$ distinct run/statement
pairs in $23$ runs. Its six-run subset has $166$ records and $152$
distinct pairs. The full audit table preserves all $800$ recovered rows,
and the curated table has $108$ entries, including the two other flagged
run examples, three intake items, the manuscript Artin passage and
Finding A06. Declared register dependencies give $201$ rows with a path
to one of $38$ curated qualitative prime-selection inputs. These paths
are recorded dependencies, not a newly certified proof graph.

For the two missing full registers, exactly $47$ A1-successor and $50$
GT identifiers are visible in the supplied summaries or flags. They
are listed separately; no unseen rows are invented. The remaining
coverage gap is therefore $390$ full register rows, together with $209$
LM7 proof blocks whose statements were recovered.

\paragraph{Finding R02 (the RH/GRH flags are false positives for conditional dependence).}
There are just two flags labelled RH/GRH. In LM8 A.69, the abbreviation
``DS/RH'' is resolved by the full proof as Deuring--Shafarevich and
Riemann--Hurwitz, not the Riemann Hypothesis. In
\texttt{S6\_SERIES\_GRIND}, row \texttt{maps:src:F.2.36}, the attached
context explicitly says that the class and unit groups were certified
with \texttt{GRH=false}. Its underlying computation was not recovered,
so this is an audit of the reported scope, not a new certification.
Neither flag establishes a formerly conditional theorem that is now
unconditional. Conversely, their absence from other rows does not show
analytic independence: the recovered LM8 proofs contain many
qualitative Chebotarev arguments.

\subsection{Qualitative prime choice and its quantitative refinements}

The finite-Galois qualitative Chebotarev theorem is unconditional:
a conjugacy class $C$ has density $|C|/|G|$. The exact source statement
checked here is Milne, \emph{Class Field Theory}, v4.03,
Chapter V, Theorem 3.23, printed p.~164 \cite{MilneCFT}.
The existing existence and density uses below therefore require neither
stipulated GRH hypothesis. A new least-prime version follows from
Finding G01 after the actual finite field and excluded primes have
been specified. Its exact analytic hypotheses are GRH for
\[
 \zeta_{\widetilde L}(s),\quad\zeta_{\mathbb Q}(s),\qquad
 \widetilde L=L\!\left(\sqrt{\operatorname{rad}
                     (2\prod_{p\in S}p)}\right).
\]
For a number-field base one uses the normal closure over $\mathbb Q$
and the rational primes below the excluded places. The stipulated
all-number-field Hecke hypothesis supplies these zeta functions as
trivial Hecke characters. Dirichlet GRH supplies this particular
argument when $L/\mathbb Q$ is abelian. No general Artin holomorphy
assumption is introduced.

\begin{longtable}{p{.22\linewidth}p{.43\linewidth}p{.26\linewidth}}
\hline
Rows & Actual finite field or predicate & Quantitative family\\
\hline
\endhead
LM8 A.16 & $M/\mathbb Q$ finite Galois containing the normalizer field;
complex-conjugation class. & $\zeta_M$, with G01 exclusions.\\[3pt]
LM8 B.20, B.22 & Normal closure of the specialized finite torsor fibre,
after the quotient, deck-constant field $K$, and Hilbert point $a$
are chosen. & Zeta of this normal closure.\\[3pt]
LM8 C.30; D.1; D.58 & Respectively
$\mathbb Q(\zeta_{12})$ with class $5$,
$\mathbb Q(\zeta_{5^r})$ with class $1$,
$\mathbb Q(\zeta_{30})$ with class $17$. & Dirichlet $L$-functions
modulo $12$, $5^r$, $30$, with finite exclusions.\\[3pt]
LM8 D.8 & $\mathbb Q(E_c[5],\mu_5)$; completely split primes.
& Zeta of this division field.\\[3pt]
LM8 D.11 & Kernel field of
$\rho_{M(-1)}\bmod\ell^{k+1}$ over the stated number field.
& Zeta of its $\mathbb Q$-normal closure.\\[3pt]
LM8 D.3, D.28 & Kernel field of
$g\mapsto(s_0(g)\bmod5,\chi_5(g)\bmod5^{r+1})$;
target $(\exp C,-1+5^r)$ in the actual Heisenberg image.
& Zeta of this actual image field.\\[3pt]
LM8 D.44, D.48, D.62, D.63, D.65 & Joint kernel field of the actual
finite Alexander action on $W_r$ and $\chi_5\bmod5^r$;
D.63 uses $r+1$. & Zeta of the joint image field.\\[3pt]
LM8 D.60, D.61, D.66 & Kernel field of the finite affine homomorphism
$g\mapsto(z(g),\rho_r(g))$ with additive coefficient module $W_r$.
& Zeta of the actual affine image field.\\[3pt]
LM8 D.67, D.69--D.71, D.73 & The same construction with finite
nonabelian coefficient group $H_r$ and its actual action;
arbitrary number field in D.70. & Zeta of the $\mathbb Q$-normal closure.\\[3pt]
LM8 D.78--D.87 & At each original finite Cartier-dual coefficient
module, the actual finite affine cocycle/action field.
& Zeta of that finite field.\\[3pt]
LM8 E.1, E.2 & $K=\mathbb Q(i)$ and $K(P)/K$ for a generator
$P$ of $E_6[12]$; degree $16$ in the source calculation.
& Finite-order Hecke $L_K(s,\chi)$, or $\zeta_{K(P)}$.\\[3pt]
LM7 A.16 & $\mathbb Q(\zeta_9,\sqrt[3]3)$;
$q\equiv4,7\pmod9$ and $3$ noncubic modulo $q$.
& Zeta of this degree-$18$ field.\\[3pt]
LM7 A.25 & $\mathbb Q(\zeta_9)$; classes $2,5$.
& Dirichlet $L(s,\chi)$ modulo $9$.\\
\bottomrule[0.4pt]
\end{longtable}

The LM7 proof sources for A.16 and A.25 are missing. Their fields are
reconstructed from the full register predicates, not from an asserted
recheck of their group presentations. Over $\mathbb Q(\zeta_3)$, the
extensions adjoining $\zeta_9$ and $\sqrt[3]3$ are independent cubic
extensions; four of the $18$ elements of the full Galois group satisfy
the conjugacy-stable A.16 predicate. This gives density $2/9$.
A.25's two classes among the six units modulo $9$ give density $1/3$.

For an additive or nonabelian cocycle, the finite homomorphism is the
actual graph $g\mapsto(z(g),\rho(g))$. A witness $\delta$ in the
action kernel with $z(\delta)\ne0$, respectively $\ne1$, gives a
nonempty conjugacy class in that actual image. One cannot replace it
by an arbitrarily chosen element of a larger formal affine group.
The recovered LM8 proofs make this distinction explicitly.

LM8 A.20 provides a separate logical warning: for a nontrivial quadratic
extension $L/K$, the inert primes have positive density and are
disjoint from the split primes, although the split set becomes
density one after pullback to $L$. This unconditional counterexample
does not refute the stronger arithmetic theorem for which a stronger
density hypothesis might suffice.

Stagewise effective prime choice does not bound a preceding Hilbert
specialization, a separating finite quotient, or the first obstruction
depth of an inverse system. LM7 A.43 and B.6 already give finite
field enumerations at each fixed depth from degree and discriminant
bounds; they expressly assert no uniform stopping rule for every
positive infinite-tower instance.

\subsection{Primitive roots, geometric existence and finite-field weights}

\paragraph{Finding R03 (fixed-prime criteria are not Artin-existence theorems).}
A1 B.3 and successor A.8, A.10, A.17 and A.35 use a hypothesis such as
$\operatorname{ord}_\ell(p)=\ell-1$ at a specified prime. This finite
condition makes $\Phi_\ell$ irreducible over $\mathbb F_p$ and its
nontrivial character values one Frobenius orbit. It does not produce
a prime satisfying the additional geometric or numerical conditions.
Successor A.12 supplies a genus-three counterexample to extending a
genus-two point-only test, even with $2$ primitive modulo $3$.

For the separate assertion of infinitely many primitive-root primes,
Hooley requires the family
\[
 K_n=\mathbb Q(\zeta_n,p^{1/n}),\qquad n\text{ squarefree},
 \qquad \zeta_{K_n}(s).
\]
The theorem was checked in \S7, printed p.~47, of the Russian
translation of his 1967 paper. The fields are generally nonabelian
over $\mathbb Q$, so the stipulated Dirichlet hypothesis does not
supply this family. The stipulated all-field Hecke hypothesis does.
Its arithmetic conclusion gives no theorem about intersection with a
curve-dependent theta-support set. The universal primitive-root theta
condition is in fact refuted by Finding A06. That example has
$p$-rank zero, and leaves the positive-$p$-rank restricted problem
unresolved. It is not a counterexample to A1.

The elementary density obstruction also remains useful: when
$\ell\equiv1\pmod8$, both $2$ and $-1$ are squares, so
$\langle2,-1\rangle$ is a proper subgroup of
$\mathbb F_\ell^\times$. A positive-density support set can therefore
be disjoint from the desired arithmetic condition. For torsion order
$\ell^a$, primitivity modulo $\ell$ also does not automatically
give primitivity modulo $\ell^a$.

A1 B.3 and successor A.11 instead use finite-field $L$-polynomials:
\[
 L(C,\chi,T)=\det(1-\operatorname{Frob}_qT
                  \mid H^1(C_{\bar k},\mathcal L_\chi)).
\]
Here $C/k$ is proper and pointed by $P\in C(k)$, $k=\mathbb F_q$,
$a:J_C(k)\twoheadrightarrow A$ is a finite prime-to-$p$ quotient,
$N=\#C(k)$, and
$N_b=\#\{R\in C(k):a([R-P])=b\}$.
Write $S_\chi=\sum_{R\in C(k)}\chi(a([R-P]))$.
The nontrivial Lang eigensheaves remain geometrically nontrivial
on $C$, since its Abel map surjects on the geometric abelian
prime-to-$p$ fundamental group. These eigensheaves have $2g-2$
eigenvalues of absolute value $\sqrt q$. The resulting pair estimate is
\[
 \sum_bN_b(N_b-1)
 \le {N^2\over|A|}
   +{|A|-1\over|A|}(2g-2)^2q-N.
\]
Deligne's \emph{Weil I}, formula (1.5.1) and Theorem (1.6), are
unconditional finite-field inputs. Neither number-field GRH hypothesis
supplies the still missing quotient or prescribed-point existence in
A1 B.5 and successor A.15. Successor A.38 explicitly distinguishes
closed-point Chebotarev from prime-order theta avoidance on an entire
nonzero cyclic orbit.

The least-prime results A1 B.27, B.30 and B.34, and successor
A.19, A.24, A.31, A.72, A.74, A.91--A.93 and A.95, concern specified
curves or explicitly bounded finite censuses. They use finite
Frobenius/Cartier calculations and finite exclusions, not analytic
least-prime theorems. A1 B.34 is a \emph{field-first} minimum, not a
least prime over arbitrary extensions. A1 B.32 and successor A.96
retain unresolved cases. A1 B.4 decides a fixed prime over all
constant extensions; it does not bound the search over all primes.
Successor A.16 rules out a bound depending only on $(p,g)$.

\subsection{Other false positives and verification scope}

LM8 A.32--A.33 use Serre's qualitative non-CM open-image theorem;
their commutator Zariski-density conclusion was already unconditional.
The recovered proofs explicitly identify
$F_a/\mathbb Q:y^2=x^3-3x+a$, for $a=6,-6$, with $j=-216$;
the later choice of a number field $L$ is a restriction of this
representation. Thus there is no unresolved descent condition.
Finding G06 applies to these actual factors, for which
$\operatorname{rad}(2N_{F_a})=6$. Under $\mathrm H_M$, the mod-$\ell$
image over $\mathbb Q$ is $\operatorname{GL}_2(\mathbb F_\ell)$ for
$\ell>964\log6+5760$, and its restricted image over $L$ has index
at most $[L:\mathbb Q]$. Full image over $L$ is not asserted.
Independently of $\mathrm H_M$, the cited Lombardo theorem over
$\mathbb Q$ and the same restriction argument give
\[
 [\operatorname{GL}_2(\widehat{\mathbb Z}):\rho_{F_a,\infty}(G_L)]
 <[L:\mathbb Q]\,\exp(10^{21483})
       \max\{1,h_F(F_a)\}^{4.8\cdot10^{10}}.
\]
Here $h_F$ is the stable Faltings height in the cited theorem's
convention. These refinements do not change the original rows' status.

LM8 A.46--A.48 use Faltings's crystalline comparison; D.25 uses
Faltings's finiteness of rational points. Neither is a finite
Faltings--Serre trace-comparison test. The flagged SA07 F.3.6 excerpt
instead refers to a Hom theorem, and its full proof is missing.

The primitive roots in LM8 A.54 are generators of the two fixed
finite fields of orders $41$ and $61$; A.55 chooses a root of unity;
D.93 uses that $2$ generates $\mathbb Z_5^\times$. D.31 uses the
Artin braid action. The phrase ``vari\'et\'e \'el\'ementaire d'Artin''
in scan 037 is geometric, not a primitive-root conjecture.
GT C.1.21 takes the least \emph{index} with
$v_\ell((J+4)!)\ge N$. Its displayed definition is not a least-prime
problem, and its full proof remains unavailable.

The retained B.23 conductor enumeration and independent Bezout verifier
were rerun successfully. The latter recomputed the division-polynomial
compositions and the complete Bezout identities at conductors $191$
and $229$, after a missing runtime dependency was supplied. This is a
rerun of an existing independent implementation, not a newly written
one. A separate new elementary implementation independently checked
the six primes, multiplicative orders and all finite differences used
in A06. These events are labelled ``computation or formal certificate'';
targeted source reading is labelled ``used and rechecked in this run''
with the precise scope stated. All events are dated 2026-11-21 and name
GPT-6 Astra Pro; no unavailable effort setting is invented.

\paragraph{Unrecovered sources.}
Exact checkpoint and register names, timestamp variants, and broad
title searches were attempted for the A1-successor and GT returns.
No matching full source was recovered, and no nested matching archive
was present in the recovered bundles. The documented Personal Context
route returned ``Personal context is disabled''; it was blocked,
not a zero-result search. The remaining concrete route is to provide
the two original checkpoint archives or their full registers and papers,
or an exact identity for a differently titled upload.
\texttt{SOURCE\_RECOVERY\_MANIFEST.tsv} retains the exact recovered
identities, creation dates, sizes and SHA-256 hashes without signed
download URLs. The ledger's coverage statement must remain
$800/1190$ registers and $591/1190$ available source blocks until those
missing sources are supplied.

\paragraph{Proof of the new scope assertions R01--R03.}
For R01, the retained inventory parses all original flag records and
all recovered CSV registers, counts distinct run/statement identifiers,
and locates the numbered source blocks. Every claimed recovery has a
source hash; the nine LM7 block identifiers are separately listed.
This proves availability counts, not correctness of every source proof.
For R02 the full A.69 text spells out both genus and rank formulas,
while the other flag literally reports the disabled GRH option.
For the finite-prime claim in R03, Frobenius acts on primitive
$\ell$th roots as multiplication by $p$, so their minimal-polynomial
degrees are exactly $\operatorname{ord}_\ell(p)$. It follows that
$\Phi_\ell$ is irreducible precisely in the stated case, without
producing such a prime. For the displayed pair estimate, with the explicitly defined Lang
eigensheaves geometrically nontrivial, character
orthogonality gives $\sum_bN_b^2=|A|^{-1}\sum_\chi|S_\chi|^2$.
The trivial character contributes $N^2$ and every nontrivial one
contributes at most $(2g-2)^2q$ by the cited finite-field trace bound.
Subtracting $\sum_bN_b=N$ gives the formula. The stronger prime
existence and geometric implications are treated, with full proofs
or exact cited hypotheses, in A02--A06 and G01--G05.
\clearpage
\part{Unchecked inputs and remaining gates}
\section{Unchecked inputs and remaining gates}
\label{sec:remaining}

The following entries are deliberately outside the checked-results
part. They are not assigned a proof status by analogy with a nearby
proved result.

\paragraph{U01: the requested Hooley-based proof.}
The strongest primitive-factor reduction is now A46: in every positive
characteristic and genus at least two it suffices to construct a pointed
odd cyclic etale cover, of degree prime to the characteristic, whose
actual primitive factor has no rational characteristic-primary torsion
over the specified finite field. A46 removes the previous elliptic
qualification and does not require geometric rank zero or a constant
deck group. Universal existence of such a cover is still unproved.

A02 proves the arithmetic consequence of M. A03 proves the exact
prime-degree geometric implication, and A11 weakens its sufficient
condition to an odd prime power $n$ with
$\langle2,-1\rangle=(\mathbb Z/n)^\times$ and an exact-order-$n$
theta point. A15 and A18 weaken this further: in characteristic two and every
genus at least two, a cyclic cover of any odd degree with primitive
factor of geometric $2$-rank zero suffices. In genus two this is
equivalent to theta meeting every primitive doubling orbit of the
defining cyclic torsion subgroup. A18 removes the exceptional
elliptic double quotient by an exact Deuring--Shafarevich and
conductor-rank calculation.
A04 refutes the inference from
two independent density statements, and A06 refutes the universal
joint primitive-root theta assertion itself. A05 proves A1 by a
different route. No claim is made here that a proof of A1 through
Hooley has been found. A strategy restricted to positive-two-rank
genus-two curves still requires a curve-dependent incidence theorem.
A07 refutes the positive-rank version over finitely generated fields
with transcendental moduli, not over finite fields. A08 shows that
no fixed finite prime list works over all finite fields; in the
ordinary case it excludes any fixed finite list of odd degrees.
A09 supplies an exact degree-$13$ rank-one witness. A13 supplies an
ordinary example of prime degree $4159$, with $2$ of index two in
the unit group. A16 supplies a rank-one example of degree $15$
whose full new factor has positive $2$-rank but whose primitive
factor has rank zero. A10 excludes order-three theta incidence on every nonordinary
genus-two base. None of
these results interchanges the finite-list and all-primes quantifiers;
in higher genus it requires nilpotence of every relevant Frobenius
orbit composite, not merely one nonordinary arrow.
A18 now permits any odd conductor with geometric primitive rank zero.
A19 makes all exponents for each specified finite odd prime support
an exact finite search. A21--A22 show that neither generic
positive rank nor any fixed finite support forces success, even
when mixed orders are allowed. A23 refutes universal success of
the older large-prime quotient inequality over constant extensions.
The remaining existence assertion must therefore allow the prime
support to depend on the particular finite-field curve, or use a
separating construction outside these sufficient tests.

A26 weakens the separate rational quotient criterion using even
augmentation. A27 nevertheless gives a rank-one curve over $\mathbf F_4$
where its numerical hypothesis fails for every extension and every
primitive-two prime. A28 recalls the recovered canonical-class criterion,
and A29 strictly improves its prime-level detection: a nonzero
$r$-primary component of $[K_C-2P]$ suffices if
$\operatorname{ord}_r(p)$ is even, with a cyclic $r^a$ cover and a
prime-to-$p$ rational top Prym. This is unconditional and needs no
prime-distribution theorem. The actual-curve strictness example in A29
shows why passage to prime powers matters even when all prime-level
canonical detectors fail. The residual points have canonical class
supported on the characteristic prime and odd primes of odd
multiplicative order, with the prime-two issue separate in odd
characteristic. In characteristic two this includes every Weierstrass
point and the purely two-primary non-Weierstrass examples of A09.
Constant extensions preserve that primary support. Other covers may
still work; no absence conclusion follows from this residual condition.

A30 refutes a proposed direct higher-genus extension of the
``every detecting quotient is good'' conclusion: a genus-three
rank-one curve has a detecting constant cyclic cubic quotient whose
relative class number is even. It does not exclude other good
quotients. The proposed unramified double-cover repair is unavailable
to the geometric prime-to-two fundamental group and does not preserve
the descending detector. A31 shows that the rank-one curve of A27
nevertheless has an exact signed-doubling theta witness of order seven.
All six of its original rational points permanently miss A29's
canonical detector, yet the geometric witness can be pointed after
a finite extension. Thus the sufficient tests have distinct limits;
their failure is not a nonexistence result.

The strongest new reduction of the characteristic-two genus-two
\emph{point-section problem} is A33: it suffices to prove the assertion
for ordinary genus-two curves with a Rosati-invariant integral
$\mathcal O_{\mathbf Q(\sqrt5)}$ action defined over the finite base field.
A cyclic quintic cover either has geometric new rank zero or gives such a
quotient; two reflection maps act on one common lifted collision.
This does not prove the separate literal primitive-cover problem.
A34 offers a further sufficient cover condition: for a constant pointed
odd cyclic cover of degree $n$, its actual rational two-primary Prym
exponent $M$ may be nontrivial if $n>(2M-1)^2$. An explicit degree-nineteen
cover after a degree-seven constant extension has $M=2$ and even relative
class number. The uniform Weil replacement for $M$ gives no improvement
on the older size test; the actual exponent or another geometric input
still has to be controlled.

A35 makes the purely two-primary canonical residual set a finite set:
all of its points lie over the field fixing geometric $J[4]$, whose degree
is at most six. For the ordinary real multiplication targets of A33 the
degree divides six. This does not treat nonzero odd-primary components
supported at primes where two has odd multiplicative order. A32 also
exhibits an ordinary twelve-dimensional new factor with odd rational
order over $\mathbf F_{256}$, losing that property exactly at extension
degrees divisible by seven. Geometric rank zero, rational unit order,
and bounded rational exponent therefore remain distinct sufficient tests.
The higher-genus and other-characteristic requirements of full A1 have
not been reduced to this genus-two real multiplication locus.


The real multiplication restriction is now verified to be an actual
proper class: the explicit ordinary curve (DRP.1) has geometric
endomorphism field with Rosati-fixed field $\mathbf Q(\sqrt2)$,
whereas the A13 curve has an integral Rosati-$\mathcal O_5$ action.
A37 closes the fixed-finite-support proposal inside this precise
smaller class: for every finite set of odd primes, including five,
there are densely many ordinary real multiplication curves with no
nonzero geometric theta point supported on those primes, with all
powers and products allowed. This does not exclude a prime depending
on the curve. A36 separately refutes the forced-five canonical
shortcut on an actual reflection target; its rational Jacobian order
is $2^{10}3^4$. The exact inhomogeneous constants have no five-primary
part there. A39 below supplies a three-primary detector on that particular target.
A universal construction for prescribed transferred points is still required.

A38 further rules out the unrestricted claim that ordinary
integral real multiplication forces a useful canonical prime at every
non-Weierstrass point: all $126$ such points of the A13 curve over
$\mathbf F_{128}$ have canonical class of exact order $4159$, and
$\operatorname{ord}_{4159}(2)=2079$ is odd. This curve still has the
geometric theta witness of A13. A40 now proves actual image membership for each prescribed geometric
point after a finite constant extension, using an inverse dihedral
construction for every ordinary real multiplication curve. Thus broad
canonical forcing fails even on actual images. A41 strengthens this to simultaneous failure at every lift and all
five reflection images of the same quintic cover, with at least eleven
split source points over $\mathbf F_{128}$. A43 proves that these
sources have two-rank one and that both of their Weierstrass points are
rational, split, and have only such bad canonical reflection images.
Thus these two qualifications do not repair the fixed-cover argument.
Whether a suitable choice of a different cover always succeeds remains open.

A39 gives a positive exact test on one earlier example: both
Weierstrass source points of A27 over $\mathbf F_{256}$ have useful
three-primary canonical detectors on the A36 quotient. A complete
$160\times155$ Riemann--Roch matrix is independently verified by
binary rank $1240$, and an explicit function identifies the image
where the canonical class is zero. Composite etale degrees $15$ or
$45$ suffice there; their geometric normal closures have odd order.
This example does not provide a universal supply theorem.


A42 makes the primary-projection alternative quantitative using one
fixed Lang cover. On an infinite collection of constant extensions, a
prescribed component has asymptotic fraction $1/\#G_S$ with an explicit
square-root error. In particular, after rationalizing a nonzero useful
odd-primary class, canonically detectable points are Zariski dense on
the Raynaud theta curve. The counting cover may have characteristic-primary
degree; it is not itself an allowable prime-to-characteristic
neighbourhood. Each resulting canonical detector is supplied separately
by A29. Density still does not put a good point into a prescribed finite
reflection fibre, and uncontrolled primary components prevent deduction
of whole-order primitive-orbit coverage. A41 and A43 show that this
distinction occurs on actual quintic transfers.


The special automorphism used in A43 is not a universal replacement
for cover selection. Its supplement shows that such an order-three
automorphism forces a nonzero order-three point on the actual Raynaud
theta divisor. A37 supplies ordinary real multiplication targets
without any such point, even geometrically. Also, varying A40's inverse
parameter changes the source curve, and its ramified degree-two
reflection map does not give a reverse admissible lifting operation.

A44 tests a controlled constant-field tower. For an abelian surface over
a field of characteristic two, every eligible odd prime occurring in
the tower of degrees $2^m$ divides the explicit nonzero integer (CT.2).
All points supported on those primes lie in the single finite field
of (CT.5), including arbitrary primary powers. At the same time,
the Lang lower bound and exact eventual two-primary growth prove that
there are infinitely many odd theta points in this tower. Thus
unbounded abundance alone does not supply unbounded eligible orders.
The finite candidate set need not be empty: (CT.15)--(CT.16) give
four exact order-three witnesses on A13. Universal success within that finite set is not settled by A44. Its
odd auxiliary-prime supplement proves that every compatible new order
has at most one new prime and a cofactor from an explicit finite list.
Poonen's bounded-index subgroup may contain multiple new primary
components, so that theorem does not supply the required prime-power
incidences. The supplement leaves their existence open.

A45 permits nonconstant geometric deck groups in the actual-exponent
criterion over the original pointed field. It proves separation whenever
$2\le g\le p$ and $n>(2M-1)^2$; in larger genus a collision forces a
separable elliptic map of degree $p^a\le g-1$, totally ramified at both
points. Primitive factors preserve their exponents at their own minimal
deck fields, but proper factors may grow at the full deck field. The
exact degree-$39$ cover has exponents $2$ and $16$, showing strictness
for this specified cover. No universal bound or supply follows.

A47 further allows a positive actual primitive exponent $M_n$ with
$n>(2^{\omega(n)}M_n-1)^2$, in the range $2\le g\le p$ or after
excluding its precise wild elliptic obstruction. For composite degree
this criterion can ignore large proper-conductor torsion, but the factor
$2^{\omega(n)}$ means that it is not uniformly stronger than A45.
A universal supply theorem for either sufficient inequality is still
missing. The literal rank-zero and rational-unit cover questions are
not settled by allowing a bounded positive exponent.

A48 supplies good unpointed twists of every genus-two cyclic cover
with constant deck group. This does not supply the specified pointing.
Exactly one twist is split above the prescribed point; every other twist
becomes isomorphic to that pointed model over an extension in which the
fibre is rational. Any original rational primitive characteristic-primary
torsion therefore persists. The affine-Frobenius calculation proves the
same obstruction for nonconstant deck groups. Varying the geometric cover
remains possible. The A47 supplement also proves that even a rounded
coarse Weil exponent bound cannot certify its numerical criterion.

A48's pullback supplement also shows why a nonordinary-cover iteration
does not repair an already failing cyclic cover: under a geometrically
connected prime-to-$p$ base change retaining its conductor, the old actual
primitive $p$-primary group injects into the new one over every common
finite field. Its geometric rank and rational exponent cannot decrease.
Raynaud's simultaneous pulled-back theta incidences do not imply
nilpotence in the larger coherent character spaces. Fresh cyclic covers
on the new base remain outside this obstruction.

The geometric A48 corollary realizes this distinction on actual curves:
after an odd Raynaud cover of degree $d=5\cdot3^a$, $13\le a\le65$,
all twelve primitive arrows of a cyclic degree-thirteen cover are singular,
while its geometric primitive two-rank lies between $12$ and $12(d-1)$.
This strengthens the evidence against the inference already refuted in A04;
it is not a new primitive-unit witness.

A49 makes the finite-quotient route exact. For an odd quotient killing
$[K_C-2P]$ whose exponent has unit group generated by $p$ and $-1$,
the point-image augmentation element is nonuniform exactly when some
cyclic quotient has a good actual primitive factor over the original
pointed field. The full odd quotient may succeed when a prime-level
size test fails. The two retained examples already have other witnesses,
so they do not establish a strict reduction of the full previous
universal existence gate. The quotient and nonuniformity must still
be supplied for each prescribed curve and point.

The variable-quotient supplement does supply nonuniformity when the
conductor is unrestricted. After fixing any finite excluded prime support
containing the canonical support, suitable constant extensions keep the
projection kernel fixed. The genus-two difference bound then supplies a
nonzero singleton coset and an empty coset. This gives one unit character
and its signed characteristic-power orbit, with pairwise coprime
conductors available. It does not supply the signed-unit condition or
control the other faithful character orbits of that conductor.

A50 checks whether a partial good packet can be used directly. A
cyclotomic norm bound forces integral projector coefficients to grow with
the required characteristic-primary precision. An exact abstract
$(\mathbf Z/8)[C_{31}]$ module has a missing ten-character signed packet
but minimum nonzero integral annihilator coefficient norm eight, too
large for the displayed degree comparison. This is not a geometric
counterexample. Neither this module calculation nor the inspected
nonabelian tower literature supplies the remaining actual cover bound.

A51 computes the exact norm-compatible point-difference orders under
prime-to-charac\-teristic refinement. The old torsion splits off, while a
new norm-kernel component remains. Norms to every proper cyclic quotient
kill the faithful product-conductor component, so the known packet
constraints do not control it by those maps. In a specified rank-one
tower, integral semilinear descent bounds the number of rational
characteristic-primary generators uniformly, but supplies no exponent
bound. Both arbitrary-exponent examples are explicitly abstract.

The full-fundamental-group point-section argument in
\cite[\S3, Theorem~3, pp.~134--135]{Faltings840} uses a level annihilating
the whole finite Jacobian group. Restricting its construction to levels
prime to the characteristic leaves precisely the characteristic-primary
residual. The inspected proof therefore does not repair this gate.
The alternative A05 proof retains its separate cuspidal/configuration
input.

The A51 supplements check two further repairs. A fixed nonzero
rational determinant in the character parameter would bound the actual
exponents, if the necessary geometric identification were proved.
Completed coefficients and finite matrix size alone do not do so.
The inspected skew-abelian tower result uses the tower prime, a growing
return product, and an integral preliminary statement whose natural-image
version fails for multiplication on an elliptic curve. Its separate
abstract results are not thereby refuted. In genus two and characteristic
two, full deck-Frobenius transitivity actually makes rational primitive
two-torsion nonzero whenever the primitive factor is ordinary. Its return
valuation remains the unknown quantity.

The coefficient hypotheses in the inspected Fourier--Mellin family
construction \cite[Theorems~1.3,1.5 and Definition~4.4]{EsnaultKerz2021}
and the unit-$L$-function result
\cite[Corollary~1.8]{EmertonKisin2001} do not supply that uniform
cross-prime determinant. In the latter result, the unit correction concerns
the $L$-function variable. It is not a rationality statement in the
varying finite-order character parameter.

For the iteration attempt, A51 proves that all geometric pointed
quintic/reflection choices have representatives over the constant extension
with degree support $\{2,3,5,13\}$, and gives an explicit common field at
each finite depth. That support does not obstruct good canonical points:
on every fixed target, a fixed Lang-cover count supplies densely many
points with a nonzero three-primary canonical component inside the same
extension. Incidence with the prescribed iterated orbit remains unproved.

The weighted-Weil supplement constructs actual ordinary principally
polarized abelian varieties with the transitive polynomial shape and one
rational characteristic-primary generator, but exponentially large exponent.
It does not realize the required cyclic deck action or Jacobian cover.
A52 permits any integral annihilator with nonzero faithful evaluation and
small enough coefficient norm. Its fixed-field lattice bound depends on
the actual ideal index $W_r$, which remains unbounded by the proved inputs.
In the fully transitive case it recovers exactly the earlier exponent test.
The actual orbit-module norm lower bound and integral trace descent exclude
an improvement from rational generator counts or inflation of an older
principal relation alone. No actual witness outside the full preceding
union of criteria, and no universal $W_r$ bound, is proved.

The final actual-point calculation gives a nonzero regular logarithmic
differential of trace zero when the compatible point-class order increases.
The two forced wild fibres cancel the degree term in the elementary
Riemann--Hurwitz lower estimate, so neither calculation bounds the order.
Finitely many good signed packets can also be preserved during further
variable-conductor supply, by a common arithmetic progression of constant
degrees. On an actual common cover the norm operators still act block by
block. Even all line quotients of $(\mathbf Z/\ell)^2$ leave a missing
signed orbit inside one rational cyclotomic block uncontrolled by the other
line norms. No universal bound against coefficient cancellation, and no
new small operator on the actual point class, is claimed.

The approaches actually tested were: fixed-base Hooley density;
Goldfeld's average over bases; Lenstra's fixed finite Chebotarev
constraints and positivity; prime-to-characteristic torsion orders
from Pop--Sa\"idi; a single singular Frobenius arrow in higher genus;
scalar-Frobenius packet counts; generic ordinary and fixed-rank
Prym theorems with finite-type specialization; a direct degree-$13$
theta certificate; prime-power character blocks; primary torsion
projections on theta; the Lang quotient and Weil bound; growth of
fixed-prime primary torsion; the composite primitive projector;
an exact order-$15$ search and certificate; the Hooley--Lenstra
family $n=3q$; the elliptic-exception rank calculation; finite
prime-support torsion via Frobenius translation orbits; the
all-three-primary obstruction; constant twists with an unspecified
split base point; the parity augmentation test; its all-extension rank-one obstruction;
inverse-character canonical detection and its prime-power refinement;
two-reflection transfer to the ordinary real multiplication locus;
actual rational Prym exponent with a degree-nineteen certificate;
the complete two-primary theta locus through geometric four-torsion;
fixed Lang covers with prescribed primary components; the exact
rank-one and two-Weierstrass analysis of the inverse family;
the attempted extension of its order-three automorphism;
finite eligible-prime support in a characteristic-power tower;
all-conductor Fourier inversion on an odd rational quotient;
nonconstant-deck geometric towers and the scope of class-number stability;
weights and principal polarization; fixed-field annihilator lattices;
trace relations and exact orbit-module norm bounds;
and the independent cuspidal route.
Their respective endpoints are recorded in A02--A52 and the
retained B.28--B.29 pointing audit. Poonen's
projection theorem supplies an order $\ell^r m$, with uncontrolled
prime factors of $m$; multiplying by $m$ moves theta, and passing to
an intermediate cover leaves proper-conductor character blocks.
Those are failures of those two strategies. A15 bypasses the
proper-conductor factors by a product projector on its stated
geometric locus, and A18 removes its geometric exception in
characteristic two. What remains is the same-subgroup orbit-coverage
condition, or the stronger single-point unit-group condition.
A17 gives density $\mathcal A/2$ to the candidate primes
$q\equiv5\pmod8$ for $n=3q$, conditionally on M, without proving
theta incidence for any particular curve. A12 gives an unconditional
curve-dependent primary-torsion test; A14 proves that enlarging the
constant field at one fixed prime cannot eventually make that
particular sufficient test succeed. Neither statement is a theorem
of nonexistence of theta torsion. The recovered Successor A.10
rational prime-cover criterion is separately distinguished in A14.
A19 now makes all primary exponents a finite exact search for
each prescribed finite odd prime support. A20 proves that even
all powers of three fail on the rank-one curve of A09, which
nevertheless has its order-thirteen witness. The prime supports
remain unbounded; no theorem proves that one of these finite
searches succeeds for every positive-rank curve.
The B.28--B.29 rational-twist alternative supplies a good split
point without prescribing it. Its good-point density estimate
does not force such a point into an arbitrarily prescribed fibre,
and extending a good twist can destroy its rational torsion
condition. These failures are recorded separately from geometric
primitive-rank-zero incidence.
The cuspidal route
settles point separation, while the scalar route refutes the
universal auxiliary existence assertion. None of the other routes
silently fills the curve-dependent incidence gap.

\paragraph{U02: the literal primitive-cover problem.}
ANABELIAN\_A1 B.24 allows the degree-one cover with convention
$P_1=J_C$. The B.23 counterexample has $p$-rank zero and is therefore
covered by that convention. The problem for positive $p$-rank
remains open in the recovered run. Similarly ANABELIAN\_A1 B.5 and
B.32 ask for prescribed-point rational class-number witnesses. A29
now settles the stated nonzero canonical-primary component cases,
while their universal existence remains open.
A bound for a prime in a number-field Frobenius class does not
produce the required covering geometry or local unit conditions.
No assertion is made that these separate questions follow from A05.

\paragraph{U03: complete source coverage.}
The two missing registers account for 390 full statements. Only
47 distinct A1-successor IDs and 50 distinct GT-tower IDs are visible
in the supplied summaries and flag excerpts. The missing LM7
proof files account for another 209 proof blocks. Exact-title,
variant-title and content searches recovered four relevant archives,
but not the two missing sources. A final prior-context recovery
attempt was unavailable because that retrieval service was disabled.
The next concrete input is the complete
\texttt{ANABELIAN\_A1\_SUCCESSOR} register/paper and the complete
\texttt{GT\_TOWER\_CHECK} register/paper, together with the omitted
LM7 proof sections. Their absence is an evidentiary gate, not a
mathematical theorem beyond GRH.

\paragraph{U04: infinite towers and unspecified finite data.}
For a recovered row with finite affine monodromy, G01 gives a
prime bound after the actual finite image, cocycle, precision,
and excluded primes are fixed. It does not bound the precision
needed to separate an infinite object, the height of a Hilbert
specialization, or a Baire-category construction. A finite test of
full adic semisimplifications is available under the hypotheses of
G03, with exact computable traces as an additional algorithmic input.
Its applicability to an unspecified row must be checked separately.

\paragraph{U05: scope of the effective open-image bounds.}
The former descent gate for LM8 A.32--A.33 is closed by their actual
models $F_{\pm6}/\mathbb Q$ and the checked refinement G06.
The conditional Mayle--Wang bound applies over $\mathbb Q$; after
restriction to the proof's number field $L$, G06 proves an index
bound $[L:\mathbb Q]$, not necessarily full image.
G-LOMBARDO is an unconditional explicit adelic index bound over
arbitrary number fields. The hypotheses must still be checked for
any other factor; CM factors have CM monodromy rather than full
$\operatorname{GL}_2$ monodromy. No corresponding general abelian-variety
bound has been asserted.

\paragraph{U06: boundaries of verification.}
The complete new proofs were read, and every included computational
certificate reports the exact execution actually performed. No
formal proof-assistant verification of this paper, reconstruction
of all earlier profinite elements, replay of every historical
computer search, or verification of Masullo's entire draft is
claimed. The check table distinguishes whole-proof reading,
independent arithmetic reimplementation, retained-verifier replay,
source-statement checks and carried background. A same-model second
or third reader is not presented as verification by a different
model or by a theorem prover.

\clearpage
\begin{thebibliography}{99}

\bibitem{Alpoge}
L.~Alp\"oge, \emph{Remark}, supplied four-page note dated 20 November
2026, especially step~20, p.~4; supplied
\texttt{ALPOGE\_Main.lean}, \texttt{ALPOGE\_Statements.lean}, and
\texttt{ALPOGE\_LEAN\_SCAN.tsv}. \url{https://alpo.ge/remark.pdf}.
Accepted as a cited theorem by the task's explicit convention.

\bibitem{Masullo}
A.~M.~Masullo, \emph{The Music of the Primes}, supplied 940-page draft,
20 November 2026. Claimed scope: \S1.2; trace formulas:
Theorems~12.5.1.1 and 13.3.1.1; unfinished intended main-statement
heading: Theorem~13.3.1.2, p.~522. All consequences in this report
are conditional on the draft's stated claim M.

\bibitem{Cote149}
A.~Grothendieck, anabelian manuscript, cote~149, supplied
transcriptions \texttt{ARCHIVE\_149\_scan033.tex} through
\texttt{ARCHIVE\_149\_scan038.tex}. No original manuscript images
were supplied or inspected in this run.

\bibitem{Intake}
\emph{INTAKE\_ANABELIAN\_149\_FINAL\_20261114}, supplied intake,
14 November 2026; especially the A.C2.5/A.C2.7/A.C2.19/A.C2.20
review, \S3.3(3) and \S5(a), and the D.30--D.31 dependency discussion.

\bibitem{A1run}
\emph{ANABELIAN\_A1}, checkpoint~03, recovered complete manuscript
\texttt{ANABELIAN\_A1\_paper.tex}, 121-row register and supporting
sources. Exact archive identity and SHA-256 hashes appear in
\texttt{SOURCE\_RECOVERY\_MANIFEST.tsv}. The proof of A1 uses
A.2--A.5, A.14--A.19 and A.30--A.33; the explicit curve is B.23.
Source labels alone are not treated as proof checks.

\bibitem{ASrun}
\emph{ANABELIAN\_SUCCESSOR}, recovered final manuscript
\texttt{ANABELIAN\_SUCCESSOR\_paper.tex} and 157-row register;
especially A.17, A.38 and the prime-to-characteristic covering
criteria. Archive identity and source hashes are retained.

\bibitem{LM7run}
\emph{LM7\_PUSH}, checkpoint~09, 13 November 2026, 218-row register.
Recovered from the user's related
\texttt{DELTA\_2-3\_reciprocity\_of\_frames\_on\_GT} bundle,
with only the proof blocks E.1--E.7, E.58 and E.65.
The complete register is available; other proof blocks are not.

\bibitem{LM8run}
\emph{LM8\_PUSH}, checkpoint~17, 16 November 2026,
304-row register and complete main proof sources. Exact section
filenames and statement IDs are recorded in the row audit and
\texttt{GT\_LEVEL\_SCOPE.tsv}. Earlier source computations are not
claimed rerun unless an event records it.

\bibitem{SGA1}
A.~Grothendieck et al., \emph{Rev\^etements etales et groupe fondamental
(SGA~1)}, Lecture Notes in Mathematics 224, Springer, 1971;
recomposed edition, Documents Mathematiques 3, SMF, 2003.
Exact Expos\'e/statement locators are given in the text.
\url{https://arxiv.org/abs/math/0206203}.

\bibitem{WeilI}
P.~Deligne, \emph{La conjecture de Weil. I}, Publications
Mathematiques de l'IHES \textbf{43} (1974), 273--307,
formula~(1.5.1), Theorem~1.6 and Frobenius convention~(1.15.1).
\url{https://www.numdam.org/item/PMIHES_1974__43__273_0.pdf}.

\bibitem{MilneLEC}
J.~S.~Milne, \emph{Lectures on Etale Cohomology}, version~2.21,
22 March 2013; Example~7.9(a)--(b), Remark~7.10,
Proposition~14.2, Theorem~14.9 and Remark~14.10.
\url{https://www.jmilne.org/math/CourseNotes/LEC.pdf}.

\bibitem{Conrad}
B.~Conrad, \emph{Chow's $K/k$-image and $K/k$-trace, and the
Lang--Neron theorem}, author version dated 17 June 2006,
Theorem~2.1. \url{https://math.stanford.edu/~conrad/papers/Kktrace.pdf}.

\bibitem{ST2009}
M.~Sa\"idi and A.~Tamagawa, \emph{A prime-to-$p$ version of
Grothendieck's anabelian conjecture for hyperbolic curves over finite
fields of characteristic $p>0$}, Publications of RIMS
\textbf{45} (2009), 135--186; Remark~1.2,
Definitions~1.3, 1.5, 1.10, Lemma~1.7 and Proposition~1.8(vi).
\url{https://doi.org/10.2977/PRIMS/1234361157}.
Primary published-version PDF checked at
\url{https://empslocal.ex.ac.uk/people/staff/ms220/Site/Publications_files/publRIMS.pdf}.

\bibitem{PopSaidi}
F.~Pop and M.~Sa\"idi, \emph{On the specialization homomorphism
of fundamental groups of curves in positive characteristic},
in \emph{Galois Groups and Fundamental Groups}, MSRI Publications
\textbf{41} (2003), 107--118; Proposition~2.1,
Propositions~2.3--2.4.
\url{https://library.slmath.org/books/Book41/files/pop.pdf}.

\bibitem{Saidi2005}
M.~Sa\"idi, \emph{On complete families of curves with a given
fundamental group in positive characteristic}, Manuscripta
Mathematica (2005), primary author PDF, Proposition~2.2 and
Theorem~2.3. \url{https://empslocal.ex.ac.uk/people/staff/ms220/Site/Publications_files/manuscripta.pdf}.
These statements concern theta avoidance and ordinarity of covers;
they do not supply the joint primitive-root theta incidence in A03.

\bibitem{Hooley}
C.~Hooley, \emph{On Artin's conjecture}, Journal f\"ur die reine
und angewandte Mathematik \textbf{225} (1967), 209--220,
\S7, concluding theorem.
\url{https://doi.org/10.1515/crll.1967.225.209}.
The original English publisher PDF was unavailable; the theorem and
field hypotheses were checked in the authorized Russian translation,
\emph{Matematika} \textbf{12}:1 (1968), 37--47,
\S7, p.~47, and in the modern numbered form below.
\url{https://www.mathnet.ru/eng/mat455}.

\bibitem{FanPollack}
K.~(S.)~Fan and P.~Pollack, \emph{Counting primes with a given
primitive root, uniformly}, arXiv:2505.05601v2, 14 September 2025,
Theorem~1.1, specialized to the fixed base~2.
\url{https://arxiv.org/abs/2505.05601}.

\bibitem{Lenstra1977}
H.~W.~Lenstra, Jr., \emph{Artin's conjecture on primes with
prescribed primitive roots}, Seminaire Delange--Pisot--Poitou,
Theorie des nombres \textbf{18}:1 (1976--1977), expose~14, 1--8;
Theorems~3.1, 4.1--4.2 and~5.1.
\url{https://www.numdam.org/item/SDPP_1976-1977__18_1_A11_0.pdf}.

\bibitem{Goldfeld}
D.~Goldfeld, \emph{Artin's conjecture on the average},
Mathematika \textbf{15} (1968), 223--226; theorem, p.~223.
\url{https://www.math.columbia.edu/~goldfeld/ArtinsConjecture.pdf}.

\bibitem{Bach1990}
E.~Bach, \emph{Explicit bounds for primality testing and related
problems}, Mathematics of Computation \textbf{55} (1990),
355--380; Theorems~2--3, p.~372, and generator consequence, p.~373;
Dirichlet scope, \S6, pp.~369--370.
\url{https://doi.org/10.1090/S0025-5718-1990-1023756-8}.

\bibitem{LLS2015}
Youness Lamzouri, Xiannan Li, and Kannan Soundararajan,
\emph{Conditional bounds for the least quadratic non-residue and related problems},
Mathematics of Computation \textbf{84} (2015), no.~295, 2391--2412.
\S1.1 explicitly distinguishes Bach's constants two and three.
\url{https://doi.org/10.1090/S0025-5718-2015-02925-1}.
For the general-subgroup bound, the corrected author version checked
is arXiv:1309.3595v3, 8~December~2016:
\url{https://arxiv.org/pdf/1309.3595v3}.
The authors' revision notice at
\url{https://arxiv.org/abs/1309.3595} identifies the corrections to
asymptotic Theorems~1.2--1.3. Theorem~1.1's displayed formula is unchanged.

\bibitem{BachSorenson}
E.~Bach and J.~Sorenson, \emph{Explicit bounds for primes in
residue classes}, Mathematics of Computation \textbf{65} (1996),
1717--1735; Theorem~5.1, p.~1729; analytic hypothesis in the
last paragraph of \S2, p.~1720; Theorem~3.2 and Corollary~3.3.
\url{https://doi.org/10.1090/S0025-5718-96-00763-6}.

\bibitem{Serre1981}
J.-P.~Serre, \emph{Quelques applications du theoreme de densite
de Chebotarev}, Publications Mathematiques de l'IHES
\textbf{54} (1981), 123--201; Proposition~6 (discriminant bound),
\S2.1, Theorem~1 (qualitative Chebotarev).
\url{https://www.numdam.org/item/PMIHES_1981__54__123_0.pdf}.

\bibitem{MayleWang}
J.~Mayle and T.~Wang, \emph{On the effective version of Serre's
open image theorem}, Bulletin of the London Mathematical Society
\textbf{56} (2024), 1399--1416; Theorems~1, 2 and~5,
Corollary~6, Lemma~7 and Proposition~12.
\url{https://doi.org/10.1112/blms.13002}.

\bibitem{MilneCFT}
J.~S.~Milne, \emph{Class Field Theory}, version~4.03,
6 August 2020; Chapter~V, Theorem~3.23, printed p.~164
(qualitative Chebotarev).
\url{https://www.jmilne.org/math/CourseNotes/CFT.pdf}.

\bibitem{Lombardo2015}
D.~Lombardo, \emph{Bounds for Serre's open image theorem for elliptic
curves over number fields}, Algebra \& Number Theory \textbf{9}
(2015), no.~10, 2347--2395; Corollary~9.3, p.~2393;
isogeny input Theorem~2.1. The alternative constants of Remark~9.4
are not used.
\url{https://doi.org/10.2140/ant.2015.9.2347}.
\url{https://msp.org/ant/2015/9-10/ant-v9-n10-p06-p.pdf}.

\bibitem{OzmanPries2019}
E.~Ozman and R.~Pries,
\emph{Ordinary and almost ordinary Prym varieties},
Asian J. Math. \textbf{23} (2019), no.~3, 455--478.
DOI: \href{https://doi.org/10.4310/AJM.2019.v23.n3.a5}{10.4310/AJM.2019.v23.n3.a5}.
Theorem~1.1(2), p.~456; Sections~8.1--8.2, p.~472;
Theorem~8.3, p.~473.

\bibitem{Poonen2005}
B.~Poonen,
\emph{Multiples of subvarieties in algebraic groups over finite fields},
Int. Math. Res. Not. 2005, no.~24, 1487--1498.
Theorem~1.8; Lemma~5.1 and the proof of Theorem~1.4, author PDF p.~6;
Lemmas~6.1--6.2, author PDF p.~7.
\url{https://math.mit.edu/~poonen/papers/multiples.pdf}.

\bibitem{Yang2021}
Y.~Yang,
\emph{Raynaud--Tamagawa theta divisors and new-ordinariness of
ramified coverings of curves},
J. Algebra \textbf{587} (2021), 263--294.
DOI: \href{https://doi.org/10.1016/j.jalgebra.2021.07.032}{10.1016/j.jalgebra.2021.07.032}.
Theorem~1.1; Definition~3.3 and Theorem~3.5;
Section~3.6.2 and Corollary~3.6.

\bibitem{PriesZhu2012}
R.~Pries and H.~J.~Zhu, \emph{The $p$-rank stratification of
Artin--Schreier curves}, Ann. Inst. Fourier \textbf{62} (2012),
707--726, Lemma~2.6, p.~712.
\url{https://www.numdam.org/item/10.5802/aif.2692.pdf}.

\bibitem{EGAIII2}
A.~Grothendieck, with J.~Dieudonn\'e, \emph{Elements de geometrie
algebrique. III: Etude cohomologique des faisceaux coherents,
seconde partie}, Publ. Math. IHES \textbf{17} (1963), 5--91.
Theorem~7.7.5(I), p.~67; Theorem~7.9.4, pp.~76--77.
\url{https://www.numdam.org/item/PMIHES_1963__17__5_0.pdf}.

\bibitem{EGAIV3}
A.~Grothendieck, with J.~Dieudonn\'e, \emph{Elements de geometrie
algebrique. IV: Etude locale des schemas et des morphismes de
schemas, troisieme partie}, Publ. Math. IHES \textbf{28} (1966),
5--255. Theorem~8.8.2(i)--(ii), p.~28; Theorem~9.7.7(ii), p.~79.
\url{https://www.numdam.org/item/PMIHES_1966__28__5_0.pdf}.

\bibitem{Stacks}
The Stacks Project, Lemma~32.10.1(1)--(3), Tag~\texttt{01ZM};
Lemma~36.32.1, Tag~\texttt{0BDN}; Lemma~30.22.1,
Tag~\texttt{07VK}. Checked 21 November 2026.
\url{https://stacks.math.columbia.edu/}.

\bibitem{Poonen2007Gonality}
B.~Poonen,
\emph{Gonality of modular curves in characteristic $p$},
Math. Res. Lett. \textbf{14} (2007), no.~4, 691--701.
DOI: \href{https://doi.org/10.4310/MRL.2007.v14.n4.a14}{10.4310/MRL.2007.v14.n4.a14}.
Proposition~2.1, Lemma~2.2, and Corollary~2.3, author manuscript p.~3.
\url{https://math.mit.edu/~poonen/papers/gonality.pdf}.

\bibitem{AnbarBeelenNguyenP}
N.~Anbar, P.~Beelen, and N.~Nguyen,
\emph{A new tower with good $p$-rank meeting Zink's bound},
Acta Arith. \textbf{177} (2017), no.~4, 347--374.
DOI: \href{https://doi.org/10.4064/aa8388-6-2016}{10.4064/aa8388-6-2016}.
Theorem~1, equation~(1), author manuscript p.~3.
\url{https://backend.orbit.dtu.dk/ws/files/123787793/p_rank_abn.pdf}.

\bibitem{MilneJV}
J.~S.~Milne, \emph{Jacobian Varieties}, author version of
12 June 2021 of the chapter in \emph{Arithmetic Geometry}
(G.~Cornell and J.~H.~Silverman, eds.), Springer, 1986, 167--212.
Section~6, Proposition~6.1, pp.~19--20;
Section~9, Proposition~9.1, p.~30, with the characteristic-primary
completion on pp.~31--32; Section~11, Theorem~11.1 and
Corollary~11.4, pp.~35--37;
Theorem~12.1 and Corollary~12.2, pp.~38--39.
\url{https://jmilne.org/math/xnotes/JVs.pdf}.

\bibitem{ScanlonVoloch1999}
T.~Scanlon and J.~F.~Voloch,
\emph{Difference algebraic subgroups of commutative algebraic
groups over finite fields},
Manuscripta Math. \textbf{99} (1999), no.~3, 329--339.
Introduction, Section~3, and reference \textnormal{[Box]}.
\url{https://math.berkeley.edu/~scanlon/papers/dag24n.pdf}.
The citation supplies literature context; the specialized
finite-support theorem in A19 is proved directly here.


\bibitem{MilneAV}
J.~S.~Milne, \emph{Abelian Varieties}, author version of
2 January 2022 of the chapter in \emph{Arithmetic Geometry},
Springer, 1986, 103--150.
Section~19, Theorem~19.1(a)--(c) and its Frobenius explanation,
pp.~40--41; Section~12, Proposition~12.9.
\url{https://jmilne.org/math/xnotes/AVs.pdf}.


\bibitem{AchterPries2008}
J.~D.~Achter and R.~Pries,
\emph{Monodromy of the $p$-rank strata of the moduli space of curves},
International Mathematics Research Notices 2008,
Art.~rnn053, 25~pp.
DOI: \url{https://doi.org/10.1093/imrn/rnn053}.
Author version:
\url{https://www.math.colostate.edu/~achter/math/prm0508.pdf}.
Theorem~4.5, author p.~14; Section~5.1 and Lemmas~5.1--5.2,
author pp.~17--18.

\bibitem{LSTX2017}
A.~Landesman, A.~A.~Swaminathan, J.~Tao, and Y.~Xu,
\emph{Lifting subgroups of symplectic groups over
$\mathbf Z/\ell\mathbf Z$},
Research in Number Theory \textbf{3} (2017), Art.~14.
DOI and published primary text:
\url{https://doi.org/10.1007/s40993-017-0078-6}.
Theorem~1 and Proposition~1(b), using the published numbering.

\bibitem{Hall2015}
J.~I.~Hall, \emph{Classical Groups and Geometry},
author notes, version 6 May 2015.
\url{https://users.math.msu.edu/users/halljo/classes/MTH913_S15/cgg_2015-05-06.pdf}.
Theorem~12.26, printed p.~173 (PDF page~179).

\bibitem{Poonen2017RPV}
B.~Poonen, \emph{Rational points on varieties},
Graduate Studies in Mathematics \textbf{186},
American Mathematical Society, Providence, 2017.
Author version: \url{https://math.mit.edu/~poonen/papers/Qpoints.pdf}.
Theorem~7.7.1(ii)--(iii), printed p.~222;
Theorem~4.3.5 and Corollary~4.4.6, printed pp.~101 and~104.

\bibitem{Stacks0245}
The Stacks Project Authors,
\emph{The Stacks Project}, Lemma~35.37.1,
Tag~\texttt{0245}, effective descent for affine morphisms.
\url{https://stacks.math.columbia.edu/tag/0245}.

\bibitem{KatzCongruence}
N.~Katz, \emph{Une formule de congruence pour la fonction
$\zeta$}, SGA~7~II, Expos\'e~XXII, Lecture Notes in
Mathematics \textbf{340}, Springer, 1973, pp.~401--438;
Proposition~2.0, Corollary~2.1, and Propositions~2.2.4--2.2.5,
pp.~406--409.
\url{https://web.math.princeton.edu/~nmk/old/congzeta.pdf}.

\bibitem{DeligneModP}
P.~Deligne, \emph{Fonctions $L$ modulo $\ell^n$ et modulo $p$},
in SGA~$4\frac12$, \emph{Cohomologie etale}, Lecture Notes in
Mathematics \textbf{569}, Springer, 1977, pp.~110--128;
\S2.1, Theorem~2.2(b), and formula~(2.4.1), p.~116.
\url{https://publications.ias.edu/sites/default/files/Number32.pdf}.

\bibitem{DelignePappas1994}
P.~Deligne and G.~Pappas,
\emph{Singularit\'es des espaces de modules de Hilbert, en les
caract\'eristiques divisant le discriminant},
Compositio Math. \textbf{90} (1994), no.~1, 59--79.
Theorem~2.2 and Sections~2.4, 2.6, 2.11--2.13.
\url{https://numdam.org/item/CM_1994__90_1_59_0.pdf}.

\bibitem{Yu2009HB}
C.-F.~Yu,
\emph{Irreducibility of the Hilbert--Blumenthal moduli spaces
with parahoric level structure},
J. Reine Angew. Math. \textbf{635} (2009), 187--211.
Theorem~1.1; full-level moduli definition immediately preceding it.
Author manuscript \texttt{arXiv:0807.0516}, pp.~1--2;
Section~6.1 and Lemma~6.1 give the corresponding
prime-adic monodromy statement.
\url{https://arxiv.org/pdf/0807.0516}.

\bibitem{Landesman2021Torelli}
A.~Landesman,
\emph{The Torelli map restricted to the hyperelliptic locus},
Trans. Amer. Math. Soc. Ser.~B \textbf{8} (2021), 354--378.
Theorem~1.2, including its genus-two assertion over
$\operatorname{Spec}\mathbf Z$.
Author manuscript \texttt{arXiv:1911.02084v3},
statement p.~1 and proof in Section~6, pp.~15--17.
\url{https://arxiv.org/pdf/1911.02084}.

\bibitem{CleryVanderGeer2022}
F.~Cl\'ery and G.~van der Geer,
\emph{Modular forms of degree 2 and curves of genus 2 in
characteristic 2},
Int. Math. Res. Not. IMRN \textbf{2022}, no.~7, 5204--5218.
Section~4, printed p.~5211, for the open genus-two Torelli image
in characteristic two.
\url{https://doi.org/10.1093/imrn/rnaa239}.

\bibitem{Stacks04GM}
The Stacks Project Authors,
\emph{The Stacks Project}, Lemma~10.153.9,
Tag~\texttt{04GM}, ``Complete local rings are Henselian,''
with the simple-root convention of Definition~10.153.1.
\url{https://stacks.math.columbia.edu/tag/04GM}.

\bibitem{Stacks00R4}
The Stacks Project Authors, \emph{The Stacks Project},
Lemma~10.128.1, Tag~\texttt{00R4} (Miracle flatness),
\url{https://stacks.math.columbia.edu/tag/00R4},
accessed 21 November 2026.

\bibitem{KumarMukamel2016}
Abhinav Kumar and Ronen E. Mukamel,
\emph{Real multiplication through explicit correspondences},
author manuscript, arXiv:1602.01924v2 (31 May 2016),
\url{https://arxiv.org/pdf/1602.01924},
Sections~2--3, pp.~5--7.

\bibitem{Crew1984}
Richard M. Crew, \emph{Etale $p$-covers in characteristic $p$},
Compositio Mathematica \textbf{52} (1984), 31--45,
Corollary~1.8, pp.~35--36,
\url{https://numdam.org/item/CM_1984__52_1_31_0.pdf}.

\bibitem{StacksRH}
The Stacks Project Authors, \emph{The Stacks Project},
Section~53.12, Riemann--Hurwitz;
Lemmas~53.12.2 and~53.12.4, Tags~\texttt{0C1D} and~\texttt{0C1F},
\url{https://stacks.math.columbia.edu/tag/0C1D},
\url{https://stacks.math.columbia.edu/tag/0C1F}.

\bibitem{Stacks023N}
The Stacks Project Authors, \emph{The Stacks Project},
Proposition~35.3.9, effective descent for modules along
faithfully flat ring maps, Tag~\texttt{023N},
\url{https://stacks.math.columbia.edu/tag/023N}.

\bibitem{RaynaudNonordinary2000}
M.~Raynaud, \emph{Rev\^etements des courbes en caract\'eristique
$p>0$ et ordinarit\'e}, Compositio Mathematica \textbf{123}
(2000), 73--88, Proposition~1(5), Theorems~2 and~17,
\url{https://doi.org/10.1023/A:1001840726893}.

\bibitem{Stacks03P8}
The Stacks Project Authors,
\emph{The Stacks Project}, Theorem~59.24.1,
Tag~\texttt{03P8} (Hilbert90 / Picard groups),
\url{https://stacks.math.columbia.edu/tag/03P8}.

\bibitem{Stacks02RT}
The Stacks Project Authors,
\emph{The Stacks Project}, Lemma~42.18.1,
Tag~\texttt{02RT} (principal divisors and pushforward),
\url{https://stacks.math.columbia.edu/tag/02RT}.

\bibitem{Faltings840}
G.~Faltings, \emph{Curves and their fundamental groups
(following Grothendieck, Tamagawa and Mochizuki)},
Seminaire Bourbaki, expose~840, Asterisque~252 (1998), 131--150,
\S3, Theorem~3 (A.~Tamagawa), pp.~134--135.
\url{https://smf.emath.fr/system/files/2017-08/smf_ast_252_131-150.pdf}.

\bibitem{ESS2022}
H.~Esnault, M.~Shusterman and V.~Srinivas,
\emph{Finite presentation of the tame fundamental group},
Selecta Mathematica (N.S.) \textbf{28} (2022), article~37.
Author manuscript:
\url{https://page.mi.fu-berlin.de/esnault/preprints/helene/141_ess.pdf}.

\bibitem{Shusterman2022}
M.~Shusterman,
\emph{Balanced presentations for fundamental groups of curves over finite fields},
Mathematical Research Letters \textbf{29} (2022), no.~4, 1251--1259,
doi:10.4310/MRL.2022.v29.n4.a13.
\url{https://arxiv.org/abs/1811.04192}.

\bibitem{KMT2017}
H.~Kodani, M.~Morishita and Y.~Terashima,
\emph{Arithmetic topology in Ihara theory},
Publications of the Research Institute for Mathematical Sciences
\textbf{53} (2017), no.~4, 629--688,
doi:10.4171/PRIMS/53-4-6.
\url{https://arxiv.org/abs/1608.07926}.

\bibitem{Asvin2023}
Asvin~G.,
\emph{On the variation of Frobenius eigenvalues in a skew-abelian Iwasawa tower},
Algebra \& Number Theory \textbf{17} (2023), no.~12, 2151--2179,
doi:10.2140/ant.2023.17.2151.
\url{https://msp.org/ant/2023/17-12/ant-v17-n12-p03-p.pdf}.

\bibitem{Sinnott1987}
W.~Sinnott,
\emph{On a theorem of L.~Washington},
Ast\'erisque \textbf{147--148} (1987), 209--224.
\url{https://www.numdam.org/item/AST_1987__147-148__209_0.pdf}.

\bibitem{EsnaultKerz2021}
H.~Esnault and M.~Kerz,
\emph{Etale cohomology of rank one $\ell$-adic local systems in positive characteristic},
Selecta Mathematica~\textbf{27} (2021), article~58,
doi:10.1007/s00029-021-00681-y.
\url{https://page.mi.fu-berlin.de/esnault/preprints/helene/134_esn_ker_hl.pdf}.

\bibitem{EmertonKisin2001}
M.~Emerton and M.~Kisin,
\emph{Unit $L$-functions and a conjecture of Katz},
Annals of Mathematics~\textbf{153} (2001), no.~2, 329--354,
doi:10.2307/2661344.
\url{https://www.math.uchicago.edu/~emerton/pdffiles/katz.pdf}.

\bibitem{Waterhouse1969AV}
W.~C.~Waterhouse,
\emph{Abelian varieties over finite fields},
Annales Scientifiques de l'\'Ecole Normale Sup\'erieure (4)
\textbf{2} (1969), no.~4, 521--560,
doi:10.24033/asens.1183.
\url{https://www.numdam.org/article/ASENS_1969_4_2_4_521_0.pdf}.

\bibitem{Diem2001Weil}
C.~Diem,
\emph{A Study on Theoretical and Practical Aspects of Weil-Restrictions
of Varieties}, doctoral dissertation, Universit\"at-Gesamthochschule
Essen, February 2001, especially Propositions~1.21--1.22, pp.~26--28.
\url{https://www.esaga.uni-due.de/f/ulrich.goertz/iem-archiv/dissertation_diem.pdf}.

\end{thebibliography}

\end{document}
